\documentclass[11pt,reqno]{amsart}
\usepackage[T1]{fontenc}
\usepackage{lmodern,microtype,amsmath,amssymb,mathtools}
\usepackage[hmargin=1.08in,vmargin=1.00in]{geometry}
\usepackage{xcolor}
\usepackage[colorlinks=true,linkcolor=blue!45!black,citecolor=blue!45!black,urlcolor=blue,bookmarksnumbered]{hyperref}
\newcommand{\codefile}[1]{\href{https://github.com/sebastienpalcoux/grothendieck-ring-classification/blob/d5750105bb5758f3ad38d2e5464dfcc12c019816/supplements/rank6/sphericality/#1}{\nolinkurl{#1}}}
\newcommand{\codecheck}[2]{\href{https://github.com/sebastienpalcoux/grothendieck-ring-classification/blob/d5750105bb5758f3ad38d2e5464dfcc12c019816/supplements/rank6/sphericality/#1}{#2}}
\newtheorem{theorem}{Theorem}[section]
\newtheorem{proposition}[theorem]{Proposition}
\newtheorem{lemma}[theorem]{Lemma}
\newtheorem{corollary}[theorem]{Corollary}
\theoremstyle{remark}\newtheorem{remark}[theorem]{Remark}
\numberwithin{equation}{section}
\setcounter{tocdepth}{1}\setlength{\emergencystretch}{1.1em}

\hypersetup{pdftitle={Spherical structures on fusion categories of rank six},pdfauthor={},bookmarksdepth=3}
\title[Spherical structures in rank six]{Spherical structures on fusion categories of rank six}
\date{}
\subjclass[2020]{18M20, 16T05}
\keywords{Fusion category, spherical structure, formal codegree, central induction}
\begin{document}
\raggedbottom
\begin{abstract}
Every complex fusion category of rank six admits a spherical structure.
For a fusion category of arbitrary rank, central induction places
the formal codegrees of its sphericalization in the field generated
by its Frobenius--Perron dimension. It also compares the ordinary
and spherical dimensions on a common family of objects. Equality
on these objects gives a fusion subcategory on which the spherical
structure descends. In rank six, bounds on the Fourier coefficients
of the reciprocal codegrees and integral trace identities for the dimension ratio
restrict the remaining dimension fields. The last quadratic case uses a
Morita dual.
\end{abstract}

\maketitle
{\centering\small\itshape
AI assisted with the research and preparation of this manuscript.
Complete human verification is pending.
\par}

\section{Introduction}
A spherical structure is a monoidal identification of each object with
its double dual for which the left and right traces agree~\cite[\S4.7]{EGNO}.
We prove the following existence theorem. The rank counts isomorphism
classes of simple objects.

\begin{theorem}\label{thm:universal}
Every complex fusion category of rank six admits a spherical structure.
\end{theorem}

In particular, every such category is pivotal. Pseudo-unitarity alone
cannot explain the theorem. Schopieray exhibited a rank-six category
whose fusion rules admit no pseudo-unitary categorification
\cite[Example~1 and Theorem~1]{SchopNP}.

\subsection*{Proof of Theorem~\ref{thm:universal}}
Write $D=\operatorname{FPdim}(\mathcal C)$ and $E=\dim(\mathcal C)$.
The intrinsic dimension $E$ is the sum of the squared norms of the
simples. If $E=D$, pseudo-unitarity supplies a canonical spherical
structure~\cite[Proposition~8.23]{ENO}. It suffices to obtain this
equality after Galois conjugation. Indeed $\mathcal C$ has an algebraic
model, the fusion rules determine $D$, and a spherical structure
transports back~\cite[Theorem~2.30 and the following paragraph]{ENO}.

There is also a direct way to obtain a spherical structure. Each
simple of $\mathcal C$ has two lifts to its sphericalization
$\mathcal S$. If a choice of one lift per simple spans a tensor
subcategory, the forgetful functor identifies it with $\mathcal C$.
Central induction places all ordinary and signed codegrees in
$K=\mathbb Q(D)$. The two dimension characters have codegrees $D$
and $E$. Evaluating them on the same induction object gives, for
every $\sigma\in\operatorname{Gal}(K/\mathbb Q)$,
the inequality of Lemma~\ref{lem:central-witness}
\[
       \frac{E}{\sigma(E)}\le\frac{D}{\sigma(D)}.
\]
Equality means that the two dimensions agree on every constituent of
the corresponding induction object. If $E/D$ lies in a proper
subfield, these constituents generate enough dimensions to recover
the larger field (Theorem~\ref{thm:proper-ratio}). The remaining
cases use bounds on the Fourier coefficients of the conjugates of
$1/D$ and $1/E$ (Lemma~\ref{lem:sextic-matrix-compression}).

The Perron character maps $K_0(\mathcal C)_{\mathbb Q}$ onto a field
containing $K$, so $[K:\mathbb Q]\le6$.
In degree one $D$ is integral and pseudo-unitarity follows
from~\cite[Proposition~8.24]{ENO}. Theorems~\ref{thm:odd-degree-descent}
and~\ref{thm:sextic-all} give sphericality in degrees three, five and
six. In degrees two and four, Theorems~\ref{thm:full-quadratic-reduction}
and~\ref{thm:quartic-final} leave four fields.
\begin{itemize}
\item In $\mathbb Q(\sqrt2)$, passage to bimodules over
$\mathbf1\oplus g$ for an invertible $g$ of order two produces two
additional codegrees. A modular bound excludes an irrational pair.
For a rational pair, the sign bimodule exchanges simples in pairs,
forcing even codegrees. A congruence for odd-index Pell numbers
excludes the last possibility
(Theorem~\ref{thm:silver-morita-end}).
\item In $\mathbb Q(\zeta_{16})^+$ and $\mathbb Q(\sqrt2,\sqrt3)$,
two characters of codegree four account for half the unit coefficient.
The resulting integral moments on the quartic factor have incompatible
residues (Theorem~\ref{thm:quartic-four-complete}).
\item In $\mathbb Q(\sqrt2,\sqrt5)$, the second Frobenius--Schur
indicator makes the ordinary and signed induction columns agree modulo
four in a suitable basis. Applying this congruence to the columns and
their squares eliminates the remaining pairs of codegrees
(Theorem~\ref{thm:forty-closed}).\qed
\end{itemize}
\smallskip\noindent
\emph{Finite calculations.} Each computation links to its code or
certificate. For example, if $z_i\ge0$, $\sum_i z_i=1$ and $Az\le b$,
a row $w\ge0$ with $\min_i(wA)_i>wb$ gives a contradiction by taking
the weighted average with weights $z_i$. These certificates use
rational arithmetic. Interval checks round outwards. The
\codecheck{README.md}{instructions} give the software requirements;
\codecheck{verify.py}{\texttt{python verify.py}} replays the checks.

\tableofcontents
\section{Two characters of the sphericalization}

Write $\varepsilon$ for the central sign of the sphericalization.
The two lifts of a simple $X$ have opposite spherical dimensions. Choose
the lift $\widetilde X$ with positive dimension $x_X$. Its square is
the squared norm of $X$~\cite[Theorem~2.3 and Remark~3.1]{ENO}. With
$d_X=\operatorname{FPdim}(X)$ we have $x_X\le d_X$,
$E=\sum_X x_X^2$ and $D=\sum_X d_X^2$. In the quotients
\[
 R=K_0(\mathcal S)/([\varepsilon]-1)\simeq K_0(\mathcal C),\qquad
 T=K_0(\mathcal S)/([\varepsilon]+1),
\]
the classes $b_X=[\widetilde X]$ give bases indexed by the same simples.
The coefficient of $b_X$ in a product is the sum of the two lift
multiplicities in $R$ and their difference in $T$. In particular
$R/2R\simeq T/2T$. We call characters of the two quotients ordinary
and signed, respectively~\cite[\S5.1]{Ostrik3}.

\subsection{Central induction}
Let $F:Z(\mathcal S)\to\mathcal S$ be the forgetful functor, with right
adjoint $I$. In either quotient the standard trace $\tau$ extracts
the coefficient of the unit. On the complexified algebra,
$\tau(a)=\sum_V\operatorname{Tr}(a|V)/f_V$, summed over the simple
modules with their formal codegrees~\cite[\S2]{OstrikCodegrees}.
For a linear character $\chi$ we have
$f_\chi=\sum_X|\chi(b_X)|^2$. In the commutative case
$\sum_\chi 1/f_\chi=1$. We call the sum over a character orbit
its \emph{reciprocal mass}. On each summand cut out by
$(1\pm[\varepsilon])/2$, the unit trace of $K_0(\mathcal S)$
restricts to $\tau/2$. Thus $V$ has codegree $2f_V$ in $\mathcal S$.

For a simple module $V$ write $\chi_V$ for its central character
(the scalar action on $V$), and $A_V$ for its label in the center.
We also write $A_\chi$ when $V$ is the linear character $\chi$.
By Ostrik's induction theorem~\cite[Example~2.9 and Theorems~2.5 and~2.13]{Ostrik3},
$A_V$ occurs $\dim V$ times in $I(\mathbf1)$ and has twist one.
With $S_{1,1}=1/(2E)$ we have
\[
      \dim A_V=E/f_V,\qquad S_{A_V,Y}=\chi_V(FY)/(2f_V).
\]
Thus symmetry of $S$ becomes
\begin{equation}\label{eq:induction-reciprocity}
       \frac{\chi_V(F A_W)}{f_V}=\frac{\chi_W(F A_V)}{f_W}.
\end{equation}
The central lift $J$ of $\varepsilon$ has dimension and twist $-1$.
Since $I(\varepsilon)=JI(\mathbf1)$, the simple summands of
$I(\varepsilon)$ have negative dimension and those of $I(\mathbf1)$
have positive dimension. Adjunction now gives
$[F A_V:\mathbf1]=\dim V$ and $[F A_V:\varepsilon]=0$.

\begin{lemma}[Simultaneous induction column]\label{lem:central-witness}
The spherical dimension field $M$ of $\mathcal S$ is abelian and
contains the codegrees of both quotients. For $\sigma\in\operatorname{Gal}(M/\mathbb Q)$
let $W=F A_{\sigma\circ\dim_{\mathcal S}}$. Its central characters and its two dimensions are
\[
  \chi_V(W)=\frac{f_V}{\sigma(f_V)},\qquad
  \dim W=\frac{E}{\sigma(E)},\quad
  \operatorname{FPdim}W=\frac{D}{\sigma(D)} .
\]
All these values are totally positive algebraic units. The formal
codegrees of both quotients lie in $\mathbb Q(D)$.
\end{lemma}
\begin{proof}
The dimensions are cyclotomic integers~\cite[Theorem~8.51]{ENO},
and $E/f_V=\dim F A_V$ puts $f_V$ in $M$. Reciprocity
\eqref{eq:induction-reciprocity} gives
$\chi_V(W)/f_V=\sigma(E/f_V)/\sigma(E)$.
The codegrees $2f_V$ are totally positive, and their principal
ideals are Galois invariant by Ostrik's $d$-number theorem~\cite{OstrikCodegrees}. If $\sigma$ fixes $D$, then $W$ has Perron dimension
one and contains $\mathbf1$, so $W=\mathbf1$. The first identity
now shows that $\sigma$ fixes every $f_V$. Conversely $D$ itself
is a codegree.
\end{proof}

In particular $\alpha=E/D=\dim A_{\operatorname{FPdim}}$ is a totally
positive algebraic integer. Since $\dim W\le\operatorname{FPdim}W$,
\begin{equation}\label{eq:least-ratio}
 \alpha\le\sigma(\alpha)\qquad(\sigma\in\operatorname{Gal}(M/\mathbb Q)).
\end{equation}
If $\alpha<1$, some conjugate exceeds one by integrality of the norm.
We often pass to a Galois conjugate maximizing $E$. Then $1/E$
and $1/D$ are both least at the distinguished embedding, the latter
by Perron domination.

\begin{lemma}[Comparison of the two quotients]\label{lem:column-comparison}
For any object $Y$ of $\mathcal S$ its ordinary and signed classes
satisfy $\|[Y]_+\|^2-\|[Y]_-\|^2\in4\mathbb Z_{\ge0}$.
Every linear character satisfies $|\chi(Y)|\le\operatorname{FPdim}Y$.
In particular $f_V/\sigma(f_V)\le D/\sigma(D)$ for every $V$.
\end{lemma}
\begin{proof}
If the two lifts of $X$ occur $m_X^+$ and $m_X^-$ times in $Y$,
the norm difference is $4\sum_X m_X^+m_X^-$. Every eigenvalue of
$[Y]$ on a simple module occurs in the regular representation, whose
spectral radius on $[Y]$ is $\operatorname{FPdim}Y$. This applies
in particular to the scalars by which $[W]$ acts.
\end{proof}

\subsection{Galois translates and their moments}
\begin{corollary}[Integral Galois translates]\label{cor:galois-columns}
In either quotient $A=R,T$, let $V_\sigma$
and $U_\sigma$ be the classes in $A$ of the forgotten induction
labels of $\sigma\circ\operatorname{FPdim}$ and
$\sigma\circ\dim_{\mathcal S}$. Here $\sigma$ acts on a common
Galois splitting field of the characters. Write $V=V_1$. Then
$U_\sigma$ is a unit of $A$, and
\[
 \chi(U_\sigma^{-1}V_\sigma)=\sigma(\chi(V))
 \qquad(\chi\text{ a central character of }A).
\]
Every $\chi(V)$ is a totally positive algebraic integer.

In the signed quotient let $M$ be the field of the spherical dimension
character, $x=E/D$, $s=\operatorname{Tr}_{M/\mathbb Q}(1/D)$ and
$t=\operatorname{Tr}_{M/\mathbb Q}(1/E)$. If $\sigma$ fixes every
codegree outside the dimension orbit, then
\[
    \operatorname{Tr}_{M/\mathbb Q}(\sigma(x)/E)\in s+\mathbb Z_{\ge0}.
\]
If all the complementary codegrees are rational, then $s\le t$.
At equality the complementary characters take value one on $V$
and on every $U_\sigma^{-1}V_\sigma$.
\end{corollary}
\begin{proof}
The determinant of multiplication by $U_\sigma$ is the product of
$f_V/\sigma(f_V)$ with multiplicities $(\dim V)^2$. Galois conjugation
permutes modules of each dimension, so this product is one.
Reciprocity~\eqref{eq:induction-reciprocity} gives
\[
  \frac{\chi(V_\sigma)}{f_\chi}
   =\sigma\!\left(\frac{\operatorname{FPdim}F(A_\chi)}{D}\right)
   =\sigma\!\left(\frac{\chi(V)}{f_\chi}\right).
\]
This proves the first formula. At $\sigma=1$ it gives $\chi(V)>0$.
Galois conjugation permutes the characters, so all conjugates are
positive, and integrality follows from $V\in A$.

In the signed quotient the dimension orbit contributes $s$ to the
unit coefficient of $V$. Every other contribution is positive, so
$s\le1$. If $\sigma$ fixes the complementary codegrees, its action
permutes the complementary values with their trace weights (module dimension divided by codegree) unchanged. The translated column therefore has
unit coefficient $1-s+\operatorname{Tr}_{M/\mathbb Q}(\sigma(x)/E)$.
This is an integer and the trace is positive, giving the inclusion.
If the complementary codegrees are rational, each complementary
orbit has constant weight. By the arithmetic--geometric mean
inequality its average value on $V$ is at least one, with equality
only if all values are one. Summing gives $1-s\ge1-t$.
\end{proof}

\begin{lemma}[Moments at equal reciprocal mass]\label{lem:mass-moments}
Suppose all complementary signed codegrees are rational and that the
two masses in Corollary~\ref{cor:galois-columns} satisfy
$s=t=u/v$ in lowest terms. Put $x=E/D$, $G=\operatorname{Gal}(M/\mathbb Q)$
and $L(y)=\operatorname{Tr}_{M/\mathbb Q}(y/E)/t$. Consider the monomials
in $g(x)$ and $E/g(E)$, $g\in G$, allowing negative powers of the
units $E/g(E)$. Each such $y$ satisfies
\[
 L(y)\in 1+t^{-1}\mathbb Z_{\ge0},\qquad
 \prod_{g\in G}(X-g(y))\equiv (X-1)^{|G|}\pmod v .
\]
In particular all these moments are positive odd integers at $t=1/2$.
\end{lemma}
\begin{proof}
Each generator has an integral lift with value one at every
complementary character. For $g(x)$ take the translate of $V$ in
Corollary~\ref{cor:galois-columns}; for $E/g(E)$ take $U_g$, an integral
unit. Consequently for any integral polynomial $P$ in these generators
and the indicated inverses, the unit coefficient of its lift gives
\[
  \operatorname{Tr}_{M/\mathbb Q}(P/E)-tP(1,\ldots,1)\in\mathbb Z .
\]
For a monomial, $P(1,\ldots,1)=1$ and the trace is positive. For the
second assertion note that $g(E)/gh(E)$ is a quotient of two of the
allowed units. All conjugates of $y$ are thus monomials as well.
Their $k$th elementary symmetric function $a_k$ is a rational integer
and its lift has complementary value $\binom{|G|}{k}$. Applying the
trace identity to this lift gives $t(a_k-\binom{|G|}{k})\in\mathbb Z$.
\end{proof}

\subsection{Discriminants and trace orders}
The relative discriminant gives a lower bound for the trace of a
totally positive generator.

\begin{lemma}\label{lem:relative-trace-discriminant}
Let $K/F$ be a totally real quadratic extension with $[F:\mathbb Q]=m$,
and relative discriminant $\mathfrak d_{K/F}$. Put $\Delta=N_{F/\mathbb Q}(\mathfrak d_{K/F})$. For a totally
positive $x\in\mathcal O_K$ with $K=F(x)$,
\[
 \operatorname{Tr}_{K/\mathbb Q}(x)\ \ge\
   m\sqrt{\Delta^{1/m}+4N_{K/\mathbb Q}(x)^{1/m}} .
\]
\end{lemma}
\begin{proof}
Pair the conjugates $a_i,b_i$ of $x$ above each embedding of $F$.
The discriminant of $\mathcal O_F[x]$ is divisible by
$\mathfrak d_{K/F}$, so $\prod_i(a_i-b_i)^2\ge\Delta$.
The geometric mean is superadditive on positive vectors by H\"older's
inequality. Apply this to $(a_i+b_i)^2=(a_i-b_i)^2+4a_i b_i$, then
use the arithmetic--geometric mean inequality on the $a_i+b_i$.
\end{proof}

We shall also need the ordinary trace dual of an order. If a commutative
quotient $A=R$ or $T$ has rational algebra $\prod_i K_i$, choose one
character in each field factor and write $f_i$ for its codegree.
Then $\tau=\sum_i\operatorname{Tr}_{K_i/\mathbb Q}(-/f_i)$. The
pairing $\tau(ab)$ is unimodular: the dual classes of the simple
basis elements form its dual basis.

\begin{lemma}[The conductor of a trace order]\label{lem:trace-order-conductor}
Let $\mathcal O=\prod_i\mathcal O_{K_i}$ and $\mathfrak D_i$ be
the different of $K_i$. Then
\[
 (A:\mathcal O)=\prod_i f_i\mathfrak D_i^{-1},\qquad
 [\mathcal O:A]^2\,|\operatorname{disc}\mathcal O|
        =\prod_i N_{K_i/\mathbb Q}(f_i).
\]
The conductor identity holds after completion at any prime. A local
field factor on which its codegree is a unit splits off as an
unramified maximal order. For a set $I$ of factors,
$A\cap\prod_{i\in I}K_i$ is the $\tau$-dual of the projection of
$A$ to those factors, where the intersection uses zero in the other
coordinates. In particular an isolated pair of rational $p$-adic axes
with codegrees $g,h$ has order
$\{(x,y)\in\mathbb Z_p^2:x\equiv y\pmod{p^v}\}$ with
$v=v_p(g)=v_p(h)$. Its rows satisfy $p^{2v}\mid g+h$ when $g,h$
are rational integers.
\end{lemma}
\begin{proof}
Write $B^\#$ for the dual of a lattice with respect to $\tau(ab)$.
Since $A=A^\#$ and $A\mathcal O=\mathcal O$, one has
$a\mathcal O\subset A$ exactly when $\tau(a\mathcal O)\subset\mathbb Z$.
This gives $(A:\mathcal O)=\mathcal O^\#$, whose $i$th component is
$f_i\mathfrak D_i^{-1}$. When $f_i$ is a local unit, containment of
the conductor in $\mathcal O$ forces $\mathfrak D_i$ to be a unit
there as well. Taking determinants of the trace form gives the index
identity. Finally an element supported on $I$ pairs integrally with
$A$ exactly when it pairs integrally with its projection. On two rational axes
the diagonal copy of $\mathbb Z_p$ lies in the order and the supported
coordinate ideals are $g\mathbb Z_p$ and $h\mathbb Z_p$. Subtracting a
diagonal element identifies these ideals and gives the congruence.
Pairing the row difference with its dual gives $g+h$ by orthogonality.
\end{proof}

\begin{lemma}\label{lem:silver-local-partner}
Let $B$ be a commutative ordinary or signed based order, and let a
rational character have formal codegree $f$ with $v_p(f)=1$.
Its completed local factor is
$\mathbb Z_p\times_{\mathbb F_p}\mathcal O_L$, with $L/\mathbb Q_p$
totally ramified. In particular an unramified quadratic branch cannot
be its partner.
\end{lemma}
\begin{proof}
Let $A$ be the completed local factor, $e$ the rational idempotent,
and $B'$ the projection of $A$ to the other branches.
Lemma~\ref{lem:trace-order-conductor} gives
$A\cap\mathbb Q_pe=p\mathbb Z_pe$. Thus the rational coordinate
modulo $p$ defines a surjection $B'\to\mathbb F_p$ and
$A=\mathbb Z_p\times_{\mathbb F_p}B'$.

As $A$ is finite free over $\mathbb Z_p$, its canonical module is
$\operatorname{Hom}_{\mathbb Z_p}(A,\mathbb Z_p)$. The perfect
pairing identifies this module with $A$, so $A$ is Gorenstein.
The Gorenstein fiber-product theorem
\cite[Fact~2.6 and Corollary~2.7]{NTV} now makes $B'$ a discrete
valuation ring. Its residue field is $\mathbb F_p$, hence its
fraction field is totally ramified over $\mathbb Q_p$.
\end{proof}

\subsection{Dimension fields and spherical sections}
\begin{lemma}\label{lem:odd-field-equality}
With trivial universal grading the ordinary and signed dimension
character fields both equal $K=\mathbb Q(D)$.
\end{lemma}
\begin{proof}
For the ordinary field this is the dimensional grading of Gannon and
Schopieray~\cite[\S1.2, (13)]{GSFields}. For the signed field, which
contains $K$ by Lemma~\ref{lem:central-witness}, let $\sigma$ fix $D$.
The label $B$ of $\sigma\circ\dim_{\mathcal S}$ forgets to $\mathbf1$,
so its half-braiding is a monoidal automorphism $\gamma$ of the
identity. On the central sign $J$ its normalized $S$-row is
\[
  \gamma_\varepsilon\dim\varepsilon
  =S_{B,J}/S_{B,\mathbf1}=\sigma(\dim\varepsilon).
\]
The first equality takes the trace of the double braiding, and the
second is the induction character formula. Hence $\gamma_\varepsilon=1$,
so the scalars on the two lifts of each simple agree. They define a
grading character of $\mathcal C$. Triviality of the universal grading
gives $\gamma=1$, hence $B=\mathbf1$. Injectivity of the labeling
in~\cite[Theorem~2.13]{Ostrik3} gives
$\sigma\circ\dim_{\mathcal S}=\dim_{\mathcal S}$.
\end{proof}

\begin{lemma}[A spherical section from a character]\label{lem:section-signs}
Let $\chi:T\to\mathbb R$ be a character. The simples $X$ of $\mathcal C$
for which $|\chi(\widetilde X)|=\operatorname{FPdim}(X)$ form a fusion
subcategory admitting a spherical structure. In particular, if
$f_\chi=D$, the whole category has a spherical structure.
\end{lemma}
\begin{proof}
Write $d_X=\operatorname{FPdim}(X)$. The spectral radius of the
fusion matrix of either lift gives $|\chi(\widetilde X)|\le d_X$.
Duality conjugates the character value, so equality is preserved by
duals. For each simple attaining equality take the unique lift with
value $d_X$. A product of two chosen lifts has character value equal
to its Perron dimension. Equality in the triangle inequality forces
every summand to be a chosen lift as well. These lifts therefore give
a spherical section over the indicated subcategory. If $f_\chi=D$,
the sum of the nonnegative differences
$d_X^2-|\chi(\widetilde X)|^2$ vanishes.
\end{proof}

For a fusion category $\mathcal A$, let $K_{\mathcal A}$ be the field
generated by the Frobenius--Perron dimensions of its simple objects.

\begin{lemma}\label{lem:field-corank}
Let $\mathcal A\subseteq\mathcal C$ be a fusion subcategory with
Perron field $F=K_{\mathcal A}$. On each indecomposable component
of the left action of $\mathcal A$, the rational span of the simple Perron
dimensions is one-dimensional over $F$. If $x_A=d_A$ for every simple
$A\in\mathcal A$, the ordinary and signed actions of these positive
lifts agree and $x_X/d_X$ is constant on each component.

In particular a proper $\mathcal A$ in rank $n$ satisfies
\[
 [F:\mathbb Q]\le\min\{\operatorname{rank}(\mathcal A),
                         n-\operatorname{rank}(\mathcal A)\}.
\]
In a faithful grading each homogeneous component has rank at least
$[\mathbb Q(D):\mathbb Q]$.
\end{lemma}
\begin{proof}
On a component let $P$ be the action matrix of the sum of the simples
of $\mathcal A$, with the dimension column as right eigenvector.
Its simple Perron root is $r=\sum_A d_A\in F$. Solving $(P-rI)y=0$
over $F$ shows that the ratios of the coordinates are in $F$.
Their rational span is stable under multiplication by each $d_A$,
so is an $F$-line.

Under the equality hypothesis the signed action matrix $P'$ satisfies
$P' x=rx$ and $P'\le P$ entrywise. Pair $(P-P')x$ with a positive
left Perron vector. The result is zero; positivity of $x$ gives
$P'=P$, also term by term for each simple of $\mathcal A$.
Simplicity of the eigenspace gives the constant ratio.
The bounds follow by counting basis elements. For a $G$-grading
take the neutral component as $\mathcal A$ and use
$D=|G|\operatorname{FPdim}(\mathcal A)$~\cite[Proposition~8.20]{ENO}.
\end{proof}

\begin{lemma}[The regular idempotent of the integral subcategory]\label{lem:integral-regular}
Let $\mathcal A$ be the full subcategory on the simples of integral
Frobenius--Perron dimension and put $n=\operatorname{FPdim}(\mathcal A)$.
In both $R\otimes\mathbb Q$ and $T\otimes\mathbb Q$ it has a regular
idempotent of unit coefficient $1/n$. A linear character takes value
one on it precisely when its restriction to $\mathcal A$ is the
Perron character, and zero otherwise.
\end{lemma}
\begin{proof}
Closure of $\mathcal A$ is the dimension-subfield criterion
of~\cite[\S1]{GSFields} applied to $\mathbb Q$. Its integral global
Perron dimension gives pseudo-unitarity~\cite[Propositions~8.23--8.24]{ENO},
so its positive lifts give a section by Lemma~\ref{lem:section-signs}.

In either quotient take
\[
 e_{\mathcal A}=\frac1n\sum_{X\in\operatorname{Irr}(\mathcal A)}
 d_X[\widetilde X].
\]
The regular-object identity $[\widetilde Y] e_{\mathcal A}=
d_Y e_{\mathcal A}$ for $Y\in\mathcal A$ gives the stated
evaluations. The coefficient at the unit is $1/n$.
\end{proof}

\subsection{When the dimension ratio lies in a proper subfield}
\begin{theorem}\label{thm:proper-ratio}
Let $K=\mathbb Q(D)$, $L=\mathbb Q(E/D)$ and $H=\operatorname{Gal}(K/L)$.
If $E<D$ and $H\ne1$, let $\mathcal A$ be the proper fusion subcategory
of simples on which $x_X=d_X$. It is pseudo-unitary. The algebraic integer
\[
       \delta=\sum_{\sigma\in H}\frac D{\sigma(D)}
\]
belongs to $K_{\mathcal A}$, and
\[
    K=L(\delta),\qquad D=\delta w,\quad E=\delta z\quad(w,z\in L),
    \qquad \operatorname{Tr}_{K/L}(1/\delta)=1 .
\]
In particular $H$ acts faithfully on $\mathbb Q(\delta)$ and
$|H|\le[\mathbb Q(\delta):\mathbb Q]\le\lfloor\operatorname{rank}(\mathcal C)/2\rfloor$.
\end{theorem}
\begin{proof}
Lemma~\ref{lem:section-signs} identifies the positive lifts of
$\mathcal A$ with a spherical section of Perron dimensions.
For $\sigma\in H$, extend $\sigma$ to the signed dimension field.
The two dimensions of $W_\sigma$ in Lemma~\ref{lem:central-witness}
agree. Thus $W_\sigma$ is supported on $\mathcal A$, so its Perron
dimension $D/\sigma(D)$ belongs to $K_{\mathcal A}$. Their sum is
$\delta=D\operatorname{Tr}_{K/L}(1/D)$, so $w=D/\delta$ and
$z=(E/D)w$ lie in $L$ and $\operatorname{Tr}_{K/L}(1/\delta)=1$.
Since $K=L(D)$, this also gives $K=L(\delta)$. The fields are
abelian, so $H$ acts faithfully on $\mathbb Q(\delta)$, and
Lemma~\ref{lem:field-corank} gives the bound.
\end{proof}

\begin{corollary}\label{cor:cyclic-prime-power-proper-ratio}
Let $n$ be the rank of $\mathcal C$. Suppose $K=\mathbb Q(D)$ is cyclic of prime-power degree $d>n/2$. If $E/D$ belongs to a proper subfield of $K$, then
$E=D$ and $\mathcal C$ is pseudo-unitary.
In particular this holds for a cyclic quartic dimension field in
rank six, with arbitrary complementary formal codegrees.
\end{corollary}
\begin{proof}
If $E<D$, Theorem~\ref{thm:proper-ratio} writes $K$ as the compositum
of two proper subfields. In a cyclic prime-power extension every
proper subfield lies in the same maximal proper subfield.
\end{proof}

\subsection{Fourier bounds from the modular representation}
The following projection estimate is the only input from congruence
representations. We use $U(r)=\left(\begin{smallmatrix}1&r\\0&1\end{smallmatrix}\right)$
and $s=\left(\begin{smallmatrix}0&-1\\1&0\end{smallmatrix}\right)$.

\begin{lemma}[Modular projection]\label{lem:mass-fourier-verification}
Let $z$ be a unit vector in a unitary representation $\pi$ of
$\mathrm{SL}_2(\mathbb Z/N)$. Suppose $U(r)$ fixes $z$ for $r=1$
or $2$ and the diagonal subgroup acts on $z$ by a character of
conductor $c$. For the projection $P_r$ onto the $U(r)$-fixed
space one has $\|P_r\pi(s)z\|^2\le(c,r^2)/c$.
\end{lemma}
\begin{proof}
Put $A=\mathbb Z/N$ and $L(b)=U(b)^{\mathsf t}$. The squared norm is
the average of $\langle z,\pi(L(b))z\rangle$ over $b\in rA$.
Group the terms by the ideal $bA$. If $u\in A^\times$ and
$u-1\in rbA$, choose $x\in rA$ with $u=1+bx$. Then
\[
 U(x)L(b)=L(b/u)\operatorname{diag}(u,u^{-1})U(x/u).
\]
On the entire stratum $b\mapsto b/u$ is a permutation. Its sum
therefore vanishes unless the diagonal character is trivial on all
$u\equiv1\pmod{rbA}$, or equivalently $c\mid rb$. These strata
occupy a proportion $(c,r^2)/c$ of $rA$, and each term has modulus
at most one~\cite[Lemma~4.1 and Proposition~4.3]{Class6Rank5}.
\end{proof}

Suppose the ordinary and signed dimension characters have the same
field $K$, and write $G=\operatorname{Gal}(K/\mathbb Q)$ and $n=|G|$.
For a nontrivial character $\eta$ of $G$ of conductor $c$, set
\[
  \widehat a=\sum_{g\in G}\eta(g)g(a),\qquad
  P=\widehat{1/D},\qquad Q=\widehat{1/E}.
\]
Every linear codegree lies in $K$ by Lemma~\ref{lem:central-witness}.
To retain contributions from the complementary labels, put
\[
  O_\eta=\frac1n\sum_{\chi\ {\rm ordinary}}|\widehat{1/f_\chi}|^2,
  \qquad
  T_\eta=\frac1n\sum_{\chi\ {\rm signed}}|\widehat{1/f_\chi}|^2,
\]
where the sums run over linear characters. The dimension orbits
contribute $|P|^2$ and $|Q|^2$, respectively.

\begin{lemma}[Simultaneous Fourier bounds]\label{lem:sextic-matrix-compression}
For each $\eta$ as above,
\begin{equation}\label{eq:orbit-projection}
  cO_\eta\le(c,2),\qquad cT_\eta\le(c,4),\qquad
  c(2O_\eta+T_\eta)\le4 .
\end{equation}
The two dimension characters also satisfy
\[
 c|Q|^2\le(2-cO_\eta)(2-c|P|^2),\qquad
 \sqrt c\,|Q|+c|P|^2\le2.
\]
\end{lemma}
\begin{proof}
Write $B_g,A_g$ for the labels of the conjugates of the signed and
ordinary dimension characters. Their translates by $J$ give two more
disjoint orbits. Only the $JB_g$ have twist $-1$; the other three
orbits have twist one. In the unitary center representation
\cite[Proposition~2.12]{ENO}, define $b=n^{-1/2}\sum_g\eta(g)e_{B_g}$
and $a=n^{-1/2}\sum_g\eta(g)e_{A_g}$. Each is a diagonal
eigenvector of conductor $c$, as are its $J$-translates: the Galois
symmetry of the modular data permutes the dimension rows by conjugation
with positive sign, and $S_{J,Y}/S_{1,Y}$ is rational
\cite[Theorem~II and Remark~4.5]{DLN}.

Reciprocity gives $S_{B_g,A_\chi}=1/(2g(f_\chi))$.
Multiplication of a source label by $J$ preserves its $S$-entry at
an ordinary target and reverses it at a signed target. Multiplication
of the target by $J$ preserves the modulus. Thus in the twist-one
projection of $Sb$, the ordinary labels and their translates
contribute $O_\eta/2$ and the signed labels contribute $T_\eta/4$.
For $(b+Jb)/\sqrt2$ and $(b-Jb)/\sqrt2$ the twist-square
projection instead retains contributions $O_\eta$ and $T_\eta$,
respectively. Lemma~\ref{lem:mass-fourier-verification} now gives
\eqref{eq:orbit-projection}, using $cO_\eta\le2$ from the mixed
bound when $4\mid c$.

For the last bounds use the twist-one source vectors $b$ and
$(a-Ja)/\sqrt2$. Project their images onto the span of the
conjugate-weighted $B$-vector and all the ordinary linear labels and their
$J$-translates. The second image vanishes at the ordinary targets.
After unitary changes of target coordinates, the resulting matrix
has the form
\[
   \begin{pmatrix}Q/2&P/\sqrt2\\ \sqrt{O_\eta/2}&0\end{pmatrix}
\]
up to conjugation of the first row. Indeed the $B_g,B_h$ and
$B_g,A_h$ entries are $1/(2gh(E))$ and $1/(2gh(D))$.
Lemma~\ref{lem:mass-fourier-verification} bounds the norm of this
matrix $M$ by $1/\sqrt c$. The determinant of $I-cM^*M$ gives the
first inequality. Since $|P|^2\le O_\eta\le2/c$, taking square
roots gives the second.
\end{proof}

When $K$ is real quadratic, $O_\eta=\Delta_R^2$ and $T_\eta=\Delta_T^2$,
where each $\Delta^2$ sums $(1/f-1/f')^2$ over the conjugate pairs,
with multiplicity. We shall use \eqref{eq:orbit-projection} in the form
\begin{equation}\label{eq:quadratic-fourier-bounds}
 \Delta_R^2\le(c,2)/c,\qquad \Delta_T^2\le(c,4)/c,\qquad
 2\Delta_R^2+\Delta_T^2\le4/c.
\end{equation}

For a quartic dimension orbit with a quadratic complementary pair
$f,f'$ in the ordinary sector, take $\eta$ cutting out
$\mathbb Q(f)$. Each embedding of that subfield occurs twice in $G$.
Hence $O_\eta=|P|^2+2(1/f-1/f')^2$; the signed formula is identical
with $T_\eta,Q$. For example, the roots of $X^2-15X+25$ have
squared reciprocal gap $1/5$. In conductor five they contribute
$2/5$ to $O_\eta$, already exceeding the bound $1/5$.

\section{Odd-degree dimension fields}
\label{sec:odd-degree-descent}
\begin{theorem}\label{thm:odd-degree-descent}
Let $\mathcal C$ have rank six and put $D=\operatorname{FPdim}(\mathcal C)$.
If $[\mathbb Q(D):\mathbb Q]$ is odd, then $\mathcal C$ has a spherical
structure. In degrees one and five a Galois conjugate is pseudo-unitary.
In degree three the same conclusion holds unless the signed algebra is
$\mathbb Q(\zeta_9)^+\times\mathbb Q(\zeta_9)^+$ and its dimension codegree is nine.
\end{theorem}

Write $K=\mathbb Q(D)$, $p=1/D$, $q=1/E$, $\alpha=E/D$,
$s=\operatorname{Tr}_K(p)$ and $t=\operatorname{Tr}_K(q)$.
Suppose first that $K$ has degree three or five and the universal
grading is trivial. By Lemma~\ref{lem:odd-field-equality}, $K$ is a
factor of both character algebras. Their complements have dimension
at most three, so both algebras are commutative. All codegrees lie in
the cyclic field $K$ by Lemma~\ref{lem:central-witness}.

\begin{lemma}[A cyclic rearrangement]\label{lem:short-orbit-rearrangement}
Let $u,v$ be real vectors indexed cyclically by two or three indices.
If both have a strict minimum at zero, then some cyclic shift $j$
satisfies $\sum_i u_i v_{i+j}<\sum_i u_i v_i$.
\end{lemma}
\begin{proof}
Subtracting the respective minima changes every sum by the same
constant. For two indices the shifted sum is now zero. For three
write the vectors as $(0,a,b)$ and $(0,c,d)$. The three sums are
$ac+bd$, $ad$ and $bc$, and $\min(ad,bc)\le\sqrt{acbd}<ac+bd$.
\end{proof}

\begin{lemma}\label{lem:odd-shift}
Suppose the ordinary and signed algebras are commutative and their
dimension characters both generate a cubic field $K=\mathbb Q(D)$.
If every complementary signed codegree is rational, a Galois conjugate
is pseudo-unitary. The same conclusion holds for $T_{\mathbb Q}=K\times K$
if $(1,-1)\in T$.
\end{lemma}
\begin{proof}
Put $h_\sigma=\operatorname{Tr}(q\sigma(\alpha))$, so $h_1=s$.
For rational complementary codegrees Corollary~\ref{cor:galois-columns}
gives $s\le t$ and $h_\sigma\ge s$.

For the second assertion write $g=(1,-1)$. The positive involution of
$K^2$ is trivial. Thus $g$ has squared norm one and is a signed basis
element distinct from $\pm1$. Its unit coefficient $2t-1$ vanishes.
Moreover multiplication by $1+g$ doubles the first coordinate and
kills the second. Applied to the integral translates of the Perron
column in Corollary~\ref{cor:galois-columns}, this shows
$2h_\sigma\in\mathbb Z_{>0}$. Since the ordinary dimension orbit
is proper, $0<s<1$, whence $s=1/2=t$ and again $h_\sigma\ge s$.

If $E$ is rational, $\operatorname{Tr}(\alpha)=Es\le Et=3$.
Total positivity and the integral norm force $\alpha=1$.
Otherwise maximize $E$ by transport. A rational $\alpha$ is a positive
integer at most one. If $\alpha$ is irrational, $q$ and $\alpha$
both have a strict minimum at the distinguished embedding, by
\eqref{eq:least-ratio}. Lemma~\ref{lem:short-orbit-rearrangement}
gives $h_\sigma<h_1=s$, a contradiction. Thus $E=D$ after transport.
\end{proof}

All nontrivial characters of the cubic and quintic fields have odd
conductor $c$. By \eqref{eq:orbit-projection}, in either sector the
sum of squared Fourier coefficients of the reciprocal codegrees,
one for each full orbit, is at most $1/c$.

We shall use the same integral coordinates at conductors $7,11,13$.
For an odd prime $\ell=2m+1$, put $b_j=\zeta_\ell^j+\zeta_\ell^{-j}$
for $1\le j\le m$. Their trace Gram matrix is $\ell I-2J$, where
$J$ is the all-ones matrix, by the identity $b_ib_j=b_{i+j}+b_{i-j}$.
Its determinant $\ell^{m-1}$ is the discriminant of
$\mathbb Q(\zeta_\ell)^+$, so the $b_j$ form an integral basis.
Writing a totally positive integer $z=-\sum c_jb_j$ of trace $A$
gives
\begin{equation}\label{eq:odd-cycle-coordinates}
 \sum c_j=A,\qquad
 \ell c_j=\operatorname{Tr}((2-b_j)z)\in(0,4A).
\end{equation}
These coordinates at eleven are also used in \cite[\S8.2]{Class6Rank5}.

\subsection{Two cubic orbits}
Put $P(x)=\sum_i x_i^2-\sum_{i<j}x_ix_j$ for triples, and write
$\nu(f)=P(1/f)/\operatorname{Tr}(1/f)^2$ for a totally positive cubic
element, using its three embeddings.
\begin{proposition}\label{prop:odd-cubic-arithmetic}
Suppose $K$ is cubic, $T_{\mathbb Q}=K^2$, and the complementary signed
codegree $H$ is irrational. Then the category has a pseudo-unitary
Galois conjugate at conductor seven. The only other possibility has
conductor nine and $E=9$, and the category is spherical there.
\end{proposition}
\begin{proof}
\textit{Bounding the conductor.} Write $w=1-t$. The cubic field has
discriminant $c^2$. The character Gram determinant of the dimension
orbit is $N(E)$. Factoring the character matrix through an integral
basis of $K$ gives $c^2\mid N(E)$. Thus rational $E$ is divisible
by $c$ (a cubic conductor is squarefree away from a possible factor
$9$). The preceding Fourier bound becomes
\[
        t^2\nu(E)+w^2\nu(H)\le1/c,\qquad P(p)\le1/c .
\]

Suppose $c\ge13$. Some irrational $f\in\{E,H\}$ has $\nu(f)\le2/c$:
if $E$ is irrational use $t^2+w^2\ge1/2$, and otherwise use
$w\ge1-3/c>1/\sqrt2$. Normalize the reciprocal triple of $f$ to
$z_0+z_1+z_2=1$ and put $v=\sqrt{P(z)}\le\sqrt{2/13}<2/5$.
The unit $z_1/z_0$ and its inverse have integral traces $A,B$, with
$A+B+3=\sum_i1/z_i$. At fixed $v$ the pair sum of $z$ is fixed,
and the product is least at $(1-2v,1+v,1+v)/3$, up to permutation
(the extremum has two equal entries). Thus
\[
  A+B+3\le\frac{9(1-v)}{(1-2v)(1+v)}<\frac{135}{7}.
\]
Up to inversion $4\le A\le B$ and $A+B\le16$. The polynomial
$X^3-AX^2+BX-1$ is irreducible with positive square discriminant.
In this range the only triples $(A,B,\sqrt\Delta)$ are $(5,6,7)$
and $(6,9,9)$, by substitution in
$\Delta=A^2B^2-4(A^3+B^3)+18AB-27$. But $c\mid\sqrt\Delta$.
Thus only the conductors seven and nine remain.

\smallskip\textit{The gap at seven and nine.} There are sharper gaps in $K_p=\mathbb Q(\zeta_p)^+$ for $p=7,9$.
Scale $1/f$ by a positive rational to clear the common ideal
valuations off the ramified prime and reduce its ramified exponent
modulo three. The resulting integer has norm $C=1,\ell,\ell^2$,
with $\ell=7,3$, respectively. Write its polynomial as
$X^3-AX^2+BX-C$. Then $\nu(f)=1-3B/A^2$, $C\mid A^3$ and $C^2\mid B^3$.
If $\nu(f)<1/4$, then $m=4B-A^2$ is positive and is $0$ or $3$
modulo four. Schur's inequality for the three distinct positive
roots gives $Am<9C$. For $C=1,9,49$ the divisibilities give
$m\ge3,27,147$, respectively, impossible since $A>3C^{1/3}$.
At $C=3$ we have $3\mid A,B$ and $m\ge3$, so $A=6$; now
$12\mid m$, again impossible. At $C=7$ one has $7\mid A,B,m$,
so $A=7$, $B=14$.

Thus $\nu(f)\ge1/4$ except for $f=n\sigma(F)$ at seven, where
\[
 F^3-14F^2+49F-49=0,\qquad
 \nu(F)=1/7,\qquad \operatorname{Tr}(1/F)=1 .
\]
The rational scalar $n$ is integral since $N(F)=49$.

\smallskip\textit{Conductor nine.} Both codegrees cannot be irrational, by
$(t^2+w^2)/4\ge1/8$. Thus $9\mid E$ and $(1-3/E)^2/4\le1/9$ give
$E=9$, $w=2/3$, $\nu(H)=1/4$. Writing $C=\operatorname{N}(H)$ gives
$C\mid9^3$ by integrality of $9/H$, and its polynomial is
$X^3-(C/9)X^2+(2C/3)X-C$. The positive discriminant forces $C>162$,
hence $C=3^5$ or $3^6$. Since the order discriminant $9^3 C$
differs from $81^2$ by a square factor, $C=3^6$.
The ordinary estimate $s^2/4\le1/9$ gives
$\operatorname{Tr}(\alpha)=9s\le6$. Its norm divides $9^3$ and is
at most eight. The trace ideal of $K_9$ is $3\mathbb Z$ by its integral power
basis in $\zeta_9+\zeta_9^{-1}$. Thus the trace of the irrational
$\alpha$ is six, and its middle coefficient $B'=(36-\operatorname{Tr}
(\alpha^2))/2$ is divisible by three. The arithmetic--geometric
mean inequalities give $3<B'<12$. At $B'=6$ the discriminants for
norms $1,3$ are $189,-459$, respectively, neither a positive square.
Thus $\alpha$ has polynomial $X^3-6X^2+9X-C'$ with $C'=1,3$.

Here $s=2/3$, so a second irrational ordinary orbit would require
$1/9\ge((2/3)^2+(1/3)^2)/4$. The three rational ordinary
complementary codegrees divide nine and their reciprocal sum is one
third, so all equal nine. If $C'=3$, the ordinary order discriminant
$\operatorname{N}(D)9^3=3^{11}$ is nonsquare. Every simple is self-dual
because $T_{\mathbb Q}=K^2$ is totally real. Hence the ordinary
complement consists of totally real abelian field factors of total
degree three. No real quadratic field ramifies only at three and
each abelian cubic discriminant is a square, a contradiction.
If $C'=1$, $D=9/\alpha$ satisfies the same cubic polynomial as $H$.
There is therefore a real signed character of codegree $D$.
Lemma~\ref{lem:section-signs} gives a spherical section.

\smallskip\textit{Conductor seven.} If $E$ is rational, $7\mid E$ gives $w\ge4/7$. The exceptional
orbit has mass $1/n$ and cannot give this $w\in[4/7,1)$. Thus
$w^2/4\le1/7$, forcing $E=7$. Then $7/H$ is a totally positive integer of trace four.
In \eqref{eq:odd-cycle-coordinates} its coefficients must be
$2,1,1$ in some order. Hence it is a conjugate of $1-b_1$,
which is negative since $2\cos(2\pi/7)>1$.

Suppose one orbit has the exceptional codegree $n\sigma(F)$. Its mass
is $1/n$. If the other is exceptional too, their masses sum to one
only at $n=2$. Otherwise $1/(7n^2)+(1-1/n)^2/4\le1/7$, so $n=2,3$.
For $\theta^3-\theta^2-2\theta+1=0$ take $F=1+\theta+2\theta^2$.
In the integral basis $(1,\theta,\theta^2-1)$ the trace form
weighted by $1/F$ has Gram matrix $I$
(\codecheck{odd_degree_checks/cubic_seven_lattice_check.py}{trace calculation}).
The coordinate matrix $Y$ of the six character values satisfies
$YY^{\mathsf t}=nI$, with first column $(1,0,0)^{\mathsf t}$.
For $n=3$ choose column signs and order so the first row is
$(1,1,1,0,0,0)$. Each other row has the form $(0,a,-a,w)$ with
$2a^2+\|w\|^2=3$. If both $a$ vanish their dot product is odd;
if just one vanishes it is $\pm1$; otherwise it is $\pm2$ plus
$0$ or $\pm1$. Thus they cannot be orthogonal. For $n=2$ the first
row, twice the orbit idempotent, is $1\pm b_j$. Hence
$(1,-1)\in T$, and Lemma~\ref{lem:odd-shift} applies.

It remains to handle $\nu(E),\nu(H)\ge1/4$. Order
$q_0\le q_1\le q_2$ by maximizing $E$ under transport. Since $P(q)\ge(q_2-q_1)^2$,
\[
 1-3q_1\le(1-t)+(q_2-q_1)
 \le\sqrt{5\big((1-t)^2/4+P(q)\big)}\le\sqrt{5/7}.
\]
Thus $q_2\ge q_1>1/20$. Using $\alpha_0\le1$ and
$\sum q_i\alpha_i=s<1$ gives $\operatorname{Tr}(\alpha)<21$.
In the power basis used by the
\codecheck{odd_degree_checks/cubic_ratio_bound.py}{ratio enumeration}
the corresponding trace bound gives coordinate bounds $20,10,10$
by the inverse trace Gram diagonal $(1,2/7,2/7)$. Positivity,
$\alpha_0\le1$ and the $d$-number divisibilities leave fifteen
irrational ratio polynomials.

For each polynomial order its roots $a_0<a_1<a_2$ and put $p_i=a_iq_i$,
where $q_0$ is least. They must satisfy
\[
 P(q)+(1-t)^2/4\le1/7,\quad P(p)\le1/7,\quad
 P(q)\ge t^2/4,\quad \textstyle\sum p_i<1.
\]
The \codecheck{odd_degree_checks/cubic_ratio_fourier.py}{quadratic-form certificates}
exclude thirteen polynomials and give the bounds in the second column
below. They use nonnegative combinations of the four inequalities and
$q_i\ge q_0$. For $X^3-7X^2+14X-7$ first pair the larger of $q_1,q_2$
with $a_1$, which minimizes $P(p)$, and split at $q_2/q_0=4,9$.
\[
 \begin{array}{c|c|c}
 \text{polynomial of }\alpha&q_0&\text{candidates before }P(p)\le1/7\\\hline
 X^3-5X^2+6X-1&>1/2000&7\\
 X^3-6X^2+5X-1&>1/44&0
 \end{array}
\]

In either case $\operatorname{Tr}(E)<6000$. Put
$\delta=1-\theta+\theta^2$, of norm seven. The unit signatures of
$-1,\theta,\sigma(\theta)$ span all signatures, so Dirichlet's theorem
makes every totally positive unit a square. Removing rational ideal
valuations therefore writes $E=n\delta^j u^2$ with $n\ge1$ integral,
$0\le j<3$, and $u$ a unit. The
\codecheck{odd_degree_checks/cubic_height_check.py}{ellipsoid enumeration}
tests $n\operatorname{Tr}(\delta^ju^2)<6000$, $49\mid N(E)$, the
bounds on $q$, and the minimum conditions on $p_0,q_0$. It gives the
last column of the table. The smallest $P(p)$ among the seven remaining
values is $9/49>1/7$.
\end{proof}

\subsection{The quintic field and its rational complement}
\begin{proposition}\label{prop:odd-quintic-arithmetic}
If the ordinary and signed dimension character fields both have
degree five and equal $K=\mathbb Q(D)$, then a Galois conjugate of
the category is pseudo-unitary.
\end{proposition}
\begin{proof}
The two rational algebras are $K\times\mathbb Q$. Write $f,g$ for
their ordinary and signed rational codegrees, respectively, so that
$s=1-1/f$ and $t=1-1/g$. The rational signed character evaluated at
the ordinary Frobenius--Perron induction column gives a positive
integer $r$. Its unit coefficient yields $1=s+r/g$, hence $g=rf$.

Let $A$ be the induction label of the rational signed character.
On the signed dimension labels its row is constant $1/(2g)$.
The unit coefficient of $F(A)$ determines its remaining signed
entry $S_{A,A}=(g-5)/(2g)$. The six entries and their $J$-translates
have squared norm $[5+(g-5)^2]/(2g^2)$, greater than one at $g=2$.
Thus $g\ge3$ by unitarity and $t>0$. A nonidentity invertible would
exchange $1,h$ and fix the other four simples, since the kernel of
the dimension map has rank one. Its positive lift has the same
permutation action, producing a signed character $(1,-1,0,0,0,0)$
of codegree two. Hence all nonunit simples have dimension greater
than one, and Perron orthogonality excludes $f=2$.

If $g=3$ then $f=g=3$ and $r=1$. Now $\operatorname{FPdim}(FA)=D/f$
by symmetry, so its ordinary row likewise has entries $1/(2f)$ five
times and $(f-5)/(2f)$ once. Each sector with its translates accounts
for half the squared norm. In $(STS)_{A,A}$ the signed translates
cancel by their opposite twists and the ordinary ones contribute $1/2$.
But $A$ has twist one, so $(STS)_{A,A}=S_{A,A}=-1/3$. Thus $g\ge4$.

The field discriminant is $c^4$ by the conductor--discriminant
formula. Write $B_T,B_R$ for the projections of the two orders onto
$K$. Lemma~\ref{lem:trace-order-conductor} gives
$T\cap\mathbb Q=g\mathbb Z$, hence $[B_T\times\mathbb Z:T]=g$;
the ordinary index is similarly $f$. Its discriminant identity gives
\[
 N(E)=gc^4[\mathcal O_K:B_T]^2,\qquad
 N(D)=fc^4[\mathcal O_K:B_R]^2.
\]
In particular $rN(\alpha)$ is a square.

Assume no conjugate is pseudo-unitary and maximize $E$ by transport.
Then $q_0,\alpha_0$ are their respective minima and $0<\alpha_0<1$.
If $E$ were rational, $s\le t$ would give
$\operatorname{Tr}(\alpha)\le5$, hence $\alpha=1$. Also $\alpha$
cannot be rational, so it has degree five. Write its polynomial
$M(X)=X^5-AX^4+BX^3-CX^2+HX-W$, where $A>5$.
The integral columns of Corollary~\ref{cor:galois-columns}
are $(\sigma(\alpha),r)$. Taking
elementary symmetric polynomials and subtracting the corresponding
scalar integers gives
\[
 M(X)\equiv(X-r)^5\pmod g.
\]
Taking unit coefficients gives
$\operatorname{Tr}(q\sigma(\alpha))=s+n_\sigma$ for nonnegative
integers $n_\sigma$, with $n_1=0$ and $\sum n_\sigma=N=At-5s$.
Fourier convolution, Parseval, and $|\widehat q(\eta)|^2\le1/c$
yield
\[
 N^2/4\le5\textstyle\sum n_\sigma^2-N^2
 \le(5\operatorname{Tr}(\alpha^2)-A^2)/c,
 \qquad cN^2\le16A^2-40B.
\]

Let $i\ne0$. Fourier inversion and
$\sum_{k=1}^4|1+\zeta_5^k|=2\sqrt5$ give
\[
 q_0\ge(t-4/\sqrt c)/5,\qquad
 q_0+q_i\ge (2/5)(t-\sqrt{5/c}).
\]
If these two lower bounds are at least $m,a>0$ respectively and
$m\le a/2$, then $\prod_iq_i\ge\min\{m(a-m)^4,(a/2)^5\}$.
Indeed the minimum of $x(a-x)^4$ on $[m,a/2]$ occurs at an endpoint.
For $c\ge31$, one may take $m=1/220$, $a=2/15$, and both products
exceed $1/(4\cdot31^4)$, contradicting the norm formula. The only
smaller conductors are $11,25$. At $c=25$, $g\ge6$ gives
$m=1/150$, $a=23/150$, again contradicting $\prod q_i\le1/(6\cdot25^4)$.
Thus at $25$ one has $f=g=4$ or $5$. At either conductor $q_i\ge
(t-\sqrt{5/c})/5$ for every $i\ne0$, so
\[
  A<1+5s/(t-\sqrt{5/c}).
\]
Using $\sqrt{5/11}<27/40$ and $\sqrt{1/5}<9/20$ bounds $A$ by $50$
at eleven and by $13,12$ at twenty-five for $g=4,5$.

At eleven \eqref{eq:odd-cycle-coordinates} gives
$1\le c_j<4A/11$, $\sum c_j=A\le50$. Retain totally positive
$d$-numbers with least conjugate below one. For each polynomial
test $r\mid A$, $rW$ square, $g\mid M(X)-(X-r)^5$ coefficientwise,
and $g/r=f\in\mathbb Z_{\ge3}$. The \codecheck{odd_degree_checks/higher_degree_quintic.py}{exact enumeration}
leaves two profiles, both violating the moment inequality.

\[
\begin{array}{rrrrrr|rr}
 A&B&W&f&g&N&11N^2&16A^2-40B\\\hline
 16&76&1&11&11&10&1100&1056\\
 18&112&32&4&8&12&1584&704
\end{array}
\]

At twenty-five the prime five is totally ramified, hence $5\mid A$.
The trace bounds and $A\equiv5\pmod g$ give $f=g=5$ and $A=10$.
The congruence $M(X)\equiv(X-1)^5\pmod5$ makes $\alpha$ congruent
to one at the ramified prime. We show directly that $\alpha$ cannot
be totally positive.

For each of the five cosets $C_i$ of $\langle7\rangle$ in
$(\mathbb Z/25)^\times$ put $b_i=\sum_{j\in C_i}\zeta_{25}^j$.
These periods sum to zero, and their trace Gram matrix is $25I-5J$.
Indeed each coset reduces bijectively to $(\mathbb Z/5)^\times$.
Among the sixteen sums of exponents from two cosets, four vanish
modulo five; these vanish modulo twenty-five exactly when the cosets
coincide. The cyclotomic traces give the stated matrix.
Consequently one and any four periods have Gram determinant
$5^8=\operatorname{disc}(K)$ and form an integral basis.
Write $\alpha=2+\sum c_i b_i$. Since $b_i\equiv4$ at the ramified
prime, $\sum c_i\equiv1\pmod5$. Thus
\[
 V=\sum_{i<j}(c_i-c_j)^2=5\sum c_i^2-(\sum c_i)^2
 \quad\text{is even and }4\pmod5.
\]
The Gram matrix gives $\operatorname{Tr}(\alpha^2)=20+5V$.
The ten pair products have geometric mean $N(\alpha)^{2/5}\ge1$,
so $20+5V\le10^2-20$. Hence $V=4$. The five coordinates are
equal except for a single entry differing by one, and their sum
modulo five fixes the sign. Thus $\alpha=2+b_i$. But these elements
have a conjugate $2+2\cos(18\pi/25)+2\cos(24\pi/25)<0$.

\end{proof}

\begin{proof}[Proof of Theorem~\ref{thm:odd-degree-descent}]
Rational $D$ gives pseudo-unitarity by \cite[Proposition~8.24]{ENO}.
In a nontrivial universal grading Lemma~\ref{lem:field-corank} gives
rank at least $[K:\mathbb Q]$ in each component. Thus $K$ is cubic
and there are two rank-three components. The neutral component
has dimension $D/2$ and trivial universal grading by the same argument.
Both its rational character algebras are $K$ by
Lemma~\ref{lem:odd-field-equality}, so Lemma~\ref{lem:odd-shift}
gives a pseudo-unitary conjugate. This property passes to the whole
extension by \cite[Proposition~2.8]{BJ}.

For trivial grading and cubic $K$, an irrational complementary codegree generates $K$,
and its row then accounts for the entire dimension-three complement
of the signed algebra. Thus Lemma~\ref{lem:odd-shift} and
Proposition~\ref{prop:odd-cubic-arithmetic} exhaust the possibilities.
The quintic case is Proposition~\ref{prop:odd-quintic-arithmetic}.
\end{proof}

\section{Sphericality in full degree}
\label{sec:sextic-all}

\begin{theorem}\label{thm:sextic-all}
Let $\mathcal C$ be a complex fusion category of rank six. If
$[\mathbb Q(\operatorname{FPdim}\mathcal C):\mathbb Q]=6$, then a
Galois conjugate of $\mathcal C$ is pseudo-unitary. In particular,
$\mathcal C$ admits a spherical structure.
\end{theorem}

\begin{lemma}[A quadratic norm obstruction]\label{lem:quadratic-trace-transfer}
Let a real quadratic field $F=\mathbb Q(\sqrt d)$, with $d$ squarefree,
act by self-adjoint operators on a rational quadratic space $V$ with an
orthonormal rational basis. If its dimension over $F$ is odd, then
$-1\in N_{F/\mathbb Q}(F^\times)$. In particular no prime
$3\pmod4$ divides $d$.
\end{lemma}
\begin{proof}
Self-adjointness writes the pairing as $\operatorname{Tr}_{F/\mathbb Q}B$
for a nondegenerate symmetric $F$-bilinear form: define $B(u,v)$ by
$\operatorname{Tr}(aB(u,v))=\langle au,v\rangle$ for $a\in F$.
Diagonalizing $B$ gives rational discriminant
$d^{\dim_F V} N_{F/\mathbb Q}(\det B)$ modulo squares. This
discriminant is one, so $d$ is a norm; rational squares are norms.
Dividing by $N(\sqrt d)=-d$ gives a norm of $-1$. Then $d$ is a
sum of two rational squares, so its squarefree part has no prime
divisor $3\pmod4$.
\end{proof}

\subsection{The two reciprocal dimension vectors}

Put $D=\operatorname{FPdim}\mathcal C$, $K=\mathbb Q(D)$, and
transport the category so that its intrinsic dimension $E$ is largest among
its conjugates. The ordinary rational algebra is $K$, since its
dimension is six and its FP image contains $D$. The signed dimension
field contains $E$ and the induction dimension $E/D$, so the signed
rational algebra is also $K$. Every ordinary simple is self-dual by
injectivity of the FP map. Cyclotomicity of the codegrees and total
reality show that $G=\operatorname{Gal}(K/\mathbb Q)=\langle g\rangle$
is cyclic of order six. Write $p_i=g^i(1/D)$, $q_i=g^i(1/E)$,
$x_i=p_i/q_i$, and $x=x_0$. Both $p,q$ are positive probability
vectors, and $p_0,q_0$ are their respective minima. Also $x_0\le1$
and $x_0\le x_i$ by \eqref{eq:least-ratio}.

For the two orders $R,T\subset K$, Lemma~\ref{lem:trace-order-conductor}
gives
\[
 N(D)=\Delta[\mathcal O_K:R]^2,\qquad
 N(E)=\Delta[\mathcal O_K:T]^2,\qquad \Delta=\operatorname{disc}(K).
\]
Thus the totally positive integer $x=\dim A_{\rm FP}$ has square
norm and invariant ideal, being a quotient of $d$-numbers.
Lemma~\ref{lem:mass-moments} gives positive integers
\[
 h_j=\sum_iq_i x_{i+j},\qquad h_0=1,\qquad
 \sum_j h_j=\operatorname{Tr}(x).
 \tag{S1}
\]

It suffices to show that $p$ and $q$ are conjugate, since their least
entries are $1/D$ and $1/E$. First we eliminate every proper subfield
with the rearrangement used in the cubic case.

\begin{lemma}\label{lem:sextic-proper-ratio}
If $x$ belongs to a proper subfield of $K$, then $E=D$.
\end{lemma}
\begin{proof}
A rational $x$ is a positive integer at most one. Otherwise let
$L=\mathbb Q(x)$ and suppose $E<D$. The factorization in
Theorem~\ref{thm:proper-ratio} writes $E=\delta z$ with
$z\in L$, $F=\mathbb Q(\delta)\subseteq K_{\mathcal A}\subseteq K$,
$K=FL$, and $\operatorname{Tr}_{K/L}(1/\delta)=1$. Since
$[F:\mathbb Q]\le3$ and the group fixing $L$ acts faithfully on $F$,
the degrees of $F,L$ are two and three in some order. Hence
$\operatorname{Gal}(K/F)$ induces all embeddings of $L$. As
$1/E$ is least at the distinguished embedding and $\delta$ is fixed
by this group, $u=1/z$ is least there too.

Taking the relative trace in (S1) gives $\operatorname{Tr}_L(u)=1$
and positive integers $\operatorname{Tr}_L(u\rho(x))$ for all
$\rho\in\operatorname{Gal}(L/\mathbb Q)$, equal to one at the identity.
If $u$ is rational then $\operatorname{Tr}_L(x)=[L:\mathbb Q]$,
forcing $x=1$ by the norm. Otherwise $u$ and $x$ have strict minima
at the identity, as $[L:\mathbb Q]$ is prime. Lemma~\ref{lem:short-orbit-rearrangement}
gives a positive integer smaller than one.
\end{proof}

\subsection{A finite list of conductors}

Let $a,b,c$ be the conductors of the quadratic, cubic and faithful
characters of $G$, respectively. Then $a\ge5$ is a fundamental
discriminant. The primitive cubic conductor $b\ge7$ is a product of
distinct primes $1\pmod3$, possibly multiplied by $9$. These
conditions follow by taking the order-three quotients of the unit
groups modulo prime powers. Characters of coprime orders do not
cancel on a local unit subgroup, so $c=\operatorname{lcm}(a,b)$.
The conductor--discriminant formula gives $\Delta=ab^2c^2$. 
For the transforms $P_k=\sum_i p_i e^{2\pi\mathrm i ki/6}$ and
$Q_k=\sum_i q_i e^{2\pi\mathrm i ki/6}$, the conductors are
$f_1=c$, $f_2=b$, $f_3=a$. We use the following inequalities from
Lemma~\ref{lem:sextic-matrix-compression}, with $f=f_k$:
\[
 f|P_k|^2\le(f,2),\quad f|Q_k|^2\le(f,4),\quad
 f(2|P_k|^2+|Q_k|^2)\le4,\quad
 \sqrt f\,|Q_k|+f|P_k|^2\le2.\tag{S2}
\]

Write $d=(a,2)$. Fourier inversion bounds each $p_i$ below by
$m=(1-\sqrt{d/a}-2/\sqrt b-2\sqrt{d/c})/6$. If $m>0$, the product
of six such coordinates of sum one is at least $m^5(1-5m)$, by
successively moving all excess over $m$ into one coordinate.
Consequently $\Delta m^5(1-5m)\le1$. For $a\ge578$ we have
$a/d,c/d\ge289$, $2/\sqrt7<19/25$ and $m>9/850$; substitution of
$\Delta\ge49\cdot578^3$ contradicts the inequality. For $b\ge121$,
$d/a\le1/4$, $c/d\ge b$ give $m\ge1/44$, and the same contradiction
follows from $\Delta\ge5\cdot121^4$. The quadratic subfield acts by self-adjoint multiplication on the
ordinary algebra $K$, of dimension three over that subfield.
Lemma~\ref{lem:quadratic-trace-transfer} therefore excludes $a$ with
a prime factor $3\pmod4$.

The 3009 candidate pairs in the resulting rectangular bounds give
94 after the product test and 41 after the norm test:
\[
\begin{array}{c|l}
 a& b\\\hline
 5,8&7,9,13,19,31,37,43\\
 13,17&7,9,13\\
 29,37,40,41&7,9\\
 65,104&7,13\\
 53,61,73,85,89,97,101,109,136&7
\end{array}
\]
Each listed $b$ has a unique cubic field $K_b$. The
\codecheck{sextic_complement/sextic_conductor_sieve.py}{conductor calculation}
uses rational upper bounds for $1/\sqrt n$.

\subsection{A height bound for the dimension ratio}
Set $r_f=\sqrt{(f,4)/f}$. Fourier inversion, applied to $q_0+q_i$
and using $2q_i\ge q_0+q_i$, gives the respective lower bounds
\[
 6q_{\pm1}\ge1-r_b-\sqrt3r_c,\qquad
 6q_{\pm2}\ge1-r_a-r_b-r_c,\qquad 6q_3\ge1-2r_b .
\]
For even $a$ we also use
$24q_{\pm2}\ge4-2\sqrt{6/a}-6/\sqrt b-4\sqrt{3/c}$.
Indeed $2p_0+q_0+q_{\pm2}\le4q_{\pm2}$; in its Fourier expansion,
use the mixed bound at $k=1,3$, and the separate bounds at $k=2$.
If the resulting common lower bound is $\beta>0$, then
$\operatorname{Tr}(x)\le1+(1-p_0)/\beta<1+1/\beta$.
This bounds the trace in 38 fields. For $(8,13)$ the
\codecheck{sextic_complement/sextic_matrix_bound.py}{Fourier certificates}
give $10^4(q_1,q_2,q_3)>(638,63,742)$ coordinatewise. They take
nonnegative rational combinations of the minimum conditions and
supporting half-spaces of (S2), using
$2\operatorname{Re}(\bar z\sqrt f P_k+\bar u\sqrt f Q_k)\le2+|z|^2$
for $|u|=1$. Reflection covers $q_4,q_5$, so $\operatorname{Tr}(x)\le160$.
Only $(8,7)$ and $(8,9)$ still require a height bound.

\begin{lemma}\label{lem:sextic-two-column-moment}
For all $i,j\in\mathbb Z/6$ one has
\begin{align}
 p_0 q_{i+j}&\le p_i q_j,\label{eq:sextic-effective-columns}\\
 \sum_l \frac{p_{l+i}^2}{p_l}&\ge
 \sum_l\frac{q_{l+i}^2}{q_l},\label{eq:sextic-effective-grams}\\
 k_j=\sum_l\frac{p_l p_{l-j}}{q_l}&\in\mathbb Z_{>0}.
 \label{eq:sextic-effective-integer}
\end{align}
\end{lemma}
\begin{proof}
Apply Lemma~\ref{lem:column-comparison} to $W_{g^i}$, whose ordinary
Perron value is $p_i/p_0$ and whose signed conjugate values are
$q_{i+j}/q_j$. Perron domination gives \eqref{eq:sextic-effective-columns};
the squared-norm comparison in the two trace forms gives
\eqref{eq:sextic-effective-grams}. For \eqref{eq:sextic-effective-integer}
apply Lemma~\ref{lem:mass-moments} to
$x\,g^{-j}(x)\,g^{-j}(q)/q$ and use $p=xq$.
\end{proof}

\begin{proposition}\label{thm:sextic-effective-height}
If $(a,b)=(8,7)$ or $(8,9)$, then
$\operatorname{Tr}_{K/\mathbb Q}(E/D)\le200$.
\end{proposition}
\begin{proof}
We show that $q_j>1/200$ for $j=2,4$. It suffices to take $j=2$
by replacing $g$ by $g^{-1}$. Suppose $q_2\le1/200$; then
$p_0,q_0,q_2\le1/200$. For either $v=p,q$, Fourier inversion
and the separate bounds in Lemma~\ref{lem:sextic-matrix-compression}
give
\begin{align*}
 v_0+v_{\pm1}&\ge(1-1/\sqrt b-\sqrt{3/(2b)})/3,\qquad
 v_0+v_3\ge(1-2/\sqrt b)/3,\\
 v_0+v_2+v_4&\ge(1-1/\sqrt2)/2.
\end{align*}
For instance the first multiplier vanishes at the quadratic
frequency and has magnitude $1,\sqrt3$ at the cubic and sextic
frequencies after pairing conjugates. As $v_0$ is minimal, the odd
coordinates exceed $1/45$; also $q_4>1/50$.

For the remaining ordinary coordinates, put
$d=1/\sqrt{8b}+1/\sqrt b+1/\sqrt8<87/100$.
The absolute multipliers in $3(v_0+v_{\pm2})-1$ are all one at
the three paired frequencies (with the quadratic frequency counted
once). Thus $\sum_{k=1}^3 |Q_k|\ge1-3(q_0+q_2)\ge97/100$.
Cauchy's inequality and the mixed bound yield
\[
 \sum_{k=1}^3|P_k|
 \le\sqrt{d\bigl(2d-\sum_{k=1}^3|Q_k|\bigr)}<83/100.
\]
Applying inversion again to $p_0+p_{\pm2}$ proves $p_2,p_4>1/50$.

Set $\lambda=p_0/q_0$, $\mu=p_0/q_2$. Then $0<\mu\le\lambda\le1$.
The integer in~\eqref{eq:sextic-effective-integer} is
\begin{equation}\label{eq:sextic-effective-objective}
 k_2=\lambda p_4+\mu p_2+
         \sum_{i\in\{1,3,4,5\}}p_i p_{i-2}/q_i .
\end{equation}
For $1\le i\le5$, the column and squared-norm inequalities imply
\begin{gather}
 \lambda q_i\le p_i,\qquad \mu q_{i+2}\le p_i,\nonumber\\
 \lambda q_i^2+\mu q_{i+2}^2
 \le p_i^2+p_0\sum_{l=1}^5 p_{l+i}^2/p_l .
 \label{eq:sextic-effective-ab}
\end{gather}
These inequalities contradict $k_2\ge1$ by the
\codecheck{sextic_height_bound.cpp}{exact interval calculation}.
Its domain is the product of the two probability simplices with
$p_0,q_0,q_2\in[0,1/200]$ and all remaining coordinates in
$[1/50,1]$. Use the weaker conductor triple $(56,7,8)$ for both
fields. On each dyadic box, (S2) supplies twelve supporting
half-spaces per Fourier disk; \eqref{eq:sextic-effective-ab}
bounds the first two terms of \eqref{eq:sextic-effective-objective}
and outward interval division bounds the other four. The subdivision visits $8\,312\,559$ boxes. Every terminal box
violates a constraint or gives $k_2<1$.

Thus $q_{\pm2}>1/200$. The pair bounds also give
$2q_i\ge q_0+q_i>2/45$ at each odd $i$, independently of $q_2$.
Since $x_0\le1$,
$\operatorname{Tr}(x)<1+200(1-p_0)<201$.
The trace is integral, as claimed.
\end{proof}

\subsection{The arithmetic conclusion}
It remains to test the ratios $x$ that generate $K$. A relative
discriminant lower bound first eliminates ten conductor pairs.
Apply Lemma~\ref{lem:relative-trace-discriminant} to $K/K_b$.
Its relative discriminant has norm $\delta=ac^2/b^2$ by the tower
formula, so $\operatorname{Tr}(x)\ge3\sqrt{\delta^{1/3}+4}$.

The upper
bounds of $18$ at $(37,9),(41,9)$, and $25$ at $(a,7)$ for
$a=73,85,89,97,101,104,109,136$, contradict this inequality.

\paragraph{Listing the ratios.}
We treat $(13,13)$ separately below. For the other thirty fields we
list the possible $x$ up to conjugacy. In an integral basis $e_i$ let
$H_{ij}=\operatorname{Tr}(e_ie_j)$ and $t_i=\operatorname{Tr}(e_i)$.
At each allowed trace $A$ the coordinate vector $v$ of $x$ satisfies
\[
 t^{\mathsf T}v=A,\qquad v^{\mathsf T}Hv\le A^2.
\]
Positivity gives the second inequality. Rational square completion
enumerates this finite ellipsoid. Retain the $x$ that generate $K$, have positive conjugates with least
conjugate at most one, have square norm, and satisfy the $d$-number
criterion for their minimal polynomial.
For $X^d+a_1X^{d-1}+\cdots+a_d$, the last test is
$a_d^j\mid a_j^d$ for every $j$~\cite{OstrikCodegrees}.

Outside the five fields below the
\codecheck{sextic_complement/sextic_lattice.py}{trace-lattice enumeration}
uses a basis of Gram determinant $ab^2c^2$, hence of the maximal order.
At coprime conductors it is the tensor product basis. For $a=65,104$
with $b=13$ take invariants of the diagonal involution on
$\mathbb Q(\zeta_{13})^+\mathbb Q(\sqrt m)$ with $m=5,2$ respectively.
Of $4\,903\,735$ vectors, 35 remain, distributed as $27,7,1$ at
$(5,13),(13,7),(29,7)$.

In the five larger searches write $x=n\pi^j u^2$, with $n\ge1$ integral,
$0\le j<3$ and $u$ a unit. The $d$-number property gives equal
valuations over each prime, and the square norm makes the valuation
even at the quadratic ramified prime. After removing rational factors
it suffices to remove the remaining valuation at the cubic ramified
prime using a totally positive generator $\pi$ of that prime in $K_b$.
Full unit signature rank makes the residual unit a square. The
\codecheck{sextic_complement/positive_dnumber_squares.py}{unit enumeration}
checks the generators and signatures, lists units in the ellipsoid
$\operatorname{Tr}(\pi^j u^2)\le A_{\max}$, and tests their multiples.

\[
\begin{array}{c|r|r|r}
 (a,b)&A_{\max}&\text{lattice vectors}&\text{primitive ratios}\\\hline
 (5,7)&1037&12\,273\,947&8007\\
 (5,9)&86&5207&128\\
 (8,7)&200&43\,959&284\\
 (8,9)&200&33\,357&206\\
 (8,13)&160&5973&46
\end{array}
\]

\paragraph{Testing the ratios.}
Place the least embedding of $x$ at zero. For the 35 ratios above
and the 128 at $(5,9)$, the
\codecheck{sextic_complement/sextic_trace_systems.py}{trace-system check}
lists the positive integral compositions $h$ of $A$ with $h_0=1$
from (S1). In an integral field basis these are rational linear
equations on $q$. The three quadratic bounds in (S2) have rational
matrices, since
$\cos(2\pi kj/6)\in\mathbb Q$; substituting $p=xq$ and completing
squares on the affine solution space excludes every composition.
Parseval first bounds $\sum h_i^2$ from
$|\widehat h_k|^2=|Q_k|^2|\widehat x_k|^2$.

For the other $8543$ ratios impose $p_i=x_iq_i$, the normalization
and minimum conditions, $h_j\ge1$ and rational supporting
half-spaces of (S2). The
\codecheck{sextic_complement/verify_complement.py}{complementary-field verifier}
and the \codecheck{sextic_certificates/verify_sextic.py}{verifier for $(8,7),(8,9)$}
verify a rational relaxation $A'q\le d'$ and supplied nonnegative
integer weights $w$ with $\min_i(w^{\mathsf T}A')_i>w^{\mathsf T}d'$.
This is impossible for a probability vector.

\paragraph{The real cyclotomic field of conductor thirteen.}
Finally at $(13,13)$, Parseval bounds both second moments by $3/13$.
The pair bounds give $q_{\pm1}>1/25$, $q_{\pm2}>1/36$, $q_3>1/14$.
The six Fourier coefficients of $(0,25,36,14,36,25)$ are
$(136,-25,-47,8,-47,-25)$. Hence
$A<1+(136+152/\sqrt{13})/6<31$. In the basis $v_j=2\cos(2\pi j/13)$, \eqref{eq:odd-cycle-coordinates}
gives $x=-\sum d_jv_j$ with $1\le d_j\le9$ and $\sum d_j=A\le30$.
The
\codecheck{sextic_complement/sextic13_descent.py}{conductor-thirteen calculation}
leaves fifteen primitive ratios, eleven units and two each of norms
$64,169$. For every one of them $x\mathcal O_K\subseteq S=\mathbb Z[Gx]$.
To check this containment, begin with the integer span of $1,Gx$ in
the period basis and adjoin products by $Gx$ twice. The resulting
lattice contains all six columns of multiplication by $x$.

Since $\operatorname{Tr}(q S)\subseteq\mathbb Z$ by
Lemma~\ref{lem:mass-moments}, $p=xq$ is in the ordinary trace dual
of $\mathcal O_K$. Any such vector $z=-\sum_j n_jv_j/13$ with
$\operatorname{Tr}(z)=1$ has integral $n_j=2-\operatorname{Tr}(zv_j)$.
If $\operatorname{Tr}(z^2)\le3/13$, the Gram matrix gives
$\sum n_j=13$ and $\sum n_j^2\le29$. The only possibility is five
twos and one three. Thus $p=(2-v_1)/13$, because $p$ is least at
the distinguished embedding. For unit $x$ we have $S=\mathcal O_K$
and the same argument gives $q=p$. Otherwise the calculated
$\operatorname{Tr}(p/x)$ is $1/2$ at norm $64$ and $7/13$ at norm
$169$, whereas $q$ has trace one.

\begin{proof}[Conclusion of the proof of Theorem~\ref{thm:sextic-all}]
Lemma~\ref{lem:sextic-proper-ratio} gives $E=D$ for a proper-subfield
ratio. The conductor and trace arguments exclude every primitive ratio.
The chosen conjugate is therefore pseudo-unitary.
\end{proof}

\section{Conductor bounds in degree two}
\label{sec:quadratic-global-reduction}

The first restriction comes from the action of the neutral component on the other homogeneous components. Once the universal grading is trivial, we can compare the two quadratic dimension characters directly. The conductor bounds at the end of the section leave a small list of discriminants.
\subsection{Nontrivial gradings in the quadratic case}
\begin{proposition}\label{prop:quadratic-graded-direct}
Let $\mathcal C$ be a complex fusion category of rank six with a nontrivial faithful finite-group grading. If $\mathbb Q(\operatorname{FPdim}(\mathcal C))$ is quadratic, then a Galois conjugate of $\mathcal C$ is pseudo-unitary.
\end{proposition}
\begin{proof}
Let $\mathcal B$ be the neutral component. Its Perron dimension
$D/|G|$ generates the same quadratic field as $D$
\cite[Proposition~8.20]{ENO}. By Lemma~\ref{lem:field-corank}
every component has rank at least two, so $\mathcal B$ has rank at
most four. It suffices to make $\mathcal B$ pseudo-unitary by
transport, since this property passes to graded extensions
\cite[Proposition~2.8]{BJ}.

At ranks two and three $\mathcal B$ is spherical by the rank-two
classification \cite{Ostrik2} or by Ostrik's quadratic rank-three
theorem \cite[Proposition~3.1 and Theorem~5.1]{Ostrik3}. In rank
three the non-self-dual ring is pointed, and the two conjugates of
$D/|G|$ give the required quadratic factor in the self-dual case.
The ordinary algebra has a unique nonrational character orbit.
It contains both dimension characters, since the spherical dimension
field contains $D/|G|$ by central induction. Thus $\mathcal B$ has
a pseudo-unitary conjugate.

At rank four consult the unconditional list
\cite[Corollary~13.1 and Table~1]{Rank4Sphere}. Among the rings with
quadratic global Frobenius--Perron dimension, all have only one value
of simple Frobenius--Perron dimension greater than one, apart from the
Haagerup ring $\mathcal H_1$ and the square of the Fibonacci ring. In the
one-value case Schopieray gives a pseudo-unitary conjugate of
$\mathcal B$ \cite[Lemma~3.1]{SchopExt}. This leaves two rings.

In either case the two exterior simples carry a two-dimensional
representation of $K_0(\mathcal B)$ by nonnegative integral matrices.
A self-dual object acts by a symmetric matrix. The Haagerup ring
has a self-dual simple of dimension $(1+\sqrt{13})/2$
\cite[Table~2]{HaagGS}. Its matrix would have trace one and
determinant $-3$, since its eigenvalue has minimal polynomial
$x^2-x-3$. Nonnegative integral diagonal entries summing to one
have product zero, forcing the off-diagonal entry to have square three.

For the Fibonacci square choose a spherical section of $\mathcal B$
by~\cite[Theorem~12.3]{Rank4Sphere}. Its two Fibonacci generators
$P,Q$ act by commuting matrices on each two-dimensional exterior
module, ordinary or signed. Each satisfies the irreducible equation
$M^2=M+I$. The centralizer of the $P$-matrix is therefore the
quadratic field it generates: the $Q$-matrix is $P$ or $I-P$,
writing $P$ also for its matrix. In the ordinary module the common
Perron eigenvector gives the same positive eigenvalue, hence $Q=P$.
In the signed module $Q\equiv P\pmod2$ now excludes $I-P$.
Its nonzero dimension functional gives $\dim P=\dim Q=z$ with
$z^2=z+1$. Thus $\dim(\mathcal B)=(1+z^2)^2$ is conjugate to
$\operatorname{FPdim}(\mathcal B)$.
\end{proof}

\subsection{The two character algebras}
Suppose the signed dimension field $M$ is quadratic. Write a prime
for its nonidentity embedding and $s,t$ for the masses in
Corollary~\ref{cor:galois-columns}.
At the transport maximizing $E$, $x=E/D$ and $1/E$ are both least
at the distinguished embedding. Thus
\begin{equation}\label{eq:quadratic-trace-order}
       2s=t\operatorname{Tr}(x)+(1/E'-1/E)(x'-x)
          \ge t\operatorname{Tr}(x),
 \qquad s=\operatorname{Tr}(1/D),\quad t=\operatorname{Tr}(1/E).
\end{equation}
If $E<D$, the totally positive algebraic integer $x$ has trace at least three
by its integral norm; at even discriminant the trace is even, hence
at least four.

\begin{proposition}\label{prop:quadratic-full-character}
Suppose $K=\mathbb Q(D)$ has degree two and no Galois conjugate of $\mathcal C$ is pseudo-unitary. Its universal grading is trivial. Both quotient algebras are commutative, both dimension characters have full field $K$, and at least one signed codegree outside the spherical dimension orbit is irrational.
\end{proposition}
\begin{proof}
Proposition~\ref{prop:quadratic-graded-direct} settles the graded case. If the universal grading is trivial, Lemma~\ref{lem:odd-field-equality} identifies both dimension-character fields with $K$.

Maximize $E$ by transport. The quadratic trace identity
\eqref{eq:quadratic-trace-order} gives $s\ge3t/2>t$ for the two
dimension-orbit masses.

If every complementary signed codegree were rational,
Corollary~\ref{cor:galois-columns} would give $s\le t$. Thus at
least one is irrational. A noncommutative signed algebra in rank
six would have complexification $\mathbb C^2\times M_2(\mathbb C)$,
with its two scalar blocks in the dimension orbit. Its only other
codegree would be rational by Galois invariance.

If instead $R_{\mathbb C}=\mathbb C^2\times M_2(\mathbb C)$, let $h$
be its rational matrix-block codegree. The Casimir $\sum b_i b_i^*$
has eigenvalue $2h$ on the matrix block and Perron value $D$, so
$D\ge2h$. From $1/D+1/D'+2/h=1$ we get
$1/D'-1/D\ge1-3/h$. The square of this gap is at most
$(c,2)/c\le1/4$ by~\eqref{eq:quadratic-fourier-bounds}, so $h\le6$.

Since $R/2R\simeq T/2T$ is commutative, the alternating tensor
$C_{ijk}=\tau([b_i,b_j]b_k)$ has even entries, where $\tau$ extracts
the coefficient of the unit. The sum of their squares is $12h$.
Indeed, expanding $\sum_{i,j}\|[b_i,b_j]\|^2$ gives twice the
difference of the Casimir square trace $\sum_V n_V^3f_V$ and the
trace of $x\mapsto\sum_i b_ixb_i^*$, equal to $\sum_V n_Vf_V$
by matrix units. Only the two-dimensional block contributes.

The regular Perron element annihilates every commutator, so
$\sum_k d_k C_{ijk}=0$ for every pair $i,j$. Take a triple with
$C_{ijk}\ne0$. Each of its three pairs must belong to another
nonzero triple, because every $d_k$ is positive. These three further
triples are distinct. Each unordered triple contributes at least
$6\cdot2^2$ to the sum of squares. Thus $12h\ge4\cdot24$,
contradicting $h\le6$.
\end{proof}

\begin{lemma}[Quadratic character rows]\label{lem:quadratic-row-difference}
Let $a,a'$ be conjugate linear character rows of either quotient
with quadratic character field of discriminant $\Delta$ and
codegrees $f,f'$. Write $a_X-a'_X=b_X\sqrt\Delta$. Then
\begin{equation}\label{eq:quad-row-trace}
      b_X\in\mathbb Z,\qquad f+f'=|\Delta|\sum_X b_X^2 .
\end{equation}
If neither row is the positive dimension row of its quotient, then
$f+f'\ge2|\Delta|$. If $\Delta<0$, their common codegree is
$|\Delta|\sum_{\{X,X^*\}}b_X^2$, summed over the unordered
pairs with $X\ne X^*$.
\end{lemma}
\begin{proof}
The integral basis of the quadratic field gives $b_X\in\mathbb Z$.
Schur orthogonality gives $\|a-a'\|^2=f+f'$.
When both rows are orthogonal to the positive dimension row,
their difference has at least two nonzero entries. For $\Delta<0$,
complex conjugation acts by duality, so $b_{X^*}=-b_X$.
\end{proof}

\begin{lemma}\label{lem:fixed-quadratic-mass}
Let $f,f'$ be conjugate totally positive quadratic $d$-numbers
and write $1/f+1/f'=a/b$ in lowest terms. For some positive integer $m$,
\[
 ff'=b^2m,\qquad f+f'=abm,\qquad
 (1/f-1/f')^2=\frac{a^2-4/m}{b^2}.
\]
If an odd prime $p$ divides $a$, then $f/f'\equiv-1$ at
every prime over $p$.
\end{lemma}
\begin{proof}
With $N=ff'$ and $L=f+f'$, the unit $f/f'$ gives
$L^2/N=2+f/f'+f'/f\in\mathbb Z$. Since $L/N=a/b$, we get
$b^2\mid N$. The three identities follow.
For an odd prime $p\mid a$, the identity
$1+f/f'=fa/b$ gives $f/f'\equiv-1$ at each prime above $p$.
Thus $p$ cannot divide $a$ if the least totally positive unit has
odd residue order at one such prime.
\end{proof}

In particular every signed codegree belongs to $K$, since $E/f=\dim F(A_f)$ lies in $K$. In the four complementary dimensions an irrational codegree therefore has its full character field of degree two or four. A quartic character field containing $K$ gives two copies of the same codegree pair. If an irrational orbit uses only two characters, its field is $K$ and the remaining two codegrees are another such pair or two rational integers. The last two characters in the rational-codegree case can be rational or can be conjugate over a separate quadratic field. The spherical dimension codegree $E$ may be rational even though its full character field is $K$.

\subsection{A conductor test for quadratic pairs}
For a codegree $f$ in the real quadratic field of discriminant $c$,
write $t_f=1/f+1/f'$ and $\delta_f=1/f'-1/f$. Let
$U=(a+b\sqrt c)/2>1$ be the least totally positive unit.
For irrational $f$, the larger of $f/f'$ and $f'/f$ is a positive
power of $U$, by the $d$-number property. Consequently
\begin{equation}\label{eq:r6-pair-gap}
       \delta_f^2\ge\lambda t_f^2,\qquad
       \lambda=\left(\frac{U-1}{U+1}\right)^2=\frac{a-2}{a+2}.
\end{equation}

Suppose the signed dimension pair and $k$ further irrational pairs
have total reciprocal mass $L$. Put $q=(c,4)/c$ and $y_+=\max(y,0)$.
Counting each pair with its multiplicity, \eqref{eq:quadratic-fourier-bounds}
and Cauchy's inequality give
\begin{equation}\label{eq:quadratic-conductor-test}
 \begin{cases}
    \lambda L^2/(k+1)\le q,&E\notin\mathbb Q,\\
    \lambda (L-2(c,2)/c)_+^2/k\le q,&E\in\mathbb Q .
 \end{cases}
\end{equation}
In the second case orthogonality of the two dimension rows gives
$2E=\sum_i(x_i-x_i')^2\in c\mathbb Z$; hence their mass $2/E$ is at
most $2(c,2)/c$.

\begin{proposition}[Three quadratic codegree pairs]
\label{prop:r6-three-quadratic-pairs}
Let the universal grading be trivial and $K=\mathbb Q(D)$ quadratic.
If $T$ is commutative with codegrees $E,E',H,H',J,J'$, where $H,J$
are irrational, then $c=5$ or $8$. The pairs may coincide.
\end{proposition}
\begin{proof}
By Lemma~\ref{lem:odd-field-equality} the dimension character field is
$K$, and the complementary codegrees generate $K$. Thus each field
factor of $T_{\mathbb Q}$ contains a copy of $K$. Their diagonal action
is self-adjoint for $\tau(uv^*)$, since involution on characters
acts by complex conjugation. As the signed basis is orthonormal and
$\dim_K T_{\mathbb Q}=3$, Lemma~\ref{lem:quadratic-trace-transfer}
excludes prime divisors $3\pmod4$ of $c$. An odd $c>5$ is at least
$13$. An even $c>8$ is at least $40$: it has the form $8d$ for an
odd squarefree $d>1$ with all prime factors $1\pmod4$.

Here $k=2$ and $L=1$, so
\begin{equation}\label{eq:r6-three-pair-fourier}
   t_E+t_H+t_J=1,\qquad
   \delta_E^2+\delta_H^2+\delta_J^2\le (c,4)/c .
\end{equation}
For odd $c\ge13$, $a^2-cb^2=4$ gives $a\ge5$, hence
$\lambda\ge3/7$. Both left sides of \eqref{eq:quadratic-conductor-test}
exceed $1/13$ already at $c=13$. For even $c\ge40$, use
$\lambda\ge1/2$; the corresponding left sides exceed $4/40$
already at $40$. They are nondecreasing with $c$ at either parity.
\end{proof}

\subsection{Separation of quadratic codegrees}
\begin{lemma}[Separation of the quadratic pairs]\label{lem:quadratic-mass-gap}
Let $K$ be a real quadratic field whose discriminant $c$ appears
in the table below. Suppose both dimension characters generate $K$,
both character algebras are commutative, and outside the dimension pair $E,E'$, the signed codegrees
consist of an irrational pair $H,H'$ and either a second copy of this
pair or two rational integers. Assume there is no pseudo-unitary
Galois conjugate. Transport to $E\ge E'$ and label $H>H'$.
Write $s=\operatorname{Tr}(1/D)$, $t=\operatorname{Tr}(1/E)$,
$w=\operatorname{Tr}(1/H)$.
In the table $U=(A+b\sqrt c)/2$ is the least totally positive unit
greater than one, $V=(1+U)/d$, and $h$ generates the trace ideal
$\operatorname{Tr}(\mathcal O_K/V)$.
\[
\begin{array}{c|rrrrrrrrr}
 c&5&8&12&13&21&24&28&40&60\\ \hline
 A&3&6&4&11&5&10&16&38&8\\
 b&1&2&1&3&1&2&3&6&1\\
 d&1&2&1&1&1&2&3&2&1\\
 h&1&2&1&1&1&2&1&2&1
\end{array}
\]
Then
\[
   \frac{s}{t}\ge\frac{A}{\max(2,d)},\qquad D/D'\ge U^2,
\]
with the sole possible exception $c=28$, $E=E'$ and
$D/D'=H/H'=U$. For $d=1$, equality in the mass bound forces
$E=E'$ and $D=EU$. For $c\ne28$ we also have $H/H'\ge U^2$.
\end{lemma}
\begin{proof}
The units in the table follow by taking the least positive $b$ for
which $a^2-cb^2=4$ has an integral solution with
$(a+b\sqrt c)/2\in\mathcal O_K$. In the standard integral basis $V$ has
coprime coordinates, $V/V'=U$, and $V>2$. Evaluating the trace
on that basis gives the last line of the table. The
\codecheck{quadratic_trace_bounds.py}{quadratic trace calculations}
also verify its entries. Every totally positive
integral $d$-number, oriented with its larger conjugate first, has the
form $mU^k$ or $mVU^k$, with $k\ge0$ and $m$ a positive integer.
Indeed its quotient by its conjugate is a power of $U$. Dividing by
the indicated element of the same quotient gives a rational scalar,
integral by primitivity.

Apply this observation to $\alpha=E/D$. It is a totally positive
algebraic integer by the induction formula, and its ratio to its
conjugate is a unit by Lemma~\ref{lem:central-witness}.
By~\eqref{eq:least-ratio}, $\alpha<1<\alpha'$.
Putting $r=E/E'\ge1$ gives the identity
\[
 \frac{s}{t}=\frac{\alpha+r\alpha'}{1+r}
       \ge\tfrac12\operatorname{Tr}(\alpha)>1 .
\]
If $\alpha=mU^{-k}$ or $mV'U^{-k}$ with $k\ge1$, its trace is at
least $A$, with equality only for $\alpha=U^{-1}$. Indeed
$VU^k>2U>A$ in the second family. Equality in the displayed
average requires $r=1$. In either family $D/D'=r\alpha'/\alpha\ge U^2$.

There remains $\alpha=mV'<1$, which requires $d>1$.
If $r>1$, then $r\ge U$ and the weighted average gives
\[
    s/t\ge m(V'+UV)/(1+U)=mA/d,\qquad D/D'=rU\ge U^2 .
\]
If $r=1$, then $D/D'=U$. The simultaneous column of
Lemma~\ref{lem:central-witness} has Perron dimension $D/D'$ and
an eigenvalue $H/H'$. Thus $U\le H/H'\le D/D'$, forcing $H/H'=U$.

We now exclude $H/H'=U$ for $c\ne28$. On the ordinary Perron induction
column the signed dimension pair contributes $s>t$, and every
rational-codegree orbit contributes at least its reciprocal mass.
Indeed its eigenvalues are totally positive algebraic integers by
Proposition~\ref{prop:quadratic-full-character}, hence their average
is at least one. The eigenvalue on a complementary $H$-row is
$\beta=H\operatorname{FPdim}(FA_H)/D\in\mathcal O_K$, even when
its character field is quartic. Comparing unit coefficients with the
reciprocal identity therefore gives
$\operatorname{Tr}(\beta/H)<\operatorname{Tr}(1/H)$ on at least one
conjugate pair.
If $H/H'=U$, then $H=nV$. For $c\ne28$ the last line of the table
gives $\operatorname{Tr}(\beta/V)\ge d=\operatorname{Tr}(1/V)$,
a contradiction.
\end{proof}

\subsection{Repeated pairs at conductors five and eight}
\begin{proposition}[A repeated complementary pair over $\mathbb Q(\sqrt5)$]
\label{prop:r6-repeated-pair-five}
Under the hypotheses of Proposition~\ref{prop:r6-three-quadratic-pairs},
suppose that $K=\mathbb Q(\sqrt5)$ and $J=H$.
Then $\mathcal C$ has a pseudo-unitary Galois conjugate.
In particular this settles the quartic complementary character-field
case over $\mathbb Q(\sqrt5)$.
\end{proposition}
\begin{proof}
Suppose otherwise, and pass to a Galois conjugate for which $E\ge E'$.
Put
\[
 u=\frac{3+\sqrt5}{2},\qquad v=\frac{5+\sqrt5}{2},\qquad
 t=\operatorname{Tr}(1/E),\quad w=\operatorname{Tr}(1/H),\quad
 s=\operatorname{Tr}(1/D).
\]
Then $t+2w=1$. Lemma~\ref{lem:quadratic-mass-gap} gives
$H/H'\ge u^2$, so
\[
 \delta_H^2\ge\frac59w^2.
\]
Because $\delta_E=t-2/E\ge0$, the parity bound and Cauchy's inequality
now imply
\[
 \frac15\ge\delta_E^2+2\delta_H^2
 \ge\frac5{23}(\delta_E+2w)^2
 =\frac5{23}(1-2/E)^2.
\]
It follows that $E\le5(5+\sqrt{23})<49$.

Orthogonality of the signed dimension row to its conjugate gives
$\operatorname{Tr}(E)=\sum_i(x_i-x_i')^2\in5\mathbb Z$.
Write $E=mP$, $H=nQ$ using the two primitive families of
Lemma~\ref{lem:quadratic-mass-gap}, and put $B_P=\operatorname{Tr}(1/P)$.
Integrality of $E/H$ gives $\ell=m/n\in\mathbb Z$; if $P$ is a
unit and $Q$ is in the $v$ family, also $5\mid\ell$. The mass
identity reads
\[
     m=B_P+2\ell B_Q,\qquad \ell\mid B_P.
\]
Here $B_Q\le22$ since $m\le E<49$ and $m\le45$ for rational $P=1$
(the other $P$ give $m<19$). The gap bound gives $Q=u^j$ or
$vu^j$ with $j\ge1$; the corresponding traces up to $22$ are
$3,7,18$ and $2,5,13$, by the recurrence $B_{j+2}=3B_{j+1}-B_j$.
For $P=u^k$ with $k\ge2$, $m>B_P\ge7$ contradicts $E<49$;
similarly $k\ge2$ gives $m>B_P\ge5$ for $vu^k$.
Thus for each $P=1,u,v,vu$, take the divisors $\ell$ of $B_P$
and the listed $Q$. Keeping $E=mP<49$, $5\mid m\operatorname{Tr}(P)$
and $H=(m/\ell)Q$ integral gives
\[
\begin{array}{c|cccc}
 (E,H)&(5v,5vu)&(7v,7u)&(11v,11vu^2)&(30,15u^2)\\ \hline
 \delta_E^2+2\delta_H^2&33/125&51/245&243/605&2/5
\end{array}
\]

Every value exceeds $1/5$.
\end{proof}

\begin{theorem}[The repeated complementary pair at conductor eight]
\label{thm:r6-repeated-pair-eight}
Under the hypotheses of Proposition~\ref{prop:r6-three-quadratic-pairs},
suppose $K=\mathbb Q(\sqrt2)$ and $J=H$. Then a Galois conjugate of
$\mathcal C$ is pseudo-unitary.
\end{theorem}
\begin{proof}
Assume the contrary. Write $u=3+2\sqrt2$, $v=2+\sqrt2$. In the notation
of Lemma~\ref{lem:quadratic-mass-gap}, $t+2w=1$. Write $E=mP,H=nQ$,
where $P,Q$ are primitive of the form $u^i$ or $vu^i$, $i\ge0$.
The traces $B_P=\operatorname{Tr}(1/P)$ are all $2\pmod4$: in each
family they satisfy the recurrence with coefficient six and start with
$(2,6)$ or $(2,10)$. Integrality of $E/H$ implies $n\mid m$, so
$m=B_P+2(m/n)B_Q\equiv2\pmod4$. Since $8\mid\operatorname{Tr}(E)$
by row orthogonality, $E=mvu^k$. The same argument at $D$,
using $E/D$ integral, gives $D=Mvu^\ell$, with $M\mid m$,
$M\equiv2\pmod4$ and $\ell>k$. Consecutive terms of $(2,10,58,\ldots)$
increase by a factor at least five, so $s\ge5t$.

Every field factor of $T_{\mathbb Q}$ contains $K$, since its codegree
is irrational outside the dimension orbit. The discriminant tower
formula therefore gives $v_2(\operatorname{disc}\mathcal O)\ge
3v_2(\operatorname{disc}K)=9$ for the maximal order $\mathcal O$.
On the other hand
\[
 [\mathcal O:T]^2|\operatorname{disc}\mathcal O|
       =N(E)N(H)^2,\qquad
 v_2(N(E))=3,\qquad v_2(N(H))\le3.
\]
Here $N(Q)=1$ or $2$ and $n\mid m$. Equality forces $T$ to be
maximal at two and $H=nvu^j$. Lemma~\ref{lem:quadratic-mass-gap}
gives $j\ge1$.

Every rational ordinary codegree $g$ is even: otherwise
Lemma~\ref{lem:trace-order-conductor} splits off an unramified factor
at two. But $R/2R\simeq T/2T$ has no reduced factor, since $T$ is
maximal at two and every field factor contains $\sqrt2$, nonzero
and square-zero in its maximal order modulo two. As $E/g$ is
integral, primitivity gives $g\mid m$, hence $g\equiv2\pmod4$.

If $g=2$, the reciprocal identity forces the character to be
rational. Its row has two nonzero entries $1,\pm1$, giving a
primitive idempotent $(1\pm X)/2$ and hence $X^2=1$ for a
nonunit simple $X$. Its positive lift is likewise an involution in
$T$. Its coordinates in the maximal order are all $\pm1$, so its
class is the unit modulo two by maximality. This contradicts the
linear independence of the distinguished basis in $T/2T$.

The formulas for $D$ and $H$ give $D/D',H/H'\ge u^3$. Their squared
reciprocal gaps are at least $49/50$ times their squared masses.
Suppose first $s<1/5$. Then $w=(1-t)/2>12/25$. For any irrational
ordinary pair $G,G'$ the mixed Fourier bound gives
\[
 \delta_G^2\le\tfrac14-\delta_H^2
       <\tfrac14-\tfrac{49}{50}(12/25)^2<\tfrac1{18}.
\]
All ordinary codegrees lie in $K$, and every irrational pair has
ratio at least $u$, hence squared gap at least half its squared mass.
It therefore contributes less than $1/3$, counting each pair with its
multiplicity. Each rational codegree contributes at most $1/6$.
The four complementary characters together contribute at most $2/3$,
contradicting $s<1/5$.

Thus $s\ge1/5$. Keeping just the $D$- and $H$-pairs in the mixed
Fourier bound gives
\[
 \tfrac12\ge2\delta_D^2+2\delta_H^2
  \ge\tfrac{49}{25}\bigl(s^2+(1-s/5)^2/4\bigr)
  \ge\tfrac{8281}{15625}>\tfrac12 .
\]
The last quadratic is increasing for $s\ge1/5$.
\end{proof}
\subsection{Two irrational pairs and two rational codegrees}
\begin{theorem}[No signed codegree two in the quadratic profile]
\label{thm:r6-quadratic-no-two}
Suppose the ordinary and signed character algebras of a rank-six
fusion category are commutative, both dimension-character fields
are the same quadratic field $K$, and the signed codegrees are
\[
 E,E',H,H',g_1,g_2,
\]
where the $H,H'$ rows form a conjugate pair and the $g_i$-rows are rational. Then $g_1,g_2\ne2$.
\end{theorem}
\begin{proof}
For the induction label $P$ of such a character $\psi$, set
$a_B=4S_{P,B}=\psi(FB)\in\mathbb Z$, $v_B=a_B\dim B$ and
$p=\operatorname{Tr}(1/D)$. Orthogonality to $JP$ gives squared
norm eight in each monodromy sector. On a $J$-pair the entries are opposite; its two
twists are opposite in the signed sector and equal in the ordinary one. Orthogonality
to $1,J$ and the modular identity at $(P,P),(P,1)$ give
\[
 \sum_{B\text{ in either sector}}v_B=0,\qquad
 \sum_B a_B^2\theta_B=4a_P,\qquad \sum_B v_B\theta_B=2E .
\]
The signed $J$-pairs cancel from both twist sums.

In each sector the six induction entries sum to two by the
unit-coefficient formula. On the dimension pair the entries are equal
positive integers by reciprocity: the two are conjugate, and lie in
$\mathbb Z$ since $P$ is rational. Their squared norm, including the
translates, is at most eight, so both are one. The four remaining
entries thus sum to zero with squared norm at most two. In the signed
sector the entries on the conjugate pair agree, so they vanish, and
the rational entries are $b,-b$ with $a_P=b$. In the ordinary sector
a nonzero pair $1,-1$ would exhaust the norm at twist one, forcing
$4b=8$. Thus the ordinary dimension labels and their translates give
four units at twist one, and $4b\ge0$ by taking real parts.
If $b=0$, all remaining ordinary entries in the support have twist
$-1$. The dimension labels and their translates have $v$-weight
$2Ep$, whence the two identities for $v$ give $4Ep=2E$. But the
Perron column has contributions $p$ from the signed dimension pair,
$1/2$ from $P$, and positive terms from the other rows.
Thus $b=1$, as in~\cite[proof of Theorem~6.2]{Class6Rank5}.

The remaining ordinary squared norm four consists of two sign pairs
of entries $\pm1$. Their twists are $\theta,-\theta$ by the first
modular identity. Their $v$-weights $V_1,V_2$ satisfy
\[
 V_1+V_2=-2Ep,\qquad \theta(V_1-V_2)=2E(1-p).
\]
Since the weights are real, $\theta=\pm1$. The twist-$-1$ pair
has weight $-E$; its positive-dimensional representative $Q$
has $a_Q=-1$ and $\dim Q=E/2$.

Now $\delta_B=4S_{Q,B}$ is an algebraic integer since $S_{Q,1}=1/4$.
The modular Galois permutations fix the rational row of $P$, preserve
both monodromy sectors (apply Galois to $SP_J=\operatorname{diag}(-M_B)S$,
where $M_B$ is monodromy with $J$), and preserve twist $-1$ by the
Galois twist rule. Hence they preserve the unique ordinary twist-$-1$
pair in the support of $P$, namely $Q,JQ$. Its entries in the row of
$P$ force positive sign at $Q$ and negative sign at $JQ$. Each
Galois automorphism thus fixes the signed-sector entries of $\delta$
and either fixes or negates all its ordinary-sector entries.

The entry $z$ of $\delta$ at the ordinary Perron label satisfies
$Dz/2=\operatorname{FPdim}(Q)>0$, so $z\in K$. Its conjugate at
the other Perron label is $\pm z$, and $z^2$ is a positive integer.
These labels and their translates give $4z^2\le8$, the ordinary-sector
squared norm. If $z^2=2$ they exhaust that sector at twist one; all
its entries are real, being rational multiples of $z$. If $z=1$,
the whole row is integral. Neither $Q$ nor $JQ$ belongs to
$I(\mathbf1)$ because both have twist $-1$. The ordinary unit
coefficient of $FQ$ therefore vanishes, so the ordinary induction
entries sum to zero. Outside the Perron pair their sum is $-2$ and
their squared sum at most two. Exactly two are $-1$, and again the
ordinary sector is exhausted at twist one. In either case
$\delta_Q=0$, while signed $J$-pairs cancel from the modular identity
at $(Q,Q)$. This gives $4\delta_Q=\sum_B\delta_B^2\theta_B=8$.
\end{proof}
\begin{proposition}[Conductors with two rational complementary codegrees]
\label{prop:r6-two-pair-conductors}
Suppose a rank-six category with trivial universal grading has no
pseudo-unitary Galois conjugate and has commutative signed codegrees
\[
 E,E',H,H',g_1,g_2,
\]
where $H$ is irrational and $g_1,g_2$ are rational integers.
Then its quadratic conductor belongs to
\[
 \{5,8,12,13,21,24,28,40,60\}.
\]
\end{proposition}
\begin{proof}
Theorem~\ref{thm:r6-quadratic-no-two} excludes a rational character
of codegree two; a conjugate pair of that codegree would consume
the full mass. Thus $g_1,g_2\ge3$ and
$L=t_E+t_H=1-1/g_1-1/g_2\ge1/3$. In
\eqref{eq:quadratic-conductor-test} take $k=1$.
Since $a^2-cb^2=4$, we have
$c\lambda=(\sqrt{c+4/b^2}-2/b)^2\ge(\sqrt{c+4}-2)^2$.
For $L\ge1/3$ the test gives $c<35$ at odd conductors and $c<106$
at even ones. Indeed the irrational case bounds $c$ by
$18+12\sqrt2$ and $72+24\sqrt2$, respectively; the rational case
gives the smaller bounds $30$ and $79$.

Testing \eqref{eq:quadratic-conductor-test} on the fundamental discriminants
in these ranges with the exact least unit leaves
\[
 5,8,12,13,17,21,24,28,40,44,56,60
 \qquad\text{(\codecheck{quadratic_conductor_check.py}{least-unit sieve}).}
\]

For $c=17,44,56$ the prime three is inert. We claim $L\ge1/2$.
Only a codegree-three character needs consideration. If it belonged
to a real quadratic orbit, squared row differences would give
$3\ge\operatorname{disc}(F)\ge5$. In an imaginary orbit the sole
dual pair gives $3=|\operatorname{disc}(F)|b^2$, hence
$F=\mathbb Q(\sqrt{-3})$. The signed discriminant of $K^2\times F$
would then make three a rational norm from $K$, impossible at an
inert prime.

The codegree-three character is therefore rational.
Lemma~\ref{lem:silver-local-partner} at three pairs it with the other rational
character: the two copies of $K$ are unramified quadratic fields
locally and cannot be its totally ramified partner.
Their rows $\rho,\eta$ are consequently congruent modulo three.
The row $\rho$ has the unit entry one, two entries $\pm1$, and
all other entries zero.
On this support $\rho_i\eta_i\equiv1\pmod3$ and the three products sum
to zero. One is therefore at most $-2$, so $\|\eta\|^2\ge 1+1+4=6$.
The two rows have total reciprocal mass at most $1/3+1/6=1/2$.

Thus $L\ge1/2$. At $17,44,56$ the least-unit gaps are
$16/17,9/11,7/8$, respectively. Neither inequality in
\eqref{eq:quadratic-conductor-test} holds at $L=1/2$.
\end{proof}
\begin{theorem}[Four quadratic conductors]\label{thm:r6-quadratic-three-fields}
Suppose $K=\mathbb Q(D)$ has discriminant $13,21,40$ or $60$. A rank-six fusion category with signed
codegrees $E,E',H,H',g_1,g_2$, with $H$ irrational and $g_1,g_2$
rational, is spherical. At $40$ and $60$ it has a pseudo-unitary
Galois conjugate.
\end{theorem}
\begin{proof}
Assume no Galois conjugate is pseudo-unitary.
Proposition~\ref{prop:quadratic-full-character} gives the hypotheses
of Lemma~\ref{lem:quadratic-mass-gap}. A pair with ratio at least
$U^2$ has squared reciprocal gap at least $1-4/A^2$ times its
squared reciprocal mass. The ordinary and signed projection bounds
of Lemmas~\ref{lem:mass-fourier-verification}
and~\ref{lem:sextic-matrix-compression} give respectively
$\delta_D^2\le(c,2)/c$ and $\delta_H^2\le(c,4)/c$,
where $\delta_f=|1/f'-1/f|$. Together with the mass bound these give
\[
\begin{array}{c|ccc}
 c & s\text{ at most} & w\text{ at most} & t\text{ at most}\\ \hline
 13 & 11/39 & 11/39 & 2/39\\
 21 & 5/21 & 5/21 & 2/21\\
 40 & 19/(60\sqrt2) & 19/60 & 1/(60\sqrt2)\\
 60 & 4/(15\sqrt2) & 4/15 & 1/(15\sqrt2)
\end{array}
\]
The rational codegrees are at least three by the reciprocal identity
and Theorem~\ref{thm:r6-quadratic-no-two}. Hence $t+w\ge1/3$.
The last two lines contradict this inequality. In the first two
lines equality holds throughout, giving $E=E'$, $D=EU$ and
$H/H'=U^2$. Thus $H=nU$, and $w=s$ gives $H=D$. The real signed
character of this codegree gives a spherical structure by
Lemma~\ref{lem:section-signs}.
\end{proof}

\section{The remaining quadratic profiles}
\label{sec:upstream-24-28}
For the conductor calculations assume that $R$ and $T$ are commutative,
their ordinary and signed dimension characters generate the same real
quadratic field $K=\mathbb Q(D)$, and the signed codegrees are
$E,E',H,H',g_1,g_2$. Here $H$ is irrational and $g_1,g_2$ are
rational integers at least three. Use the embeddings with
$E\ge E'$ after Galois transport and $H>H'$ by labelling its two
rows. Put $t=\operatorname{Tr}(1/E)$, $w=\operatorname{Tr}(1/H)$,
$s=\operatorname{Tr}(1/D)$, and
$\delta_f=|1/f'-1/f|$. The induction formula gives integral ratios
$E/f$ and the simultaneous-column unit gives
$H/H'\le D/D'$.
If the characters of common rational codegree $g$ form a quadratic
pair over $F$, Lemma~\ref{lem:quadratic-row-difference} gives
$g\ge\operatorname{disc}(F)$ in the real case and
$g=|\operatorname{disc}(F)|b^2$ in the imaginary case, since the
involution has exactly one nontrivial two-cycle.

We will use one divisibility for the fields $K=\mathbb Q(\sqrt d)$,
$d=3,6,7$. In all three the second primitive family has $V=3+\sqrt d$.
For a quadratic row of codegree $mV^eU^k$, one has
\begin{equation}\label{eq:three-quadratic-scalars}
              2d/\gcd(3^e,d)\mid m .
\end{equation}
Indeed write $V^eU^k=A+B\sqrt d$ in primitive coordinates.
The squared distance of the two rows gives $2d\mid mA$;
expanding the sum of squares gives $2\mid mB$. Thus $m$ is even.
The rational coordinate of $U^k$ is prime to $d$ by its norm,
so $\gcd(A,d)=\gcd(3^e,d)$. For even $d=6$, $A$ is odd.
These facts give the assertion, including for a pair of distinct
characters of rational codegree.

\begin{theorem}\label{thm:upstream-24-28}
Under these hypotheses, if $K=\mathbb Q(\sqrt6)$ or
$K=\mathbb Q(\sqrt7)$, a Galois conjugate of the category
is pseudo-unitary.
\end{theorem}

\begin{lemma}\label{lem:integral-three-supported}
Suppose the signed codegrees are $E,E',H,H',3,3$, and the last two
characters have absolute value at most one on each basis element.
If $s>t$ and a simple object has integral Frobenius--Perron dimension,
it is invertible.
\end{lemma}
\begin{proof}
Apply Lemma~\ref{lem:integral-regular} to the subcategory
$\mathcal A$ of integral-dimensional simples.
Both ordinary Perron conjugates take the value one, giving a lower
bound $s$. If $\mathcal A$ contains a simple of dimension at least two,
the two short signed characters give zero. The dimension pair gives
$t$, and the $H$-pair gives either zero or its full mass, since dimensions
on $\mathcal A$ are rational. Thus the unit coefficient is $t$ or
$t+w=1/3$. The first is too small, the second would give
$\operatorname{FPdim}(\mathcal A)=3$, which cannot contain a simple
of dimension at least two.
\end{proof}

\subsection{Two rational characters of codegree three}
\label{sec:two-threes}

\begin{theorem}\label{thm:two-threes-exclusion}
Suppose $\mathcal C$ has rank six, $R$ is commutative, and
\[
 T_{\mathbb Q}\simeq K\times K\times\mathbb Q\times\mathbb Q,
\]
where $K$ is a real quadratic field.
Suppose the signed spherical dimension character is in the first
$K$-orbit and the Frobenius--Perron character of $R$ has field $K$.
If the two rational characters of $T$ both have formal codegree three, then $3$ splits in $K$. Moreover, writing $E=\dim(\mathcal C)$ and $D=\operatorname{FPdim}(\mathcal C)$, one has
\[
 E\equiv D\equiv-3\pmod{9\mathcal O_K}.
\]
\end{theorem}

We record the estimate for a rational signed row of codegree three
that will also be used at $(3,8)$. Write $A$ for its induction label
and put $M=6S$. The row $M_{A,*}$ is integral by the central-character
formula. Its two signed dimension entries are one. Its two ordinary
Perron entries are equal positive integers, say $b$, by central
symmetry, Galois conjugation and positivity of the Perron dimension.
At every induction label $C$ the two unit coefficients give
\begin{equation}\label{eq:sector-row-sum-33}
 \sum_{V\ {\rm signed}}2S_{A_V,C}
 =\sum_{V\ {\rm ordinary}}2S_{A_V,C}=1 .
\end{equation}

For a real vector $v$ supported on such rows, put $y=Mv$ and
$z_W=S_{J,W}/S_{1,W}$. The two sign sectors of $y$ each have
squared norm $18\|v\|^2$: $v\perp Jv$ and $(SJv)_W=z_W(Sv)_W$.
All entries of $y$ are real. In the twist sum the negative sector
cancels in $J$-pairs, since $\theta_{JW}=z_W\theta_W$. The
twist identity $S\Theta S=\Theta^{-1}S\Theta^{-1}$ therefore gives
\[
 6v^tMv\ \ge\ 2\!\sum_{z_W=\theta_W=1}y_W^2-18\|v\|^2.
\]
For $v=e_A$, put $r=M_{A,A}$. The two signed dimension entries
and their $J$-partners give $r^2\le7$. The four ordinary Perron
entries, including their translates, give $6r\ge8b^2-18$, forcing
$b=1$. The remaining four ordinary induction entries now sum to one,
so their squared norm is at least one. Including their translates
improves the bound to $6r\ge-6$.

Apply these facts to the two rational three-rows $A,B$. At the
two labels $C,C'$ of the complementary quadratic pair the entries
in each row agree by Galois symmetry; call them $a,b$. Write
$\left(\begin{smallmatrix}r&z\\z&s\end{smallmatrix}\right)$ for the
$A,B$ block of $M$. We claim
\begin{equation}\label{eq:two-threes-vanishing}
               a=b=0 .
\end{equation}
The signed row sums and squared norms give $r+z=1-2a$,
$s+z=1-2b$ and
$2a^2+r^2+z^2,\ 2b^2+s^2+z^2\le7$. Thus $a,b\in\{-1,0,1\}$.
For $v=e_A+e_B$, the four ordinary Perron entries of $Mv$ are two,
and the other four ordinary induction entries sum to two. With their
translates the twist-one contribution is at least $16+4=20$.
The twist inequality gives $6(r+2z+s)\ge4$, so $a+b\le0$.
The signed norm gives
\[
 8+2(a+b)^2+(1-2a)^2+(1-2b)^2\le18 .
\]
This leaves $(a,b)=(0,0)$ or $\{a,b\}=\{-1,1\}$. In the latter
case the row sums and bounds force one diagonal entry to be $-2$,
contrary to $r,s\ge-1$. This proves the claim.

Now $r=s=1-z$. By symmetry the signed class of $F(A-B)$ vanishes
on both quadratic factors and has rational values $r-z,z-r$.
Its simple coefficients are $(r-z)(\rho_i-\eta_i)/3$.
The two rational rows have three entries $\pm1$, with supports
intersecting in precisely two positions by orthogonality. Their
difference has an entry $\pm1$, so $3\mid r-z$ and $r,z\equiv-1\bmod3$.

\begin{proof}[Proof of Theorem~\ref{thm:two-threes-exclusion}]
Put $T_3=T\otimes\mathbb Z_3$. Its decomposition into complete local rings partitions the field factors of $T_3\otimes\mathbb Q_3$. The two rational characters have distinct reductions, since their supports differ and all nonzero entries are $\pm1$. They belong to different local factors of $T_3$. Neither factor consists of its rational axis alone. Otherwise its identity would be the primitive idempotent of that character, whose unit coefficient $1/3$ is not in $\mathbb Z_3$.

By symmetry of the central character formula and \eqref{eq:two-threes-vanishing}, the signed class of $F(A)$ has value zero on the second quadratic field. Its values on the two rational characters are $r,z\equiv-1\pmod3$. A branch of that quadratic field cannot lie in either rational character's local factor, since an element of a local ring cannot have residue zero on one branch and nonzero on another. Each rational character therefore needs a different branch of the first quadratic field. This proves that $3$ splits.

At either prime of $K$ above three, the spherical dimension character evaluates $F(A)$ to $E/3$. Its residue must equal that of the rational character in the same local factor, namely $-1$. Similarly, the spherical dimension character evaluates $F(A_D)$ to $E/D$, whereas both rational characters take the value one. Thus $E/3\equiv-1$ and $E/D\equiv1$ at both primes. Since $3$ splits, these congruences give the result.
\end{proof}

In particular this configuration is impossible over $\mathbb Q(\sqrt2)$, $\mathbb Q(\sqrt3)$ and $\mathbb Q(\sqrt5)$. The prime three is ramified in the second field and inert in the other two.

\subsection{Conductor twenty-four}

\begin{proposition}[The complementary profiles at conductor twenty-four]
\label{prop:r6-twenty-four-complements}
At conductor twenty-four, if no Galois conjugate is pseudo-unitary,
the last two characters form an Eisenstein pair of codegree three.
\end{proposition}

\begin{proof}
Put $U=5+2\sqrt6$, $V=3+\sqrt6$. The scalars in the unit and
$V$ families are divisible by twelve and four, respectively,
by \eqref{eq:three-quadratic-scalars}.

Lemma~\ref{lem:quadratic-mass-gap} gives $s\ge5t$ and squared gaps
at least $24/25$ times the squared masses for $H$ and $D$. The mixed bound gives
\[
 \frac16\ge\frac{24}{25}(w^2+2s^2)
       \ge\frac{16}{17}(w+t)^2.
\]
In particular $L=t+w<7/15$.

As $1/g_1+1/g_2>8/15$, the only possibilities are $(3,3)$ and
$(3,4)$. Two rational threes would force three to split by
Theorem~\ref{thm:two-threes-exclusion}, whereas it ramifies.
A quadratic orbit must have codegree three: its field is
$\mathbb Q(\sqrt{-3})$ by Lemma~\ref{lem:quadratic-row-difference}.

It remains to exclude two rational rows of codegrees three and four.

Here $t+w=5/12$. If $H=nU$, then $w=10/n<5/12$ and $12\mid n$,
so $n\ge36$ and $t\ge5/36$. But $s\ge5t$ gives
$2\delta_D^2\ge48t^2>1/6$. Every other allowed $H$ has
$H/H'\ge U^3$ and $\delta_H^2\ge(242/243)w^2$.
The \codecheck{quadratic_24_28_checks/quadratic_twenty_four_rational_audit.py}{mixed bound}
then gives
\[
 \frac16\ge48t^2+\frac{242}{243}w^2
 \ge\frac{5808}{5953}(t+w)^2=\frac{3025}{17859}>\frac16.\qedhere
\]
\end{proof}

\begin{lemma}[A cyclic stabilizer obstruction]\label{lem:r6-fixed-cyclic-dimensions}
Suppose the invertibles of a fusion category form $C_3$ and fix every
noninvertible simple on both sides. Its simple Frobenius--Perron
dimensions cannot generate $K=\mathbb Q(\sqrt6)$.
\end{lemma}
\begin{proof}
Let $g$ generate $C_3$. A fixed simple gives a one-simple module
category over $\langle g\rangle$, so the associator on this group is
cohomologically trivial. Write $B=\mathbf1\oplus g\oplus g^2$
for its group algebra. Each fixed simple $X$ admits left and right
$B$-actions separately, since $H^2(C_3,\mathbb C^\times)=0$.
In particular $X^*\otimes_B X$ has Perron dimension $d_X^2/3$,
an algebraic integer~\cite[Exercises~7.16.8--7.16.9]{EGNO}. Write
$d_X=a_X+b_X\sqrt6$ with $a_X,b_X\in\mathbb Z$. The rational
coordinate of $d_X^2/3$ gives $3\mid a_X$.

Let $l_X:gX\to X$ and $r_X:Xg\to X$ give the two actions.
Their interchange on $gXg$ is a scalar $\kappa_X$ with
$\kappa_X^3=1$, since the actions are cyclic. If $\kappa_X\ne1$,
take any noninvertible simple $Z$, put $Y=X^*$ and act on
$\operatorname{Hom}(Z,XY)$ by
\[
 A(f)=(l_X\otimes1)(1_g\otimes f)l_Z^{-1},\qquad
 T(f)=(r_X\otimes1)(1\otimes l_Y^{-1})f .
\]
The composites differ only by the interchange on $gXg$, so
$AT=\kappa_X^{\pm1}TA$. Both operators are invertible; determinants
give $3\mid N_{XY}^Z$. Every invertible occurs once in $XX^*$ by
rigidity. Taking rational coordinates in
$d_X^2=3+\sum_Z N_{XY}^Z d_Z$ gives $6b_X^2\equiv3\pmod9$,
which is impossible.

Thus $\kappa_X=1$ for every fixed $X$. The actions define a
$B$-bimodule on $X$, of Perron dimension $d_X/3$ in the Morita
dual by the same dimension formula. Hence $d_X\in3\mathcal O_K$.
Orthogonality of the distinct Perron character and its conjugate
gives $3+\sum_X N_{K/\mathbb Q}(d_X)=0$, with the sum over the
noninvertibles. Reduction modulo nine gives the contradiction.
\end{proof}

\begin{theorem}[The nonreal-three branch at conductor twenty-four]
\label{thm:r6-twenty-four-imaginary}
Let $K=\mathbb Q(\sqrt6)$ and suppose the signed rational algebra is
\[
 T_{\mathbb Q}=K\times K\times\mathbb Q(\sqrt{-3}),
\]
with codegrees $E,E',H,H',3,3$, where $H$ is irrational.
Then $\mathcal C$ has a pseudo-unitary Galois conjugate.
\end{theorem}
\begin{proof}
Assume the contrary. Lemma~\ref{lem:quadratic-mass-gap} gives
$D/D'\ge U^2$, where $U=5+2\sqrt6$.

\smallskip\noindent\emph{The ordinary nonreal pair has codegree three.}
There is exactly one dual pair $X,X^*$ among the basis simples.
The nonreal signed row of norm three has the form
\[
 (1,\omega,\omega^2,0,0,0),\qquad \omega^3=1,\quad\omega\ne1.
\]
Indeed its nonreal entries occupy $X,X^*$ and exhaust the norm
remaining after the unit entry.
Orthogonality to the positive dimension row selects $\omega$,
rather than $-\omega$, and gives $x_X=1$.
Every real signed row then takes the value one on $X,X^*$.
In particular any actual nontrivial invertible must be one of
these two basis objects; a self-dual invertible is impossible.

At two the odd-codegree branch is isolated and maximal unramified
by Lemma~\ref{lem:trace-order-conductor}. Its residue field is
$\mathbb F_4$, with nontrivial involution.
The remaining generic factors are the two ramified quadratic
fields $K\otimes\mathbb Q_2$; their completed local factors have
ranks two or four and residue field $\mathbb F_2$.
The parity isomorphism $R/2R\simeq T/2T$ preserves the involution.
Thus the ordinary nonreal quadratic factor is unramified and
maximal at two, and has odd rational codegree $g$.
Its nonreal row is nontrivial modulo two at $X$.
Writing its value there as $(a+b\sqrt{\operatorname{disc}(F)})/2$
gives $b$ odd and
\[
 g=|\operatorname{disc}(F)|b^2\equiv3\pmod8.
\]
Here the row identity uses the unique dual pair; an imaginary
quadratic field unramified and nonsplit at two has discriminant
$5$ modulo eight.

Any further odd rational codegree would split off an unramified
factor by its orbit idempotent and the perfect trace pairing. Such a
factor has trivial involution and is absent from the dyadic decomposition.
Thus the remaining rational ordinary codegrees are even and
at least four.
A rational codegree two would force a self-dual invertible; a
real complementary quadratic pair of codegree two is excluded
by the complementary-pair discriminant bound.

Suppose $g\ge11$.
Every ordinary codegree lies in $K$ by central induction.
If the other real ordinary pair has irrational codegree, its
character field is $K$ and its mass together with the FP mass
is at least $9/11$.
The least squared gap in $K$ is $2/3$, so the ordinary Fourier
bound would give
\[
 1/12\ge\delta_D^2+\delta_G^2
 \ge\tfrac13(9/11)^2=27/121,
\]
which is impossible.
Otherwise its two rational codegrees contribute at most $1/2$,
and $s\ge7/22$.
But then $2\delta_D^2\ge(48/25)(7/22)^2>1/6$,
since $D$ has index at least two.

Consequently $g=3$.

The ordinary nonreal row is now the same three-supported row.
Its orthogonality to the FP row gives $d_X=1$.
Thus $1,X,X^*$ form the actual group $C_3$ of invertibles.
Multiplication by $X$ has spectrum $1,1,1,1,\omega,\omega^2$;
as a positive permutation it consists of the three-cycle on
these invertibles and fixes all three remaining simples.

Commutativity gives the same right action. Lemma~\ref{lem:r6-fixed-cyclic-dimensions}
now gives the contradiction.
\end{proof}

\subsection{Conductor twenty-eight}

\begin{proposition}\label{prop:r6-quadratic-twenty-eight}
If $K=\mathbb Q(\sqrt7)$ and no Galois conjugate is pseudo-unitary,
the short rows are rational of codegree three, and
$E=mU^k$, $H=nU^j$, $D=MU^\ell$, where $U=8+3\sqrt7$,
$k,j\ge1$, $\ell>k$, $M\mid m$, and
$E/D\equiv1\pmod{3\mathcal O_K}$.
\end{proposition}
\begin{proof}
Put $V=3+\sqrt7$, so $V/V'=U$ and $N(V)=2$.
The primitive families are those of Lemma~\ref{lem:quadratic-mass-gap}.

The scalars in both primitive families are divisible by fourteen
by \eqref{eq:three-quadratic-scalars}.

Since $g_1,g_2\ge3$, $t+w\ge1/3$.
For the short pair $nV$, its squared gap is $(7/9)w^2$;
every other irrational pair has squared gap at least
$(63/64)w^2$.
We use the ordinary, signed and mixed bounds
\[
 \delta_D^2\le\frac1{14},\qquad
 \delta_H^2\le\frac17,\qquad
 \delta_E^2+\delta_H^2+2\delta_D^2\le\frac17.
\]

Put $A_i=\operatorname{Tr}(U^i)$, so $A_0=2$, $A_1=16$ and
$A_{i+2}=16A_{i+1}-A_i$. We first exclude $H=nV$. Its eigenvalue
$\beta=a+b\sqrt7$ on the Perron induction column is totally positive
and integral. Hence
\[
 \operatorname{Tr}(\beta/V)=3a-7b\ge2\sqrt{N(\beta)/2}\ge\sqrt2,
\]
so this integral trace is at least two. The dimension pair contributes
$s$ on that column, and the rational-codegree orbits contribute at
least their reciprocal mass. Since $\operatorname{Tr}(1/V)=3$, the
available deficit is therefore
\begin{equation}\label{eq:r6-twenty-eight-short}
                       s-t\le w/3 .
\end{equation}
By \eqref{eq:quadratic-trace-order}, $s\ge2t$, so $t\le w/3$.
But $14\mid n$ gives $t+w\le4w/3\le2/7<1/3$, a contradiction.

Now $H/H'\ge U^2$. This also rules out the exceptional ordinary pair,
so $s/t\ge16/3$ and $D/D'\ge U^2$. The separate Fourier bounds give
\[
 t+w\le\frac1{14\sqrt2}+\frac8{21}<\frac7{15}.
\]
The only possible integer codegree pairs are $(3,3)$ and $(3,4)$.
The latter pair is rational and requires $12N(EH)$ to be a square,
whereas $N(EH)$ has squareclass one or two.

For $(3,3)$, a quadratic character field would be
$\mathbb Q(\sqrt{-3})$.
Indeed orthogonality of a real quadratic pair gives
$\operatorname{disc}(F)\mid2g$, and no positive quadratic
discriminant divides six.  An imaginary pair here has
$g=|\operatorname{disc}(F)|b^2$.
The order discriminant would then require $N(EH)$ to have
squareclass three, again impossible.
Thus both characters are rational, and their order discriminant
forces $E,H$ to belong to the same primitive family.

For later use, Theorem~\ref{thm:two-threes-exclusion} gives
$E\equiv D\equiv-3\pmod{9\mathcal O_K}$. Neither can be in the
$V$-family: dividing by three, the two residues modulo the primes
above three must both be $-1$, whereas $V=3+\sqrt7$ has opposite
residues and $U\equiv-1$ at both primes. Hence $E,H,D$ belong to the
unit family in the rational $(3,3)$ case.

We have reduced all cases to rational codegrees $3,3$ and
$E=mU^k,H=nU^j$.
Integrality of $E/H$ gives $q=m/n\in\mathbb Z_{>0}$, and the
reciprocal identity gives
\[
 m=3A_k+3qA_j,\qquad q\mid3A_k.
\]
Since $A_r\equiv2\pmod7$ and $7\mid m$, this forces
$q\equiv-1\pmod7$.
If $E$ were rational, then $A_k=2$, so $q\mid6$ and $q=6$.
But $n=3A_j+1$ would be odd, contradicting $14\mid n$.
Thus $k\ge1$, while $j\ge1$ follows from irrationality of $H$.

The same congruence yields $E/D\equiv1\pmod{3\mathcal O_K}$, and
integrality and $E<D$ give $M\mid m$, $\ell>k$.
\end{proof}

\begin{lemma}[Two rational-three rows and the dyadic factors]
\label{lem:r6-two-threes-dyadic}
Suppose a rank-six fusion category has commutative ordinary and
signed character algebras, their dimension-character fields equal
one quadratic field $K$ ramified at two, and
$T_{\mathbb Q}=K^2\times\mathbb Q^2$, with the two rational
signed characters of codegree three.
Then every simple is self-dual and the unit is the only invertible
simple. The ordinary order at two has exactly two local factors
of rank one, and every residue field is $\mathbb F_2$.
Its rational-odd-codegree character orbits have total degree two.
If these are two rational rows, their codegrees are three or seven
modulo eight. Every other rational ordinary codegree is even,
and no ordinary codegree is two.
\end{lemma}
\begin{proof}
Total reality of the signed characters makes every simple self-dual.
The two norm-three rational rows have supports intersecting in exactly
two positions, the unit and one nonunit $X$, where their signs are
opposite. An invertible must be nonzero on both rows. If $X$ were
invertible, its positive dimension would be one, so in the row with
$X$-entry $-1$ orthogonality would force the remaining supported
simple to have dimension zero.

At two the two short idempotents have odd denominators and split off
two copies of $\mathbb Z_2$. The remaining generic algebra is $K_2^2$
with $K_2$ ramified quadratic. Thus there are exactly two rank-one
factors, and all residue fields are $\mathbb F_2$. The isomorphism
modulo two transfers these properties to the ordinary ring: the
idempotents lift uniquely on completion and the rank of each factor
is its dimension after reduction.

An ordinary orbit of odd rational codegree is isolated, maximal and
unramified at two by Lemma~\ref{lem:trace-order-conductor}. With all residue fields $\mathbb F_2$ this
accounts for one rank-one factor per character. Conversely a rank-one
factor cannot have irrational codegree in $K$, since $K$ does not
embed in $\mathbb Q_2$, and its denominator must be a unit by
perfectness. Thus exactly two characters have odd rational codegree.
If rational, their reductions coincide up to permutation with the
two short signed rows. Three odd entries give squared norm $3$ or
$7\pmod8$. Finally a rational row of norm two would have idempotent
$(1\pm X)/2$, implying $X^2=1$. A nonrational orbit of norm two
would exhaust the reciprocal mass.
\end{proof}

\begin{lemma}[Integral dimensions from short rows]
\label{lem:r6-three-three-small}
Let $R$ be the commutative Grothendieck ring of a rank-six fusion
category in which every simple is self-dual and no nonunit simple
is invertible.  If $R$ has two rational characters of codegree
three and also has a rational character of codegree four or six,
then some noninvertible simple has integral Frobenius--Perron dimension.
\end{lemma}
\begin{proof}
The two rows have supports intersecting in the unit and one other
simple, with opposite signs there.
After reordering, positivity of FP dimensions gives
\[
 \rho=(1,1,-1,0,0,0),\qquad
 \eta=(1,-1,0,1,0,0),
\]
on the basis $1,X,Y,Z,P,Q$, with
$d_Y=d_X+1$ and $d_Z=d_X-1$.
Every other rational row is
\[
 (1,x,1+x,x-1,u,v),\qquad
 f=3+3x^2+u^2+v^2.
\]
For $f=6$, the only integer possibilities are
$x=\pm1,u=v=0$.
Orthogonality to the FP row gives $3\pm3d_X=0$,
which is impossible or makes $X$ invertible.

For $f=4$, $x=0$ and $(u,v)$ has squared norm one.
Orthogonality to the Perron row gives $d_P=3$ or $d_Q=3$.
\end{proof}

\begin{theorem}[Descent at conductor twenty-eight]
\label{thm:r6-quadratic-twenty-eight}
Suppose $K=\mathbb Q(\sqrt7)$.  Then $\mathcal C$ has a pseudo-unitary Galois conjugate;
in particular it is spherical.
\end{theorem}
\begin{proof}
Assume that there is no pseudo-unitary Galois conjugate, and apply Proposition~\ref{prop:r6-quadratic-twenty-eight}.
We have
\[
 T_{\mathbb Q}=K^2\times\mathbb Q^2,
 \qquad K=\mathbb Q(\sqrt7),\qquad U=8+3\sqrt7,
\]
with signed codegrees $E,E',H,H',3,3$, where
\[
 E=mU^k,\quad H=nU^j,\quad D=MU^\ell,
 \qquad k,j\ge1,\quad \ell>k,
\]
and $E/D\equiv1\pmod{3\mathcal O_K}$.
By Lemma~\ref{lem:r6-two-threes-dyadic} there are no nontrivial
invertibles, and all characters are real. Exactly two ordinary
characters have odd rational codegree, and all other rational
codegrees are even and at least four. If the odd characters are rational,
their codegrees satisfy
\begin{equation}\label{eq:r6-twenty-eight-odd-codes}
                        a,b\equiv3,7\pmod8.
\end{equation}
If they are conjugate over a quadratic field $F$, then two splits in
$F$. Thus $\operatorname{disc}(F)\ge17$, and their common codegree
is at least $\operatorname{disc}(F)$ by the row difference bound.

Put $t=\operatorname{Tr}(1/E)$, $s=\operatorname{Tr}(1/D)$,
$w=1/3-t$, $\lambda=63/64$. The congruence for $E/D$ gives
$(m/M)(-1)^{k-\ell}\equiv1\pmod3$. Thus $m/M\ge2$ if
$\ell=k+1$. The trace ratio $A_{k+1}/A_k$ is at least $127/8$ for
$k\ge1$, so in all cases $s/t\ge127/4$.
The mixed Fourier bound, with all three pairs in the unit family, gives
\begin{equation}\label{eq:r6-twenty-eight-strong-mass}
       \lambda\bigl(t^2+(1/3-t)^2+2s^2\bigr)\le1/7.
\end{equation}
In fact $s<3/22$. The ordinary bound gives $s<1/3$, so the
left side decreases with $t$ on $0\le t\le4s/127$.
Substituting $t=4s/127$ gives a quadratic in $s$ which is
increasing for $s\ge3/22$ and already exceeds $1/7$ there.

An irrational complementary ordinary codegree generates $K$, and its
character field has no rank-one factor at two. Since two places must
be reserved for odd codegrees, it gives a single quadratic pair $G,G'$.
Unless the two odd codegrees are both three, their total mass is at most
$10/21$, by \eqref{eq:r6-twenty-eight-odd-codes} or the bound $2/17$
for a nonrational pair. The masses of the $D$- and $G$-pairs then sum
to at least $11/21$. Their gaps contradict the ordinary bound since
$(7/18)(11/21)^2>1/14$. If both odd codegrees are three, their
characters are rational and the discriminant forces $N(G)$ square.
All pairs are then in the unit family. The ordinary and signed pair
masses each sum to $1/3$, and Cauchy's inequality in the mixed bound
gives $1/7\ge\lambda/6=21/128$. Thus all complementary codegrees are rational.

Their total mass exceeds $6/7$. This excludes a nonrational odd pair,
which contributes at most $2/17$ and leaves at most $1/2$.
If the two even characters form one quadratic orbit of codegree $g$,
then $g\ge\operatorname{disc}(F)\ge5$. For $(a,b)\ne(3,3)$ the
mass bound gives $2/g>6/7-10/21$, hence $g=4$, impossible. For
$(a,b)=(3,3)$ it gives $g<21/2$, whereas the discriminant squareclass
of $28\operatorname{disc}(F)$ forces $F=K$ and $g\ge28$.
Therefore all four complementary rows are rational.

Write the odd codegrees as $a\le b$ and the even ones as $c\le d$.
Their product has squareclass seven and their reciprocal sum exceeds
$19/22$. Thus $a=3$. If $b\ge7$, the mass bound gives $c=4$.
At $d=4$ the squareclass would require $b=21h^2$, impossible since
$b\equiv3\pmod4$. At $d\ge6$ the mass bound forces $b=7,d=6$,
whose product with $3,4$ has squareclass fourteen.
Thus $a=b=3$. Lemmas~\ref{lem:r6-three-three-small}
and~\ref{lem:integral-three-supported} give $c,d\ge8$.
One is a multiple of seven and therefore at least fourteen. But then
\[
 1-s\le 2/3+1/8+1/14=145/168<19/22.
\]
\end{proof}

\begin{proof}[Proof of Theorem~\ref{thm:upstream-24-28}]
At conductor twenty-four Proposition~\ref{prop:r6-twenty-four-complements}
reduces to the Eisenstein pair, excluded by
Theorem~\ref{thm:r6-twenty-four-imaginary}.
Theorem~\ref{thm:r6-quadratic-twenty-eight} treats conductor twenty-eight.
\end{proof}

\subsection{Conductor twelve}

\begin{theorem}\label{thm:quadratic-twelve-upstream}
Let both dimension characters of a rank-six fusion category have
field $K=\mathbb Q(\sqrt3)$ and both character algebras be commutative.
Suppose the signed codegrees are $E,E',H,H',g,h$, where $H$ is
irrational and $g,h$ are rational and at least three. Then the
category is spherical, or the last two characters are rational with
unordered pair $(3,3),(3,8),(3,12)$ or $(4,4)$, or they form a
Gaussian pair of codegree four.
\end{theorem}
\begin{proof}
Transport to $E\ge E'$ and order $H>H'$. Assume nonsphericity.
Write $t=\operatorname{Tr}(1/E)$, $w=\operatorname{Tr}(1/H)$,
$s=\operatorname{Tr}(1/D)$ and $L=t+w=1-1/g-1/h$. Use the two families $mU^k,mVU^k$ of Lemma~\ref{lem:quadratic-mass-gap},
with $U=2+\sqrt3$, $V=1+U$.
Equation~\eqref{eq:three-quadratic-scalars}
gives $6\mid m$ in the unit family and $2\mid m$ in the $V$ family.

The primitive reciprocal traces are $A_i,B_i$,
with initial pairs $(2,4),(1,3)$ and recurrence $z_{i+2}=4z_{i+1}-z_i$.
The mixed conductor-twelve bound reads
$\delta_E^2+\delta_H^2+2\delta_D^2\le1/3$, where
$\delta_f=1/f'-1/f$; it also allows twice the squared gaps of any
additional ordinary pairs. Lemma~\ref{lem:quadratic-mass-gap} gives $s\ge2t$ and
$D/D'\ge H/H'\ge U^2$. Thus $\delta_D^2\ge3s^2/4$ and $s\le\sqrt2/3$.

If $H=nU$, its gap squared is $12/n^2<1/3$, so $n\ge12$ and
$w\le1/3$. This gives $L\le(2+\sqrt2)/6$. In all other cases
$\delta_H^2\ge25w^2/27$, so $6t^2+25w^2/27\le1/3$.
Consequently $L^2\le187/450$ in every case. For two rational
characters the order-discriminant identity makes $gh$ square or
six times a square. With $g\le h$, the mass bound gives $g\le5$
and, for $g=3,4,5$, respectively $h\le45,9,6$. The resulting
list is
\[
 (3,3),(3,8),(3,12),(3,18),(3,27),(3,32),
 (4,4),(4,6),(4,9),(5,5).
\]
At five the last rational pair would have to occupy a single local
factor by Lemma~\ref{lem:silver-local-partner}, since five is inert
in $K$. Its trace $2/5$ is not integral. If the last two characters
instead form one quadratic orbit, the common codegree is at most
five and the absolute squareclass of the field discriminant is one
or six. The real possibility $\mathbb Q(\sqrt6)$ requires codegree
at least its discriminant by the quadratic row-difference test.
In the imaginary case the unique dual pair gives $h=|\operatorname{disc}F|b^2$
by conjugate-row orthogonality. Among $\mathbb Q(i),\mathbb Q(\sqrt{-6})$
only the Gaussian pair of norm four remains.

Consider the five rational pairs which still require exclusion.
The order-discriminant identity says the primitive families of
$E,H$ agree if $gh$ is a square, and differ otherwise. First let
$E=m$ be rational. In $(3,27),(4,9)$ the other orbit is $H=nU^j$.
For $j=1$, $t\ge17/27-1/3$ contradicts $6t^2\le1/3$. For
$j\ge2$ minimization of $6t^2+(48/49)(L-t)^2$ gives
$(16/19)L^2>1/3$. In $(3,18),(3,32),(4,6)$ put $H=nVU^j$.
Integrality of $E/H$ gives $m=6dn$ with $d$ integral. For
$L=P/R$ in lowest terms,
\[
 n=(RB_j+R/(3d))/P.
\]
Since $n$ is even, the respective possibilities for $d$ are
$1,3$; $1,2,4,8,16$; and $1,2$. The values $B_1=3,B_2=11$ give no integral $n$. For $B_j=B\ge41$,
the mass equation gives $n\le(B+1/3)/L$. In all three cases
$L\ge7/12$, hence
\[
 \delta_H^2\ge(7/12)^2\frac{B^2-2/3}{(B+1/3)^2}>1/3.
\]
The last inequality amounts to $3B^2-96B-114>0$.

This finishes the rational-$E$ alternatives which require exclusion.

Assume now $E$ irrational. From $E/D$ integral and $D>E$ the
same-family exponent increases by at least one; from the $V$ family
to the unit family it increases by at least two since $V>U$, and the
reverse change multiplies the scalar quotient by at least six and
increases the exponent at least once since $V<6$. The reciprocal
trace sequences give, respectively for unit-family and $V$-family
$E$,
\[
 \delta_E^2+2\delta_D^2\ge(99/4)t^2\quad\hbox{or}\quad17t^2.
\]
Indeed in the former case $s/t\ge7/2$ and the squared gap fractions
are at least $3/4,48/49$; in the latter $s/t\ge3$ and they are at
least $1/3,25/27$. The case $H=nU$ contradicts $17t^2\le1/3$
since $t\ge7/12-1/3=1/4$. Every remaining $H$ has squared gap
fraction at least $48/49$ or $25/27$ according to its primitive
family. Minimizing the resulting two terms in $t$ gives the following
coefficients multiplying $L^2$:
\[
\begin{array}{c|cc}
 &H\text{ unit family}&H\text{ in the }V\text{ family}\\\hline
 E\text{ unit family}&1584/1681&2475/2773\\
 E\text{ in the }V\text{ family}&816/881&425/484
\end{array}
\]
For $(3,27),(4,9)$ the diagonal entries exceed $1/(3L^2)$ since
$L\ge17/27$; for $(3,32)$ both off-diagonal entries do likewise
at $L=61/96$. For $(3,18)$ the lower left entry suffices at
$L=11/18$. In the other orientation put $E=mU^k$, $H=nVU^j$.
At three the three-axis requires one totally ramified partner; the
eighteen-axis cannot join it as $1/3+1/18$ is nonintegral. The
simultaneous unit excludes the $H$ branch, whose value reduces to
minus one. Thus the three-axis and $E$ form the local factor,
implying $1/3+A_k/m\in\mathbb Z_3$. Since $3\nmid A_k$, we have
$3\parallel m$. Yet $m=6dn$ and the mass equation is
$11m=18(A_k+6d B_j)$, a contradiction.

It remains to exclude $(4,6)$. Here $L=7/12$. A nonunit invertible
would give $t\pm w\pm1/4\pm1/6=0$ by its signed unit coefficient.
Only $t=1/12$ is consistent with $17t^2\le1/3$, and then
$17t^2+25w^2/27>1/3$. The two short signed rows have different supports modulo two.
Indeed their squared norms are four and six and their unit entry is one,
so they have respectively four and either two or six odd entries.
They belong to different dyadic factors. Each needs a partner,
by the nonintegral trace of its rational idempotent. Only the two
ramified quadratic branches are available, so the two factors each
have rank three and residue field $\mathbb F_2$. This description
modulo two carries over to $R$. Every ordinary rational codegree
is even: an odd codegree would split off an unramified factor,
necessarily of rank one at each prime above two. Codegree two would
give an invertible.

All simple objects are self-dual, since $T_{\mathbb Q}$ is totally real.
By Lemma~\ref{lem:real-four-six}, each ordinary four-row is rational.
It reduces to the unique residue character with support of size four.
Two such rows would give a nonunit invertible: for their idempotents
$e_1,e_2$, the common support makes $1-2(e_1+e_2)$ integral.
This element has norm one, zero unit coefficient and Perron value one.
Thus among the ordinary characters of rational codegree at most one
has codegree four; the rest have codegree at least six.

Put $L=7/12$ as above. The ratio bounds already proved give
$s\ge3t$, $\delta_D^2\ge(25/27)s^2$ and
$\delta_H^2\ge(25/27)(L-t)^2$, since $H=nU$ has been excluded.
If $s\ge1/4$, the mixed bound is at least
\[
      \frac{25}{27}\bigl((L-s/3)^2+2s^2\bigr)
      \ge\frac{25}{72}>\frac13 ,
\]
the quadratic being increasing on $[1/4,\sqrt2/3]$. Otherwise
the ordinary complementary codegrees cannot all be rational, as
their total mass would be at most $1/4+3/6=3/4$. If there is one
irrational pair its mass is at least $1-s-1/4-1/6=L-s$.
If there are two, their masses sum to $1-s$. Each squared gap is
at least a third of its squared mass, and
$(1-s)^2\ge2(L-s)^2$ for $0\le s\le1/4$.
Thus the additional term in the mixed bound is at least
$2(L-s)^2/3$. Since $L-t>1/2$, we obtain
\[
 \frac13\ge\frac{25}{108}+\frac{50}{27}s^2+\frac23(L-s)^2
       \ge\frac{25}{108}+\frac{25}{51}L^2>\frac13 .
\]
This also excludes $(4,6)$.
\end{proof}

\subsection{Reduction to the conductor-eight profile}

\begin{theorem}\label{thm:full-quadratic-reduction}
Let $\mathcal C$ be a complex fusion category of rank six whose global Frobenius--Perron dimension generates a quadratic field. If $\mathcal C$ does not admit a spherical structure, this field is $K=\mathbb Q(\sqrt2)$ and its signed character table has formal codegrees
\[
 E,E',H,H',4,4,\qquad E,H\notin\mathbb Q.
\]
The two short characters are rational or constitute a Gaussian pair. For a Gaussian pair or two rational rows with the same support, the category has a nonidentity invertible with two fixed points and two free orbits, the ordinary negative rows match the signed short rows, and the other ordinary rows form two irrational quadratic pairs. In this case $\operatorname{Tr}(1/E)<1/1000$.
\end{theorem}
\begin{proof}
By Proposition~\ref{prop:quadratic-full-character} both algebras
are commutative, both dimension characters generate $K$, and some
complementary signed codegree is irrational. Remove the two characters
of the dimension orbit. If all four remaining codegrees are irrational,
Proposition~\ref{prop:r6-three-quadratic-pairs} leaves conductors five
and eight. A repeated pair is excluded by
Proposition~\ref{prop:r6-repeated-pair-five} and
Theorem~\ref{thm:r6-repeated-pair-eight}. This includes a single
quartic character field. Two distinct pairs give the algebra $K^3$.
At conductor five Lemma~\ref{lem:force-golden-short} produces the
pair excluded by Theorem~\ref{thm:q5-three-short}. At conductor eight
use Theorem~\ref{thm:silver-three-complete}.

There remain one irrational pair and two rational codegrees. Both
rational codegrees are at least three. The reciprocal identity excludes
one and also a conjugate pair of common codegree two.
Theorem~\ref{thm:r6-quadratic-no-two} excludes a rational character of
codegree two. Proposition~\ref{prop:r6-two-pair-conductors} now leaves
the following conductors:
\[
\begin{array}{c|l}
 c & \text{exclusion or remaining profile} \\ \hline
 13,21,40,60 & \text{Theorem~\ref{thm:r6-quadratic-three-fields}}\\
 24,28 & \text{Theorem~\ref{thm:upstream-24-28}}\\
 5 & \text{Theorems~\ref{thm:golden-two-pairs} and~\ref{thm:golden-two-rational-E}}\\
 12 & (3,3),\ (3,8),\ (3,12),\ (4,4)
                       \quad\text{(Theorem~\ref{thm:quadratic-twelve-upstream})}\\
 8 & (4,4)\quad\text{with irrational }E .
\end{array}
\]
At conductor twelve, $(3,3)$ would force three to split by
Theorem~\ref{thm:two-threes-exclusion}. The asymmetric pairs are
excluded by Theorems~\ref{thm:twelve-3-8} and~\ref{thm:twelve-3-12}.
For $(4,4)$ use Theorem~\ref{thm:four-four-distinct} for rational
characters with distinct support, and
Theorem~\ref{thm:folded-twelve-completion} for the other possibilities.

For the last line of the table,
Proposition~\ref{prop:silver-two-pair-finite} gives the finite list.
Theorem~\ref{thm:silver-large-pairs} and
Propositions~\ref{prop:silver-four-nine} and~\ref{prop:silver-four-sixteen}
leave $(4,4)$, and Proposition~\ref{prop:silver-four-rational-E}
excludes rational $E$. For the matched characters the asserted
conclusions follow from Propositions~\ref{prop:folded-silver-small-mass}
and~\ref{prop:folded-silver-actual}, together with
Corollaries~\ref{cor:folded-silver-matching}
and~\ref{prop:folded-silver-ordinary-irrational}.
\end{proof}

\section{The quadratic field of conductor five}
\label{sec:quadratic-five-cases}

In the split case with two complementary irrational codegree pairs,
the projection estimate forces one pair of reciprocal mass $1/2$.
The other arrangement treated here has two rational complementary codegrees. Their contributions to a single induction
column impose a congruence. We use the same elementary arithmetic for
both arrangements and for the two small conductors treated below.

\begin{lemma}[Quadratic codegrees and their reciprocal traces]
\label{lem:small-quadratic-recurrence}
For the three fields in the table, let $u$ be the least totally positive
unit greater than one and let $v$ be as indicated. An oriented totally
positive integral $d$-number $f\ge f'$ has the form
$f=mP_{e,i}$, where $P_{e,i}=v^eu^i$,
$m\in\mathbb Z_{>0}$, $e\in\{0,1\}$ and $i\ge0$.
Write $B_{e,i}=\operatorname{Tr}(1/P_{e,i})$. For fixed $e$, these
integers satisfy $B_{e,i+2}=\operatorname{Tr}(u)B_{e,i+1}-B_{e,i}$.
Their initial values are
\[
\begin{array}{c|c|c|c|c}
\operatorname{disc}K & u &v &(B_{0,0},B_{0,1})&(B_{1,0},B_{1,1})\\\hline
5 &(3+\sqrt5)/2 &(5+\sqrt5)/2 &(2,3)&(1,2)\\
8 &3+2\sqrt2 &2+\sqrt2 &(2,6)&(2,10)\\
12&2+\sqrt3 &3+\sqrt3 &(2,4)&(1,3)
\end{array}
\]
For any shift $r$, the determinant of two consecutive terms
$B_{e,k}B_{f,k+r+1}-B_{e,k+1}B_{f,k+r}$ is independent of $k$.
\end{lemma}
\begin{proof}
The unit $f/f'$ equals $u^{2i+e}$. Since $v/v'=u$, division by
$P_{e,i}$ gives a positive rational number. The coordinates of $v^e$
in the integral basis are primitive, and multiplication by the unit
$u^i$ preserves primitivity. This proves that the rational factor is
an integer. The trace recurrence follows from the minimal polynomial
of $u$, and its companion matrix has determinant one.
\end{proof}

\begin{lemma}[Relative indices in a quadratic recurrence]\label{lem:relative-quadratic-indices}
In the notation of Lemma~\ref{lem:small-quadratic-recurrence} write
$E=mP_{e,k}$, $H=nP_{f,j}$ and $D=MP_{a,\ell}>E$. Set $b=N(v)$,
$c=m/M$, $h=m/n$, $r=\ell-k$ and $z=j-k$, with $c,h$ positive
integers. Write $t,w,s$ for the reciprocal traces of $E,H,D$ and
$L=t+w$. Then $s/t>c^2/(2b)$. If $h\mid qB_{e,k},qcB_{a,\ell}$
for an integer $q>0$, then
\[
 h\mid qcW,\qquad
 W=W_{e,a}(r)=B_{e,0}B_{a,r+1}-B_{e,1}B_{a,r}\ne0 .
\]
For a lower bound $I$ on $k,j,\ell$ put $C_g=(v/b)^g$,
$x=h(C_f/C_e)u^z$, $y=c(C_a/C_e)u^r$. With
$0<\eta\le(1-u^{-2I})/(1+u^{-2I})$ the three reciprocal gaps obey
\[
 (\delta_E,\delta_H,\delta_D)\ge
 \frac{\eta L}{1+x}(1,x,y),\qquad
 \frac{\eta Ly}{1+x}\le s\le\frac{Ly}{\eta(1+x)} .
\]
For fixed relative indices, all three trace sequences modulo any
integer are exhausted by iterating one period.
\end{lemma}
\begin{proof}
Writing $P=P_{g,i}$ gives
$P/b^g\le B_{g,i}\le2P/b^g$ for $i\ge0$. Since
$P_{a,\ell}>cP_{e,k}$, this bounds $s/t=cB_{a,\ell}/B_{e,k}$
as claimed. Take the constant determinant at $k$ for the divisibility.
Its vanishing would make the two sequences proportional. Their
coefficients of $u^i,u^{-i}$ would then give $2r+a-e=0$, hence
$r=0,a=e$, contrary to $D>E$. The exact trace and gap numerators
are $C_g u^i(1\pm u^{-2i-g})$. Dividing by
$L=(B_{e,k}+hB_{f,j})/m$ gives the stated bounds. Finally the
recurrence is invertible on consecutive pairs modulo every integer,
so the initial triple of pairs recurs.
\end{proof}

\subsection{A forced pair of reciprocal mass one half}
\begin{lemma}\label{lem:force-golden-short}
Assume that $R$ is commutative, both dimension character fields equal $K=\mathbb Q(\sqrt5)$, and $T_{\mathbb Q}\simeq K^3$, with irrational codegrees on both complementary factors. If no Galois conjugate is pseudo-unitary, one complementary pair is $5\pm\sqrt5$.
\end{lemma}
\begin{proof}
Put $u=(3+\sqrt5)/2$ and $v=(5+\sqrt5)/2$.
Orient each pair by its larger conjugate, and transport the category
so that $E$ has this orientation. Every codegree is $mu^k$ or
$mvu^k$, with $m\ge1$ and $k\ge0$; irrational unit-family
codegrees have $k\ge1$.
Write $t_A=\operatorname{Tr}(1/A)$ and
$\delta_A=1/A'-1/A$.
The separate signed Fourier bound is
\[
 t_E+t_H+t_J=1,\qquad
 \delta_E^2+\delta_H^2+\delta_J^2\le\frac15.
\]
Equation~\eqref{eq:quadratic-trace-order} gives $s=\operatorname{Tr}(1/D)\ge3t_E/2$.
The ordinary Perron column has unit coefficient one; the dimension
orbit contributes $s>t_E$, so one complementary orbit contributes
less than its reciprocal mass.

The shortest orbit $nv$ admits no such deficiency, since
$\operatorname{Tr}(\beta/v)$ is a positive integer for every totally
positive integral $\beta\in K$. Thus the complementary pairs
cannot both be shortest orbits. Every other irrational pair satisfies
$\delta_A^2\ge(5/9)t_A^2$.

We first show that a complementary orbit is $2v$ or $3v$.
If both complementary orbits are long, Cauchy's inequality gives
\[
 \frac15\ge\frac5{23}(1-2/E)^2,
\]
and hence $E<49$. If the shortest orbit is $nv$ with $n\ge4$,
put $w=1/n\le1/4$ and let $z$ be the other complementary mass.
Since $\delta_E+z=1-w-2/E$, the same inequality gives
\[
 \frac15\ge\frac5{14}(1-w-2/E)^2+\frac{w^2}{5}.
\]
If $E\ge80$, the right side is at least
$\frac5{14}(39/40-w)^2+w^2/5$. This function decreases on
$0\le w\le1/4$ and has value $897/4480>1/5$ at $w=1/4$.
Therefore $E<80$ unless a complementary orbit is $2v$ or $3v$.

For $E<80$, the squared distance between the two dimension rows
gives $5\mid\operatorname{Tr}(E)$. Enumerate $E=mP_{e,k}<80$
with this divisibility and $t_E<1$. There are 49 choices and $m\le75$.
The scalars for $H,J$ divide $m$, by integrality of $E/H,E/J$.
Their exponents are at most four: larger indices have reciprocal
trace greater than 75 before scaling. Of these possibilities, exactly
ten satisfy the mass identity, the signed Fourier bound and the
order-discriminant test $N(EHJ)/125\in(\mathbb Z_{>0})^2$.
All ten contain $2v$ or $3v$ (\codefile{short_orbit_check.py}).

We next exclude the shortest orbit $H=3v$ whenever the other
complementary orbit is long. The discriminant condition forces
$E$ and $J$ to belong to the same primitive family.
Suppose both belong to the unit family. The integral central unit
of Lemma~\ref{lem:central-witness} for the nontrivial automorphism of $K$
has values $u^{2k},u,u^{2j}$ on the three factors.
At three the field $K$ is unramified quadratic and $u^2=-1$ in
its residue field. The middle factor therefore has residual
polynomial $x^2+1$, while the other two have residual values
$1$ or $-1$. Polynomial Chinese remaindering isolates the middle
factor by an idempotent in $T\otimes\mathbb Z_3$.
Its coefficient of the unit is $\operatorname{Tr}(1/H)=1/3$,
contradicting integrality of that coefficient.

Suppose instead $E=mvu^k$ and $J=nvu^j$, with $j\ge1$.
Integrality gives $3\mid m$ and $n\mid m$.
Put $t=t_E$ and $z=t_J$, so $t+z=2/3$.
If $k\ge1$, the two long $v$-family pairs have squared-gap
ratios at least $4/5$. Consequently
\[
 \frac15\ge\frac45(t^2+z^2)+\frac1{45}\ge\frac15.
\]
Equality requires $k=j=1$ and $t=z=1/3$, hence
$E=J=6vu$. An ordinary unit-family $D=Mu^\ell$ would require
$5\mid M$ and $M\mid6$, which is impossible. In the other
family $D=Mvu^\ell$, the inequalities $D>E$ and $M\mid6$
force $\ell\ge2$. They give
$\delta_D^2\ge121/(5M^2)\ge121/180>1/5$, contradicting the
separate ordinary Fourier bound.
If $k=0$, then $t=1/m$ and
\[
 \frac15\ge\frac15t^2+\frac45(2/3-t)^2+\frac1{45}.
\]
Thus $t\ge(8-2\sqrt6)/15>1/5$, so $m<5$ and $3\mid m$
give $m=3$. But $t_J=1/3$ then implies
$n=3\operatorname{Tr}(1/(vu^j))>3$, contradicting $n\mid m$.
This proves that the forced shortest orbit is $2v$.
\end{proof}
\subsection{Three quadratic character pairs}

Write $\varphi=(1+\sqrt5)/2$, $K=\mathbb Q(\varphi)$, $u=\varphi^2$
and $v=2+\varphi$. Write $E=\dim(\mathcal C)$ and
$D=\operatorname{FPdim}(\mathcal C)$.
\begin{theorem}\label{thm:q5-three-short}
Let $\mathcal C$ have rank six, with ordinary commutative Grothendieck
algebra and signed rational algebra $K^3$. Suppose the ordinary
Frobenius--Perron character has field $K$, the spherical dimension
character belongs to the first factor, and a second factor has
codegrees $2v,2v'$. If the last factor has irrational codegrees,
then $\mathcal C$ is spherical.
\end{theorem}
\begin{proof}
Set $H=2v$ and write $G,G'$ for the last pair, oriented by $G>G'$.
For $f\in K$ write $t_f=\operatorname{Tr}(1/f)$ and
$\delta_f=1/f'-1/f$. We record two conductor-five projection bounds
\[
 \delta_E^2+\delta_G^2+\delta_H^2\le1/5,\qquad
 \delta_D^2+1/20\le1/5.
\]
The first bound is the signed inequality in
\eqref{eq:quadratic-fourier-bounds}. For the second project the
ordinary difference $(e_P-e_{P'}-e_{JP}+e_{JP'})/2$ onto the signed targets, as in Lemma~\ref{lem:sextic-matrix-compression}.
The two signed dimension labels and their translates contribute
$\delta_D^2$.

Put $\beta=\chi_H(F(P))$ for the signed character of codegree $H$.
Both conjugates of $\beta$ are positive, by symmetry and the
positivity of the ordinary FP character on actual objects.
The contribution of this pair to the signed unit coefficient of
$F(P)$ is $\operatorname{Tr}(\beta/v)/2<1$.
As $\operatorname{Tr}((a+b\varphi)/v)=a$, we obtain
$\beta=1+b\varphi$ with $b=0$ or $1$. At $A_H$ the two rows
$P,P'$ have entries $\beta/(2H)$ and its conjugate by the
ordinary character formula. The four coordinates at $A_H,A_{H'}$
and their translates contribute
$|\beta/H-\beta'/H'|^2=1/20$. In both bounds the full
$U(2)$-fixed projection has squared norm at most $1/5$.

Assume there is no pseudo-unitary conjugate and maximize $E$.
Write $t=t_E$, $w=t_G$ and $s=t_D$. As in
Lemma~\ref{lem:force-golden-short}, $s\ge3t/2$. Since $t+w=1/2$
and the $H$-pair contributes $1/2$ to the Perron induction column,
the $G$-pair contributes less than $w$. This excludes $G=nv$:
$\operatorname{Tr}(\gamma/v)$ is a positive integer for totally
positive integral $\gamma$.

For $e\in\{0,1\}$ set $P_{e,i}=v^e u^i$ and
$B_{e,i}=\operatorname{Tr}(1/P_{e,i})$, with $i\in\mathbb Z$.
Apply Lemma~\ref{lem:small-quadratic-recurrence} to $E,G,D$.
The trace discriminant of $T$ is $N(EHG)$ and differs from $5^3$
by a square factor. As $N(H)=20$, $E$ and $G$ thus have the same
primitive family. Write
\[
 E=mP_{e,k},\qquad G=nP_{e,j},\qquad D=MP_{f,\ell},
 \qquad h=m/n,\quad a=m/M.
\]
Here $k,\ell\ge0$, $j\ge1$, and $h,a$ are positive integers by
integrality of $E/G$ and $E/D$. If $f>e$, then $5\mid a$.
Orthogonality to the conjugate of any $K$-row of codegree $z$
gives $5\mid\operatorname{Tr}(z)$ by taking the squared norm
of the row difference. Thus a scalar attached to the unit family
is divisible by five, since $u\equiv-1\pmod{\sqrt5}$.
Lemma~\ref{lem:column-comparison} gives
$2\ell+f\ge2j+e$ by comparing $D/D'$ with $G/G'$.

Set $B=B_{e,k}$, $L=B_{e,j}$, $C=B_{f,\ell}$. The mass equations
provide an integral congruence. The signed $G$-orbit value at
the FP column is $\operatorname{Tr}(\gamma/G)=1/2-s$ for
$\gamma\in\mathcal O_K$. Multiplying by $n$ yields an integer:
$1/P_{e,j}$ belongs to the codifferent, since $u$ is a unit and
$\operatorname{Tr}(\mathcal O_K/v)\subset\mathbb Z$. Likewise
$n(1/2-t)=L$. It follows that
\begin{equation}\label{eq:golden-congruences}
 m=2(B+hL),\qquad h\mid2B,\qquad h\mid aC-B.
\end{equation}
This second divisibility of the scalar ratio will bound it uniformly.

We claim that $D=G$. First suppose $k\ge2$. The initial values
and the recurrence give
\[
 B_{e,i+1}/B_{e,i}\ge5/2\quad(i\ge2),\qquad
 B_{1,i+1}\ge B_{0,i},\qquad B_{0,i+2}\ge7B_{1,i}.
\]
Indeed each difference satisfies the same recurrence and starts
with nonnegative, nondecreasing values at the indicated index.
In the same family $D>E$ forces $r=\ell-k\ge1$. In passing from
the unit family to the $v$ family, $5\mid a$ and $v<5$ again give
$r\ge1$. In the opposite direction $u<v$ gives $r\ge2$.
Thus the squared gap fraction of $D$ is at least $45/49$, and
the ordinary projection bound gives $s\le7/\sqrt{300}<5/12$.

Also $t\ge1/10$. Indeed if $j\ge2$, both signed gap fractions
are at least $45/49$, giving
$t^2+(1/2-t)^2\le49/300$, impossible for $t\le1/10$.
If $j=1$, then $w=3/n\le3/10$ in the unit family ($5\mid n$),
and $w=2/n\le2/5$ in the other family, since $w<1/2$.
It follows that $s/t<5$. The trace comparisons above now force
$f=e$, $a=r=1$ in $s/t=aB_{f,k+r}/B_{e,k}$.

The determinant $W_{e,e}(1)$ of
Lemma~\ref{lem:relative-quadratic-indices} is five in the unit
family and one in the other family. By \eqref{eq:golden-congruences}
we have $h\mid2W_{e,e}(1)$ and $h\mid2B$; since
$5\nmid B_{0,k}$, in either family $h\le2$.
Dominance gives $j\le k+1$. If $j\le k$, then
$t=B/(2(B+hL))\ge1/6$ and $s\ge5t/2\ge5/12$.
If $j=k+1$ and $h=2$, then $w\ge5/12$ and
$\delta_G^2\ge(45/49)(5/12)^2>3/20$.
Hence $h=1$ and $j=k+1$, proving the claim.

It remains to check $k=0,1$, where $B\le3$.
For $j\ge1$ the squared gap fraction of $G$ is at least $4/5$,
except when $(e,j)=(0,1)$, where $w\le3/10$.
Thus $w\le\sqrt3/4$ and $t>1/16$. Hence $m<16B\le48$,
$hL<7B\le21$, and $j\le3$. Dominance gives
$\delta_D^2\ge5s^2/9$, so $C=Ms<27$ and $\ell\le3$.
For transparency the remaining scalar checks are as follows.
The middle column tests $h\mid2B$, $m=2(B+hL)<16B$,
the family divisibilities and $\delta_E^2+\delta_G^2\le3/20$.
The last tests $a\mid m$, $h\mid aC-B$, the family divisibilities,
$D>E$, dominance and $\delta_D^2\le3/20$.
\[
\begin{array}{c|l|c}
(e,k)&(j,h)\text{ after the signed test}&(j,h,f,\ell,a)\\\hline
(0,0)&(1,1)&(1,1,0,1,1)\\
(0,1)&(2,1)&(2,1,0,2,1)\\
(1,0)&(1,1),(1,2)&(1,1,1,1,1)\\
(1,1)&(1,1),(2,1),(1,2),(1,4)&(2,1,1,2,1)
\end{array}
\]
These substitutions also give $D=G$
(\codecheck{golden_orbit_check.py}{arithmetic check}).
The real signed character of this codegree gives a spherical
structure by Lemma~\ref{lem:section-signs}.
\end{proof}
\subsection{Two irrational pairs and two rational codegrees}

Let $K=\mathbb Q(\sqrt5)$, $u=(3+\sqrt5)/2$ and $v=(5+\sqrt5)/2$. Write $R$ and $T$ for the ordinary and signed rings.
\begin{theorem}\label{thm:golden-two-pairs}
Suppose a complex rank-six fusion category has commutative ordinary and signed character algebras. Assume the ordinary Frobenius--Perron and signed spherical dimension characters have field $K$. If the signed codegrees are $E,E',H,H',a,b$, with $E$ the intrinsic global dimension, $E,H$ irrational in $K$, and $a,b$ rational integers at least three, then the category admits a spherical structure.
\end{theorem}
\begin{proof}
Write $D$ for the Frobenius--Perron dimension and $g'$ for conjugation in $K$. Order $a\le b$ and $H>H'$. If there is no pseudo-unitary Galois conjugate, transport to an embedding with $E>E'$; then $E<D$ and $\alpha=E/D$ is a totally positive integer of $K$ with $\alpha<1<\alpha'$. Put
\[
 t=\operatorname{Tr}(1/E),\quad w=\operatorname{Tr}(1/H),\quad s=\operatorname{Tr}(1/D),\quad L=1-1/a-1/b=t+w,
\]
and write $\delta_g=1/g'-1/g$ for $g\ge g'$. At conductor five \eqref{eq:quadratic-fourier-bounds} gives
\begin{equation}\label{eq:two-golden-fourier}
 \delta_E^2+\delta_H^2\le 1/5,\qquad \delta_D^2\le1/5.
\end{equation}
The central class attached to conjugation has ordinary FP value $D/D'$ and signed $H$-value $H/H'$ by the induction character formula. Thus $D/D'\ge H/H'>1$. Also $\operatorname{Tr}(\alpha)\ge3$ and $\alpha'-\alpha\ge\sqrt5$, whence
\begin{equation}\label{eq:two-golden-rearrange}
 s=\tfrac12\big(t\operatorname{Tr}(\alpha)+\delta_E(\alpha'-\alpha)\big)\ge\tfrac32 t,
 \qquad \delta_D=s-2/D\ge ut-(2+\sqrt5)/E.
\end{equation}

Use $P_{e,k}=v^e u^k$, $B_{e,k}=\operatorname{Tr}(1/P_{e,k})$ and the $d$-number parametrization from Theorem~\ref{thm:q5-three-short}. Write
\[
 E=mP_{e,k},\qquad H=nP_{f,j},\qquad D=MP_{d,\ell},
\]
with $e,f,d\in\{0,1\}$ and nonnegative indices. The four integral ratios $E/H$, $E/D$, $E/a$ and $E/b$ give $a,b,n,M\mid m$, with $5n\mid m$ if $f>e$ and $5M\mid m$ if $d>e$. Every unit-family scalar is divisible by five, by orthogonality of the two rows. For $q=\operatorname{lcm}(a,b)$ the two last signed characters contribute $Z/q$ to the unit coefficient of the ordinary FP induction object. Here $Z$ is an integer between $q/a+q/b$ and $q-1$: at rational characters the evaluations are positive integers, and for a quadratic orbit of common codegree the trace of the totally positive integral evaluation is at least two. At the first pair the contribution is $s$; at the second it is $\operatorname{Tr}(\beta/H)$, with $\beta$ totally positive integral. Thus
\begin{equation}\label{eq:two-golden-column}
 0<1-Z/q-s=\operatorname{Tr}(\beta/H)<w,\qquad n(1-Z/q-s)\in\mathbb Z.
\end{equation}
Indeed $s>t$ and $1/P_{f,j}$ lies in the codifferent. This excludes $H=nv$, since $\operatorname{Tr}(\beta/v)$ is a positive integer and $\operatorname{Tr}(1/v)=1$. Thus $2j+f\ge2$.

For exponent $2j+f=2$ we have $w=3/n$, with $5\mid n$. The value $n=5$ exhausts the Fourier bound while $\delta_E>0$, so $w\le3/10$. For exponent three we have $w=2/n$ and $\delta_H^2=4w^2/5$; likewise $n>4$, hence $w\le2/5$. For exponent at least four, $H/H'\ge u^4$ gives $\delta_H^2\ge45w^2/49$. Consequently $w\le7/15$ in all cases. Except at exponent two, $\delta_H^2\ge4w^2/5$.

\smallskip\emph{Large rational codegrees.}
Suppose $1/a+1/b\le9/25$, so $L\ge16/25$. If $E\ge1000$, \eqref{eq:two-golden-rearrange} implies $t<173/1000$, and then $w>467/1000$. This contradicts $w\le7/15$. Hence $E<1000$. Now $m<1000$, $k,j\le7$ by $E\ge u^k$, $B_{f,j}<n\le m$, and $B_{f,8}\ge1597$. Test $1\le m<1000$, $0\le k\le7$, $1\le j\le7$, $n,a,b\mid m$ against $E<1000$, the scalar divisibilities, $L=t+w\ge16/25$, and the signed inequality in \eqref{eq:two-golden-fourier}. The three surviving rows $(e,k,m;f,j,n;a,b)$ are
\[
 (0,2,20;0,1,10;4,10),\quad (1,1,6;1,1,6;6,6),\quad (1,2,20;1,1,5;4,10).
\]
For distinct rational axes the absolute regular trace discriminant $N(EH)ab$ is a square, the two quadratic field discriminants multiplying to $25$. This excludes both rows with $a=4,b=10$. In the middle row the unit family for $D$ would require $5\mid M\mid6$. In the other family $D>E$ forces $\ell\ge2$, and then $\delta_D^2\ge(25-4/5)/6^2>1/5$.

\smallskip\emph{Small rational codegrees.}
We may now assume $1/a+1/b>9/25$, so either $a=3$ and $b\le37$, or $a=4$ and $b\le9$, or $a=5$ and $b\le6$. For two rational rows the discriminant condition requires $ab$ to have squareclass $1$ or $5$. If $a=3$, its partner at three is the rational $b$-axis by
Lemma~\ref{lem:silver-local-partner}, since $K$ is inert there.
The congruent rows have squared distance $3+b$, divisible by nine.

If the complement is quadratic, $a=b$ and its field $F$ has $|\operatorname{disc}(F)|$ of squareclass $1$ or $5$ since the absolute order discriminant is $N(EH)a^2$. Moreover $|\operatorname{disc}(F)|\mid2a$ by taking the squared norm of the difference of the two conjugate rows. This excludes $a=3$; for $a=4$ it gives $F=\mathbb Q(i)$. The complete list of profiles is therefore
\[
 (3,15),\quad(4,4),\quad(4,5),\quad(4,9),\quad (5,5).
\]
In particular $L\ge1/2$. The discriminant forces $e+f$ odd for $(3,15),(4,5)$ and even for $(4,4),(4,9)$.

\smallskip\emph{The finite recurrence.}
Put $h=m/n$, $c=m/M$, $r=\ell-k$, $z=j-k$ and write $B_0=B_{e,k}$, $B_1=B_{f,j}$, $B_2=B_{d,\ell}$. We have
\begin{equation}\label{eq:golden-two-integrality}
 m=(B_0+hB_1)/L,\qquad m(q-Z)-qcB_2\in qh\mathbb Z_{>0}.
\end{equation}
The second assertion is \eqref{eq:two-golden-column}. Thus $h\mid qB_0,qcB_2$. From $w\le7/15$, $L\ge1/2$ and
$\delta_D^2\ge5s^2/9$ we have $t\ge1/30$, $s\le3/5$ and
$s/t\le18$. Lemma~\ref{lem:relative-quadratic-indices} gives
$c\le13$, $1/v\le u^r<180v$ and $-1\le r\le6$. The initial
pairs give $0<|W|\le720$ on the allowed range, so $h\mid qcW$
and $h\le36\cdot13\cdot720$. Since $1/4\le L/t-1=hB_1/B_0<30$
and $u^z/3<B_1/B_0<3u^z$, we may take $-15\le z\le4$.

The \codecheck{golden_two_pairs_check.py}{finite recurrence check}
enumerates these relative indices with the family divisibilities,
$2z+f\le2r+d$ and $v^{d-e}u^r>c$. Equality
$(f,z,h)=(d,r,c)$ gives a spherical section since $D=H$.
For all other tuples and $0\le k<36$, test $2k+e\ge1$,
$j,\ell\ge1$, the scalar divisibilities,
\eqref{eq:golden-two-integrality} and \eqref{eq:two-golden-fourier}.
None survives.

For $k\ge36$ take $I=21$, $\eta=999999/1000000$ and $x,y$ of
Lemma~\ref{lem:relative-quadratic-indices}. The two Fourier tests
and positivity of the $H$ contribution require
\[
 5\eta^2 L^2(1+x^2)\le(1+x)^2,\qquad
 5\eta^2 L^2 y^2\le(1+x)^2,\qquad
 1-Z/q>\eta Ly/(1+x).
\]
The first two inequalities leave $0,16,6,0,0$ tuples for the five
profiles in order. Writing $L=N/P$ in lowest terms, the modulus
$N\operatorname{lcm}(qh,5h,5c,q)$ determines integrality and the
required residues of $m=P(B_0+hB_1)/N$. The largest recurrence
periods are $120$ and $750$; no state satisfies the divisibilities
and \eqref{eq:golden-two-integrality} for an allowed $Z$.
\end{proof}

\subsection{Rational intrinsic dimension}
\begin{theorem}\label{thm:golden-two-rational-E}
Theorem~\ref{thm:golden-two-pairs} also holds when $E$ is rational and the signed dimension character still has field $K$.
\end{theorem}
\begin{proof}
Use the notation of that theorem, with $E=m$ a multiple of five and $t=2/m$. Assume $D>E$. The bound $\delta_H^2\le1/5$ now allows equality. As before, $H$ has exponent $2j+f\ge2$. The exponent-two and exponent-three families still give $w\le7/15$ except when $H=5u$ or $H=4vu$, of masses $3/5$ and $1/2$, respectively. Hence $w\le3/5$ always. If $L\ge16/25$ then $t\ge1/25$, giving $m\le50$. Otherwise the same discriminant and local argument gives the five previously listed $(a,b)$ and $L\ge1/2$. Unless $H$ is exceptional, $t\ge1/30$, so $m\le60$. At the two exceptional values, the only possibilities for $L$ are $1/2,11/20,3/5,23/36$. Direct substitution in $m=2/(L-w)$, requiring a positive integer, gives $m\le40$. Thus $m\le60$ in every case.

For $m=5,10,\ldots,60$ all scalar divisors and the rational codegrees
divide $m$, and $j,\ell\le4$ since both primitive traces at index
five exceed 60. Put aside $D=H$, which gives a spherical section. Test $D>m$,
$D/D'\ge H/H'$, $s\ge3/m$, the mass equation $2/m+w=L$,
the discriminant and local restrictions, and the two Fourier
bounds of the preceding theorem. One further projection, as in
Theorem~\ref{thm:q5-three-short}, gives
\[
 \delta_D^2+n^{-2}\bigl(\operatorname{Tr}(\gamma/v^f)^2
               -4N(\gamma)/5^f\bigr)\le1/5,\qquad
 \gamma=\beta/u^j=x+y\varphi\in\mathcal O_K .
\]
Here $t_0=n(1-s-Z/q)>0$ equals $2x+y$ if $f=0$ and $x$ if
$f=1$; positive norm gives $|y|\le2t_0$. This leaves seven
numerical profiles.

To reconstruct the character rows for each profile, write a
quadratic row of codegree $c+d\varphi$ as $A+B\varphi$.
Orthogonality to its conjugate gives three finite sphere conditions:
\[
 B^2=(2c+d)/5,\quad A^2=c-B^2,\quad 2A\cdot B=d-B^2,
 \qquad A_0=1,\ B_0=0 .
\]
Enumerate the two mutually orthogonal signed pairs with nonvanishing
dimension row, and the Perron rows with $d_i\ge\max(1,|d_i'|)$.
The signed pairs leave a plane with projector $P$. For unequal
complementary codegrees $g,h$, a rational row $r$ of norm $g$
determines the other row as $h(P_{0i}-r_i/g)_i$. For equal codegrees
their means are $C_i=gP_{0i}/2$ and their centered products are
$gP_{ij}/2-C_iC_j$. These determine the contributions to the signed
table for either real or complex conjugate rows.
Keep integral tables $T$ dominated by the Perron row and satisfying
\[
 d_i d_j-\sum_k|T_{ij}^k|d_k\in2\sum_{k\ne0}\mathbb Z_{\ge0}d_k .
\]
The \codecheck{golden_rational_two_check.py}{row computation} leaves
only $(E,H,g,h,D)=(20,4vu,5,5,20u^2)$ and
$(60,15u^2,4,4,60u^3)$, each with two signed matches.

For $E=60$, put $d=1+2\varphi$, so $d^2=1+4d$. Order the
simples as $1,g,Y,A,B,V$ with dimensions
$1,1,4d,d,d,2+6d$. In both signed matches the coefficient of $V$
in $Y^2$ is $-2$. The unique simple $Y$ of its dimension is
self-dual and fixed by $g$. Hence $1,g$ occur once each in its
ordinary square. If $l$ denotes the multiplicity of $V$, comparison
of constant coordinates in the basis $1,d$ gives
$16=2+2l$, since $\operatorname{FPdim}(Y^2)=16+64d$.
Thus $l=7$, contradicting the parity of the signed coefficient.

For $E=20$ the Perron row is
$d=(1,\varphi,\varphi,1+3\varphi,1+\varphi,3+4\varphi)$.
Let $W$ be the simultaneous column of Lemma~\ref{lem:central-witness}
for the nontrivial automorphism of $K$. In the signed order it is
$(1,1,1,1,1,2)$ in both tables
(\codecheck{golden_witness_parity.py}{witness calculation}).
Indeed $W$ acts as $f/f'$ on a row of codegree $f$, so character
orthogonality recovers its coefficients. More explicitly, the
rational codegrees give eigenvalue one, and the $H$-pair has
$H=12+16\varphi$, $N(H)=80$. Thus the coefficient of $b_i$ in
$W-1$ is
$(1/H'-1/H)(\chi_H(b_i^*)-\chi_H'(b_i^*))$, the
$\varphi$-coefficient of $\chi_H(b_i^*)$.
Since the two orders agree modulo two, evaluating this vector on
$d$ must give $D/D'$ modulo $2\mathcal O_K$. The evaluation is
$9+14\varphi$, whereas $D/D'=u^4=13+21\varphi$.
\end{proof}

\section{The quadratic field of conductor twelve}
\label{sec:quadratic-twelve-cases}
The two rational codegrees left by the conductor reduction are
$(3,3)$, $(3,8)$, $(3,12)$ and $(4,4)$.
Theorem~\ref{thm:two-threes-exclusion} excludes the first pair since three ramifies in $K$.
For the middle two, a projection estimate bounds the reciprocal
mass of the Perron orbit. The last pair is governed by the supports
of its two rows and reduction at two.

\subsection{Rational codegrees three and twelve}
\label{sec:twelve-3-12}

\begin{lemma}\label{lem:real-four-six}
Let $R$ be the commutative Grothendieck ring of a complex fusion
category of rank six. A real character of formal codegree four or six
is rational.
\end{lemma}
\begin{proof}
Let $F$ be the field of a real row of codegree $h=4$ or $6$.
It is totally real and abelian by cyclotomicity
\cite[Theorem~8.51]{ENO}. On its factor the coefficient-of-unit
pairing is $\operatorname{Tr}_F(xy)/h$, so $h\mathfrak D_F^{-1}$
pairs integrally with $R$. Perfectness puts it in
$R\cap F\subseteq\mathcal O_F$. Hence
\[
      v_p(\operatorname{disc} F)\le [F:\mathbb Q]v_p(h).
\]

For $h=4$ the field lies in the real cyclotomic $2$-power tower.
This tower is cyclic, with unique quartic subfield
$\mathbb Q(\zeta_{16})^+$ of discriminant $8\cdot16^2=2^{11}$.
The tower discriminant formula would give
$v_2(\operatorname{disc}F)/[F:\mathbb Q]\ge11/4>2$ if
$[F:\mathbb Q]\ge4$. Thus $F\subseteq\mathbb Q(\sqrt2)$.

For $h=6$ the conductor--discriminant formula bounds the average
conductor exponent by one at both two and three. At two, each
ramified character has exponent at least two, so nontrivial inertia
has order two and every ramified character has exponent two.
At three, wild inertia would give exponent at least two for at least
two thirds of the characters. Thus the conductor divides twelve,
and $F\subseteq\mathbb Q(\zeta_{12})^+=\mathbb Q(\sqrt3)$.

In either quadratic alternative write $c=8$ or $12$ respectively.
By row orthogonality, the squared distance between the two
conjugate rows is $2h=c\sum_i b_i^2$, where $b_i$ are the integral
coefficients of their differences in units of $\sqrt c$.
Since $c=2h$, exactly one coordinate differs. This contradicts
orthogonality of their difference to the positive Perron row.
Indeed the orbit cannot contain that row: $D=h\le6$ would force
every simple to be invertible, and the Perron row would be rational.
\end{proof}

\begin{theorem}\label{thm:twelve-3-12}
Let $\mathcal C$ be a complex rank-six fusion category with commutative ordinary and signed Grothendieck algebras. Suppose its ordinary Frobenius--Perron and signed spherical dimension characters have field $K=\mathbb Q(\sqrt3)$. Suppose the signed rational algebra is $K\times K\times\mathbb Q\times\mathbb Q$, with rational characters of formal codegrees three and twelve. Then $\mathcal C$ admits a spherical structure.
\end{theorem}

\begin{proof}
Use $U=2+\sqrt3$, $V=3+\sqrt3$, $B_{e,i}$ from
Lemma~\ref{lem:small-quadratic-recurrence}, and $A_i=B_{0,i}$.
Assume nonsphericity, orient $E\ge E'$, $H\ge H'$, and write
$t,w,s$ for the reciprocal traces of $E,H,D$. Then $t+w=7/12$.
For $\delta_f=1/f'-1/f$ we use
\begin{equation}\label{eq:twelve-312-mixed}
 \sum_{\rm signed\ pairs}\delta_f^2+
 2\sum_{\rm ordinary\ pairs}\delta_f^2\le1/3 .
\end{equation}

At three the rational rows have distinct residues, since their
squared distance is $15$. Each needs another branch and there
are just two, both ramified quadratic. The simultaneous unit has
value one on the rational rows, while on a branch of type $VU^i$
its value $U^{2i+1}$ reduces to $-1$. Thus $E=mU^k$, $H=nU^j$.
The identity of each local factor has integral trace
$1/a+\operatorname{Tr}(1/f)$, with $a=3$ or $12$ and
$f=mU^k$ or $nU^j$. Since $3\nmid A_i$, the scalar has
three-adic valuation one. Multiply the trace by three and reduce
modulo the ramified prime: $3/a+2/(f/3)=0$ in $\mathbb F_3$.
Consequently $m/3\equiv(-1)^k$ and $n/3\equiv(-1)^j\bmod3$.

Write $D=MV^dU^\ell$, $c=m/M$, $h=m/n$, $r=\ell-k$, $z=j-k$.
Integrality gives $c,h\in\mathbb Z_{>0}$ and $6\mid c$ if $d=1$.
Also $r\ge1$ by $D>E$, and $\ell\ge j$ by $D/D'\ge H/H'$.
The complementary pair is irrational since otherwise
Corollary~\ref{cor:galois-columns} would give $s\le t$, contrary to
\eqref{eq:quadratic-trace-order}. In particular $j\ge1$.

We claim $j\ge3$ and $t>1/180$. For $j=1$, $6\mid n$ by
\eqref{eq:three-quadratic-scalars} and $\delta_H^2=12/n^2<1/3$,
so $w\le1/3$. Then $s\ge2t\ge1/2$ by
\eqref{eq:quadratic-trace-order}, contrary to
$2\delta_D^2\ge3s^2/2$. For $j\ge2$, the signed dominance and
the recurrence give $\delta_D^2\ge48s^2/49$, $s/t\ge7/2$:
in the unit family use $A_{k+1}/A_k\ge7/2$ for $k\ge1$ and
$\ell\ge2$ for $k=0$; in the other use $6B_{1,k+1}\ge9A_k$.
If $j=2$, then $192/n^2<1/3$ gives $n\ge30$ and $t\ge7/60$.
The mixed bound fails since $24t^2+(48/49)(7/12-t)^2>1/3$
on this interval. For $j\ge3$, $\delta_H^2\ge675w^2/676$
gives $w<26/45$, proving the claim.

Only relative indices remain. The two unit coefficients give
\[
 7m=12(A_k+hA_j),\qquad
 m(12-Z)-12cB_{d,\ell}\in12h\mathbb Z,
\]
where $Z/12$ is the rational complementary contribution on the
Perron induction class; the $H$ contribution belongs to
$\mathbb Z/n$. Since $6h\mid m$, we obtain
$h\mid2A_k,2cB_{d,\ell}$, so
$h\mid2cW_{0,d}(r)$ by Lemma~\ref{lem:relative-quadratic-indices}.
Also $h\equiv(-1)^z\bmod3$. Since $2\delta_D^2\ge(96/49)s^2$, we have $s<1/2$ and $s/t<90$.
The same lemma gives $c\le32$. From $U^r\le2(6/V)^d s/t$ we get $r\le4$, hence
$|W_{0,d}(r)|\le672$, $h\le43008$. Finally $1<hA_j/A_k<104$ and
$U^z/2\le A_j/A_k\le2U^z$ give $-11\le z\le4$.

At $D=H$ apply Lemma~\ref{lem:section-signs}. For the other
tuples in these bounds the
\codecheck{twelve_three_twelve.py}{exact recurrence check}
uses the mass identity, scalar divisibilities and
\eqref{eq:twelve-312-mixed} at $k=0,1,2$; none survives.
For $k\ge3$ all three indices are at least three. Put
$x=hU^z$, $y=c(V/6)^dU^r$, $\eta=999/1000$.
Lemma~\ref{lem:relative-quadratic-indices} and the mixed bound give
\[
 147\eta^2(1+x^2+2y^2)\le144(1+x)^2 .
\]
Only $(d,r,c,h,z)=(0,1,1,8,1)$ survives. But then
$7\mid A_k+8A_{k+1}$, whose residues modulo seven cycle through
$6,4,3,1,1,3,4,6$. This final contradiction proves the theorem.
\end{proof}

\subsection{Rational codegrees three and eight}

\begin{theorem}\label{thm:twelve-3-8}
The conclusion of Theorem~\ref{thm:twelve-3-12} also holds when the two rational signed codegrees are three and eight.
\end{theorem}

\begin{proof}
Use $K,U,V,E,H,D,t,w,s,A_i,B_i$ as in the preceding theorem and again assume nonsphericality. Here $t+w=L=13/24$. Again \eqref{eq:quadratic-trace-order} gives $s>t$, so $H$ is irrational. The discriminant of the signed maximal order is a square. The signed trace discriminant has squareclass $24N(E)N(H)$, hence exactly one of $E,H$ belongs to the family of multiples of $VU^j$.

\emph{The local factor of the codegree-three row.}
Denote its induction label by $A$ and the eight-label by $B$.
As in the proof of Theorem~\ref{thm:two-threes-exclusion}, the row
of $M=6S$ at $A$ has entries one on both dimension pairs.
Let $a$ be its common entry on the $H$-pair and put
$r=M_{A,A}$, $z=M_{A,B}$. The signed row sum and norm give
\[
 r+z=1-2a,\qquad 2a^2+r^2+z^2\le7 .
\]
The rational eight-row takes the integral value $8z/3$ on $F(A)$.
Thus $z=0$ and $(a,r)$ is $(0,1)$ or $(1,-1)$.

At three the rational eight-axis splits off by its idempotent. The central element of simultaneous character values $f/f'$ separates the $V$-family branch, with residue $-1$, from the rational three-axis. Since the latter cannot split off by itself, its local factor consists of that axis and the unit-family branch, of codegree $f=m_0U^i$. The integral trace of its identity is $1/3+A_i/m_0$. As $3\nmid A_i$, it follows that $3\parallel m_0$. Writing $\mathfrak p$ for the prime of $K$ over three, multiplication by three and reduction further give $1+2/(f/3)=0$ in its residue field. Hence $f/3\equiv1\pmod{\mathfrak p}$.
If $f=H$, the signed values of $F(A)$ on the two paired branches are $r$ and $aH/3$, so $r\equiv a\pmod3$. Both possibilities above contradict this. Thus $f=E$. The signed character values of $F(A_D)$ on the paired branches are $1$ and $E/D$, so $E/D\equiv1\pmod{\mathfrak p}$. This forces $D$ to be in the unit family like $E$, since a $V$-family element has odd valuation at $\mathfrak p$. We have therefore proved
\[
 E=mU^k,\quad H=nVU^j,\quad D=MU^\ell,
 \qquad 3\parallel m,\quad 6n\mid m,\quad M\mid m .
\]
Here $E/H$ and $E/D$ are integral by induction. The coefficient of $\sqrt3$ in $VU^j$ is odd, while that of a sum of squares in $\mathbb Z[\sqrt3]$ is even; thus $n$ is even. We have $\ell\ge k+1$ because $D>E$, and $\ell\ge j+1$ because $D/D'\ge H/H'=U^{2j+1}$.
Lemma~\ref{lem:quadratic-mass-gap} gives $H/H'\ge U^2$, hence $j\ge1$.

Now $s/t\ge7/2$: for $k\ge1$ this follows from $A_{k+1}/A_k\ge7/2$, and when $k=0$ use $\ell\ge2$. If $j=1$, the mixed inequality gives $25/(3n^2)=\delta_H^2<1/3$, so $n>5$. Since $3\nmid n$ and $2\mid n$, this gives $n\ge8$, hence $w\le3/8$, $t\ge1/6$. But $2\delta_D^2\ge(96/49)s^2\ge24t^2$, a contradiction. We have obtained
\begin{equation}\label{eq:38-gaps}
 j\ge2,\qquad \delta_H^2\ge(361/363)w^2,\quad
 \delta_D^2\ge(48/49)s^2,\quad
 s\le7/(12\sqrt2)<5/12 .
\end{equation}

\emph{Invertibles and dyadic local factors.}
On a positive lift of an invertible all character values are signs. A nonunit invertible's zero unit coefficient gives
$t+\lambda w+e/3+h_0/8=0$ with signs $\lambda,e,h_0$. As $L>1/3+1/8$, one needs $\lambda=-1$. The four possible $t$ are $1/24,1/6,3/8,1/2$, of which only $1/24$ satisfies $24t^2\le1/3$. It forces $e=h_0=1$. Thus multiplication by the positive lift acts as a permutation with eigenvalue $-1$ exactly on the two $H$-rows. It is an involution since every ordinary simple is self-dual and the positive lift respects duality. Two transpositions occur, one interchanging the unit and this invertible. Accordingly the $H$-row has the form $(1,-1,z,-z,0,0)$; orthogonality to its conjugate gives $N(z)=-1$. That norm equation is impossible in $\mathbb Z[\sqrt3]$ modulo three. Thus all nonunit FP dimensions exceed one.

At two the rational three-axis is isolated. The other axes constitute a single local factor. Indeed, the eight-axis cannot be isolated, having nonintegral coefficient trace $1/8$. An isolated $H$-axis would have regular trace discriminant $N(H)$ of odd valuation, impossible since the field discriminant has even valuation. For the $E$-axis note that
\[
  13m/24=A_k+(m/n)B_j\in\mathbb Z ,
\]
so $24\mid m$. Any order on an isolated $K$-axis has the form $A=\mathbb Z_2+2^v\sqrt3\,\mathbb Z_2$ for $v\ge0$, since projection to the irrational coordinate gives a submodule of $\mathbb Z_2$. Its trace dual is $(2^{v+1}\sqrt3)^{-1}A$, by direct computation or the derivative formula, and is $E^{-1}A$ by perfectness. Thus $mU^k/(2^{v+1}\sqrt3)\in A^\times$, giving $v=v_2(m)-1\ge2$. Writing $U^k=c_k+d_k\sqrt3$ now gives $4\mid c_k$, impossible since $c_0=1,c_1=2,c_{k+2}=4c_{k+1}-c_k$. This establishes the claimed local decomposition.

Transfer the factor ranks across $R/2R\simeq T/2T$.
As in Lemma~\ref{lem:r6-two-threes-dyadic}, the rank-one factor is
the unique odd rational-codegree row, of norm $a\equiv3\bmod4$. All even-codegree rational rows of $R$ reduce to the unique residue character of the other local factor, the reduction of the eight-row of $T$. A norm-eight integral vector of length six with unit coordinate one has exactly four odd entries. Hence these rational ordinary codegrees are divisible by four. Two rational four-rows would have the same support. If $e_1,e_2$ are their primitive idempotents, $u=1-2(e_1+e_2)$ is integral by the row idempotent formulas and the common support. It satisfies $u^2=1$, has unit coefficient zero and ordinary FP value one. Since $\|u\|^2=1$, it is a basis element of FP dimension one, a contradiction. Thus there is at most one rational ordinary four-row.

\emph{The ordinary spectrum.}
If an ordinary complementary pair had irrational codegrees of mass $z$, one would have $z\ge5/12-s>0$ and squared gap at least $z^2/3$ as before. By \eqref{eq:twelve-312-mixed} and \eqref{eq:38-gaps},
\[
 1/3\ \ge\ (361/363)(13/24-2s/7)^2
          +(96/49)s^2+(2/3)(5/12-s)^2 .
\]
The minimum exceeds $1/3$ by $222697/41601600$, a contradiction. Also $s<1/4$, since the sum of the first two indicated terms is increasing on $s\ge1/4$ and at $1/4$ exceeds $1/3$ by $92425/10245312$.

Thus the four ordinary complementary codegrees are rational.
Lemma~\ref{lem:real-four-six} makes the four-rows rational and
excludes quadratic six-orbits.

Since $1-s>3/4$ and at most one complementary character has codegree four, the odd codegree must be $a=3$: for $a\ge7$ the sum is at most $1/7+1/4+2/6<3/4$. A four-row is necessary, since all three even rows would otherwise
have norms at least eight: Lemma~\ref{lem:real-four-six} makes a
six-row integral, whereas integral even rows here have norm divisible
by four.

If the two remaining characters form a quadratic orbit of common rational codegree $g$, the regular trace discriminant has squareclass $N(D)\cdot3\cdot4g^2\sim3$, while the maximal order has squareclass $3\operatorname{disc}(F)$, impossible for a real quadratic field $F$. Thus all four characters are rational. Write their codegrees as $3,4,c,d$, where $8\le c\le d$ are multiples of four. The strict inequality $1/3+1/4+1/c+1/d>3/4$ forces $c=8$ and $d<24$. The discriminant squareclass requires $8d$ to be a square, so $d=8$.

\emph{The four ordinary rows.} The four-row and both eight-rows
have the same four odd positions $S$. An eight-row is $\pm1$ on $S$
with just one further entry $\pm2$. The latter entries must be in
different positions: otherwise orthogonality forces the two rows
to agree on $S$ (their unit entries are one) and have opposite even
entries. Their difference would be supported at one position, contrary
to orthogonality to the Perron row.

Negate basis elements, keeping the unit fixed, so the four-row is
one on $S$ and both even entries are positive. This preserves the
integrality of the multiplication table. On $S$ the three rows are
now pairwise orthogonal sign vectors. Permuting the three nonunit
positions gives the first three rows below. The three-row meets $S$
in the unit and exactly one other position, by orthogonality to the
four-row. Up to interchanging the eight-rows its third entry lies at
the first even position. Orthogonality then fixes the entire row:
\[
 \begin{pmatrix}
 1& 1& 1& 1&0&0\\
 1& 1&-1&-1&2&0\\
 1&-1& 1&-1&0&2\\
 1& 0&-1& 0&-1&0
 \end{pmatrix}.
\]
The Perron row in this signed basis must consequently be
$(1,y-2,y,1-2y,1-y,-1-y)$.
Put $v=2y-1$. Its squared norm is $D=2(v^2+3)$, and its
orthogonality to the conjugate row gives $vv'=-3$.
At the fifth basis element, of Perron value $(1-v)/2$, the unit
coefficient of its cube is the integer
\[
 \operatorname{Tr}_K\!\left(\frac{(1-v)^3}{8D}\right)
          +\frac{2^3}{8}-\frac13
       =\frac{14-\operatorname{Tr}_K(v)}{16}.
\]
The equality follows from $v'=-3/v$. Write $v=A+B\sqrt3$.
Then $A\equiv7\bmod8$; the norm equation gives
$3B^2=A^2+3\equiv4\bmod16$, which is impossible.
\end{proof}

\subsection{Two rational four-rows with different supports}

The same calculation at two serves both remaining conductors. We use
the notation $u,v,B_{e,k}$ of Lemma~\ref{lem:small-quadratic-recurrence}.

\begin{proposition}\label{prop:twelve-four-four-dyadic}
Suppose $T_{\mathbb Q}=K^2\times\mathbb Q^2$, with $K$ of discriminant
$c=8$ or $12$. Suppose the two rational characters have codegree four
and different supports. At two the order has two local factors of
rank three, each containing one rational and one quadratic branch.
For either quadratic branch write $f=mv^eu^k$ for its oriented
codegree and $\mathbb Z_2+2^b\mathcal O_{K,2}$ for its projected order.
Then
\[
\begin{array}{c|c|c|c}
 c & e & v_2(m)& b\\ \hline
 8&1&3&1\\
 12&0&3\ (k\text{ even}),\ 4\ (k\text{ odd})&v_2(m)-2 .
\end{array}
\]
At conductor twelve, $3\mid m$ as well.
\end{proposition}
\begin{proof}
Each rational row has exactly four entries $\pm1$, counting its
unit entry. Different supports distinguish the residue characters.
Neither rational branch can form a local factor by itself, since
its idempotent has unit coefficient $1/4$. Both quadratic branches
are ramified, so each rational branch takes precisely one.

Write $A$ for one local factor and $B$ for its quadratic projection.
Perfectness of the coefficient-of-unit pairing gives
$A\cap(\mathbb Q_2\times0)=4\mathbb Z_2\times0$:
the rational projection is $\mathbb Z_2$ and its trace weight is $1/4$.
Thus $[\mathbb Z_2\times B:A]=4$. Computing the regular discriminant
first by index and then by formal codegrees gives
\[
 4+v_2(c)+2b=2+v_2(N(f)).
\]
At conductor eight this forces $e=1$, at conductor twelve $e=0$,
and in both cases $v_2(m)=2+b$. The local identity has integral
unit coefficient $1/4+B_{e,k}/m$, so $b=v_2(B_{e,k})$.
The trace recurrence gives the table. At conductor twelve the
quadratic row-difference test also gives $3\mid m$ as in
Theorem~\ref{thm:twelve-3-12}.
\end{proof}

\begin{lemma}\label{lem:two-four-supports}
In Proposition~\ref{prop:twelve-four-four-dyadic}, let $t,w$ be
the reciprocal masses of the signed quadratic branches. A nonunit
invertible forces $t=w=1/4$. In the ordinary ring every integral
character has norm divisible by four, and no orbit has odd common
rational codegree. If the unit is the only invertible, there are
at most two ordinary rational four-rows.
\end{lemma}
\begin{proof}
The signed four-rows intersect in precisely two coordinates: their
intersection contains the unit, has even size by orthogonality, and
has size less than four. At the other coordinate their signs are
opposite. An invertible must belong to both supports. All signed
characters are real, so at an invertible each quadratic pair has
a common sign. The zero unit coefficient of a nonunit invertible
now gives $\pm t\pm w=0$, and $t+w=1/2$.

The parity isomorphism identifies the two residue characters and
the ranks of their local factors in the ordinary and signed orders.
Every integral ordinary row therefore has four odd entries. An orbit of odd rational codegree is an isolated maximal
unramified factor by Lemma~\ref{lem:trace-order-conductor}. Its residue fields are all $\mathbb F_2$, so it would
have a local factor of rank one. Finally two rational four-rows
with equal support give a nonunit invertible by
$1-2(e_1+e_2)$, as in Theorem~\ref{thm:twelve-3-8}. There are only
two residue characters, hence only two possible supports.
\end{proof}

\begin{theorem}\label{thm:four-four-distinct}
Under the hypotheses of Proposition~\ref{prop:twelve-four-four-dyadic},
suppose the ordinary Frobenius--Perron field is also
$K=\mathbb Q(\sqrt3)$ and the ordinary character algebra is
commutative. Then the category is spherical.
\end{theorem}
\begin{proof}
Use $U,V,A_k,B_k,t,w,s$ as in Theorem~\ref{thm:twelve-3-12}.
Suppose the category is nonspherical and choose $E\ge E'$. Write
$E=mU^k$ and $H=nU^j$ for the quadratic signed codegrees by
Proposition~\ref{prop:twelve-four-four-dyadic}. We have $t+w=1/2$.
If $H$ is rational the positivity argument on the ordinary
Frobenius--Perron induction label in Theorem~\ref{thm:twelve-3-12}
gives $s\le t$, contrary to strict nonspherical dominance. Thus
$j\ge1$. The integrality of $E/D$ and $D>E$ give
$s\ge2t$, as in that theorem. If $j=1$, then $48\mid n$ by the
dyadic proposition. Hence $w\le1/12$, $s\ge5/6$ and
$2\delta_D^2\ge(3/2)s^2>1/3$, contradicting the mixed projection
bound. Therefore $j\ge2$. Central dominance
$D/D'\ge H/H'$ and the same unit computations as in
Theorem~\ref{thm:twelve-3-12} give
\[
 s\ge7t/2,\qquad \delta_H^2\ge48w^2/49,\qquad
 \delta_D^2\ge48s^2/49.
\]
Put $F(s)=(48/49)(1/2-2s/7)^2+(96/49)s^2$. Then the mixed
bound gives $F(s)\le1/3$. This quadratic is increasing on
$[7/24,1]$ and $F(7/24)>1/3$, hence $s<7/24$.

By Lemma~\ref{lem:two-four-supports}, a nonunit invertible would give
$t=w=1/4$, contradicting $s\ge7t/2$ and $s<7/24$. Hence there are
at most two rational four-rows. Integral rows have norm divisible by
four, and odd common rational codegrees are excluded.

Lemma~\ref{lem:real-four-six} also excludes a real quadratic
orbit of common codegree four or six.

Consider an irrational ordinary complementary pair of mass $z$.
Every irrational codegree lies in $K$ and its squared gap is at
least $z^2/3$. If there are two such pairs counting multiplicity,
their combined mixed contribution is at least $(1-s)^2/3$, but
$F(s)+(1-s)^2/3\ge8268/17089>1/3$. With just one pair, the
remaining two rows of rational codegree have reciprocal mass at most
$3/8$, unless both are rational four-rows. Indeed an integral row
has norm divisible by four and a quadratic orbit of common norm
has even codegree at least eight. Without two four-rows we get
$F(s)+(2/3)(5/8-s)^2\ge7617/19490>1/3$. In the two-four-row case
the two distinct dyadic local factors each contain one rational
row and one ramified quadratic branch $K$. In the factor containing
the irrational complementary pair, the regular trace discriminant
has valuation $2+v_2(N(h))$, where $h$ is its oriented codegree.
The maximal order has discriminant valuation two, so $v_2(N(h))$
is even. Thus $h$ belongs to the $U$-family, for which the squared
gap is at least $3z^2/4$. Since $z=1/2-s$, we obtain the final
contradiction $F(s)+(3/2)(1/2-s)^2\ge204/515>1/3$.
All these displayed lower bounds follow by completing the square.

The four ordinary complementary codegrees are now all rational,
even, with reciprocal sum $1-s>17/24$. At least one is four,
since four rows of codegree at least six contribute at most $2/3$.
Every four-row is rational by Lemma~\ref{lem:real-four-six}.

If there were just one four-row, the mass would require a six-row.
It would be integral by Lemma~\ref{lem:real-four-six}, contradicting
divisibility of its norm by four.

There are consequently two rational four-rows. If the remaining
two rows are rational their codegrees are multiples of four at
least eight, and their reciprocal sum exceeds $5/24$. This forces
both codegrees to be eight. The product of the codegrees has
squareclass $N(D)$, either one or six, whereas the character algebra
$K\times\mathbb Q^4$ has discriminant squareclass three. If the
remaining rows form a quadratic orbit their common even codegree
is at least eight and the same reciprocal sum forces eight. Its
quadratic discriminant $c$ divides $16$ and satisfies $2c\le16$,
so the field is $\mathbb Q(\sqrt2)$. Each of the two rational
four-rows lies in a different rank-three dyadic factor. The factor
containing $\mathbb Q_2(\sqrt2)$ has regular discriminant valuation
$v_2(4\cdot8^2)=8$, while its maximal order has valuation three.
These valuations must differ by twice an index valuation. This
excludes the last possibility.
\end{proof}

\subsection{An involution fixing a near-group subcategory}
\begin{proposition}\label{prop:folded-e6-descent}
Let $\mathcal C$ be a rank-six complex fusion category whose
sphericalization has commutative Grothendieck ring. Suppose all
ordinary Frobenius--Perron dimensions lie in $K=\mathbb Q(\sqrt3)$.
Suppose a nonidentity invertible $g$ is an involution with action
$(1,g)(X,Y)$ on the simples and fixed points $P,Q$.
If $1,g,P$ form a fusion subcategory with $P^2=1+g+2P$,
then $\mathcal C$ is spherical.
\end{proposition}
\begin{proof}
Write $p=1+\sqrt3$, $x=\operatorname{FPdim}(X)
=\operatorname{FPdim}(Y)$ and $q=\operatorname{FPdim}(Q)$.
The span of $X,Y,Q$ is a module for $1,g,P$ by reciprocity.
Since $P=P^*$ and $gP=P$, the matrix of $P$ on this span is
\[
 \begin{pmatrix}a&a&c\\a&a&c\\c&c&d\end{pmatrix}
 \quad (a,c,d\in\mathbb Z_{\ge0}).
\]
On the symmetric subspace its eigenvalues solve
$z^2-2z-2=0$. This quadratic is irreducible over $\mathbb Q$, so
the trace and determinant there give $2a+d=2$ and $2ad-2c^2=-2$.
Thus $c=1$ and $(a,d)$ equals $(0,2)$ or $(1,0)$.

In the first case $PX=PY=Q$, so $q=px$ and $Q$ is self-dual:
$X^*$ belongs to $\{X,Y\}$, since duality permutes the
$g$-orbits and fixes $P$ and the unit. The multiplicities of
$1,g,P$ in $XX^*$ are $1,0,0$ by reciprocity. Therefore
$x^2=1+(r+sp)x$ for nonnegative integers $r,s$. This makes $x$
an algebraic unit. No element of $\mathbb Z[\sqrt3]$ has norm
$-1$, modulo three. Thus $x-1/x=x-x'$ has rational trace zero.
The nonnegative coefficients of $r+sp$ force $r=s=0$ and $x=1$.
The four objects $1,g,X,Y$ are precisely the invertibles.

In the second case $PQ=X+Y$, so $x=pq/2$. Here again $Q=Q^*$.
The multiplicities of $1,g,P$ in $Q^2$ are $1,1,0$. The equality
$gQ=Q$ gives equal multiplicities on $X,Y$. It follows that
\[
 q-2/q=h p+k,\qquad h,k\in\mathbb Z_{\ge0}.
\]
In particular $q$ divides two in $\mathbb Z[\sqrt3]$, and
$q>1$ since $Q$ has nontrivial stabilizer. The ramification
$(2)=(1+\sqrt3)^2$ as ideals and the positive unit group
$\langle U\rangle$, $U=2+\sqrt3$, give
$q=U^n,(1+\sqrt3)U^n$ or $2U^n$, with $n\ge0$. For the
unit option $q-2/q$ has negative trace. For the middle option
$N(q)=-2$, so $q-2/q=q+q'$ is rational and $h=0$.
For the last option write $U^n=A+B\sqrt3$. Then
$h=3B$ and $k=A-3B$. For $n\ge1$ we have $B\ge1$ and
$A^2=3B^2+1<9B^2$, which is impossible. At $n=0$ again $h=0$.
Therefore $Q^2=1+g+kQ$ and $1,g,Q$ form a fusion subcategory.

It remains to lift the category. The subcategory on $1,g,P$ has
a spherical tensor section of the sphericalization by the corank-one
sphericity theorem \cite[Theorem~4.5]{DCW}.
The second subcategory, on $1,g,X,Y$ in the first case and
$1,g,Q$ in the second case, likewise has a spherical tensor
section, positive on its invertible objects in the pointed case.
All these sections choose the lift of $g$ of dimension one:
in the rank-three subcategories it fixes a simple of nonzero
spherical dimension. These two subcategories generate $\mathcal C$, by $PX=Q$ or
$PQ=X+Y$, and their lifted images have the same intersection
$\langle g\rangle$. For commuting Grothendieck subrings the regular
Perron sums satisfy
\[
 R_{\mathcal A} R_{\mathcal B}
  =\operatorname{FPdim}(\mathcal A\cap\mathcal B)\,
    R_{\mathcal A\vee\mathcal B}.
\]
Indeed their product is a common Perron vector for multiplication
by the simples of $\mathcal A,\mathcal B$, and the unit coefficient
determines its scalar. Apply this identity downstairs and upstairs,
where commutativity is a hypothesis. The lifted join has the same
Perron dimension as $\mathcal C$ and maps dominantly to it, hence
equivalently. Its spherical structure descends.
\end{proof}

\subsection{The folded codegree-four pairs at conductor twelve}
\begin{lemma}\label{lem:folded-twelve-mass}
Suppose the ordinary and signed character algebras are commutative
of rank six, their respective Frobenius--Perron and spherical
dimension fields are $K=\mathbb Q(\sqrt3)$, and the signed
rational algebra has two $K$ factors. Assume the other two
signed characters have common codegree four and are either rational
with the same support, or Gaussian conjugates. If no Galois
conjugate is pseudo-unitary, write $t=\operatorname{Tr}_K(1/E)$ and
$s=\operatorname{Tr}_K(1/D)$ and orient $E\ge E'$. Then
$t>1/128$. Moreover, either $s>2/9$ or the ordinary complementary
irrational pairs contribute at least $(2/3)(1/2-s)^2$ to the
mixed squared-gap bound.
\end{lemma}
\begin{proof}
Write $w=1/2-t$ for the other signed mass, of codegree $H$.
The two short signed rows can be written
$(1,-1,\pm z,\mp z,0,0)$ on $1,g,X,Y,P,Q$, with $z=1$ or $i$.
For rational rows with equal supports this follows by adding and
subtracting their two orthogonality relations with the positive row.
For the Gaussian pair, $X,Y$ are the unique dual pair, and norm
four and positivity force the stated row. In either case the four
real signed rows take value one at $g$ and agree at $X,Y$. Thus
$g^2=1$ in the signed ring. Its multiplication matrix is integral
orthogonal and preserves the strictly positive spherical dimension
vector, so it is a permutation matrix. Its minus-one space has
dimension two and contains the two differences $1-g,X-Y$.
Consequently its permutation is $(1,g)(X,Y)$.

Equation~\eqref{eq:quadratic-trace-order} gives $s\ge2t$.
The complementary four-rows contribute at least $1/2$ on the
ordinary Perron column. Thus $H$ is irrational by
Corollary~\ref{cor:galois-columns}; and $H/H'\ge U^2$ by
Lemma~\ref{lem:quadratic-mass-gap}, where $U=2+\sqrt3$ and
$V=3+\sqrt3$. Central dominance gives $D/D'\ge H/H'$,

so $\delta_D^2\ge3s^2/4$ and
$s<1/2$ by the mixed inequality. If $X$ or $Y$ were invertible
its spherical value would be one, its $H$-values equal signs,
and the four-row pair would contribute zero to its unit coefficient.
Then $t=w=1/4$, contradicting $s<1/2$. Every nonunit invertible
must have nonzero value on the short rows, so $g$ is the only
possible one.

At two, every signed branch is ramified or rational, with residue
field $\mathbb F_2$. The two short rational rows have the same
reduction, so their rational branches share a dyadic local
factor. Their joint projected order consists of pairs congruent
modulo two: the unit and the value pair $(1,-1)$ at $X$ generate
this order. They cannot form an isolated local factor, whose
regular discriminant would have valuation $v_2(4^2)=4$ by the
perfect codegree trace, whereas their projected order has
discriminant valuation two. In the Gaussian case the projection
is the full $\mathbb Z_2[i]$, again of discriminant valuation
two and so not isolated. There are therefore at most two dyadic
local factors, all of rank at least two and residue degree one.
This carries over to the ordinary order via the parity of the
structure constants and lifting idempotents. Lemma~\ref{lem:trace-order-conductor} excludes odd rational
codegrees: their unramified maximal factors with residue field
$\mathbb F_2$ would have rank one.

All rational-codegree ordinary rows have norm at least four. In the
absence of nonunit invertibles, a rational two-row is impossible
by the idempotent formula, and an orbit of common codegree two
cannot be nonrational (each nonunit coordinate is a totally real
or complex algebraic integer of house at most one, hence zero
or a root of unity; there is room for just one nonzero entry,
and duality and orthogonality to the FP row force it to be $-1$).
If $g$ is invertible its ordinary action agrees with the signed
permutation on positive lifts. The two characters nontrivial at $g$
have rows $(1,-1,\lambda,-\lambda,0,0)$. In the Gaussian case
$Y=X^*$, so $\lambda$ is nonzero purely imaginary. In the real
case all simples are self-dual, so $\lambda$ is real. Evaluating
$X^2$ gives $\lambda^2=\pm1+n\lambda$ with integral $n$, the
constant being minus one in the Gaussian case and plus one in the
real case by reciprocity. The first case forces $n=0,\lambda=\pm i$.
In the second, if $n\ne0$ the roots are quadratic conjugates
of norm $-1$ and the codegrees $4+2n\lambda$ lie in $K$ by
central induction, which is impossible modulo three. Thus
$\lambda=\pm1$. Every remaining ordinary row takes value one
on $g$ and cannot be a two-row by orthogonality to FP.

If there is one irrational ordinary pair of mass $z$, the other
two rational-codegree rows contribute at most $1/2$. Hence
$z\ge1/2-s$ and the extra mixed contribution is at least
$(2/3)(1/2-s)^2$. In the case of two irrational pairs (including
multiplicities) their combined contribution is at least
$(1-s)^2/3\ge(2/3)(1/2-s)^2$. Suppose henceforth that all
complementary codegrees are rational and $s\le2/9$. Their
reciprocal sum is at least $7/9$, so at least two rows have
codegree four. Real four-rows and six-rows are rational by Lemma~\ref{lem:real-four-six}.

Any nonreal ordinary row has imaginary quadratic character field $F$: there is only one nonreal pair, and character fields are abelian. A nonreal pair of common codegree $h$ has
$h=|\operatorname{disc}(F)|b^2$, $b\in\mathbb Z_{>0}$, by
orthogonality to its conjugate, since $X,Y$ are the unique
dual pair. Hence it cannot have codegree six. At codegree four
its field is Gaussian and its row has exactly one nonunit real
entry, equal to $-1$ on an invertible by orthogonality to FP.
This forces $g$ to be an actual invertible by the preceding
discussion.

In the absence of a nonunit invertible there are therefore at
least two integral ordinary four-rows. They have different
supports, since the usual idempotent $1-2(e_1+e_2)$ from two
four-rows with equal support would be a nonunit invertible.
Their distinct supports exhaust the possible dyadic residue
rows, so all integral ordinary rows have norm divisible by
four, and there are at most two ordinary four-rows. The remaining
rows have even norms at least eight: the only other possibility,
a nonrational pair of common codegree six, was just excluded.
Their total mass is at most $3/4$, a contradiction.

It remains to consider actual $g$. The two $g$-negative rows are
the matched four-rows above. The other ordinary rows are all real
and have shape $(1,1,z,z,b,c)$. If neither remaining row has
codegree four, the reciprocal sum requires $(6,6)$ or $(6,8)$
for their two codegrees. Both are integral rows by the quadratic
orbit exclusion at six. At $(6,6)$ the product of the ordinary
codegrees has absolute squareclass $N(D)$, one or six. The
maximal order has absolute discriminant squareclass three since
the matched anti-invariant algebra is rational split or Gaussian.
At $(6,8)$ one would have $s=5/24$. By Lemma~\ref{lem:fixed-quadratic-mass}, $U^3\equiv1\pmod5$ excludes this numerator, since
no power $D/D'$ of $U$ can equal $-1$ modulo five as
$1+D/D'=Ds$ would require.

If there is a third four-row it has shape
$(1,1,0,0,-1,1)$ after relabelling. The alternative supported at
$X,Y$ would force these objects invertible, and the two tail
signs cannot be equal since both fixed objects have dimension
greater than one. Thus $d_P=d_Q+2$. The last row must be integral
and has shape $(1,1,z,z,b+2,b)$ by orthogonality. Its even
codegree $h\le36$ has squareclass two or three by the same
discriminant comparison, hence $h=8,12,18$ or $32$.
Since $h=4+2z^2+2(b+1)^2$, the last two would represent seven
or fourteen as a sum of two squares. At $h=12$, either
$z=\pm2,b=-1$ forces $d_X=1$, or $z=0,b=1,-3$ contradicts
$d_P=d_Q+2$ and $d_Q>1$. At $h=8$ the two possibilities give
the FP rows $(1,1,x+1,x+1,x,x-2)$ and
$(1,1,x-1,x-1,x+2,x)$. Orthogonality with the conjugate yields
$N(x)=-2$, and $D=4(x^2+2)=4x(x-x')$. At the fixed simple of
dimension $x$ the coefficient of its own simple in its square is
$\operatorname{Tr}(x^3/D)\pm3/4=(\operatorname{Tr}(x)\pm3)/4$.
This contradicts the even trace in $\mathbb Z[\sqrt3]$. Thus
$s>2/9$ whenever the ordinary complement has only rational codegrees.

We now bound $w$. If $2j+e=2$, the mixed inequality excludes
$n=6$ by $\delta_H^2=12/n^2$ and the unit-family multiple-of-six
condition, so $w=4/n\le1/3$. For $2j+e=3$ we have
$\delta_H^2=25/(3n^2)$ and $2\mid n$ by the even irrational
coefficient of a sum of squares in $K$. Since $w=3/n<1/2$,
$n\ge8$ and $w\le3/8$. For exponent at least four, both $H$
and $D$ have squared gap at least $48/49$ times their squared
reciprocal mass. If $s>2/9$, the mixed inequality gives
$w^2<49/144-2(2/9)^2<(63/128)^2$. With an extra irrational
ordinary contribution it gives
$\delta_H^2\le1/3-\min_s[(3/2)s^2+(2/3)(1/2-s)^2]=17/78$.
Thus again $w^2\le(49/48)(17/78)<(63/128)^2$.
In every case $t=1/2-w>1/128$.
\end{proof}

\begin{theorem}\label{thm:folded-twelve-completion}
Suppose a rank-six complex fusion category has commutative ordinary and signed character algebras, and both dimension characters have field $K=\mathbb Q(\sqrt3)$. Suppose its signed rational character algebra is $K^2\times F$, with the dimension character in the first quadratic factor. Each of the two complementary characters of $F$ has formal codegree four. If $F=\mathbb Q(i)$, or if $F=\mathbb Q^2$ and its characters have the same support, then the category admits a spherical structure.
\end{theorem}
\begin{proof}
Suppose otherwise. Use $U=2+\sqrt3$, $V=3+\sqrt3$, and $E,H,D,t,w,s$ as in Theorem~\ref{thm:twelve-3-12}, with $E\ge E'$ and $H\ge H'$. We have $t+w=1/2$, $D>E$, $s\ge2t$, and $D/D'\ge H/H'$. The inequality $s\ge2t$ also follows from the trace rearrangement since $\operatorname{Tr}(E/D)\ge4$ for a totally positive nonrational integer of $K$ with least conjugate less than one. The other signed $K$-factor has irrational codegree: if $H$ were rational, all complementary signed orbits would contribute at least their reciprocal mass on the ordinary FP induction column, contradicting $s>t$.

By Lemma~\ref{lem:folded-twelve-mass} the only possible nonunit invertible is the second coordinate $g$ in the short rows, whose action, if it is invertible, has two free orbits and two fixed simples.

We may take $D/D'>H/H'$. For if equality holds, the real signed $H$-character attains the ordinary Perron bound on the central conjugation witness, whose FP dimension $D/D'$ is irrational. Its equality subcategory is spherical. If proper it has field $K$ (it contains that witness) and rank
at most four by Lemma~\ref{lem:field-corank}. The pivotal
classifications~\cite{Ostrik2,Ostrik3,DPP} leave the rank-three $E_6$ ring, the $C_3$ Tambara--Yamagami ring and the rank-four $C_3$ near-group ring of parameter six. Our invertible-object restriction excludes both $C_3$ cases. The $E_6$ case requires the actual involution $g$ and implies sphericality by Proposition~\ref{prop:folded-e6-descent}. If the equality subcategory is the whole category the spherical section already exists.

Write $P_{e,i}=V^eU^i$ and $B_{e,i}=\operatorname{Tr}(1/P_{e,i})$, with $e\in\{0,1\}$. The $d$-number and row-orthogonality calculations from Theorem~\ref{thm:twelve-3-12} give unit-family scalars divisible by six and $V$-family scalars divisible by two (the coefficient of $\sqrt3$ in a sum of squares is even, and that in $VU^i$ is odd). The absolute order discriminant is $16 N(EH)$; the absolute maximal-order discriminant is a square, also for $F=\mathbb Q(i)$. Hence $E,H$ are in the same family. Put
\[
 E=mP_{e,k},\qquad H=nP_{e,j},\qquad D=MP_{d,\ell},\qquad h=m/n,\quad c=m/M .
\]
The quantities $h,c$ are positive integers, and $6\mid c$ if $d>e$, by integrality of $E/H,E/D$. Moreover $4\mid m$ by $E/4\in\mathcal O_K$. At the FP induction label the remaining two signed codegrees give $Z/4$ with $Z=2$ or $3$, by integral positivity or by taking the trace of a totally positive pair. Its $H$-mass is positive, strictly smaller than $w$, and belongs to $\mathbb Z/n$, since $\operatorname{Tr}((a+b\sqrt3)/V)=a-b$. In particular $H\ne nV$ and $j\ge1$, so $D/D'\ge H/H'\ge U^2$. 

\smallskip\emph{A uniform recurrence calculation.}
We use the preceding bound $t\ge1/128$ and the alternatives $s>2/9$ or an additional ordinary contribution at least $(2/3)(1/2-s)^2$ in the mixed Fourier inequality. Set $r=\ell-k$ and $z=j-k$, and abbreviate $b_0=B_{e,k}$, $b_1=B_{e,j}$, $b_2=B_{d,\ell}$. The two column mass identities give
\begin{equation}\label{eq:folded-twelve-h}
 m=2(b_0+h b_1),\qquad m(4-Z)-4cb_2\in4h\mathbb Z_{>0}.
\end{equation}
In particular $h\mid4b_0,4cb_2$. The ordinary Fourier bound and
$D/D'\ge U^2$ give $s\le\sqrt2/3$, hence $t\le s/2<1/4$ and
$s/t<61$. Lemma~\ref{lem:relative-quadratic-indices} gives $c\le27$,
$1/V\le U^r<12\cdot61V$, so $-1\le r\le6$. The initial
recurrence pairs give $0<|W|\le9360$ on the allowed range, with
$h\mid4cW$. Finally $1<hb_1/b_0=1/(2t)-1<64$ and
$U^z/2\le b_1/b_0\le2U^z$ give $-14\le z\le3$.

The \codecheck{folded_twelve_check.py}{relative-index check}
tests these tuples with $6\mid c$ if $d>e$, $V^{d-e}U^r>c$ and
$2z+e<2r+d$. At $0\le k<28$, no tuple with $j,\ell\ge1$ satisfies
\eqref{eq:folded-twelve-h}, the scalar divisibilities, $s\ge2t$
and both Fourier conditions
\[
 G=\delta_E^2+\delta_H^2+2\delta_D^2\le1/3,\qquad
 s>2/9\quad\text{or}\quad G+\tfrac23(1/2-s)^2\le1/3.
\]
For $k\ge28$ take $I=14$ and $\eta=999999/1000000$ in
Lemma~\ref{lem:relative-quadratic-indices}. The resulting bounds
for the gaps and $s$, with $s<1/2$, leave
\[
 (e,d,r,c,h,z)=(0,0,1,1,6,0),\ (0,0,1,1,24,-1),\
 (0,0,1,3,24,0),\ (1,1,1,3,24,0).
\]
Modulo $\operatorname{lcm}(4h,6h,6c)$ their periods are
$12,24,24,144$. None satisfies \eqref{eq:folded-twelve-h} and the
scalar divisibilities.
\end{proof}

\section{The quadratic field of conductor eight}
\label{sec:quadratic-eight-cases}
When every complementary signed codegree is irrational, the two
projection estimates again force one orbit of reciprocal mass $1/2$.
When the complement has rational codegrees, the same estimates
leave a pair of codegree four. Its two rows determine the induction
objects needed for the final descent.

\subsection{Three quadratic character pairs}

Let $K=\mathbb Q(\sqrt2)$, $u=3+2\sqrt2$ and $v=2+\sqrt2$.

\begin{lemma}[Three ordinary four-rows]\label{lem:silver-three-four-rows}
Suppose the ordinary algebra is commutative with Perron field $K$
and the signed rational algebra is $K^3$. If its four ordinary
complementary codegrees are even and at least four, then
\[
 s=\operatorname{Tr}(1/D)<21/100
 \quad\Longrightarrow\quad
 \text{three ordinary codegrees are four}.
\]
They force an actual involution fixing precisely two simples, and a
signed $v$-family pair of reciprocal mass $1/2$. If $s<1/6$, the
ordinary codegrees are $4,4,4,8$ and $D=8A_\ell u^\ell$, where
$A_\ell=u^\ell+u^{-\ell}$.
\end{lemma}
\begin{proof}
Sorting the denominators with reciprocal sum $1-s>79/100$ leaves
$(4,4,4,d)$, $(4,4,6,6)$ or $(4,4,6,8)$, with largest denominator at most $24$.
The signed algebra is totally real, so all simples are self-dual.
All four rows are rational by Lemma~\ref{lem:real-four-six} and
Galois invariance. Their product times $N(D)$ has squareclass
two, with $N(D)$ a square or twice a square. This excludes
$(4,4,6,8)$. At $(4,4,6,6)$ the two six-rows are distinct
modulo three by squared distance twelve. Each needs a partner,
whereas the four-axes split off and three is inert in $K$.
Finally at $s<1/6$ we get $4<d<12$, and the same squareclass
test gives $d=8$, $N(D)$ square and $s=1/8$.

Each rational codegree-four row has four entries $\pm1$, including
the unit. By orthogonality the two-element zero sets of two such
rows coincide or are disjoint. There are three such sets on five
nonunit coordinates, so two coincide. Their sum has support the
unit and a nonunit $g$, with coefficients $2,-2$ by orthogonality
to the Perron row. Thus $\operatorname{FPdim}(g)=1$ and
$1-g$ is twice the sum of the two primitive idempotents. Hence $g$
is an involution with a two-dimensional minus-eigenspace. Consequently $g$ swaps $1,g$ and $X,gX$ for
one additional simple $X$, and fixes the other two. The same is true
on the positive signed basis. The two signed characters with $g=-1$
therefore form a full $K$-orbit and have rows
$(1,-1,q,-q,0,0)$. On the two-dimensional space with basis
$1-g,X-gX$, signed multiplication by $X$ has matrix
$\left(\begin{smallmatrix}0&1\\1&d_0\end{smallmatrix}\right)$ by
self-duality and symmetry. Thus $qq'=-1$; $q$ cannot be rational
since the characters are distinct and their values on $X$ determine
them. The quotient of their codegrees $2(1+q^2)$ is $q^2$ up to
inversion, an odd power of $u$, since the norm-negative units of
$K$ are the signed odd powers of $1+\sqrt2$. This orbit is therefore
in the $v$-family and has reciprocal mass $1/2$.
\end{proof}

\begin{theorem}\label{thm:silver-finite-orbit}
Suppose $\mathcal C$ has rank six, its ordinary character algebra
is commutative, its signed rational algebra is $K^3$, and the
signed dimension and ordinary Frobenius--Perron characters both have field
$K$. Suppose both complementary signed codegrees are irrational.
Then $\mathcal C$ is spherical, or its two complementary orbits
can be represented by
\[
 H=8+4\sqrt2,\qquad G=nv^e u^j,\qquad
 e\in\{0,1\},\quad j\ge2,\quad n\in\mathbb Z_{>0}.
\]
In the latter case $s=\operatorname{Tr}(1/D)<1/3$.
\end{theorem}
\begin{proof}
Argue under the assumption that $\mathcal C$ is nonspherical and
orient $E\ge E'$ by Galois transport. As in the preceding section
write $t=\operatorname{Tr}(1/E)$, $s=\operatorname{Tr}(1/D)$, and
$\delta_f=1/f'-1/f$. The orbit masses are $t,w,z$ on the signed
side, with $t+w+z=1$. All codegrees lie in $K$ by the dimension
formula $E/f=\dim F(A_f)$. For a conjugate ordinary pair with
irrational codegrees $L,L'$ include one term $\delta_L^2$, counting
repeated pairs separately. The modular projection inequality is
\begin{equation}\label{eq:silver-mixed}
 \sum_{\rm signed\ pairs}\delta_f^2+
   2\sum_{\rm ordinary\ pairs}\delta_f^2\le1/2.
\end{equation}
This is \eqref{eq:quadratic-fourier-bounds} at discriminant eight. Lemma~\ref{lem:central-witness} gives $D/D'\ge f/f'$ by
spectral domination. Equation~\eqref{eq:quadratic-trace-order} gives $s\ge2t$.

Write $f=mv^eu^k$ as in Lemma~\ref{lem:small-quadratic-recurrence},
and put $A_k=B_{0,k}$, $B_k=B_{1,k}$. The squared distance between
a quadratic row and its conjugate gives $8\mid\operatorname{Tr}(f)$.
Since $\operatorname{Tr}(u^k)\equiv2\pmod4$ and
$\operatorname{Tr}(vu^k)\equiv4\pmod8$, we have $4\mid m$ for
$e=0$ and $2\mid m$ for $e=1$. Both reciprocal trace sequences
are $2\pmod4$.

Suppose first that both complementary signed exponents are at least
two. Their squared gap ratios $\delta_f^2/\operatorname{Tr}(1/f)^2$
and that of $D$ are then at least $c=288/289$, since $f/f'\ge u^4$.
We also have $s/t\ge10/3$. In the same primitive family the
ordinary exponent increases by at least one; the successive trace
ratios are at least $5$ except $A_1/A_0=3$, which is excluded by
$D/D'\ge u^4$. From unit-family $E=mu^k$ to $D=Mvu^\ell$
the integer $m/M$ is even and $\ell\ge k$, giving
$s/t\ge2B_k/A_k\ge10/3$ for $k\ge1$. From $E=mvu^k$ to
$D=Mu^\ell$ one has $\ell\ge k+1$ and
$s/t\ge A_{k+1}/B_k\ge17/5$ for $k\ge1$.
In either cross-family case $k=0$ and
$D/D'\ge u^4$ give the same asserted bound. Hence, if $s\ge1/6$,
the left side of \eqref{eq:silver-mixed} is at least
\[
 c\bigl((1-3s/10)^2/2+2s^2\bigr)>1/2.
\]
The quadratic is increasing for $s\ge1/6$ and exceeds $1/2$
already at that endpoint. Thus $s<1/6$ and $t<1/20$.

Every rational ordinary codegree is even and at least four.
Indeed the completed signed order at two has residue fields
$\mathbb F_2$ and has no local factor of rank one, because its
normalization has three ramified quadratic branches. Both facts
transfer by $R/2R\simeq T/2T$, lifting idempotents over
$\mathbb Z_2$. By Lemma~\ref{lem:trace-order-conductor}, an orbit of
odd rational codegree splits off maximally and unramified at two.
Its residue fields have degree one, which would supply a rank-one factor.

Codegree two on at least
two characters exhausts the reciprocal mass. A single character of
codegree two has row $(1,-1,0,\ldots,0)$ by orthogonality to
positive dimensions; call the second basis element $g$. The
idempotent $(1-g)/2$ is primitive, implying $g^2=1$ and a
one-dimensional minus-eigenspace for multiplication by $g$.
The ordinary and positive signed lifts of an actual invertible
act by the same permutation, since the positive signed dimension
row forces all its nonzero signs. This would give a minus-eigenspace
of odd dimension in $K^3$, which is impossible. Here and below all
simples are self-dual since $T_{\mathbb Q}$ is totally real.

If exactly two of the four remaining ordinary codegrees are
irrational, their pair has mass at least $1/2-s>1/3$. Its
contribution to the left side of \eqref{eq:silver-mixed} is at
least $(1/2-s)^2$, since the minimal quadratic squared gap ratio
is $1/2$. With four irrational codegrees their contribution is
at least $(1-s)^2/2$. Either contribution exceeds $1/9$, whereas
the two signed complementary orbits already supply
$c(1-t)^2/2>c(19/20)^2/2$. This contradicts \eqref{eq:silver-mixed}.

Lemma~\ref{lem:silver-three-four-rows} now gives $s=1/8$,
$D=Mu^\ell$ with $M=8A_\ell$, and a signed $v$-family orbit of
mass $1/2$.

If $s/t\ge5$, the signed complementary orbits and $D$ already
contribute at least
$c((1-1/40)^2/2+2/64)>1/2$. Thus $s/t<5$.
The preceding exponent comparisons with unit-family $D$ force
$E=mvu^k$ with $m=M$, $\ell=k+1$: a larger scalar quotient
or exponent gap gives $s/t\ge 6$.
The signed-order discriminant squareclass requires an odd number
of $v$-families among the three signed factors. Since both $E$
and the half-mass orbit are of that type and $t<1/20$, the remaining
factor is $G=nvu^p$ as well. Set $q_0=m/n\in\mathbb Z_{>0}$. Its
mass and the half-mass identity give
\[
 q_0B_p=4A_{k+1}-B_k,\qquad q_0\mid m=8A_{k+1}.
\]
Since $B_p,B_k\equiv2\pmod4$, $q_0$ is odd and divides
$\gcd(A_{k+1},B_k)$. This gcd divides
$|B_kA_{k+2}-B_{k+1}A_{k+1}|=|2\cdot34-10\cdot6|=8$.
Thus $q_0=1$ and $B_p>B_k$, so $p>k$. On the other hand
$D/D'\ge G/G'$ gives $2\ell\ge2p+1$, equivalently $p\le k$.
This contradiction proves that a signed complementary exponent is
at most one.

Choose such an orbit $H$ with mass $w$, giving preference to $nv$
if it occurs. The two complementary orbits cannot both have this
form: their contributions at the ordinary FP induction column
are at least their reciprocal masses since for integral totally
positive $b$ one has $\operatorname{Tr}(b/v)\in2\mathbb Z_{>0}$.
This contradicts $s>t$. The other orbit $G$ thus satisfies
$G/G'\ge u^2$. The same holds for $D$, giving squared gap ratios
at least $8/9$. Dropping all other terms from \eqref{eq:silver-mixed}
yields
\[
 \tfrac12\ge\tfrac89\bigl((1-w-t)^2+8t^2\bigr)
 \ge\tfrac{64}{81}(1-w)^2 .
\]
Hence $w>1/5$. An exponent comparison improves $s\ge2t$ to
$s\ge3t$: with $D/D'\ge u^2$ the only two additional exponent-zero
exceptions to the comparisons above are $A_1/A_0=A_1/B_0=3$.
Retaining the gap of $H$ yields
\[
 \delta_H^2+\tfrac{16}{19}(1-w)^2\le1/2
\]
after minimizing over $t$. The possible forms of $H$ are $nv,nu,nvu$,
with $w=2/n,6/n,10/n$, respectively. Applying the scalar
divisibilities, $w>1/5$, and this last inequality gives
\[
\begin{array}{c|l}
 H & n\\ \hline
 nv&4,6\\
 nu&12,16,20\\
 nvu&18,20,22,24,26,28,30,32 .
\end{array}
\]

If both complementary exponents are at most one, our choice of $H$
leaves $G=nu$ or $nvu$. Since $w\le5/9$ and
$t\le s/3\le1/(4\sqrt2)$, its mass exceeds $1/10$ and $n<100$.
The bound $\delta_H^2+\delta_G^2+16t^2\le1/2$ leaves eight
unordered pairs, all with $t\ge1/36$. Their complementary scalars
have least common multiple $N\le220$. Both $Nt$ and $Ns$ are
integral by the signed mass and FP induction identities, since
$\operatorname{Tr}(\mathcal O_K/(v^eu^j))\subseteq2\mathbb Z$.

Write $E=mv^eu^k$, $D=Mv^fu^\ell$, $a=m/M$, $r=\ell-k$.
Now $s/t\le36\cdot3/(4\sqrt2)<20$ and
Lemma~\ref{lem:relative-quadratic-indices} gives $a^2<4s/t<80$.
The leading coefficients of the primitive traces give $s/t\ge au^r/v$,
so $0\le r\le2$. Because $m$ divides both $NB_{e,k}$
and $NaB_{f,\ell}$, it divides $NaW_{e,f}(r)$. Here
$0<|W_{e,f}(r)|\le328$, whence $m\le220\cdot8\cdot328$
and $B_{e,k}=mt$ forces $k\le6$.

In these finite ranges the divisibility and Fourier comparisons
leave exactly
\[
\begin{array}{c|c|c|c}
(H,G)& E&D&s\\\hline
(12u,22vu)&132u&66vu&5/33\\
(16u,4v)&16&16u&3/8\\
(4v,22vu)&44v&44vu&5/22\\
(20vu,22vu)&220vu&220u^2&17/110
\end{array}
\]
The file \codefile{silver_small_check.py} verifies these reductions by
integer arithmetic.

The second row of the last table has $D$ as a signed codegree,
so gives sphericality. In the other three rows the available
remaining upper bound for the contributions of the ordinary complementary orbits is, respectively,
$31/1089,17/242,7/2420$. Each is less than $(1/2-s)^2$,
the minimum contribution of an irrational ordinary pair; two such
pairs give an even larger bound. Thus the ordinary complements
are rational. The masses $5/33$ and $17/110$ are both less than $1/6$ and
different from $1/8$, contrary to Lemma~\ref{lem:silver-three-four-rows}.

For $s=5/22$ the
rational ordinary codegrees divide $44$ by integrality of $E/f$,
so the only possible tuple is $(4,4,4,44)$; its discriminant
squareclass is $22$, whereas the maximal order gives two.
This excludes all eight pairs.

The other complementary orbit $G$ consequently has exponent at
least two. We again have $s/t\ge10/3$, and $\delta_G^2/z^2$,
$\delta_D^2/s^2$ at least $c=288/289$. In particular
\[
 L(s):=\delta_H^2+c\bigl((1-w-3s/10)^2+2s^2\bigr)\le1/2 .
\]
Minimizing the quadratic over $s\ge0$ gives
$\delta_H^2+(57600/60401)(1-w)^2\le1/2$, reducing the preceding
table to $4v,6v,12u,16u,18vu,20vu,22vu,24vu$.
For all of these except $4v$, the ordinary complementary codegrees
are rational. Indeed $2cs^2+\delta_H^2\le1/2$ and
$\delta_H^2\ge1/18$ give $s<1/2$. An irrational ordinary complement
then requires $L(s)+(1/2-s)^2\le1/2$, since two quadratic pairs
give at least $(1-s)^2/2\ge(1/2-s)^2$.
For the seven values the minimum of this last quadratic exceeds
$29/50$. These comparisons are direct substitutions of $w$ and
$\delta_H^2$.

For $16u$ and the four $nvu$ cases, $L(s)>1/2$ on $s\ge1/6$:
the quadratic is increasing from its value at $1/6$.
Thus $s=1/8$ and Lemma~\ref{lem:silver-three-four-rows}
supplies a signed $v$-family half-mass orbit. Since $t\le3/80$, the masses
of $18vu,22vu,24vu$ preclude a half-mass orbit. At $H=16u$,
$L(1/8)>1/2$. At $H=20vu$ we have $w=1/2$ and
$\delta_H^2=49/200$. If $s/t\ge5$, then the mixed sum is at least
$49/200+c((19/40)^2+2/64)>1/2$. Thus $s/t<5$ and the preceding
Casoratian argument for $s=1/8$ applies verbatim to the other
$v$-family orbit, contradicting its exponent bound.

It remains to rule out $H=6v$ and $12u$. In these two cases
$L(21/100)>1/2$ and $L$ is increasing thereafter, hence $s<21/100$.
Lemma~\ref{lem:silver-three-four-rows} now forces a signed
$v$-family orbit of mass $1/2$. For $H=6v$ none of the three
masses is $1/2$, since $w=1/3$ and $t\le3s/10$.
For $H=12u$ the half-mass orbit is in the unit family.

This leaves $H=4v$, of squared gap $1/8$. The long signed
pair has mass $1/2-t$. Using $s\ge10t/3$ in the mixed bound gives
\[
 \frac12\ge\frac18+\frac{288}{289}
          \bigl((1/2-3s/10)^2+2s^2\bigr).
\]
The right side increases for $s\ge1/3$ and already exceeds $1/2$
at $1/3$. Thus $s<1/3$.
\end{proof}

\begin{lemma}\label{lem:silver-saturation}
In the conductor-eight setting, let a real signed complementary
character of codegree $U$ satisfy $U/U'=D/D'>1$. Then $\mathcal C$
is spherical or contains a spherical rank-four subcategory with
Perron dimensions $(1,1,1+\sqrt2,1+\sqrt2)$. Its nonunit
invertible $g$ has positive lift in every spherical section of
this subcategory. If the signed algebra is $K^3$, then
$\mathcal C$ is spherical.
\end{lemma}
\begin{proof}
The character attains the Perron bound on the conjugate dimension
induction object $Z=F(B_\tau)$. By Lemma~\ref{lem:section-signs}
and equality in the triangle inequality, its constituents generate
a spherical subcategory. Its Perron field is $K$, since it contains
$Z$ of irrational Perron dimension $D/D'$. If proper, it has rank
at most four by Lemma~\ref{lem:field-corank}. The spherical
classifications leave the Ising ring and the asserted rank-four
rings~\cite{Ostrik2,Ostrik3,DPP}. In the Ising ring the rational
character $(1,-1,0)$ takes $Z$ to one: it extends to the full algebra,
where the values on $Z$ are positive units. Since the unit coefficient
is one, $Z=1+bX$, $b\ge0$. Its Perron dimension $1+b\sqrt2$ has
norm one, forcing $b=0$, a contradiction.

In rank four, $g$ interchanges the other two nonunit simples $Y,Z_0$.
If a spherical section gave $g$ dimension $-1$, their dimensions
would be opposite. This is impossible when $Y^*=Z_0$. In the
self-dual ring the two characters with this sign are
$(1,-1,1,-1)$ and $(1,-1,-1,1)$, so the intrinsic dimension
would be four. Its quotient by the Perron dimension $8+4\sqrt2$
has norm $1/2$, contrary to integrality by induction in the
spherical center.

Finally suppose the signed algebra is $K^3$. Then every simple
is self-dual. The two displayed characters extend to the signed
algebra, and each restriction occurs on a full quadratic orbit.
Thus the positive lift of $g$ would have at least four negative
eigenvalues. It permutes six basis elements and can have at most
three transpositions.
\end{proof}

\begin{proposition}\label{prop:silver-ordinary-rational}
In any nonspherical case of Theorem~\ref{thm:silver-finite-orbit},
all four ordinary complementary codegrees are rational integers.
\end{proposition}
\begin{proof}
Use the orbit $H=4v$, and write $E=mv^eu^k$, $G=nv^eu^j$,
$D=Mv^fu^\ell$. The primitive families of $E,G$ agree by the
discriminant. The inequalities $2\ell+f\ge2j+e$ and
Lemma~\ref{lem:silver-saturation} give strict inequality.
Set $a=m/M$, $h=m/n$, $r=\ell-k$, $b=j-k$.
The positive FP induction column has $H$-orbit contribution
$\operatorname{Tr}(\beta/v)/4=1/2$: its numerator is a positive
even integer strictly smaller than four. Put $C=B_{f,\ell}$,
$B=B_{e,k}$, $L=B_{e,j}$. Since
$\operatorname{Tr}(\mathcal O_K/(v^eu^j))\subset\mathbb Z$,
the $G$-orbit contributions on this and the unit column give
$n(1/2-s),n(1/2-t)\in\mathbb Z$. Thus
\[
  m=2(B+hL),\quad h\mid2B,\quad h\mid aC-B.
\]
In particular $8\nmid h$.

If an ordinary complementary codegree is irrational, its additional
mixed contribution is at least $(1/2-s)^2$ as above. We claim $s/t<5$.
Otherwise the mixed bound and $288/289>49/50$ would require
\[
 \frac18+\frac{49}{50}\bigl((1/2-s/5)^2+2s^2\bigr)
          +(1/2-s)^2\le\frac12 .
\]
The quadratic on the left has minimum $1/2+1/1304$.

Now $s/t=aB_{f,k+r}/B_{e,k}$. The only choices with $s/t<5$ are
\[
       (e,f,r,a)=(0,1,0,2)\quad\text{or}\quad(1,0,1,1).
\]
Indeed $v^{f-e}u^r>a$ by $D>E$. In the same family one must advance
at least one index, and the successive trace ratios are at least
five except $B_{0,1}/B_{0,0}=3$, which cannot occur since $j\ge2$
and $2j+e<2(k+r)+f$. From the unit family to the $v$-family,
$a$ is even, $B_{1,k}\ge B_{0,k}$, and advancing one index gives
a ratio at least six including $a$. Without an advance $k\ge j\ge2$
and $B_{1,k}/B_{0,k}\ge5/3$, forcing $a=2$. In the opposite direction
$r\ge1$, with $B_{0,k+1}/B_{1,k}\ge3$, so $r=a=1$.

Both choices give $j\le k$ and $|W_{e,f}(r)|=8$. By
Lemma~\ref{lem:relative-quadratic-indices}, $h\mid16a$; since
$8\nmid h$, we have $h\le4$. Thus
$t=B/[2(B+hL)]\ge1/10$ because $L=B_{e,j}\le B$.
But $s\ge10t/3$ then contradicts $s<1/3$.
\end{proof}

\begin{lemma}\label{lem:silver-no-invertible-mass}
Suppose the only invertible is the unit, the Perron field is
$K=\mathbb Q(\sqrt2)$, and the four ordinary complementary codegrees
are even integers at least four. Assume their four-rows and six-rows
are rational. Then $s\ge1/4$. If the dyadic factors have ranks six or four and two,
all with residue field $\mathbb F_2$, then $s\ge5/14$.
\end{lemma}
\begin{proof}
Three four-rows force an invertible. Above mass $3/4$ the only
profiles are therefore $(4,4,6,d)$, $d=6,8,10$. At the inert
prime three a six-row needs a totally ramified partner by
Lemma~\ref{lem:silver-local-partner}. The prime-to-three axes split
off, and two six-rows are incongruent by squared distance twelve.
This proves $s\ge1/4$.

Now impose the dyadic hypotheses and suppose the mass exceeds $9/14$.
Without two four-rows
the sorted norms must be $(6,6,6,6)$, $(4,6,6,d)$ with $d\le16$,
or $(4,6,8,8)$. In each case the same argument at three leaves too few partners for the six-rows.

Two four-rows have different supports, since a matched pair gives
an invertible. They exhaust the dyadic residue characters, so all
integral rows have norm divisible by four. In particular there are
exactly two four-rows: their distinct supports, after deletion of
the unit, have disjoint two-element complements in five positions.
A four-row in the rank-two factor needs a rational local partner,
as the Perron field is ramified. If the last two rows formed a
quadratic field $F$, their common codegree would contribute a square
to the order discriminant. Thus $2|\operatorname{disc}(F)|$ has
the squareclass of $N(D)$, which is one or two. The only possible
fields are $\mathbb Q(i),\mathbb Q(\sqrt2),\mathbb Q(\sqrt{-2})$,
none split at two. Hence the two last rows are rational.

The four-row's partner has norm $d\equiv12\pmod{16}$ by
Lemma~\ref{lem:trace-order-conductor}. At $d=12$ the difference
of the two rows is divisible by four and has squared norm sixteen.
Its sole nonzero entry contradicts orthogonality to the Perron row.
Thus $d\ge28$. Since $1/2+1/28+1/12<9/14$, the remaining norm
is eight; then $d<56$ gives $d=28$ or $44$. Neither
$4^2\cdot8d N(D)$ has squareclass two, as required by the ordinary
algebra $K\times\mathbb Q^4$.
\end{proof}

\begin{proposition}\label{prop:silver-invertible}
In any nonspherical case of Theorem~\ref{thm:silver-finite-orbit},
there is exactly one nontrivial invertible object $g$. Its
permutation on the simple classes has two transpositions and two
fixed points.
\end{proposition}
\begin{proof}
Write the entries of the $H=4v$ signed character as $a_i+b_i\sqrt2$.
The codegree and conjugate-row orthogonality give
$\sum a_i^2=4$, $\sum b_i^2=2$, $\sum a_i b_i=2$. Since the unit
entry is one, the three other nonzero $a_i$ are signs. The two
nonzero $b_i$ equal the corresponding $a_i$. Orthogonality of
both rows to the positive dimension character now forces their
shape, up to permutation,
\[
 (1,-1,1+\sqrt2,-(1+\sqrt2),0,0).
\]
Every character outside this pair takes the same values on the
first two entries, and equal values on the next two. Thus the
second simple $g$ satisfies $g^2=1$ in $T$, and signed left
multiplication by $g$ is the two-transposition permutation. Indeed
this matrix is symmetric and integral with square one, and the
positive dimension character rules out negative permutation signs.
If there is any nontrivial ordinary invertible, its positive lift
takes value $+1$ on the spherical dimension orbit, and the orbits where it
takes value $-1$ must have total reciprocal mass $1/2$ by its zero
unit coefficient. Only the $H$-orbit has this property outside the
physical orbit. Therefore that invertible must be $g$.

Suppose $g$ is not invertible. All four ordinary complementary
codegrees are even and at least four, and their four- and six-rows
are rational by Lemma~\ref{lem:real-four-six}.

The algebra $T/2T$ has a local factor of rank at least four.
Indeed the ideal generated by $g-1$ has basis the two disjoint
pair differences, and is square-zero. It is indecomposable: its
integral version maps under the $H$-character isomorphically onto
$2\mathbb Z[\sqrt2]$, so reduction gives the regular module of
the local algebra $\mathcal O_K/2\mathcal O_K$. The two differences
remain independent after reduction. Consequently the ideal lies in
one local factor. Its coefficient pairing vanishes there, whereas
the pairing on the entire factor is perfect. The rank is therefore
at least four. All these statements about local ranks transfer
to $R$ by parity.

The signed ramified quadratic branches exclude rank-one factors;
all residue fields are $\mathbb F_2$. Lemma~\ref{lem:silver-no-invertible-mass}
now contradicts $s<1/3$ from Theorem~\ref{thm:silver-finite-orbit}.
\end{proof}

\begin{lemma}\label{lem:silver-common-ordinary-mass}
In the conductor-eight setting, suppose the ordinary algebra is
commutative, its Perron field is
$K=\mathbb Q(\sqrt2)$, and its only nonunit invertible acts by
$(1,g)(X,Y)$ with fixed simples $P,Q$. Suppose both negative rows
have codegree four and common support, with values $\pm1$ or $\pm i$
at $X,Y$. At two assume the order is local or has factors of ranks
four and two, with both negative rows in the rank-four factor. If the
two remaining ordinary characters outside the Perron pair are real and
have even rational codegrees at least four, then
$\operatorname{Tr}(1/D)\ge5/14$.
\end{lemma}
\begin{proof}
Write $s=\operatorname{Tr}(1/D)$ and $a\le b$ for the last two codegrees.
The two positive rows have the form $(1,1,z,z,p,q)$. If $s<5/14$, then
\begin{equation}\label{eq:silver-ordinary-final-denoms}
       1/a+1/b=1/2-s>1/7.
\end{equation}
There is no third four-row. Such a row would be rational by Lemma~\ref{lem:real-four-six}. In the shape $(1,1,z,z,p,q)$ it cannot have $z=\pm1$ by orthogonality to FP and $d_X>1$. Its support thus differs from that of the matched short pair. The latter pair cannot occupy a rank-two dyadic factor, by the dyadic hypothesis. The third four therefore lies in the rank-two factor with one rational partner of codegree $b$; the FP field is ramified at two and cannot supply a rank-one branch. The two rows in this factor are congruent modulo four by perfect trace duality and projection to its two axes. Thus $b\equiv12\pmod{16}$ by the squared row distance. The absolute discriminant identity gives $b$ of squareclass one or two, neither compatible with that residue modulo sixteen.

Thus $a=6,8,10$ or $12$. For distinct $a,b$ the two folded rational-codegree rows are rational and
\begin{equation}\label{eq:silver-folded-squareclass}
 abN(D)\text{ has squareclass two}.
\end{equation}
The short split algebra or Gaussian field contributes square absolute discriminant, while $K$ contributes squareclass two. If $a=b$ on a quadratic field $F$, its positive discriminant $c$ satisfies $c\mid2a$, $c\le a$ by the real row-difference test, and $N(D)$ has squareclass $2c$.

At $a=6$ both rows are rational by Lemma~\ref{lem:real-four-six}. The first has norm two modulo four and hence represents a different dyadic residue row from that of the matched four-pair. It lies in the rank-two factor and must pair with the other rational row, so $b\equiv2\pmod4$. At three the inert $K$-branch cannot pair with the six-row by Lemma~\ref{lem:silver-local-partner}; the codegree-four summand splits off. Hence the rational rows are congruent modulo three and $9\mid6+b$. But \eqref{eq:silver-folded-squareclass} forces $b=6z^2$ or $3z^2$. The first fails the modulo-nine test and $b$ even in the second case requires $4\mid b$.

At $a=10$, \eqref{eq:silver-ordinary-final-denoms} gives $10\le b\le22$. If both rows are rational, the same argument at the inert prime five gives $25\mid10+b$, impossible for even $b$ in this range. If they form a quadratic orbit, $c\mid20$ and $c\le10$ force $c=5$, inconsistent with $N(D)$ having squareclass $2c$. At $a=12$, the same mass bound gives $b=12,14,16$, with $14,16$ excluded by \eqref{eq:silver-folded-squareclass}. Two rational twelve-rows would have to pair at three by the same local argument, whereas nine does not divide their squared distance $24$. For a quadratic pair $c=8$ or $12$. Only $8$ is compatible with the discriminant, and then $D=Mu^\ell$ has square norm. But $s=1/3$ gives $M=3(u^\ell+u^{-\ell})$, of two-adic valuation one, whereas $4\mid M$.

Finally let $a=8$. A quadratic pair requires $b=8,c=8$. For its two rows write the first as $(1,1,z,z,p,q)$ with $z,p,q\in\mathbb Z[\sqrt2]$. Their squared difference $16$ gives
\[
 2[\sqrt2](z)^2+[\sqrt2](p)^2+[\sqrt2](q)^2=2,
\]
where each bracket denotes the coefficient of $\sqrt2$ before squaring. If $[\sqrt2](z)\ne0$ the row difference is supported on $X,Y$ with equal entries, contrary to FP orthogonality. Otherwise the last two coefficients are opposite signs and $d_P=d_Q$. Write $p=A+\sqrt2,q=C-\sqrt2$ after conjugation. The rational norm equation gives $A=C$, $z\in\mathbb Z$ and $2z^2+2A^2=2$. FP orthogonality would require $d_X=1$ or $d_P=1$.

If the eight-row is rational it has form $(1,1,z,z,2e,0)$ after interchanging $P,Q$, where $z,e$ are signs. Its reduction agrees with the two negative rows. In a dyadic decomposition of ranks four and two the negative pair and the eight-row occupy three dimensions in the rank-four factor. The last denominator is rational, so that factor must also contain its axis, leaving the FP branch alone. Its idempotent has dyadically integral unit coefficient $s=3/8-1/b$. Hence $b=8B$ with $B\equiv3\pmod8$. By \eqref{eq:silver-ordinary-final-denoms}, $b<56$, so $b=24$, excluded by \eqref{eq:silver-folded-squareclass}. We can thus assume the ordinary dyadic order local. Its remaining row reduces to the same short row, so $4\mid b$. The discriminant identity and $8\le b<56$ give $b=8,16,32,36$. Write its coordinates as $(1,1,z',z',-e(1+zz'),C)$, using orthogonality to the eight-row. Its norm is
\[
 b=3+3(zz')^2+2zz'+C^2,
\]
where $z'$ is odd by reduction at two. At $b=16,32$ there is no solution. At $b=36$ the only possibility $zz'=3,C=0$ gives $d_X=3/5$ by the two FP orthogonality equations. At $b=8$ both rows cannot be nonzero at the same fixed simple, since their two equations would force $d_X=1$. The other possibility has $zz'=-1,C=\pm2$. The two fixed dimensions must then be $x-1,x+1$ with $x=d_X=d_Y$. Conjugate-row orthogonality gives $N(x)=-1$ and $D=4(x^2+1)=4x(x-x')$. At the fixed simple of dimension $x+1$, the two eight-rows contribute $\pm1$ to its self-square coefficient and the negative rows contribute zero. The FP pair contributes
\[
 \operatorname{Tr}((x+1)^3/D)=(\operatorname{Tr}(x)+4)/4,
\]
since $1/x=-x'$. Writing $x=A+B\sqrt2$, the equation $A^2-2B^2=-1$ makes $A$ odd. Thus this coefficient is nonintegral. All rational alternatives are excluded.
\end{proof}

\begin{theorem}\label{thm:silver-three-complete}
Let $K=\mathbb Q(\sqrt2)$. Suppose a rank-six fusion category has
commutative ordinary character algebra, signed rational algebra $K^3$,
and both its ordinary Frobenius--Perron and spherical dimension
characters have field $K$. If each complementary signed character
has irrational formal codegree, the category admits a spherical
structure.
\end{theorem}
\begin{proof}
Assume otherwise. By Theorem~\ref{thm:silver-finite-orbit} and
Propositions~\ref{prop:silver-ordinary-rational} and
\ref{prop:silver-invertible}, the signed short orbit has codegree
$4v=8+4\sqrt2$, all four ordinary complementary codegrees are
rational even integers, and the unique nontrivial invertible $g$
acts on simples by $(1,g)(X,Y)$ with two fixed points $P,Q$.
The two $g$-negative ordinary characters have rows
$(1,-1,y,-y,0,0)$. On the span of $1-g,X-Y$ multiplication by
$X$ has matrix $\left(\begin{smallmatrix}0&1\\1&n\end{smallmatrix}\right)$,
by self-duality. Since $2(1+y^2)$ is rational for both roots and
their product is $-1$, necessarily $n=0$: otherwise $y$ is
rational, and the two rational algebraic units must be $1,-1$.
Thus both negative rows are rational four-rows. All other ordinary
rows have shape $(1,1,z,z,b,c)$.

We can now apply Lemma~\ref{lem:silver-common-ordinary-mass}.
The proof of Proposition~\ref{prop:silver-invertible} gives a dyadic
factor of rank at least four. There are no rank-one factors, since
the signed order has three ramified quadratic branches. In the
two-factor case the negative rows cannot occupy the rank-two factor:
their projection $\{(a,b):a\equiv b\pmod2\}$ has discriminant
valuation two, whereas $4\cdot4$ has valuation four. Thus $s\ge5/14$.

This contradicts $s<1/3$ from Theorem~\ref{thm:silver-finite-orbit}.
\end{proof}

\subsection{Two rational complementary codegrees}

\begin{proposition}\label{prop:silver-two-pair-finite}
Suppose that a rank-six fusion category has commutative ordinary and
signed character algebras and common dimension-character field
$K=\mathbb Q(\sqrt2)$. Suppose the signed codegrees are
$E,E',H,H',g,h$, with $H$ irrational and $g,h$ rational integers at least three.
If the category has no pseudo-unitary Galois conjugate, then either the
last two characters form a Gaussian pair with $g=h=4$, or they are
rational and, upon interchanging them,
\[
\begin{array}{c|l}
 g&h\\\hline
 4&4,8,9,16,18,25,32,36,49,50\\
 7&7 .
\end{array}
\]
\end{proposition}
\begin{proof}
Write $u=3+2\sqrt2$, $v=2+\sqrt2$ and orient $E,H$ by their
larger conjugate. Let $t=\operatorname{Tr}(1/E)$,
$w=\operatorname{Tr}(1/H)$, $s=\operatorname{Tr}(1/D)$ and
$L=t+w=1-1/g-1/h$. For $f\ge f'$ put $\delta_f=1/f'-1/f$.
From \eqref{eq:silver-mixed} and Lemmas~\ref{lem:central-witness}
and~\ref{lem:quadratic-mass-gap} we obtain
\[
 \delta_E^2+\delta_H^2+2\delta_D^2\le\tfrac12,\qquad
 D/D'\ge H/H'\ge u^2,\qquad s\ge3t .
\]
Write an oriented codegree as $nv^eu^j$ as before. Its ratio to
its conjugate is $u^{2j+e}$, and $8\mid\operatorname{Tr}(f)$
for a quadratic row by orthogonality. Thus $4\mid n$ if $e=0$,
and $2\mid n$ if $e=1$.

If $H/H'\ge u^3$, the squared gap divided by squared mass is at least
$49/50$ for both $D$ and $H$. Hence
\[
 \tfrac12\ge(49/50)(18t^2+w^2),\qquad L^2\le 475/882.
\]
If $H/H'=u^2$, write $H=nu$. Now $4\mid n$ and
$\delta_H^2=32/n^2<1/2$, so $n\ge12$, and $w=6/n\le1/2$.
The squared gap ratio is $8/9$ at $u^2$, so
$16t^2+(8/9)w^2\le1/2$. The upper bound
$w+\tfrac14\sqrt{1/2-(8/9)w^2}$ for $L$ is increasing on
$[0,1/2]$ and is less than $2/3$ at the endpoint. Thus
$L^2\le475/882$ always.

We show there is no rational signed character $\rho$ of codegree
three. Three is inert in $K$. The local partner
Lemma~\ref{lem:silver-local-partner} puts $\rho$ and the other rational
character $\eta$ of codegree $h$ in their own two-branch factor
at three, with congruent character rows modulo three. The
coefficient trace of the identity of this factor is integral over
$\mathbb Z_3$, hence $1/3+1/h\in\mathbb Z_3$ and
$h/3\equiv-1\pmod3$. The absolute discriminant of the signed
order in $K^2\times\mathbb Q^2$ gives that
$N(E)N(H)3h$ is a square. The squareclasses of the two norms are
one or two, so $h=3a^2$ or $6a^2$. The congruence forces
$h=6a^2$ with $3\nmid a$. Consequently $E,H$ belong to opposite
primitive families. The simultaneous column $W$ has components
$E/E',H/H',1,1$, allowing reversal of orientation in the second
component. The residue of $u$ has order four in $\mathbb F_9$.
Thus $W$ has a residue generating $\mathbb F_9$ on the unique
$v$-family branch, and a residue in $\mathbb F_3$ on every
other branch. The $v$-family branch splits integrally at three:
the two factors of the characteristic polynomial on the disjoint
residue sets are coprime over $\mathbb Z_3$. Its order is maximal
since it contains $W$, and $1,W$ span its residue extension.
Its coefficient functional is $\operatorname{Tr}_K(-/f)$ for
the corresponding codegree $f$; as $K$ is unramified at three,
integrality forces $f$ to be a $3$-adic unit. On the other hand
$E/3$ is an algebraic integer, the dimension of an induction label.
It follows that $H$ is the $v$-family codegree and is a $3$-adic unit.

For each signed induction label $B$ write $m_B=\rho(FB)\in\mathbb Z$.
On the signed dimension labels the two entries are one, since
$F(B_1)=1$ and the rational row takes $W$ to one. The entries on
$H,H'$ agree by rationality and Galois symmetry; call them $c$.
Symmetry gives $\chi_H(F A_\rho)=(H/3)c$, an algebraic integer,
so $3\mid c$. Write $m=m_{A_\rho}$ and $n=m_{A_\eta}$.
The same symmetry gives
$\eta(F A_\rho)=(h/3)n\equiv-n\pmod3$.
Since $\eta\equiv\rho$, this makes $m+n\equiv0\pmod3$.
But the signed unit coefficient of $F A_\rho$ is one, and
symmetry expresses it as $(2+2c+m+n)/3$. Its numerator is two
modulo three, a contradiction.

It remains to list the complements. A rational pair $g\le h$
has $g,h\ge4$, and $gh$ is a square or twice a square by
the discriminant. In
$(1-1/g-1/h)^2\le475/882$ we have $g\le7$.
For $g=4,5,6,7$ this inequality bounds $h$ by $61,15,10,8$
respectively. The squareclasses leave exactly the table
and $(5,5),(5,10),(6,6)$. At the inert prime five, the first two
pairs have integral combined trace by the local partner lemma, whereas
neither $2/5$ nor $3/10$ is integral there. In the third case use the
same lemma at the inert prime three: $1/6+1/6$ is nonintegral.
For a quadratic complementary field $F$ with a common rational
codegree $g$, the same bound gives $g\le7$; the absolute
discriminant of $F$ has squareclass one or two. A real such
quadratic field is $\mathbb Q(\sqrt2)$ and the complementary
codegree is at least its discriminant eight by character
orthogonality. In the imaginary case the single nonself-dual
pair gives $g=|\operatorname{disc}(F)|b^2$ by row orthogonality.
Hence $F=\mathbb Q(i)$ and $g=4$.
\end{proof}

\subsection{Two invertibles in a rank-four subcategory}
\begin{proposition}\label{prop:folded-silver-descent}
Suppose a complex rank-six fusion category has commutative sphericalization and all ordinary Frobenius--Perron dimensions in $K=\mathbb Q(\sqrt2)$. Suppose tensoring by an involutive invertible $g$ interchanges $1,g$ and $Y,Z$ and fixes $P,Q$. If $1,g,Y,Z$ form a fusion subcategory $\mathcal B$ admitting a spherical structure and satisfying $\operatorname{FPdim}(Y)=\operatorname{FPdim}(Z)=\xi=1+\sqrt2$, then the category admits a spherical structure.
\end{proposition}
\begin{proof}
The quadratic relation in the Grothendieck ring of $\mathcal B$ is $Y^2=\epsilon+Y+Z$, where $\epsilon$ is $1$ or $g$. Indeed exactly one invertible occurs by rigidity, and two copies of just one of $Y,Z$ are impossible. In the self-dual case they would form a rank-two fusion subcategory on that simple with product $X^2=1+2X$, excluded by the rank-two classification. In the nonself-dual case the eigenvalues of $Y$ in the $g=-1$ component are purely imaginary, since $Y^*=Z=gY$; a relation $\lambda^2=-1\pm2\lambda$ is impossible there.

Multiplication by $Y$ preserves the span of $P,Q$ by reciprocity. Its matrix is symmetric since $Y^*$ is $Y$ or $gY$ and $g$ is trivial on this span. Its minimal polynomial is $x^2-2x-1$, irreducible over $\mathbb Q$. The trace and determinant give, after interchanging $P,Q$, the matrix $\left(\begin{smallmatrix}0&1\\1&2\end{smallmatrix}\right)$. Thus $YP=Q$ and the dimensions of these objects are $p,\xi p$. In particular $P,Q$ are self-dual, since duality preserves the fixed objects and their dimensions differ. In $P^2$ reciprocity gives the multiplicities $1,1,0,0$ at $1,g,Y,Z$. Hence
\[
 p-2/p=r+s\xi,\qquad r,s\in\mathbb Z_{\ge0}.
\]
In particular $p\mid2$. Since $(2)=(\sqrt2)^2$ as ideals, and the units are $\pm\xi^k$, we have $p=\xi^k,\sqrt2\xi^k$ or $2\xi^k$ for $k\ge0$ by $p\ge1$. Also
\[
 \operatorname{FPdim}(\mathcal C)/\operatorname{FPdim}(\mathcal B)=1+p^2/2
\]
is a $d$-number in the sense that its quotient by its conjugate is a unit: both global FP dimensions have this property. Write $u=\xi^2$, the generator of the totally positive units. For $k\ge1$ the quotient of $2+u^k$ by its conjugate lies strictly between $u^{k-1}$ and $u^k$, while the quotient of $1+2u^k$ by its conjugate lies strictly between $u^k$ and $u^{k+1}$. Indeed division by $u^k$ gives $(1+2u^{-k})/(2+u^{-k})$ or its reciprocal. These exclude $p=\xi^k,2\xi^k$ for $k\ge1$. A fixed object cannot be invertible, so $p\ne1$. At $p=2$ the equation above gives $s=0$. In the remaining family $N(p)=\pm2$, so $p-2/p$ is $p+p'$ or $p-p'$; in the first case it is rational, while in the second it has trace zero. Since $r,s$ are nonnegative, either case gives $s=0$. Consequently $1,g,P$ form a rank-three fusion subcategory with $P^2=1+g+rP$. It has a spherical tensor section, by the corank-one theorem \cite[Theorem~4.5]{DCW}, and its section takes $g$ to its positive lift since $gP=P$.

By Lemma~\ref{lem:silver-saturation} and its rank-four argument,
every spherical section of $\mathcal B$ takes $g$ to its positive
lift. Thus the two lifted subcategories agree on $g$.

Their products furnish one lift of each simple: $YP=Q$, and the products of $P$ with the selected lifts of $Y,Z=gY$ agree since $gP=P$. The span is tensor-closed because the sphericalization ring is commutative and the two original lifted subcategories are tensor-closed. Forgetting gives a tensor equivalence and the spherical structure descends.
\end{proof}

\begin{proposition}\label{prop:silver-four-rational-E}
Under the hypotheses of Proposition~\ref{prop:silver-two-pair-finite}, suppose both complementary rational codegrees are four and $E$ is rational. Then the category is spherical.
\end{proposition}
\begin{proof}
The absolute discriminant of the maximal signed order is a square for the split or Gaussian complement. Thus $H=nu^j$ is in the unit family, and $4\mid n$. We have $j\ge1$ by the positive-induction argument, $h=E/n\in\mathbb Z$, and the signed reciprocal mass gives $E=4+2hA_j$, where $A_j=u^j+u^{-j}\equiv2\pmod4$. Hence $h\mid4$. Were $h$ even the last equation would give $v_2(E)=2$, contrary to $4h\mid E$. Thus $h=1$ and $n=E$. Put $w=A_j/(4+2A_j)\ge3/8$. On the positive FP induction column the two rational codegrees contribute at least $1/2$, so $s<1/2$. Write $D=Mv^d u^\ell$ with $v=2+\sqrt2$. By integrality of $E/D$, $c=E/M$ is an integer, even if $d=1$. By central dominance $\ell\ge j$. Therefore in the unit family $s/w=cA_\ell/A_j<4/3$ gives $c=1,\ell=j$; the successive trace ratios from index one exceed $4/3$. But then $D=H$, which yields a spherical section by signed domination. In the other family $s/w=c\operatorname{Tr}(1/(vu^\ell))/A_j\ge2$, since the reciprocal traces in this family equal or exceed the unit-family traces at the same index. This contradiction completes the proof.
\end{proof}

\subsection{Two short rational profiles}

\begin{proposition}\label{prop:silver-four-sixteen}
In the setting of Proposition~\ref{prop:silver-two-pair-finite}, the
signed rational codegrees $(4,16)$ force sphericality. For $(4,9)$,
a nonspherical category would have $h=6$ or $9$ and $b=0$, or
$h=36$ and $b=-1$, where $h=m/n$ is the ratio of the signed scalars
and $b=j-k$ the difference of their exponents.
\end{proposition}
\begin{proof}
Let $L=t+w$ be the total quadratic mass. We first record the
ordinary contribution to the mixed bound, writing
\[
 L=11/16\quad\text{or}\quad23/36,\qquad \kappa=1/4,
 \qquad G_0=\delta_E^2+\delta_H^2+2\delta_D^2 .
\]
In both cases
\begin{equation}\label{eq:silver-short-deficit}
 G_0\le1/2,\qquad s\ge\kappa\quad\text{or}\quad
 G_0+(1/2-s)^2\le1/2 .
\end{equation}

Indeed there is no nonunit invertible. Its signed unit coefficient
would give $t\pm w\pm1/4\pm1/q=0$, $q=16,9$. Since
$t\le\sqrt{9/32}/3$, this is impossible at $q=16$ and forces
$(t,w)=(5/36,1/2)$ at $q=9$. In the latter case
$(8/9)(w^2+2(3t)^2)>1/2$. We have used $s\ge3t$ and the
squared gap ratios at $H,D$. Lemma~\ref{lem:silver-saturation}
now gives $D/D'>H/H'$.

The dyadic orders have only degree-one residue fields. Their
rank-one factors are exactly the odd-codegree axes: the two
quadratic branches ramify, and an even rational axis cannot be
isolated. Thus at $q=16$ all rational ordinary codegrees are even
and at least four. For an all-rational complement,
Lemma~\ref{lem:silver-no-invertible-mass} gives $s\ge1/4$.

At $q=9$ there is exactly one isolated odd row. Its reduction agrees
with the signed nine-row, so its norm is $1\pmod4$ and is at least
five. For an all-rational complement with two four-rows, their
distinct reductions place them in factors of ranks two and three.
The partner of the four-row in rank two has norm $d\equiv12\pmod{16}$
by Lemma~\ref{lem:trace-order-conductor}. At $d=12$ the rational
row difference has squared norm sixteen and is divisible by four,
contrary to orthogonality to the Perron row. Hence $d\ge28$ and
$s\ge1-1/5-1/2-1/28>1/4$. With at most one four-row the only
exception to $s\ge1/4$ would be $(5,4,6,6)$. The rational six-rows
have distinct reductions at three by squared distance twelve, and
neither can pair with the inert Perron field by
Lemma~\ref{lem:silver-local-partner}. Finally an irrational ordinary
pair contributes at least $(1/2-s)^2$ when $s<\kappa$, since the
rational companions have reciprocal mass at most $1/2$.
Two irrational pairs contribute at least $(1-s)^2/2$.
This proves \eqref{eq:silver-short-deficit}.

The case $H=nu$ would give $w=6/n$, $4\mid n$, $n\ge12$.
Since $t\le\sqrt{9/32}/3$, only $q=9,n=12$ is possible, giving
the excluded pair $(t,w)=(5/36,1/2)$. Thus the squared gap
ratios at $H,D$ are at least $c=49/50$. We claim
\[
                            s/t<15 .
\]
Otherwise $G_0\ge c((L-s/15)^2+2s^2)$. On $s\ge\kappa$ this
quadratic increases from a value greater than $1/2$ for both values of $L$. Adding $(1/2-s)^2$ gives a quadratic with minimum
greater than $1/2$ in both cases, contradicting
\eqref{eq:silver-short-deficit} also for $s<\kappa$.

Write
\[
 E=mv^eu^k,\qquad H=(m/h)v^eu^{k+b},\qquad
 D=(m/a)v^fu^{k+r}.
\]
The signed discriminant gives the same family for $E,H$. Set
$B=B_{e,k}$, $G=B_{e,k+b}$, $C=B_{f,k+r}$ and $L=P/Q$ in lowest
terms. The unit and Perron columns give
\[
 Pm=Q(B+hG),\qquad h\mid QB,\ QaC.
\]
For the last divisibility multiply the Perron column by $Qm/h$:
the rational terms and the $H$-trace are then integral.
Strict dominance gives $2(k+b)+e<2(k+r)+f$.

Only six choices remain for the family and scalar of $D$:
\[
 (f,r,a)=(e,1,1),\ (e,1,2),\ (1-e,e,2-e),\qquad e=0,1 .
\]
Indeed $a<v^{f-e}u^r$ and $aC/B<15$. An advance of two indices
already gives $C/B\ge17$ in either family. At one index in the
same family the trace ratio is at least five, except at $e=k=0$,
which violates strict dominance. From $e=0$ to $f=1$ the scalar
$a$ is even, and any advance gives $aC/B\ge2B_{1,2}/B_{0,1}>15$
for $k\ge1$. At $k=0,r=1$ strict dominance would force $H=nu$.
These comparisons follow directly from the two-term trace formula.

In all six cases $b\le0$ and $|W_{e,f}(r)|$ is $8,16$ or $32$.
Lemma~\ref{lem:relative-quadratic-indices} gives $h\mid QaW$.
Since $B\equiv2\pmod4$ and $h\mid QB$, we obtain $h\mid2Q$.
If $h$ is even, $v_2(m)=1+v_2(Q)$ and $m/h$ is a multiple of
$4/(1+e)$, so $h\cdot4/(1+e)\mid2Q$. From $1<w/t=hG/B\le2hu^b$
and $2Q\le72$, it follows that $b=-2,-1$ or $0$.

These short lists can be checked at every exponent at once.
Put $x=hu^b$, $y=a(v/2)^{f-e}u^r$ and $\eta=989/1000$.
All three ratio exponents are at least three: $H=nu$ is excluded,
$k\ge k+b$ and $D/D'>H/H'$. Thus the correction to each leading
trace or gap term has relative size at most $u^{-3}$.
Since $\eta<(1-u^{-3})/(1+u^{-3})$,
Lemma~\ref{lem:relative-quadratic-indices} and its two-term formulas
give
\[
 G_0\ge (\eta L)^2\frac{1+x^2+2y^2}{(1+x)^2},\qquad
 s\le \frac{Ly}{\eta(1+x)} .
\]
Substitute these bounds in \eqref{eq:silver-short-deficit}. At
$Q=16$ the $81$ choices leave none. At $Q=36$ the $135$ choices
leave $(h,b)=(6,0),(9,0)$ for either mixed family and $(36,-1)$
for $e=1,f=0$. The \codecheck{silver_short_profiles.py}{exact comparisons}
use rational endpoints for $3+2\sqrt2$.
\end{proof}

\begin{proposition}\label{prop:silver-four-nine}
The pair of rational signed codegrees $(4,9)$ also forces sphericality.
\end{proposition}
\begin{proof}
The signed mass equation requires $23\mid B+hG$, the trace of
$\alpha=(1+h u^{-b})/(v^e u^k)$. A trace-zero element of
$\mathbb Z[\sqrt2]/23$ has norm $-2$ times a square. As two is
a square modulo $23$, this requires $-N(1+h u^{-b})$ to be a
square. For $(h,b)=(6,0),(9,0)$ the norm is the nonzero square
$(1+h)^2$. For $(36,-1)$ it is $1+36^2+6\cdot36\equiv18\pmod{23}$,
also a nonzero square. But $-1$ is a nonsquare modulo $23$.
\end{proof}

\subsection{The eight remaining codegree pairs}
\begin{theorem}\label{thm:silver-large-pairs}
A category satisfying the hypotheses of
Proposition~\ref{prop:silver-two-pair-finite} is spherical whenever
its two rational complementary characters have one of the following
pairs of formal codegrees:
\[
   (4,8),\ (4,18),\ (4,25),\ (4,32),\ (4,36),\ (4,49),\ (4,50),\ (7,7).
\]
\end{theorem}
\begin{proof}
Suppose the category is not spherical. Then it has no pseudo-unitary
Galois conjugate, since Galois transport preserves sphericality.
Write $L=1-1/g_1-1/g_2$, and keep the notation of
Proposition~\ref{prop:silver-two-pair-finite}. Except at $(4,8)$, there is no nonunit
invertible. Such an object would have order two by total reality,
and its positive lift would have dimension one. Its character values
on each field are constant signs. Since $s\le\sqrt{9/32}$ and
$s\ge3t$, the gap $w-t=L-2t$ is at least
$25/36-2\sqrt{9/32}/3>11/36$, greater than the total rational mass.
The coefficient of the unit of such a lift cannot vanish.
By Lemma~\ref{lem:silver-saturation}, equality $H/H'=D/D'$ would
supply a spherical rank-four subcategory and a nonunit invertible
$g$. This leaves only $(4,8)$. Here the spherical section takes $g$
to its positive lift, and the saturating $H$-character takes this
lift to one. Thus both signed quadratic pairs take $g$ to one.
Its two transpositions within the subcategory already span a
two-dimensional negative eigenspace, so the other two simples
are fixed and the two rational characters span that eigenspace.
They vanish on the fixed simples and restrict to the two rational
minus-characters of the self-dual rank-four ring. Both then have
codegree four, contrary to $(4,8)$. Hence $D/D'>H/H'$ in all cases.

We recall the necessary bounds on $t=\operatorname{Tr}(1/E)$. For
the six profiles with $L=1-1/g_1-1/g_2>17/24$, one has
$t\ge L-17/24$. Indeed $H=nu$ gives $w\le1/2$ as above; at
$H=nvu$ the gap $7\sqrt2/n$ leaves $n>14$, hence $w\le5/8$.
All later orbits have gap ratio at least $\sqrt{288/289}$, giving
$w\le17/24$.
For $(4,18)$ the even signed rational axes and the ramified branches
give no dyadic factor of rank one. The same holds at $(4,8)$, where
$t<1/8$ also rules out an invertible, since $w-t>3/8$. As in
Proposition~\ref{prop:silver-four-sixteen}, every rational ordinary
codegree is then even and at least four, and the ordinary four- and
six-rows are rational. Lemma~\ref{lem:silver-no-invertible-mass}
gives $s\ge1/4$ for an all-rational complement. Otherwise, for $s<1/4$
the additional ordinary contribution is at least $(1/2-s)^2$.

Suppose $t\le1/100$. At $(4,18)$ we have $w>2/3$, excluding
$H=nu$, and the squared gap ratios at $H,D$ are at least $c=49/50$.
At $(4,8)$ we have $123/200\le w<5/8$, excluding both $H=nu$
and $H=nvu$, since $10/n$ skips this interval for even $n$.
Here take $c=288/289$. In both cases $w\ge L-1/100$ and the
mixed bound would require
\[
 c(L-1/100)^2+\min\{c/8,\ c/(2(1+2c))\}\le1/2 .
\]
The two terms in the minimum bound respectively $2cs^2$ on
$s\ge1/4$ and $2cs^2+(1/2-s)^2$ on the line. Substitution of
the two values of $L,c$ proves $t>1/100$.

Write $\delta$ for the asserted positive lower bound on $t$.

In this proof $B_{e,k}=\operatorname{Tr}(1/(v^e u^k))$ and
$d_{e,k}=(1/(v'^e u^{-k})-1/(v^e u^k))/\sqrt2$.
The pairs $(B_{e,k},d_{e,k})$ start at $(2,0)$ and $(2,1)$
for $e=0,1$ and advance by $(x,y)\mapsto(3x+4y,2x+3y)$.
Write
\[
 E=m v^e u^k,\quad D=(m/a)v^f u^{k+r},\quad
 H=(m/h)v^i u^{k+b}.
\]
Here $a,h$ are positive integers by integrality of $E/D,E/H$
at both conjugate branches, with $a$ even if $f>e$ and $h$ even
if $i>e$. The squareclass of $g_1g_2$ gives $i=e$ if this
product is a square, and $i=1-e$ otherwise.
The inequalities $D>E$, $D/D'\ge E/E'$ and nonsaturation yield
\[
  a<v^{f-e}u^r,\qquad r\ge0,\qquad 2r+f>2b+i .
\]
Put $B=B_{e,k}$, $C=B_{f,k+r}$, $G=B_{i,k+b}$, and write $L=P/Q$
in lowest terms. Then
\[
  m=(B+hG)Q/P,\qquad h\mid QB,\qquad h\mid Q_0aC,
  \qquad Q_0=\operatorname{lcm}(g_1,g_2).
\]
The second follows from $h\mid m$. For the third, the signed trace
of the ordinary FP induction column is $s$ plus the two rational
terms and the trace of an integral $H$-evaluation divided by $H$.
Multiplication by $Q_0(m/h)$ makes the latter three terms integral.
Since $Q\mid Q_0$, Lemma~\ref{lem:relative-quadratic-indices} gives
$h\mid aQ_0 W_{e,f}(r)$. By the leading terms of the traces,
$s<9/16$, $t\ge\delta$, and $t\le s/3$ give
\[
 a<2\sqrt{9/(16\delta)},\qquad au^r\le4(9/(16\delta)),\qquad
 \tfrac12\le hu^b\le4/\delta .
\]
Here $\delta\ge1/600$ and $Q_0\le196$. The first two bounds give
$0\le r\le5$, hence $0<|W_{e,f}(r)|\le64952$; the third now gives
$-16\le b\le6$.

There are two local simplifications. For $(4,25)$ the two primitive
families agree. If they are both $v$-families, the residues of
$E/E',H/H'$ are not one at five since $u$ has order six modulo
five. The simultaneous column therefore splits the rational axes
from the quadratic branches. The codegree-four idempotent is
integral at five, so the codegree-25 rational factor would be
isolated, contrary to its nonintegral trace. For $(4,18)$ the
families are opposite. At three the unique $v$-family branch is
maximal, as proved in Proposition~\ref{prop:silver-two-pair-finite}.
Since $E/18$ is integral this gives $e=0$. Moreover $k$ is even:
otherwise $W$ has residue $-1$ on this branch and generates
$\mathbb F_9$ on the $v$-branch, isolating the codegree-18 axis
after the codegree-four idempotent splits.

For $k<32$ test the scalar divisibilities and mixed bound directly.
For $k\ge32$ all three indices are at least $16$. Put
$x=h(v/2)^{i-e}u^b$, $y=a(v/2)^{f-e}u^r$ and $\eta=999999/1000000$.
Lemma~\ref{lem:relative-quadratic-indices} gives the necessary bound
\[
             2\eta^2L^2(1+x^2+2y^2)\le(1+x)^2
\]
and upper bounds $L/(\eta(1+x))$, $Ly/(\eta(1+x))$ on $t,s$.
Keep $t\ge\delta$. At $(4,8)$, if $t<1/8$, require also $s\ge1/4$
or the extra $(1/2-s)^2$ contribution. The remaining conditions are
$h\mid Q_0aC$ and divisibility of $m$ by
$g_1,g_2,4/(1+e),4a/(1+f),4h/(1+i)$.
Testing one full recurrence period at each prime power and combining
the exponents by the Chinese remainder theorem excludes every tuple,
as verified in \codefile{silver_large_pairs_check.py}.
\end{proof}

\begin{proposition}\label{prop:silver-distinct-four}
Assume the setting of Proposition~\ref{prop:silver-two-pair-finite}
with two rational signed four-rows of distinct supports. If the
category has no pseudo-unitary Galois conjugate, its only invertible
object is the unit. Each signed quadratic codegree is $mvu^k$,
$8\parallel m$, and both quadratic branch projections at two
equal $\mathbb Z_2+2\sqrt2\mathbb Z_2$. If the four ordinary
complementary codegrees are all rational, $s\ge3/8$.
If $t\le1/1000$, the ordinary complement consists of two rational
four-rows of distinct supports and an irrational quadratic
codegree pair $G,G'$ having $G=nv u^j$, $j\ge1$, $8\parallel n$.
Its mass is $1/2-s$.
\end{proposition}
\begin{proof}
Proposition~\ref{prop:twelve-four-four-dyadic} gives the assertion
on the quadratic branches. Lemma~\ref{lem:two-four-supports} excludes
a nonunit invertible: $t=w=1/4$ would give $s\ge3/4$ and
$2\delta_D^2\ge(16/9)s^2>1/2$ by
Proposition~\ref{prop:silver-two-pair-finite}. The same lemma excludes
odd common rational codegrees, gives $4\mid h$ for every integral
ordinary row of norm $h$, and permits at most two rational four-rows.

Lemma~\ref{lem:real-four-six} also shows that each four-row is
rational and excludes a quadratic six-orbit.

Suppose all remaining ordinary codegrees are rational. All four-rows and six-rows are integral by Lemma~\ref{lem:real-four-six}.
An integral row has norm divisible by four, so no six-row occurs.
With at most one four-row the four masses sum to at most
$1/4+3/8=5/8$, giving $s\ge3/8$.

With two four-rows, each local factor contains one of them.
The ramified FP branch therefore lies in one with a four-row, so
the preceding dyadic calculation forces $N(D)$ to have squareclass
two. If the other two rows are rational, write their norms as $c,d$.
The global discriminant gives $cd$ a square. On the local factor
carrying these three rational branches,
$1/4+1/c+1/d\in\mathbb Z_2$. If $v_2(c)\ne v_2(d)$ this is impossible
as $c,d$ are divisible by four. At equal valuations their odd parts
are congruent modulo eight since $cd$ is a square, so
$v_2(1/c+1/d)=1-v_2(c)$. We obtain $v_2(c)=v_2(d)=3$.
Writing $c=8a,d=8b$, the same trace condition gives
$a\equiv b\equiv3\pmod4$, hence $c,d\ge24$. For a quadratic pair
of common rational codegree $c$ the same argument on $1/4+2/c$
gives $c\ge24$. Thus in the two-four-row case $s\ge5/12$.

Now assume $t\le1/1000$, so $499/1000\le w<1/2$. The short
$v$-family exponent $j=0$ for $H$ was excluded in
Proposition~\ref{prop:silver-two-pair-finite}. The next exponent
has $w=10/m$, impossible for $8\parallel m$ in this interval.
Thus $H/H'\ge u^5$ and by dominance the same holds for $D/D'$;
their squared gap ratios are at least $c_0=1681/1682$.
The all-rational case is impossible, as
$c_0((499/1000)^2+2(3/8)^2)>1/2$.
With two irrational ordinary complementary pairs their mixed
contribution is at least $(1-s)^2/2$, and
$c_0(499/1000)^2+2c_0 s^2+(1-s)^2/2>1/2$.
With a single irrational pair of mass $z$, unless the remaining
ordinary rows are two rational four-rows their reciprocal sum
is at most $3/8$. Thus $z\ge5/8-s$, and the same contradiction
follows from $c_0(499/1000)^2+2c_0s^2+(5/8-s)^2>1/2$
for $s\le5/8$ (the bound on $D$ already gives $s<5/8$).
The two four-rows pair dyadically with the two ramified
$K$ branches, so the irrational ordinary complement has
$G=nv u^j$, $8\parallel n$, by the local calculation above.
If $j=0$ and $n=8$, the projected row entries are
$a_i+2b_i\sqrt2$ with $a_i,b_i\in\mathbb Z$.
Its norm $8v$ and orthogonality to its conjugate give
$\sum b_i^2=1$. The difference from its conjugate would have just
one nonzero coordinate, contradicting orthogonality to FP.
For every other $j=0$ value we have $n\ge24$, $z=2/n\le1/12$,
$s\ge5/12$, impossible as before.
\end{proof}

\begin{corollary}\label{cor:silver-distinct-small-mass}
In the nonspherical case of Proposition~\ref{prop:silver-distinct-four}
one has $t<1/1000$ and the ordinary complementary exponent satisfies
$j\ge2$. In particular both the ordinary and signed character algebras
have two $K$ factors and two rational four-rows of distinct supports.
\end{corollary}
\begin{proof}
Write $E=mvu^k,H=nvu^j,D=Mv^d u^\ell$. Since there is no
nonunit invertible, Lemma~\ref{lem:silver-saturation} gives strict
dominance $2j+1<2\ell+d$. Suppose $t\ge1/1000$ and set
$h=m/n$, $c=m/M$, $z=j-k$, $r=\ell-k$; $c,h$ are positive
integers by central induction.

Let $B_{e,i}=\operatorname{Tr}(1/(v^eu^i))$, with recurrence
$B_{e,i+2}=6B_{e,i+1}-B_{e,i}$ and initial pairs $(2,6),(2,10)$.
At a $v$-family branch $B_{1,i}\equiv2\pmod8$. The dyadic
integrality $1/4+B_{1,i}/m\in\mathbb Z_2$ and $8\parallel m$
therefore sharpen to $m\equiv24\pmod{32}$, and likewise for $n$.
The central FP induction class on the signed pair of rational rows
contributes $Z/4$, where $Z=2$ or $3$ by total positivity.
Its contribution on the $H$-rows is $1-Z/4-s$, positive, and
$n(1-Z/4-s)\in2\mathbb Z$ since
$\operatorname{Tr}(\mathcal O_K/(vu^j))\subset2\mathbb Z$.
Consequently, writing $b=B_{1,k},g=B_{1,k+z},o=B_{d,k+r}$,
\[
 m=2(b+hg),\qquad s=co/m,\qquad
 0<m(4-Z)-4co\in8h\mathbb Z .
\]
In particular $h\mid2b$ and $h\mid4co$. Hence $h\mid4cW_{1,d}(r)$, with $W_{1,d}(r)\ne0$ by
Lemma~\ref{lem:relative-quadratic-indices}. The leading terms
give $c(v/2)^{d-1}u^r<1200$ from $s/t<600$. Together with
$v^{d-1}u^r>c$ this gives $c^2<1200\cdot2^{d-1}$, $c\le34$,
$1\le r\le4$. Now $4c|W_{1,d}(r)|<5^{14}$ and $1<w/t<500$,
so $-15\le z\le3$.

If $s<3/8$ there is
an irrational ordinary pair. Its mixed contribution is at least
$(1/2-s)^2$, or at least $(5/8-s)^2$ if $d=0$: in the latter
case the two remaining rational-codegree rows cannot both be four,
by the dyadic computation applied to $D$. With two irrational
pairs the bound $(1-s)^2/2$ is at least either displayed bound.
These expressions are used with positive parts when necessary.
For each $d,r,c$, take the positive divisors $h\equiv1\pmod4$ of
$4c|W_{1,d}(r)|$; retain $z$ satisfying strict central dominance.
For $0\le k<32$, test directly the preceding congruence, scalar
divisibilities, $t\ge1/1000$, $s\ge3t$ and the mixed inequality.
For $k\ge32$ all indices are at least $17$. The two-term formulas
of Lemma~\ref{lem:relative-quadratic-indices} bound each mass
and gap to a relative factor in $[1-10^{-6},1/(1-10^{-6})]$
of its mass at $k=32$: every exponential correction is below $5^{-34}$.
Use lower bounds on the gaps and the upper bound on $s$ in the ordinary
deficit. Neither range contains a surviving tuple among the $9908$
tested in \codefile{silver_four_distinct_check.py}. Thus $t<1/1000$.

Proposition~\ref{prop:silver-distinct-four} now applies. If the
ordinary irrational pair had $j=1$, its mass would be $z=10/n$
with $n\equiv24\pmod{32}$, so $z=5/12$ or $z\le5/28$.
For $c_0=1681/1682$ the mixed inequality would give
\[
 c_0(499/1000)^2+2c_0(1/2-z)^2+2(49/50)z^2\le1/2.
\]
The quadratic exceeds $1/2$ at $5/12$ and on $[0,5/28]$.
Thus $j\ge2$.
\end{proof}

\begin{proposition}\label{prop:folded-silver-small-mass}
Under the hypotheses of Proposition~\ref{prop:silver-two-pair-finite}, suppose that the complementary signed characters have codegree four and are Gaussian conjugates or rational with common support. If the category is not spherical, choose $E\ge E'$. Then $\operatorname{Tr}(1/E)<1/1000$.
\end{proposition}
\begin{proof}
The two rows have the forms $(1,-1,\pm1,\mp1,0,0)$ or $(1,-1,\pm i,\mp i,0,0)$, exactly as in Lemma~\ref{lem:folded-twelve-mass}. The signed involution at their second coordinate acts by the permutation $(1,g)(X,Y)$. It is the only possible nontrivial actual invertible: at $X,Y$ a putative invertible would give $t=w=1/4$ by cancellation in its signed unit coefficient, contradicting $s\ge3t$ and $s<1/2$. This last upper bound follows on the ordinary FP induction column because the short pair contributes at least $1/2$ in the signed quotient. Indeed its two evaluations sum to an even integer: the sum of their rows is twice the difference of two coordinate functionals. Positivity and the unit coefficient bound force this sum to be exactly two, so the short pair contributes precisely $1/2$.

If $D/D'=H/H'$, Lemma~\ref{lem:silver-saturation} supplies the
spherical rank-four subcategory, with nonunit invertible necessarily
$g$. Proposition~\ref{prop:folded-silver-descent} then gives a
spherical section on all of $\mathcal C$. Thus $D/D'>H/H'$.

Use $P_{e,k}=v^eu^k$ and $B_{e,k}=\operatorname{Tr}(1/P_{e,k})$, $u=3+2\sqrt2$, $v=2+\sqrt2$. As in Proposition~\ref{prop:silver-four-rational-E}, the square discriminant gives $E=mP_{e,k}$ and $H=nP_{e,j}$ in the same family. Put $D=MP_{d,\ell}$, $h=m/n$, $c=m/M$, $r=\ell-k$, $z=j-k$. Write $b=B_{e,k}$, $G=B_{e,j}$, $C=B_{d,\ell}$. The unit column and the ordinary FP column give
\[
 m=2(b+hG),\qquad (cC-b)/(2h)\in\mathbb Z .
\]
The second equality uses $n(1/2-s),nw\in2\mathbb Z$, since $\operatorname{Tr}(\mathcal O_K/(v^e u^j))\subset 2\mathbb Z$. Thus $h\mid2b,2cC$, so $h\mid2cW$, where $W=W_{e,d}(r)$ in Lemma~\ref{lem:relative-quadratic-indices}. All primitive traces are two modulo four. If $h$ is even then $v_2(m)=2$, so $h$ is odd for $e=0$ and is never divisible by four for $e=1$.

Suppose $t\ge1/1000$. Then $s/t<500$, so $c\le45$, $-1\le r\le6$ by $s/t\ge cP_{d,\ell}/(4P_{e,k})>c^2/4$ and $V^{d-e}u^r>c$, with $V=v$. The initial trace pairs $(2,6),(2,10)$ and recurrence $B_{i+2}=6B_{i+1}-B_i$ give $0<|W|<2\cdot10^6$ in this range when $v^{d-e}u^r>c$. From $2<hG/b=1/(2t)-1\le499$ and $u^z/2\le G/b\le2u^z$ we may take $-14\le z\le3$.

The same finite verification as in Theorem~\ref{thm:folded-twelve-completion} is given in \codefile{folded_silver_check.py}. For reproducibility, enumerate $(e,d,r,c,h,z)$ in the stated ranges with $h\mid2cW$, the $h$-parity conditions, $2\mid c$ for $d>e$, $v^{d-e}u^r>c$, and $2z+e<2r+d$. For $0\le k<28$ require $j,\ell\ge1$, $4\mid m$, scalar multiples of four in the unit family and of two in the other family, $2h\mid cC-b$, $3t\le s<1/2$, $t\ge1/1000$, and the mixed gap bound. At $k\ge28$ use $x=hu^z$, $y=c(v/2)^{d-e}u^r$, $T=2(1+x)$ and $\eta=999999/1000000$. The three gaps are at least $\eta/T$ times $1,x,y$ and $s\ge\eta y/T$, $t\le1/(\eta T)$ by Lemma~\ref{lem:relative-quadratic-indices}. Test the resulting necessary bounds and iterate the six recurrence residues modulo $\operatorname{lcm}(4h,4c)$ to decide the congruences. Among $228239$ parameter tuples there are nine small matches. Twenty-two tuples pass the tail inequalities, with largest congruence period $42$. After congruence tests all matches have $(e,d,r,c,h,z)=(1,1,3,1,10,2)$, with $k\ge1$ (at $k=0$ the scalar divisibility fails).

In that tuple $H,D>10^5$, $t<15/10000$, and $2905/10000<s<2906/10000$. The ratios $t=B_{1,k}/[2(B_{1,k}+10B_{1,k+2})]$ and $s=B_{1,k+3}/[2(B_{1,k}+10B_{1,k+2})]$ lie between their values at $k=1$ and their limits: each is a linear fractional function of $u^{-2k}$. These endpoint inequalities follow from $s_1=197/678$, $t_1=1/678$ and the limits $u^3/[2(1+10u^2)]$, $1/[2(1+10u^2)]$. Consequently $\delta_H^2+2\delta_D^2>(4984/10000)^2+2(2904/10000)^2>417/1000$.

All rational ordinary codegrees are even and at least four. The dyadic signed order has no rank-one factor and residue fields $\mathbb F_2$ by the argument of Lemma~\ref{lem:folded-twelve-mass}, valid also for $\mathbb Q(\sqrt2)$. The parity and perfect-trace argument there rules out odd common rational codegrees. A rational two-row requires an invertible, necessarily $g$; its value on $g$ is $-1$. If such $g$ exists its two negative rows have shape $(1,-1,\lambda,-\lambda,0,0)$, and $\lambda\ne0$ because $\lambda^2=\pm1+n\lambda$. A nonrational pair of codegree two alone exhausts the reciprocal mass.

Two ordinary complementary irrational pairs would contribute at least $(1-s)^2/2>83/1000$. For one pair of mass $z_0$, unless both rational companions have codegree four, $z_0\ge7/12-s$ and its contribution is at least $z_0^2>83/1000$. Thus both rational companions must be four and $2094/10000<z_0=1/2-s<2095/10000$. An exponent at least three in the ordinary pair contributes at least $2(49/50)z_0^2>83/1000$. The shorter possibilities have mass $2/n$ or $6/n$, neither in the displayed interval for any admissible integer scalar $n$ (the latter is a multiple of four). Finally four even rational codegrees at least four cannot have reciprocal sum in $(7094/10000,7095/10000)$. The smallest must be four, the next four or six, and the only candidate first three are $(4,4,6),(4,4,8),(4,6,6)$; substitution gives no even fourth denominator. This contradicts every ordinary decomposition.
\end{proof}

\begin{proposition}\label{prop:folded-silver-actual}
Under the hypotheses of Proposition~\ref{prop:folded-silver-small-mass},
suppose $\mathcal C$ is not spherical. Then
$s=\operatorname{Tr}(1/D)<356/1000$, and the signed involution
is an actual invertible. Its ordinary action interchanges two pairs
of simples and fixes the other two.
\end{proposition}
\begin{proof}
Put $t=\operatorname{Tr}(1/E)<1/1000$ and $w=1/2-t$.
For ratio exponents one, two and three the possible masses of the
$H$-pair are $2/n,6/n,10/n$, respectively. None belongs to
$(499/1000,1/2)$ for integral $n$. Thus both $H$ and $D$ have
squared reciprocal gap at least $288/289$ times their squared mass.
With $c_0=288/289$ the mixed bound gives
\[
      c_0\bigl((499/1000)^2+2s^2\bigr)<1/2,\qquad s<356/1000.
\]
Suppose now that the signed involution is not an actual invertible.
All ordinary rational codegrees are even and at least four as before. At two, the signed short factor cannot be isolated (in the rational case its two branches share a common reduction): on the projection its regular discriminant valuation is two, whereas its isolated perfect-trace value would be four. The two $K$-branches are ramified. Thus the only possible dyadic local ranks are six or four and two, and all residue fields are $\mathbb F_2$. The ordinary order has the same local ranks and residue fields by reduction modulo two.

Two irrational ordinary complementary pairs contribute at least $(1-s)^2/2$. With just one such pair the two smaller rational denominators cannot be $(4,4)$ or $(4,6)$. A nonrational pair of codegree four is Gaussian, forcing an actual invertible by its real part, while two rational four-rows with common support do so by their idempotents. In the other cases the two rational rows have distinct reductions (norm four and norm six cannot be congruent modulo four). Each ramified quadratic branch then has to join one of them, since neither rational axis splits off at two. This would give dyadic ranks three and three. Therefore the irrational pair has mass at least $5/8-s$ and contributes at least $(5/8-s)^2$. Both alternatives contradict the mixed bound: the minima over real $s$ of $c_0((499/1000)^2+2s^2)+(1-s)^2/2$ and $c_0((499/1000)^2+2s^2)+(5/8-s)^2$ both exceed $1/2$.

Suppose all four ordinary complementary codegrees are rational. By Lemma~\ref{lem:real-four-six}, every real four-row or six-row
is integral. A nonreal four-row forces a nonunit invertible as
above. A nonreal six-row would give $6=|\operatorname{disc}(F)|b^2$
for its imaginary quadratic character field, since there is only
one dual pair of simples. No such discriminant exists.

Lemma~\ref{lem:silver-no-invertible-mass} now gives
$s\ge5/14>356/1000$, the desired contradiction.
\end{proof}

\begin{corollary}\label{cor:folded-silver-matching}
In every nonspherical case of Proposition~\ref{prop:folded-silver-actual}, both ordinary characters taking $g$ to $-1$ have formal codegree four. Their ordinary rows agree with the signed short pair up to ordering.
\end{corollary}
\begin{proof}
They have shape $(1,-1,\lambda,-\lambda,0,0)$; in the Gaussian case $Y=X^*$ and the equation $\lambda^2=-1+a\lambda$ forces $\lambda=\pm i$. In the real case $\lambda^2=1+a\lambda$. If $a\ne0$, the two roots are conjugate quadratic units of norm $-1$. The second formal codegree $2(1+\lambda^2)=4+2a\lambda$ lies in $K$, so the roots belong to $K$. Up to signs their larger absolute value is at least $1+\sqrt2$, since every unit has absolute value a power of $1+\sqrt2$. Their total reciprocal mass is $1/2$ and the squared gap is at least $1/8$. By the previous proposition $w>499/1000$, and the squared gaps at $H,D$ exceed or equal $c_0=288/289$ times their respective squared masses. Thus
\[
 1/2\ge c_0((499/1000)^2+2s^2)+1/4,
\]
which gives $s<1/25$. The two remaining ordinary characters have mass $1/2-s>23/50$. If irrational, they contribute at least $(1/2-s)^2$ to the same mixed inequality, an impossibility. If rational, their codegrees are even and at least four and must both equal four, since $1/4+1/6<23/50$. That would leave $s=0$. Hence $a=0$ and the negative roots are $1,-1$.
\end{proof}

\begin{corollary}\label{prop:folded-silver-ordinary-irrational}
In the nonspherical alternative of Corollary~\ref{cor:folded-silver-matching},
the two remaining ordinary rows have irrational conjugate codegrees.
\end{corollary}
\begin{proof}
They take the same value at $X,Y$ and are real, since these are the
only possibly non-self-dual simples. The dyadic decomposition and
the bound $s<356/1000<5/14$ are given in the proof of
Proposition~\ref{prop:folded-silver-actual}. In the two-factor case
the negative pair cannot occupy the rank-two factor: on its projected
order the discriminant has valuation two, whereas two codegrees four
would give four. Lemma~\ref{lem:silver-common-ordinary-mass} now
excludes rational codegrees.
\end{proof}

\section{Completing the quadratic case}\label{sec:guide-pell}
The two remaining short characters have codegree four. Pell
divisibility rules out different supports, and also the case in
which the supports agree and the order is local at two. In the
remaining case, bimodules over $\mathbf1\oplus g$ supply two
new ordinary codegrees with reciprocal sum
$s=\operatorname{Tr}(1/D)$. A free involution of the simple
bimodules imposes the parity that rules out rational codegrees.

\subsection{The dyadic orders and their Pell traces}
Distinct supports and the matched local case will lead to the same
contradiction. We first determine which quadratic branches attach to
the short rows at two in the matched case.

\begin{proposition}\label{prop:folded-silver-dyadic-planes}
In the nonspherical case of Corollary~\ref{prop:folded-silver-ordinary-irrational}, write each of the two conjugate pairs of positive-character codegrees in a given ring as $m v^e u^k,n v^e u^j$, with $v=2+\sqrt2$, $u=3+2\sqrt2$. If its completed dyadic order is local then
\[
 v_2(m)=v_2(n)\in\{3,4\}\quad(e=0),\qquad
 v_2(m)=v_2(n)=3\quad(e=1).
\]
Otherwise both signed and ordinary rings have local factors of ranks four and two, both have $e=1$, and the rank-two factor in each is the maximal quadratic order. Its codegree scalar has valuation one and its character takes every simple to one modulo $\sqrt2$. The rank-four factor contains the signed physical pair or the ordinary Perron pair, respectively, and its quadratic codegree scalar has valuation two.
\end{proposition}
\begin{proof}
The two split or Gaussian negative rows have a common residual character and cannot form a dyadic factor by Proposition~\ref{prop:folded-silver-actual}. The completed local ranks agree between rings; each four-dimensional factor, when present, therefore contains the two negative rows. The squareclasses of the absolute discriminants of the rational algebras are squares, so the two positive characters belong to the same primitive family in either ring. An isolated quadratic factor is an order in $K_2=\mathbb Q_2(\sqrt2)$ and hence has discriminant valuation $3+2a$ for $a\ge0$. If its codegree is $F=n v^e u^j$, perfect trace duality gives $3+2a=v_2 N(F)=2v_2(n)+e$. Its integral idempotent has unit coefficient $\operatorname{Tr}(1/F)$, of valuation $1-v_2(n)$. Thus $a=0,e=1,v_2(n)=1$. Since the two positive pairs have total mass $1/2$, the other scalar has valuation two.

Represent one of the positive rows in its four coordinates as $(1,A+\sqrt2 B)$, where $A,B\in\mathbb Z^3$ and the first entry in each triple corresponds to the other free pair. Put $[x,y]=2x_1y_1+x_2y_2+x_3y_3$. Writing $F=c+d\sqrt2$, orthogonality to its conjugate gives
\[
 2+[A,A]=2[B,B]=c/2,\qquad [A,B]=d/2.
\]
For another pair, represented by $C,D$, the four cross equations are
\[
 2+[A,C]=[B,C]=[A,D]=[B,D]=0.
\]
An attached row, in the factor containing the negative pair, has
$A=(a,2p,2q)$ with $a$ odd. Suppose first that both pairs are
attached and write $C=(a',2p',2q')$. The identity
\[
   [B,B]=1+a^2+2(p^2+q^2)
\]
shows that $B_2,B_3$ have the same parity $t$. Put $x=p+q\pmod2$.
Reduction modulo four then gives $B_1\equiv1+x+t\pmod2$.
Write $y,t'$ for the analogous parities of the other pair.
Dividing $[B,C]=[A,D]=0$ by two gives
\[
       1+x=t(1+y),\qquad 1+y=t'(1+x)\quad\text{in }\mathbb F_2 .
\]
Thus $x=y$. If $x=1$, exactly one of $p,q$ is odd, so
$[B,B]\equiv4\pmod8$. If $x=0$, then $t=t'=1$ and $B_1$ is
even. Thus $[B,B]\equiv2\pmod8$, and $[A,B]\equiv2(p+q)=0\pmod4$.
The same conclusions hold for the second pair. Hence both rational
parts $c=4[B,B]$ have the same valuation, either three or four.
For $e=0$ this equals $v_2(m)$. For $e=1$ we have
$v_2(c)=1+v_2(m)$ and $v_2(d)=v_2(m)$. Valuation three for $c$
is impossible in that family, since $d=2[A,B]$ would be divisible
by eight.

For an isolated pair $C,D$ we have already proved $e=1$ and
$v_2(n)=1$. Thus $[D,D]$ and $[C,D]$ are odd, and
$[C,C]=2[D,D]-2\equiv0\pmod4$. The latter congruence forces
$C_1,C_2,C_3$ to have the same parity; $[C,D]$ odd forces all three
to be odd.

It remains to identify the isolated pair in each ring. The signed physical scalar is a multiple of four since $E/4$ is an induction dimension and $v^e u^k$ is primitive. Thus its pair is attached. For an ordinary fixed simple $P$ its Perron dimension $p$ satisfies $p^2/2\in\mathcal O_K$. Indeed $P$ admits a right module structure over the algebra $1\oplus g$. The pointed subcategory has trivial associator class since its action preserves the simple $P$, and a module structure on $P$ can be normalized since $H^2(C_2,\mathbb C^\times)=0$. The dimension of $P\otimes_{1\oplus g}P^*$ is $p^2/2$, by the Frobenius--Perron relative tensor product formula. Hence the rational part of $p$ is even. This rules out the isolated residual pattern at the Perron row and proves the result.
\end{proof}

\begin{theorem}\label{thm:silver-distinct-round3}
If a nonspherical conductor-eight category has two complementary
signed characters of formal codegree four, their supports agree and
the dyadic order splits into factors of ranks four and two.
\end{theorem}
\begin{proof}
Suppose that the supports differ, or that they agree and the order is
local at two. Write $u=3+2\sqrt2$, $v=2+\sqrt2$, and put
\[
 E=mv^eu^k,\quad H=nv^eu^j,\quad
 D=Mv^du^\ell,\quad G=Nv^du^i.
\]
In both cases $m,n,M,N$ are divisible by eight and $h=m/n$ is odd.
For different supports this follows from
Proposition~\ref{prop:silver-distinct-four} and
Corollary~\ref{cor:silver-distinct-small-mass}; for the local order
it follows from Proposition~\ref{prop:folded-silver-dyadic-planes}.
The same reductions give $t=\operatorname{Tr}(1/E)<1/1000$ and
$2\ell+d>2j+e$ (strict dominance over $H$).
Write $w=1/2-t$, $s=\operatorname{Tr}(1/D)$, $z=1/2-s$.
We have $s<w$: on the distinct supports the mixed
bound with $H/H',D/D'\ge u^5$ gives
$2(1681/1682)s^2\le1/2-(1681/1682)w^2$; on the matched supports
$s<356/1000$ by Proposition~\ref{prop:folded-silver-actual}.

\smallskip\noindent\emph{One common divisor.}
Put $\xi=1+\sqrt2$ and define the odd integers
\[
 C_p=\frac{\xi^p+\xi^{-p}}{d_p},\qquad
 d_p=\begin{cases}2,&p\text{ even},\\2\sqrt2,&p\text{ odd}.\end{cases}
\]
Indeed these are $Q_p$ for even $p$ and $P_p$ for odd $p$, where
$\xi^p=Q_p+P_p\sqrt2$. They are nondecreasing for $p\ge0$, since
$Q_{2n}\le P_{2n+1}=P_{2n}+Q_{2n}\le Q_{2n+2}$.
The reciprocal trace of $v^e u^k$ is $B_{e,k}=2C_{2k+e}$.
Write $K_0=2k+e$, $J=2j+e$, $L=2\ell+d$ and reduce
$h/c=h_0/a$ to lowest terms, where $c=m/M$ is integral by
central induction. Then
\[
 \frac{s}{w}=\frac{aC_L}{h_0C_J}.
\]
Consequently $M>n$ and $h_0>1$, since $L>J$ and $s<w$.

The signed mass identity $m=2(B_{e,k}+h B_{e,j})$ gives
$h\mid2B_{e,k}$. On the ordinary Perron induction column the short
signed rows contribute $Z/4$ with $Z\in\mathbb Z$, and the
$H$-orbit contributes $1-Z/4-s$. Multiplication by $n$ makes this
last trace integral, since $\operatorname{Tr}(\mathcal O_K/(v^e u^j))
\subset\mathbb Z$. Thus $h\mid4cB_{d,\ell}$. Oddness of $h_0$
gives $h_0\mid C_{K_0},C_L$, or equivalently
$u^{K_0}=u^L=-1$ modulo $h_0\mathbb Z[\sqrt2]$.
In any residue field of odd characteristic dividing $h_0$, the order of $u$
has two-adic valuation $1+v_2(K_0)=1+v_2(L)$. In particular
$K_0>0$, $d=e$, and for $g=\gcd(K_0,L)$ both $K_0/g,L/g$ are odd.
Bézout's identity gives $u^g=-1$ modulo $h_0$, hence $h_0\mid C_g$.
Now $D>E$ and $c\ge1$ imply $r=\ell-k>0$. Since $2g\mid L-K_0$,
we have $g\le r$.

\smallskip\noindent\emph{Rounding to an odd integer.}
Let $b=\ell-j>0$ and $\omega=s-1/4$. Expanding
\eqref{eq:silver-mixed} with
$\delta_f^2=\operatorname{Tr}(1/f)^2-4/N_K(f)$ gives
\[
 4\omega^2\le t-2t^2+4/m^2+4/n^2+8/M^2+8/N^2 .
\]
Here $mt=B_{e,k}$ and $nt=n/2-B_{e,j}$ are positive even
integers; moreover $M>n$. If $\omega\ne0$, then $N\omega=N/4-B_{d,i}$
is a nonzero even integer. In either case
$2\omega^2\le t+2t^2$, so $|\omega|<0.71\sqrt t$ and $s<0.273$.

Set $d_j'=u^{-2j-e}$. The closed formula for $B_{e,j}=nw$ gives
$B_{e,j}^2 d_j'\le4$, hence $d_j'\le(2t/w)^2<17t^2$.
Using the same formula on both sides of the mass ratio yields
\[
 \frac{2a u^b}{h_0}=\frac{2s}{w}\frac{1+d_j'}{1+d_\ell'},
 \qquad \left|\frac{2a u^b}{h_0}-1\right|<4\sqrt t,qquad
 tu^r=\frac{s}{c}\frac{1+d_k'}{1+d_\ell'}<1 .
\]
Indeed $d_\ell'<d_j'$ and $|2s/w-1|\le(2|\omega|+t)/w<3\sqrt t$.
In particular $h_0>u^b$, so $h_0\mid C_g\le\xi^g$ forces $g>2b$.
The even integer $2aC_g/h_0$ differs from $C_g u^{-b}$ by less
than $4u^{-b}$, since $g\le r$ and $\sqrt t\,\xi^r<1$.
But $C_g u^{-b}$ differs from the odd integer $C_{g-2b}$ by at
most $1/(2\xi)$, by the definition of $C_g$. The total error is
at most $4/u+1/(2\xi)<1$, a contradiction.
\end{proof}

\subsection{The split case and Morita duality}
In the split case the two isolated quadratic factors constrain
the induction of the unit after Morita duality. The two new linear
characters will have reciprocal codegrees summing to $s$. We first
need an ordinary lower bound on this mass.

\begin{lemma}\label{lem:split-ordinary-gap}
In the split matched case, $G/G'\ge u^3$ and $s=\operatorname{Tr}(1/D)>1/6$.
\end{lemma}
\begin{proof}
First, only $\mathbf1,g$ have integral Perron dimension. For the
integral subcategory of Perron dimension $n$, evaluate its regular
idempotent (Lemma~\ref{lem:integral-regular}) in each order. The
short characters take $g$ to $-1$, the dimension pair agrees on
the subcategory, and the last pair occurs together or not at all.
The two unit coefficients give $1/n\in\{s,1/2\}\cap\{t,1/2\}$.
Since $t<s<1/2$, we get $n=2$.

Here $G=Nvu^i$. If $i=0$, then $2/N=1/2-s$ and $s<356/1000$
give $4<N<14$. Write the row as $(1,1,x,x,y,z)$ with
$(x,y,z)=A+\sqrt2 B$. The entries of $A$ are odd by
Proposition~\ref{prop:folded-silver-dyadic-planes}. In the form
$[A,B]=2A_1B_1+A_2B_2+A_3B_3$, the norm and conjugate
orthogonality give $[A,A]=N-2$, $[A,B]=[B,B]=N/2$.
Thus $N\equiv6\pmod8$, so $N=6$ and every $A_i$ is a sign.
The last equations give $|B_1|=1$, one further sign among $B_2,B_3$,
and $A_i=B_i$ on their support. Orthogonality of their difference to the Perron row says the
signed Perron dimensions on this support sum to zero (with the
first coordinate counted twice). Orthogonality of their mean
then leaves two unit entries and a single signed fixed dimension. That dimension would be two, a contradiction.

Consequently the squared gaps of $G,H,D$ are at least $49/50$
times their squared masses. Put $w=1/2-t>499/1000$, $z=1/2-s$.
At $s\le1/6$ we have $s^2+z^2\ge5/36$. Then
$(49/50)(w^2+2s^2+2z^2)>1/2$, contradicting
\eqref{eq:silver-mixed}.
\end{proof}

The isolated signed quadratic factor now locates every copy of $g$
in the induction of the unit.

\begin{lemma}\label{prop:common-four-allocation}
In the split matched case $g$ occurs once in $F A$ for each of the
four induction labels of codegrees $H,H',G,G'$, and nowhere else.
Here $g$ denotes the positive lift of the involution.
\end{lemma}
\begin{proof}
Write $m_A=[FA:g]$. The sum $FI(\mathbf1)=\sum_Y YY^*$ contains
four copies of $g$, one for each fixed lift, and no copy of
$\varepsilon g$, which reverses the sign of the spherical dimension.
The two signed short rows have sum $(2,-2,0,0,0,0)$.
As $[FA:\mathbf1]=1$ and $[FA:\varepsilon]=0$, they contribute
$(1-m_A)/2$ to the unit coefficient of $FA$ in the signed order.

If $A$ has codegree $f$ in either sector, the signed dimension pair
contributes $\operatorname{Tr}(1/f)$, since $\dim FA=E/f$ and its
conjugate row takes the conjugate value. The remaining signed pair
is the isolated $H$-factor of
Proposition~\ref{prop:folded-silver-dyadic-planes}. Its idempotent
is integral at two. Projecting $FA$ to this factor therefore gives
\[
          (1+m_A)/2-\operatorname{Tr}(1/f)\in\mathbb Z_2 .
\]
For $f=H,H',G,G'$ the trace is $2$-adically integral. Indeed the
isolated factors in the two orders have integral idempotents, of
unit coefficients $\operatorname{Tr}(1/H)$ and $\operatorname{Tr}(1/G)$.
Each of the four multiplicities is thus odd. Their sum exhausts the
four copies of $g$.
\end{proof}

\begin{proposition}\label{prop:four-morita}
In the split matched case the Morita dual \({}_B\mathcal C_B\) for
$B=\mathbf1\oplus g$ has rank twelve and a freely acting invertible
involution. The irreducible representations of its Grothendieck algebra have
dimensions and formal codegrees as follows.
\[
\begin{array}{c|cccccc}
 \text{dimension}&1&1&2&2&1&1\\
 \text{codegree}&D&D'&G&G'&a&b .
\end{array}
\]
The pair $\{a,b\}\subset K$ is Galois-stable, $1/a+1/b=s$, and the
two corresponding labels in the center of the sphericalization occur
in the dual unit induction outside the original list.
\end{proposition}

\begin{proof}
The action on a fixed simple trivializes the associator class of
$\langle g\rangle$. Choose the group algebra multiplication on
$B$ in this gauge. The positive lift of $g$ has dimension one,
so the pivotal structure on $\langle g\rangle$ is the ordinary one.

Put $\mathcal D={}_B\mathcal S_B$ and transport its unit induction
$L$ to $Z(\mathcal S)$. For the original induction labels we claim
\begin{equation}\label{eq:four-full-center}
 [L:A_\chi]=(1+\chi(g))/2+[FA_\chi:g].
\end{equation}
Under the center equivalence $Y$ has underlying bimodule $FY\otimes B$.
A right-linear map to $B$ is a pair of maps $f_a:FY\to g^a$, $a=0,1$.
Write $m_{a,b}:g^a\otimes g^b\to g^{a+b}$ for the chosen multiplication.
Left linearity at $g$ says exactly
\[
 m_{a,1}(f_a\otimes1_g)=m_{1,a}(1_g\otimes f_a)c_{Y,g},
\]
where $c_{Y,g}:FY\otimes g\to g\otimes FY$ is the half-braiding.
Using $g^*=g$ and $\dim g=1$, this is the fixed-space equation
for the rotation $E^{(a,1)}_{Y,g}$ of \cite[Definition~2.4]{NgSchInd}.
For $Y=A_\chi$ the first space is a line. The rotation on it has trace
\[
 \sum_W S_{Y,W}[FW:g]=\chi(FI(g))/(2f_\chi)=\chi(g)
\]
by \cite[(5.4)]{NgSchInd}, since $FI(g)=gFI(\mathbf1)$
in the commutative ring and $\chi(FI(\mathbf1))=2f_\chi$.

On the second space the rotation is
precomposition with $\theta_Y^{-1}=1$
\cite[Lemma~2.7]{NgSchInd}. This proves \eqref{eq:four-full-center}.

The preceding lemma now gives four constituents of $L$ with
multiplicity one and four with multiplicity two, hence
$\operatorname{rank}(\mathcal D)=\dim\operatorname{End}(L)\ge20$.
The left and right $B$-actions make a bimodule an equivariant
object for the two translations by $g$. Thus its simple objects
on each orbit are indexed by the irreducible projective
representations of the stabilizer \cite[Corollary~2.13]{BN}.
Each of the four size-two orbits contributes two simples. Each of
the four fixed lifts contributes one or four, according as the two
stabilizer generators anticommute or commute. The central sign pairs
opposite lifts with the same cocycle, so the rank is $12,18$ or $24$.
Thus the rank is $24$ and every stabilizer representation is
one-dimensional. Each orbit downstairs has two lifts: multiplication
by the positive lift of $g$ preserves spherical dimension.
Its stabilizer maps lift uniquely on either orbit and have the same
composition scalars. Each simple bimodule downstairs therefore has
two lifts, giving rank twelve.
Twisting the left $B$-action by the automorphism $g\mapsto-g$
is tensoring on the left by the correspondingly twisted regular
bimodule. It tensors the stabilizer representation by
$(i,j)\mapsto(-1)^i$ on the left-right translation group $C_2^2$.
This sign is nontrivial on each stabilizer: these are the diagonal
$C_2$ on the two-point orbits and all of $C_2^2$ on fixed points.
Thus it changes every one-dimensional projective representation
and acts freely.

The haploid normalized special symmetric algebra $B$ gives
$\mathcal D$ its spherical structure \cite[\S3.10, after (70)]{MuleviciusOrb}.
Closing the bimodule trace with the unit and counit in
\cite[(2.16)]{KMRS} multiplies its scalar by $\dim_{\mathcal S}B=2$.
Thus
\begin{equation}\label{eq:quadratic-bimodule-trace}
        \dim_{\mathcal D}(V)=\dim_{\mathcal S}(V)/2 .
\end{equation}
It follows that the center equivalence preserves spherical dimensions
and hence twists: the underlying bimodule of $Y$ is $FY\otimes B$.
Also $\dim\mathcal D=2E$ by \cite[Theorem~2.15]{ENO}.

The central sign $J$ acts by $\varepsilon\otimes B$, which interchanges
the two lifts of each simple bimodule. Forgetting lifts respects
relative tensor products, since it preserves the separability projector.
Thus the positive quotient is $K_0({}_B\mathcal C_B)_{\mathbb C}$,
of dimension twelve. An irreducible representation of dimension $n$
and codegree $f$ there has full codegree $2f$ by summing over both
lifts. Its label $Y$ occurs in $L$ with multiplicity $n$ and has
dimension $E/f$. The scalar action of the sign is $S_{J,Y}/S_{1,Y}$,
so the old labels keep their sectors.

Lemma~\ref{prop:common-four-allocation} and \eqref{eq:four-full-center}
give multiplicities $1,1,2,2$ for the inherited ordinary labels
$D,D',G,G'$. Their matrix blocks occupy ten dimensions, so precisely
two new linear characters remain. Denote their labels by $Y_1,Y_2$
and codegrees by $a,b$. The unit coefficient gives
$1/a+1/b=1-s-2(1/2-s)=s$. The discarded pairs are Galois-stable,
hence so is $\{a,b\}$. Since $\dim Y_i=\dim_{\mathcal S}FY_i$, the dimension-character
field in Proposition~\ref{prop:quadratic-full-character} contains
$E/\dim Y_i$.
\end{proof}

\begin{theorem}\label{thm:silver-morita-end}
Every complex fusion category of rank six whose global Frobenius--Perron
dimension generates $\mathbb Q(\sqrt2)$ admits a spherical structure.
\end{theorem}
\begin{proof}
By Theorems~\ref{thm:full-quadratic-reduction}
and~\ref{thm:silver-distinct-round3}, only the split matched case
remains. Use the two new labels of Proposition~\ref{prop:four-morita},
and write $z=1/2-s$, $w=1/2-t$.

\smallskip\noindent\emph{The two remaining codegrees are rational.}
If they are irrational then they are distinct conjugate totally
positive $d$-numbers, with
\[
                  (1/a-1/b)^2\ge s^2/2
\]
by the smallest totally positive unit $u=3+2\sqrt2$.
The four labels $Y_i,JY_i$ have twist one: $Y_i$ occurs in dual
unit induction and its sign is ordinary. The original labels
occurring in the dual induction are precisely those already counted
in \eqref{eq:four-full-center}. Thus these two new labels are outside
the original list, and negative dimensions distinguish the translates
from both lists of induction labels.

At $Y_i$, the signed dimension-orbit row indexed by $\sigma$ has entry
$\sigma(\dim_{\mathcal S}FY_i)/(2\sigma(E))$, hence $1/(2\sigma(a))$
or $1/(2\sigma(b))$. Retaining these four twist-one coordinates
in the projection of \eqref{eq:silver-mixed} gives
\[
\tfrac12\ \ge\ \delta_H^2+2\delta_D^2+2\delta_G^2+2(1/a-1/b)^2 .
\]
By the gap estimates of Lemma~\ref{lem:split-ordinary-gap},
the right side is at least
\[
   (49/50)w^2+(74/25)s^2+(49/25)z^2
   \ \ge\ (49/50)(49/100)^2+\frac{74\cdot49}{100(74+49)}
   \ >\ 1/2.
\]
The middle inequality minimizes the last two terms at fixed
$s+z=1/2$. Thus $a,b$ are rational integers.

\smallskip\noindent\emph{The rational codegrees contradict the Pell norm.}
The free involution in Proposition~\ref{prop:four-morita} pairs
the entries of either linear character with equal absolute values.
Its codegree is therefore twice an algebraic integer, so $a,b$ are
even. Also $a,b\mid m$, since $E/a,E/b$ are integral and
$E=mvu^k$ with $vu^k$ primitive. In the split profile
$v_2(m)=v_2(M)=2$, $v_2(N)=1$. Writing $P_j$ for the Pell
numbers as above gives
\[
     s=2P_{2\ell+1}/M,\qquad z=2P_{2i+1}/N .
\]
Pell numbers of odd index are odd. Hence $v_2(s)=-1$ and
$v_2(z)=0$. Of the possibilities $v_2(a),v_2(b)\in\{1,2\}$,
only $(2,2)$ can give $v_2(1/a+1/b)=-1$.
Since $s>1/6$ by Lemma~\ref{lem:split-ordinary-gap}, the smaller
denominator must be four. Write the other as $4\beta$ with
$\beta$ odd. Now $z=(\beta-1)/(4\beta)$ has valuation zero,
so $\beta\equiv5\pmod8$. The reduced numerator of
$s=(\beta+1)/(4\beta)$ is $(\beta+1)/2\equiv3\pmod4$.
But it divides $P_{2\ell+1}$, whose odd prime divisors all admit
a square root of $-1$ by $Q_{2\ell+1}^2-2P_{2\ell+1}^2=-1$.
This contradiction completes the proof.
\end{proof}

\section{The dimension ratio in degree four}\label{sec:guide-quarticdesc}
Write $K=\mathbb Q(D)$ and assume $[K:\mathbb Q]=4$. The ordinary
and signed dimension characters both generate $K$: their fields contain
$D$ by Lemma~\ref{lem:central-witness} and have degree at most six. Each rational character
algebra is therefore the product of $K$ and a two-dimensional algebra.

\begin{theorem}\label{thm:all-proper-quartic}
Let $\mathcal C$ have rank six. Suppose $K=\mathbb Q(D)$ is
quartic. If $E/D$ belongs to a proper
subfield of $K$, then $\mathcal C$ has a pseudo-unitary Galois
conjugate.
\end{theorem}

\begin{proof}
A rational ratio is an integer in $(0,1]$, and
Corollary~\ref{cor:cyclic-prime-power-proper-ratio} treats cyclic $K$.
Suppose then that $K$ is biquadratic, $E<D$ and $G=\mathbb Q(E/D)$ is
quadratic. For the equality subcategory $\mathcal A$ let
$D_0=\operatorname{FPdim}(\mathcal A)$ and $F=K_{\mathcal A}$.
Theorem~\ref{thm:proper-ratio} and Lemma~\ref{lem:field-corank} give
a second quadratic field $F\ne G$ and a factorization
$D=\delta w$, $E=\delta z$ with $\delta\in F$ and $w,z\in G$.

It suffices to show $s=\operatorname{Tr}_K(1/D)\le
t=\operatorname{Tr}_K(1/E)$. Indeed, transport if necessary by
the nontrivial automorphism fixing $F$, so $z\ge z'$. This fixes
$\delta$ and leaves $D,w,s,t$ unchanged. The new ratio $x=z/w$
has $x\le x'$ by \eqref{eq:least-ratio}. The two-term identity
\eqref{eq:quadratic-trace-order}, with weights $1/z,1/z'$, gives
\begin{equation}\label{eq:fib-quartic-shift}
 \operatorname{Tr}_G(x)\le
 2\operatorname{Tr}_G(x/z)/\operatorname{Tr}_G(1/z)=2s/t\le2.
\end{equation}
Thus $x=1$ by its positive integral norm.

Over $\mathbb R$, consider in each quotient the regular idempotent of $\mathcal A$,
$e=D_0^{-1}\sum_{X\in\mathcal A}d_X b_X$. Its rank is the number of
indecomposable $\mathcal A$-components: it projects to the Perron line
in each component, and the signed action agrees with the ordinary one
by Lemma~\ref{lem:field-corank}. In either quartic factor precisely
the two embeddings fixing $F$ take $e$ to one. If there is a single
outside component, these exhaust its rank. Taking the unit coefficient
in both sectors gives
\[
 1/D_0=1/D+1/\gamma D=1/E+1/\gamma E
\]
for the nontrivial $\gamma$ fixing $F$. Tracing from $F$ gives $s=t$.

Otherwise all three components have rank two by
Lemma~\ref{lem:field-corank}. Then $\mathcal A$ is Fibonacci
by~\cite{Ostrik2}, with generator of dimension
$\varphi=(1+\sqrt5)/2$ and $D_0=2+\varphi$. For the nontrivial $\tau$
fixing $G$, the simultaneous column is supported on positive
Fibonacci lifts and has unit coefficient one. Thus
$D/\tau D=1+m\varphi$ with $m\in\mathbb Z_{>0}$, since $D$ generates
$K$. Its norm $1+m-m^2$ equals one, so $m=1$ and
$\delta=1+D/\tau D=D_0$ in Theorem~\ref{thm:proper-ratio}.

The Fibonacci generator embeds $F$ in the two-dimensional signed
complement. Its character with positive Fibonacci dimension is the
third constituent selected by $e$. If its codegree is $h$, the unit
coefficient of $e$ gives
$1/h=1/D_0-1/E-1/\gamma E=(1-t)/D_0$, since
$\operatorname{Tr}_F(1/D_0)=1$. Let $\beta\in\mathbb Z[\varphi]$
be the complementary value of the ordinary Perron induction class.
It is totally positive, and $\operatorname{Tr}_F(\beta/D_0)$ is a
positive integer: it is its unit coefficient in the Fibonacci order.
Hence the unit coefficient of the induction class gives
$1-s=(1-t)\operatorname{Tr}_F(\beta/D_0)\ge1-t$, as required.
\end{proof}

\begin{lemma}[A Galois difference in the trace dual]\label{thm:primitive-divisibility}
Let $x=E/D$ and let $\tau$ fix the character fields of the signed
complement. Then $(\tau x-x,0)$ belongs to the signed order.
If $x$ has quartic minimal polynomial $P$, then
$(\tau x-x)P'(x)/E\in\mathbb Z[x]$.
\end{lemma}
\begin{proof}
Subtract the signed class of the Perron induction label from its
integral translate in Corollary~\ref{cor:galois-columns}.
Multiplication by any polynomial in that class gives
$\operatorname{Tr}_K((\tau x-x)\mathbb Z[x]/E)\subseteq\mathbb Z$.
For $\deg r<4$, Lagrange interpolation gives
$\operatorname{Tr}_K(r(x)/P'(x))=[X^3]r(X)$. On the power basis
this coefficient pairing is integral with anti-diagonal ones and
zeros for total degree below three. Thus the trace dual is
$P'(x)^{-1}\mathbb Z[x]$ (see also~\cite{Narkiewicz}).
\end{proof}

\section{Reduction to three quartic fields}\label{sec:guide-quarticred}
\begin{theorem}\label{thm:quartic-final}
Let $\mathcal C$ be a complex fusion category of rank six, and suppose its global Frobenius--Perron dimension generates a quartic field $K$. If
\[
 K\notin\{\mathbb Q(\sqrt2,\sqrt3),\ \mathbb Q(\sqrt2,\sqrt5),\ \mathbb Q(\zeta_{16}+\zeta_{16}^{-1})\},
\]
then a Galois conjugate of $\mathcal C$ is pseudo-unitary. In particular $\mathcal C$ admits a spherical structure.
\end{theorem}

The dimension orbit accounts for four of the six characters. We first
consider the case in which the other two signed codegrees are irrational
conjugates. For rational codegrees, Lemma~\ref{lem:quartic-rational-mass-strata}
restricts the reciprocal masses, and Theorem~\ref{thm:quartic-rational-frontier}
reduces the problem to twenty-four fields. The last step combines
these two reductions.

\subsection{An irrational pair of complementary codegrees}
\label{sec:quartic-fourier-preliminaries}
Write $G=\operatorname{Gal}(K/\mathbb Q)$, $p_g=g(1/D)$,
$q_g=g(1/E)$ and $\alpha_g=p_g/q_g$. Write $s=\sum p_g$,
$t=\sum q_g$, and $\widehat y(\eta)=\sum_g\eta(g)y_g$.
The two residual signed codegrees, when distinct and irrational,
are quadratic conjugates $f,f'$. Their field $F$ lies in $K$ by
Lemma~\ref{lem:central-witness}, and $t=1-1/f-1/f'$.
We transport to maximize $E$. Then $q_1,p_1,\alpha_1$ are each
the smallest entry of their respective orbit by \eqref{eq:least-ratio}
and Perron domination. In particular $\alpha_1\le1$ and $s<1$.

\begin{lemma}\label{lem:quartic-sector-fourier}
Let $c$ be the conductor of a nontrivial character $\eta$ of $G$.
Then
\[
 |\widehat q(\eta)|^2\le \frac{\gcd(c,4)}{c},\qquad
 |\widehat p(\eta)|^2\le \frac{\gcd(c,2)}{c}.
\]
For a quadratic complementary pair $f,f'$ in either sector add
$2(1/f-1/f')^2$ to the corresponding left side at the character
cutting out its field.
\end{lemma}
\begin{proof}
In \eqref{eq:orbit-projection}, each complementary row contributes
$\frac14|2/f-2/f'|^2$ at this character. The four dimension rows
contribute $|\widehat p|^2$ or $|\widehat q|^2$.
\end{proof}
\begin{lemma}[From Fourier coefficients to a trace bound]\label{lem:quartic-pair-height}
At an involution $j$ put $M_j=|\widehat q(\eta)|$ for the nontrivial character
trivial at $j$. At a cyclic generator put $M_j=\sqrt2|\widehat q(\eta)|$
for a faithful character. If $M=\max_{j\ne1}M_j<t$, then
\[
             \operatorname{Tr}_K(E/D)<1+4s/(t-M).
\]
\end{lemma}
\begin{proof}
Fourier inversion gives $2(q_1+q_j)=t+\widehat q(\eta)$ at an
involution. At a cyclic generator with $\eta(j)=\mathrm i$, the last
term becomes $\operatorname{Re}((1-\mathrm i)\widehat q(\eta))$.
Thus each $q_j\ge(t-M_j)/4$ for $j\ne1$ by minimality of $q_1$.
Now sum $p_j/q_j$, using $p_1/q_1\le1$ and $\sum_{j\ne1}p_j<s$.
\end{proof}

\begin{lemma}\label{thm:quadratic-pair-mass}
In either sector a distinct quadratic complementary pair satisfies $1/f+1/f'\le1/2$.
\end{lemma}
\begin{proof}
The character field $F=\mathbb Q(f)$ is real quadratic, as the two
codegrees are totally positive. Put $c=\operatorname{disc}(F)$,
$L=f+f'$ and $N=ff'$. The $d$-number theorem for formal codegrees
\cite{OstrikCodegrees}, applied to $2f$, shows that $f/f'$ is a
unit, hence $L^2/N=f/f'+f'/f+2$ is an integer. The positive dimension row belongs to the quartic orbit.
Thus \eqref{eq:quad-row-trace} gives $c\mid L$ and $L\ge2c$.

Lemma~\ref{lem:quartic-sector-fourier} implies in either sector
\begin{equation}\label{eq:quad-row-fourier}
 2\left(1/f-1/f'\right)^2\le\gcd(c,4)/c .
\end{equation}
Put $u=L/N$, $h=L^2/N$ and $\kappa=(c,4)/2$. Then
$c u^2\le\kappa+4cu/L\le\kappa+2u$. If $u>1/2$, this
forces $c=5$ or $8$. Orient $f>f'$ and take the least totally
positive unit $U>1$ of the quadratic field. For some $k\ge1$,
\[
 h=2+U^k+U^{-k},\qquad u=h/L,\qquad
 (1/f-1/f')^2=u^2(1-4/h).
\]
At $c=5$, $u>1/2$ forces $h<20/3$ and hence $h=5$
by the unit table in Lemma~\ref{lem:quadratic-mass-gap}.
But $L\ge10$ then gives $u\le1/2$.
At $c=8$ put $b=2L-N$. The bound $(1/f-1/f')^2\le1/4$
gives $(N+b)^2\le N^2+16N$, so $1\le b\le7$.
Since $N\mid b^2$,
$4h=N+2b+b^2/N\le(b+1)^2\le64$.
Only $h=8$ occurs for this quadratic field. Again $L\ge16$
gives $u\le1/2$.
\end{proof}

\begin{lemma}[Short quadratic codegrees]\label{lem:quartic-short-pairs}
Let $f,f'$ be distinct quadratic conjugate codegrees complementary
to the quartic dimension orbit, in either sector. Put
$c=\operatorname{disc}\mathbb Q(f)$, $L=f+f'$, $N=ff'$ and
$u=L/N$. The table gives all possibilities in the indicated ranges,
including all ordinary fields at $u=1/2$. An empty entry in the last
column excludes that range.
\[
\begin{array}{c|c|l|l|l}
 \text{sector}&c&\text{range}&h\text{ to test}&(L,N)\\\hline
 \text{signed}&5&u>1/3&5,9,20&(10,20)\\
 &8&N\le98&8&(16,32),(24,72)\\
 &12&N\le98&6,16&(24,96),(36,81)\\
 &12&u>4/9&6,16&\\
 &13&u\ge1/4&\varnothing&\\
 &24&u\ge29/100&12,100&\\
 &28&u\ge27/100&18&\\\hline
 \text{ordinary}&8&37/100\le u\le1/2&8,36&(16,32),(96,256)\\
 &12,24&u\ge37/100&6\ (c=12)&\\
 &5&u=1/2&5&(10,20)\\
 &8&u=1/2&8&(16,32)
\end{array}
\]
\end{lemma}
\begin{proof}
The row-difference calculation \eqref{eq:quad-row-trace} applies
in either sector. Take $\kappa=(c,4)/2$ in the signed sector and
$\kappa=(c,2)/2$ in the ordinary one by Lemma~\ref{lem:quartic-sector-fourier}. As above, $h=2+U^k+U^{-k}$
and $c u^2\le\kappa+2u$. Thus $u\ge r>0$ requires
$c\le(\kappa+2r)/r^2$ and, if $c r^2>\kappa$,
$h\le4/(1-\kappa/(c r^2))$. A bound $N\le B$ gives
$h\le4+B\kappa/c$, since $(1/f-1/f')^2=(h-4)/N$.
These bounds and the unit traces of Lemma~\ref{lem:quadratic-mass-gap}
give the fourth column. For each entry take
\[
 c\mid L,\quad L\ge2c,\quad h\mid L^2,\quad
 \kappa L^2\ge c h(h-4),\quad u=h/L,\quad N=L^2/h.
\]
Use $u\le1/2$ in the signed sector. For example, at $c=24$,
$h=12$ gives $L\ge48>12/(29/100)$, while $h=100$ requires
$339<L<345$, containing no multiple of $24$.
\end{proof}

\begin{lemma}\label{lem:quadratic-compression}
Suppose the signed rational algebra in rank six is $K\times F$,
where the dimension field $K$ is quartic and the complementary
pair has distinct quadratic codegrees $f,f'\in F$. Write
$c=\operatorname{disc}(F)$, $u=1/f+1/f'$, $v=1/f-1/f'$,
$Q=\widehat{(1/E)}(\chi)$ for the character $\chi$ cutting out $F$,
and $M=\sqrt{\gcd(c,4)/c}$. There exists a trace-zero $z\in F$
such that $f(1/2-u+z)\in\mathcal O_F$ and
\[
 \left\|\begin{pmatrix}Q&\sqrt2\,v\\
                 \sqrt2\,v&2z\end{pmatrix}\right\|_{\mathrm{op}}
                \le M.
\]
In particular,
\[
 |Q|\le\sqrt2/3\quad(f,f'=8\pm4\sqrt2),\qquad
 |Q|\le\sqrt5/10\quad(f,f'=5\pm\sqrt5).
\]
\end{lemma}
\begin{proof}
Let $A,B$ be the signed induction constituents indexed by the
complementary characters. At every central simple $Y$ the center
character formula shows that $S_{A,Y}\in F$ and that conjugation
over $\mathbb Q$ sends $S_{A,Y}$ to $S_{B,Y}$. By symmetry
$S_{A,A}=S_{B,B}$ and $z=S_{A,A}-S_{A,B}$ has trace zero.
Summing the six signed induction entries in the column of $A$
gives $1/2$ by its signed unit coefficient. Four dimension labels
contribute $u$, so $S_{A,A}+S_{A,B}=1/2-u$.
The character value $2f S_{A,A}=f(1/2-u+z)$ is an algebraic
integer.

For the four dimension labels $B_\sigma$ put
$v_1=\frac12\sum_\sigma\chi(\sigma)e_{B_\sigma}$;
put $v_2=(e_A-e_B)/\sqrt2$. Their $S$-compression in this
orthonormal basis is
\[
 \begin{pmatrix}Q/2&v/\sqrt2\\v/\sqrt2&z\end{pmatrix}.
\]
Indeed $S_{B_\sigma,A}=1/(2\sigma(f))$, giving the off-diagonal
entries, and the other two entries follow from the dimension-column
formula and the definition of $z$.
Let $P_J$ be tensoring by the central sign on the simple basis.
Both $(v_i-P_Jv_i)/\sqrt2$ belong to the $T^2$-fixed subspace and
transform by the quadratic character $\chi$ under the diagonal torus:
Galois permutes the two complementary rows without signs, since
their character formulas are conjugate. Between signed induction
labels multiplication by $J$ in either entry reverses the sign
of $S$, so the compression in these two new vectors is twice
the preceding matrix. The $U(2)$-fixed projection estimate of
Lemma~\ref{lem:mass-fourier-verification}, applied to their entire
span, gives the asserted norm bound.

In both explicit pairs $u=1/2$ and $fz$ is integral. At
$f=8+4\sqrt2$, write $z=r\sqrt2$. Then $8r\in\mathbb Z$,
and the second column norm gives $|z|\le1/4$; therefore
$|z|\le\sqrt2/8$. Positive semidefiniteness of $MI$ plus and
minus the displayed symmetric matrix yields
\[
 |Q|\le M-\frac{2v^2}{M+2|z|}.
\]
Here $M=1/\sqrt2$, $v^2=1/8$, so the right side is at most
$\sqrt2/3$. At $f=5+\sqrt5$, $z=r\sqrt5$ has $10r\in\mathbb Z$,
and $v^2=1/20$. The second column norm gives $|2z|\le1/\sqrt{10}$,
so $z=0$. Thus $|Q|\le1/\sqrt5-(1/10)\sqrt5=\sqrt5/10$.
\end{proof}

\begin{proposition}\label{prop:quartic-primitive-norm-height}
Let $\mathcal C$ have rank six, let $K=\mathbb Q(D)$ be quartic,
and suppose the complementary signed codegrees are distinct
quadratic conjugates with field $F$. Make a Galois transport
so that $E$ has its largest conjugate as spherical dimension.
Assume both of the following:
\begin{enumerate}
\item Each quadratic subfield of $K$ of discriminant $8$ or $12$
is $F$ (if such a subfield exists).
\item If $K$ is cyclic, its conductor is different from $16$ and $20$.
\end{enumerate}
Then $\operatorname{Tr}_{K/\mathbb Q}(E/D)\le94$.
In particular there are only finitely many possible such dimension
fields with $E/D$ primitive.
\end{proposition}
\begin{proof}
Put $q=1/E$, $t=\operatorname{Tr}_K(q)$ and $\alpha=E/D$.
Write $f,f'$ for the complementary codegrees and put $c=\operatorname{disc}(F)$,
$L=f+f'$, $N=ff'$, $u=L/N=1-t$, and $v=1/f-1/f'$.
By Lemma~\ref{thm:quadratic-pair-mass}, $t\ge1/2$.
For the character $\eta$ cutting out $F$,
Lemma~\ref{lem:quartic-sector-fourier} gives
$|\widehat q(\eta)|^2+2v^2\le\gcd(c,4)/c$.
Since $v^2=u^2-4/N$ and $(1-u)^2+2u^2\ge2/3$,
\[
 t-|\widehat q(\eta)|\ge\tfrac12\bigl(2/3-\gcd(c,4)/c-8/N\bigr)
\]
when the right side is positive. The denominator in obtaining this
bound is $t+|\widehat q(\eta)|\le2$. For $N\ge99$ this gives a lower bound $1/12-4/99>4/94$.

For $N\le98$ and $c=8,12$ use Lemma~\ref{lem:quartic-short-pairs}
(also checked in \codefile{verify_quartic_bounds.py}).
Substituting $t=1-L/N$ and
$v^2=(L^2-4N)/N^2$ gives $t-|\widehat q|>4/94$ except at
$(8,16,32)$. There the matrix bound of
Lemma~\ref{lem:quadratic-compression} gives
$|\widehat q|\le\sqrt2/3$, so the gap exceeds $2/94$.
Moreover $s=\operatorname{Tr}_K(1/D)=1/2$ in this case.
Indeed if $\beta=m+n\sqrt2$ is the complementary value of the
signed induction class of the ordinary FP character, its unit
coefficient gives $1-s=\operatorname{Tr}(\beta/f)=(m-n)/2$.
Since $0<s<1$, equality follows.

For every other involution $j$, the quadratic character of kernel
$\{1,j\}$ has conductor $a$ distinct from $8,12$, so the separate
estimate gives $t-|\widehat q|\ge1/2-1/\sqrt5>4/94$. In a cyclic
quartic field the two adjacent pair sums satisfy
$2(q_0+q_j)\ge t-\sqrt{2\gcd(J,4)/J}$ by Fourier inversion,
where $J$ is the faithful conductor. For odd $J\ge15$ and even
$J\ge40$ this again exceeds $4/94$. The possible cyclic field
conductors below forty that are divisible by four are $16,20$:
indeed the unit groups modulo $4,8,12,24$ have exponent at most two,
and modulo $28,36$ have exponent dividing six; modulo $32$ the even
characters factor through the cyclic group of order eight, and its
order-four characters factor through modulus sixteen. (Odd or
twice-odd conductors in this range satisfy $\gcd(J,4)/J\le1/15$.)

Apart from $8\pm4\sqrt2$, all pair gaps $t-M_j$ exceed $4/94$,
using $1/2-1/\sqrt5$ also when the complementary conductor differs
from $8,12$. At $8\pm4\sqrt2$ they exceed $2/94$ and $s=1/2$.
Lemma~\ref{lem:quartic-pair-height} now gives
$\operatorname{Tr}_K(\alpha)<95$. Bounded trace bounds all elementary
symmetric functions of the positive conjugates.
\end{proof}

\begin{theorem}\label{thm:quartic-all-fields}
Let a rank-six fusion category have quartic field $K=\mathbb Q(D)$
and two distinct irrational complementary signed codegrees.
Write $a\le b\le c$ for the conductors of the three nontrivial
characters of $K$. Transport the category so that $E$ is its
largest Galois conjugate, and put $A=\operatorname{Tr}_K(E/D)$.
If $b\ge33$, then $A\le94$. If $E/D$ generates $K$, then
$A\ge 2\sqrt{\sqrt{bc/a}+4}$. 
\end{theorem}
\begin{proof}
Let $c_0$ be the discriminant of the complementary pair. When $b\ge33$
and $c_0\le32$, that pair gives the only character of conductor at
most $32$, so Proposition~\ref{prop:quartic-primitive-norm-height}
applies. When $c_0>32$ the inequality
$c_0u^2\le(c_0,4)/2+2u$ of Lemma~\ref{lem:quartic-short-pairs}
gives $u\le1/4$ (the least even discriminant here is $40$).
Every pair amplitude is at most $1/\sqrt2$. Since
$3/4-1/\sqrt2>4/94$, Lemma~\ref{lem:quartic-pair-height} gives $A\le94$.
For the quadratic subfield $F$ of least conductor, the
conductor--discriminant formula gives
$N_{F/\mathbb Q}(\mathfrak d_{K/F})=bc/a$. If $x=E/D$ generates
$K$, then $K=F(x)$ and the totally positive integer $x$ has norm at
least one. Lemma~\ref{lem:relative-trace-discriminant} gives
$A=\operatorname{Tr}_K(x)\ge2\sqrt{\sqrt{bc/a}+4}$.

\end{proof}

\begin{corollary}\label{cor:quartic-five-unbounded}
Let $F$ be the field generated by two distinct signed complementary
codegrees. If its discriminant $c_0$ is different from $5,8,12$, then
$t=\operatorname{Tr}_K(1/E)>71/100$. Under the same hypothesis, unless $c_0=24,28$ one has
$t\ge3/4$.
\end{corollary}
\begin{proof}
The bound $c_0u^2\le(c_0,4)/2+2u$ gives $u=1-t\le1/4$ for
odd $c_0\ge17$ or even $c_0\ge40$. For the remaining
$c_0=13,24,28$ use Lemma~\ref{lem:quartic-short-pairs}.
\end{proof}

\begin{lemma}\label{thm:quartic-five-to-three}
Suppose the two distinct signed complementary codegrees $f,f'$ generate
$\mathbb Q(\sqrt5)$. Then either $t=\operatorname{Tr}_K(1/E)\ge2/3$
or the pair is $5\pm\sqrt5$ and
$\operatorname{Tr}_K(1/D)=t=1/2$. In the latter case, with unnormalized
$p_i=i(1/D)$ and $q_i=i(1/E)$, the moments
\[
 h_j=\sum_i q_i\alpha_{i+j},\qquad
 K_j=\sum_i p_i p_{i+j}/q_i
\]
are at least $1/2$; $2K_j$ is odd when $j$ exchanges $f,f'$.
\end{lemma}
\begin{proof}
Lemma~\ref{lem:quartic-short-pairs} gives the stated alternative.
For the exceptional pair orient
$f=5+\sqrt5$ and write the complementary evaluation of the ordinary
FP column as $m+n(1+\sqrt5)/2$. Its contribution to the unit
coefficient is then $m/2=1-\operatorname{Tr}_K(1/D)$, which belongs
to $(0,1)$. Hence $m=1$ and both masses equal $1/2$.

Put $\phi=(1+\sqrt5)/2$ and let $\beta$ be the complementary value
of the ordinary Perron induction class $V$. Positivity at both
embeddings gives $\beta=1$ or $\phi^2=f/f'$. By
Corollary~\ref{cor:galois-columns} the integral translate of $V$ has
complementary value $j(\beta)$, and the raw induction class $V_j$
has value $f j(\beta/f)$. Thus the complementary contributions to
the unit coefficients defining $h_j$ and $K_j$ are respectively
$\operatorname{Tr}_F(j(\beta)/f)$ and
$\operatorname{Tr}_F(\beta j(\beta/f))$, where $F=\mathbb Q(\sqrt5)$.
Both are $1/2$ or $1$; the second is $1/2$ whenever $j$ exchanges
$f,f'$, since $N_F(\beta)=1$. Subtract from the integral unit
coefficients and use positivity of both moments.
\end{proof}

\begin{lemma}[Coordinates for the trace bounds]\label{lem:quartic-trace-box}
Let $x$ be a totally positive integer of trace $A$ in a quartic field $K$.
If $K$ is biquadratic with quadratic discriminants $a,b,c$, then
\[
 4x=A+i\sqrt a+j\sqrt b+k\sqrt c,\qquad
 i,j,k\in\mathbb Z,\qquad ai^2,bj^2,ck^2<A^2.
\]
If $K$ is cyclic and contains $\sqrt2$, then
$K=\mathbb Q(y_+)$, $y_\pm=\sqrt{m(2\pm\sqrt2)}$, for odd
squarefree $m>0$. Its conductor is $16m$. In the integral basis
$1,\sqrt2,y_+,y_-$ the coordinates $(r,i,j,k)$ satisfy
\[
       4r=A,\qquad 4r^2+8i^2+8m(j^2+k^2)<A^2 .
\]
\end{lemma}
\begin{proof}
In the biquadratic case the three relative traces give the integral
coefficients, and positivity bounds each signed sum of conjugates by
$A$. For a cyclic quartic the parameterization and integral basis are
those of Hudson and Williams~\cite[pp.~145--146, Theorem~(i)]{HudsonWilliams}
with $D=2$, $B=C=1$. The trace Gram matrix has diagonal $4,8,8m,8m$
and zero off-diagonal entries. Its determinant is $8(16m)^2$, so the
conductor--discriminant formula gives $16m$. Finally
$\operatorname{Tr}(x^2)<A^2$ by positivity.
\end{proof}

\begin{theorem}\label{thm:quartic-high-quadratic}
Suppose $\mathbb Q(D)=K$ is quartic and the two complementary signed codegrees are distinct and irrational. A Galois conjugate is pseudo-unitary in either of the following cases:
\begin{enumerate}
\item $K$ has no quadratic subfield of discriminant $5,8$ or $12$.
\item $K$ is biquadratic, contains $\sqrt5$, and contains neither $\sqrt2$ nor $\sqrt3$.
\end{enumerate}
\end{theorem}
\begin{proof}
By Theorem~\ref{thm:all-proper-quartic} we may take $x=E/D$ primitive after transport maximizing $E$. Write $A=\operatorname{Tr}(x)$, $s=\sum p_i<1$, $t=\sum q_i$, and $u=1-t$. We apply Lemma~\ref{lem:quartic-pair-height} with the separate Fourier
bounds of Lemma~\ref{lem:quartic-sector-fourier}.

In (1) the complementary quadratic field has discriminant at least thirteen, so Corollary~\ref{cor:quartic-five-unbounded} gives $t>71/100$. On a quadratic character take $M=1/\sqrt6$ (the smallest even fundamental discriminant above twelve is twenty-four). In a cyclic field the adjacent pair bound is at most $1/\sqrt5$: for odd $J$ one has $J\ge13$, and the even quartic conductors below forty can only be $16,20$, whose quadratic subfields are excluded here. It follows that $A\le14$ in a biquadratic field and $A\le16$ in a cyclic field.

Write $P=X^4-AX^3+BX^2-CX+H$ for the minimal polynomial. For these small heights one may enumerate the polynomials directly. Positivity gives $0<B\le3A^2/8$, $0<C\le A^3/16$, $256H\le A^4$, and the $d$-number divisibilities give $H\mid A^4$, $H^2\mid B^4$, $H^3\mid C^4$. Keep irreducible polynomials with four positive real roots, one less than one, and abelian Galois group. The cubic resolvent is
\[
 Y^3-BY^2+(AC-4H)Y-(A^2-4B)H-C^2.
\]
In the biquadratic case its roots $r$ are rational integers. The pair sums of the four roots give the three signed differences of complementary pair sums. Their squares are $A^2-4(B-r)$. The nonzero differences generate the quadratic subfields present in $x$ and hence generate $K$. In the cyclic case the discriminant of $P$ generates the quadratic subfield on taking its square root, since a generator of order four acts by an odd permutation. The following five polynomials remain. In the cyclic rows the entry in the right column is the quadratic conductor.
\[
\begin{array}{rrrr|l}
 A&B&C&H&\text{quadratic conductors}\\\hline
 9&19&9&1&13\\
 12&45&54&1&28,44,77\\
 13&40&13&1&17\\
 14&38&28&4&21,60,140\\
 15&60&72&9&13
\end{array}
\]
For each polynomial fix its least root at the identity and allow every
ordering of the other three. Impose $p_i=x_iq_i$, the minimum conditions,
$s,t\le1$, $t\ge71/100$ and the Fourier inequalities of
Lemma~\ref{lem:sextic-matrix-compression}. In the cyclic cases the odd
quadratic conductor $d$ can replace the faithful conductor: the latter
is an odd multiple of $d$ or a multiple of $4d$.

In (2) all characters except the one of conductor five have separate amplitude at most $1/\sqrt6$. If the complementary character field differs from $\mathbb Q(\sqrt5)$ then $t>71/100$. Otherwise Lemma~\ref{thm:quartic-five-to-three} gives either $t\ge16/25$ or $f,f'=5\pm\sqrt5$. In the former case $M=1/\sqrt5$ gives $A\le21$. For the exceptional pair $t=s=1/2$ and Lemma~\ref{lem:quadratic-compression} improves its own signed Fourier amplitude to $\sqrt5/10$. Taking $M=1/\sqrt6$ gives $A\le22$.

In the biquadratic coordinates of Lemma~\ref{lem:quartic-trace-box},
primitivity requires two nonzero coefficients, so the smaller
discriminant other than five is $<22^2$. The resulting 120 fields
give ten primitive ratios. Test each in both mass ranges, using the
sharper conductor-five amplitude at half mass.

Here and in the remaining quartic tests the Fourier disks and their
mixed bounds have rational supporting half-spaces. The root intervals
are rounded outwards. If the resulting system is $Bq\le d$ with
$q\ge0$, $\sum q_i\le1$, an integer vector $w\ge0$ excludes it when
\begin{equation}\label{eq:quartic-separator}
 \min\{0,\min_i(B^{\mathsf T}w)_i\}>d^{\mathsf T}w .
\end{equation}
Indeed $(B^{\mathsf T}w)^{\mathsf T}q$ is at least the left-hand side.
The \codecheck{generic_quartic/verify_quartic_fields.py}{polynomial check}
verifies the \codecheck{generic_quartic/quartic_field_duals.json}{30 orderings
in (1)}, and the \codecheck{generic_quartic/golden_biquadratic.py}{biquadratic
check} verifies the \codecheck{generic_quartic/golden_biquadratic_duals.json}{20
mass and ratio combinations in (2)}.
\end{proof}

\begin{theorem}\label{thm:quartic-silver-fields}
Let $\mathcal C$ have rank six and quartic dimension field $K=\mathbb Q(D)$,
with two distinct irrational complementary signed codegrees. If $K$
contains $\sqrt2$ or $\sqrt3$, a Galois conjugate is pseudo-unitary except possibly when $K$ is one of
\[
 \mathbb Q(\sqrt2,\sqrt3),\quad \mathbb Q(\sqrt2,\sqrt5),\quad
 \mathbb Q(\sqrt3,\sqrt5),\quad \mathbb Q(\zeta_{16})^+.
\]
\end{theorem}
\begin{proof}
Write $x=E/D$ after transport maximizing $E$, and put
$p_g=1/g(D)$, $q_g=1/g(E)$, $s=\sum p_g$, $t=\sum q_g$. By
Theorem~\ref{thm:all-proper-quartic} it suffices to exclude primitive
$x$. Let $F$ be the complementary field of discriminant $c_0$,
$u=1-t=1/f+1/f'$, $L=f+f'$, $N=ff'$.

\smallskip\noindent\emph{The trace bounds.}
Outside the four stated fields, $K$ contains just one of
$\sqrt2,\sqrt3$ and does not contain $\sqrt5$. Thus $c_0=5,24$
are impossible: $\sqrt6$ together with either $\sqrt2$ or $\sqrt3$
generates both. At $c_0=8,12$ apply
Proposition~\ref{prop:quartic-primitive-norm-height}. All other
$c_0$ except $28$ give $t\ge3/4$ by
Corollary~\ref{cor:quartic-five-unbounded}. Every pair amplitude
is at most $1/\sqrt2$, so the trace is at most $94$.
At $c_0=28$ one has $t>73/100$. The field is then
$\mathbb Q(\sqrt3,\sqrt7)$ or $\mathbb Q(\sqrt2,\sqrt7)$;
the maximal amplitudes $1/\sqrt3,1/\sqrt2$ give $A\le94,176$
respectively. For the latter field and for cyclic conductor $48$
we also use
\[
             t-|\widehat q(\chi_8)|\ge2/94.                 \tag{Q1}
\]
This follows from Proposition~\ref{prop:quartic-primitive-norm-height}
when $c_0=8$, and from $73/100-1/\sqrt2>2/94$ otherwise.
\smallskip\noindent\emph{Enumerating fields and ratios together.}
By Lemma~\ref{lem:quartic-trace-box} at least two radical coefficients
of a primitive biquadratic ratio are nonzero. The second smallest
quadratic discriminant is therefore $<A^2$. For $A\le94$, enumerate
$\mathbb Q(\sqrt d,\sqrt n)$ with $d=2,3$ and squarefree
$1<n<94^2$, including the field with cap $176$. This gives
$4\,358$ conductor triples, nine with second conductor below $33$. A real cyclic quartic cannot
contain $\sqrt3$, since the quadratic character at three is nontrivial
on $-1$ and cannot be a square. With quadratic subfield
$\mathbb Q(\sqrt2)$, the cyclic coordinates of the same lemma give
$8m<A^2$ by primitivity. There are 446 odd squarefree $m>1$ to test,
with $m=1$ one of the stated exceptions.

\smallskip\noindent\emph{The separating inequalities.} In the preceding boxes retain primitive totally positive integral
$d$-numbers with least conjugate below one. If
$P=X^4-AX^3+BX^2-CX+H$ is their polynomial, the divisibilities are
$H\mid A^4$, $H^2\mid B^4$, $H^3\mid C^4$. The
\codecheck{verify_quartic_silver_fields.py}{enumeration} leaves
662 biquadratic and 74 cyclic polynomials. For each impose
$p_g=g(x)q_g$, $1/2\le t\le1$, $s\le1$, the two coordinate minima
and the separate and mixed Fourier inequalities, with (Q1) in its two
fields. Every system has \codecheck{quartic_silver_field_weights.json}{integer
weights} satisfying \eqref{eq:quartic-separator}.
\end{proof}

\subsection{Rational complementary codegrees in quartic fields}

\begin{lemma}[Codegree three at an inert prime]\label{lem:quartic-three-induction}
If every prime of the quartic dimension field $K$ above three has
even residue degree, neither sector has a complementary character
of codegree three.
\end{lemma}
\begin{proof}
A rational row $\rho$ of codegree three needs a totally ramified
local partner by Lemma~\ref{lem:silver-local-partner}. It must be the
other rational row $\eta$, say of codegree $h$. The two rows agree
modulo three and the unit coefficient of their common local
idempotent gives $1/3+1/h\in\mathbb Z_3$, hence $h/3\equiv-1\bmod3$.

Write $m=\rho(FA_\rho)$ and $n=\rho(FA_\eta)$.
Central symmetry gives $\eta(FA_\rho)=(h/3)n$, so $m+n\equiv0\bmod3$.
The four dimension labels of the chosen sector have the same integral
value $b$ on $\rho$: it is one on signed simultaneous columns, and
on ordinary columns Galois permutes their normalized rows while
fixing $\rho$. In the ordinary case
$b=3\operatorname{FPdim}(FA_\rho)/D>0$ by symmetry. The row
$6S_{A_\rho,-}$ has squared norm $36$ and includes the vacuum entry
one besides these four entries, so $b\le2$. In either case
the unit coefficient of $FA_\rho$, computed by symmetry, gives
$4b+m+n=3$, contradicting $3\nmid b$.

A quadratic orbit of common codegree three has field
$\mathbb Q(\sqrt{-3})$ by Lemma~\ref{lem:quadratic-row-difference}.
At three both the discriminant of $K$ and any norm from $K$
have even valuation, since every residue degree is even.
The quadratic discriminant contributes one and the two codegrees
contribute two, contrary to Lemma~\ref{lem:trace-order-conductor}.
\end{proof}

\begin{lemma}\label{lem:quartic-rational-mass-strata}
In rank six suppose $K=\mathbb Q(D)$ is quartic and the two signed
complementary codegrees are rational. Put
$s=\operatorname{Tr}_K(1/D)$ and $t=\operatorname{Tr}_K(1/E)$.
Then $s\le t$, $t\ge1/2$, and either $s\ge1/2$, or $s=1/3$ with
$t\ge2/3$. The latter case requires two ordinary codegrees three.
At $s=t$ the complementary signed evaluations of every forgotten
ordinary dimension induction label are one.
\end{lemma}
\begin{proof}
Corollary~\ref{cor:galois-columns} gives $s\le t$ and the equality
assertion; the raw induction columns have the same complementary
evaluations as their integral Galois translates because the
simultaneous column evaluates at one. It remains to bound $s$ and $t$.

Neither algebra has a complementary codegree two. A signed
character of codegree two must be rational, as a second conjugate
of the same codegree would exhaust the mass. Its central row has
four entries $1/4$ on the signed dimension orbit and four
$\beta/4$ on the ordinary orbit, with $\beta\in\mathbb Z_{>0}$.
The distinct $J$-partners have the same magnitudes. Their squared
norm is already at least one, precluding any other support. Yet
the sum on all six signed induction labels is $1/2$ by the
coefficient of the unit, while the four signed dimension entries
alone sum to one. For an ordinary rational codegree two, the
orthogonal character row is $(1,-1,0,0,0,0)$ up to permutation:
its norm is two and orthogonality to FP forces the negative sign
and the second FP dimension to be one. The corresponding ordinary
basis element $g$ is $1-2e_\chi$, hence an involutive invertible
with a one-dimensional negative eigenspace. Its positive spherical
lift acts by the same permutation since it sends the positive lift
of each simple to the positive lift of its unique product. Thus in
the signed order $(1-g)/2$ is the primitive idempotent for a
rational character, with unit coefficient $1/2$, giving the
excluded signed codegree. A common ordinary codegree two already
exhausts the mass. Codegrees below two on rational rows are
excluded by positivity of the residual mass.

A pair of rational characters cannot have codegrees $(3,4)$ or
$(3,5)$, in either algebra. For the corresponding quartic
codegree $H=D$ or $E$, Lemma~\ref{lem:trace-order-conductor} gives
$N(H)gh=|\operatorname{disc}(K)|I^2$, where $I$ is the order index.

At three the quartic field is tame, and
its inertia index is at most two: tame inertia in an abelian
extension at three has order dividing two. Hence the field
discriminant has even valuation, and the $d$-number property
makes $v_3(N(H))$ a multiple of $4/e$, which is also even.
Both proposed products have odd valuation.

We record a trace inequality valid for the two complementary
sectors. For the signed complementary labels $A,B$ put $v=e_A+e_B$ and $\Phi(x)=6x^2-8x+3$.
The vector $Sv$ is real, since $v$ is invariant under duality.
For any induction label of either sector with codegree $f$,
the sum of its six signed entries is $1/2$ by adjunction and
the unit coefficient. Its four signed dimension entries sum to
$\operatorname{Tr}_K(1/f)/2$ by the center formula. Consequently
$Sv$ has value $(1-t)/2$ at each of the four signed dimension
labels and its two complementary signed values sum to $2t-1$;
the analogous ordinary values have $s$ in place of $t$.
In the complementary sum use $\sum_{\text{two }f}
\operatorname{Tr}_K(1/f)=4(1-t)$ or $4(1-s)$, respectively.
Including the $J$-partners shows that the squared norms on these
signed and ordinary coordinates are $M_s\ge\Phi(t)$ and
$M_o\ge\Phi(s)$ by Cauchy's inequality. In the quadratic form
$(Sv)^tT(Sv)$ the signed pairs cancel, the ordinary pairs
contribute $M_o$, and the remaining real part is at least
$-(2-M_s-M_o)$. The twist fixes $v$ and the center has central
charge one, so $v^tSTS v=v^tSv=2t-1$. Therefore
\[
                  \Phi(t)+2\Phi(s)\le2t+1.
\]

If the ordinary complementary codegrees are rational, reciprocal
mass greater than half thus leaves only $(3,3)$ and $s=1/3$,
also allowing a quadratic character pair of common codegree three.
The displayed inequality at $s=1/3$ reads
$(3t-2)(t-1)\le0$, so $t\ge2/3$.

Distinct quadratic ordinary codegrees have $s\ge1/2$ by
Lemma~\ref{thm:quadratic-pair-mass}. Finally $t<1/2$ would give
signed codegrees $(3,3)$, hence $s=t=1/3$ by $s\le t$, contrary
to the preceding bound $t\ge2/3$.
\end{proof}
\begin{lemma}\label{lem:quartic-inert-congruence}
Suppose an odd prime $p$ is inert in a quadratic subfield of $K$.
If two rational complementary characters in either sector have
codegrees $g,h$ of odd $p$-adic valuation, their rows are congruent
modulo $p$. In particular $p^2\mid g+h$.
\end{lemma}
\begin{proof}
Distinct residue characters belong to distinct factors of the
completed order. Neither factor can consist of a rational branch
alone, since its idempotent has nonintegral unit coefficient $1/g$
or $1/h$. Each therefore needs a branch of $K$. Their residue degrees are even
by inertness, so degree four allows exactly two unramified quadratic branches. On each
factor the unit-coefficient form is unimodular. Passing to the
regular trace multiplies its discriminant by its rational codegree
and the norm from an unramified quadratic field of the dimension
codegree. This gives odd valuation, whereas every order in an
unramified algebra has even discriminant valuation. The rows are
therefore congruent. Their squared distance is $g+h$ by orthogonality.
\end{proof}

\begin{theorem}\label{thm:quartic-rational-frontier}
Let a rank-six category have quartic dimension field $K=\mathbb Q(D)$
and rational complementary signed codegrees. A Galois conjugate is pseudo-unitary outside
the following list of twenty-four fields:
\[
\begin{array}{c|l}
 K&\text{parameters}\\\hline
 \mathbb Q(\sqrt2,\sqrt d)&d=3,5,7,11,13,15,17,19,21,23,29,33,37,41\\
 \mathbb Q(\sqrt3,\sqrt d)&d=2,5,7,10,11,13\\
 \text{cyclic quartic of conductor }J&J=16,20,48,80,112 .
\end{array}
\]
The first two rows have one common field. The cyclic fields in this
list have quadratic subfield $\mathbb Q(\sqrt2)$ except at $J=20$,
where it is $\mathbb Q(\sqrt5)$; each specification is unique.
\end{theorem}
\begin{proof}
By Theorem~\ref{thm:all-proper-quartic} take $x=E/D$ primitive at
the conjugate maximizing $E$. Put $p_i=i(1/D)$, $q_i=i(1/E)$,
$s=\sum p_i$, $t=\sum q_i$ and $A=\operatorname{Tr}(x)$. We use
Lemma~\ref{lem:quartic-rational-mass-strata}, in particular
$1/3\le s\le t<1$, $t\ge1/2$ and $s+t\ge1$. The two-orbit
compression gives $\sqrt c|Q|+c|P|^2\le2$ at every nontrivial
character of conductor $c$, together with the separate bounds
$c|P|^2\le(c,2)$ and $c|Q|^2\le(c,4)$. Here and below the
identity coordinates of $p,q,x$ are minima, with $p_0\le q_0$.

If $\sqrt2\in K$ then $s\ge1/2$: three is inert in $\mathbb Q(\sqrt2)$,
so Lemma~\ref{lem:quartic-three-induction} excludes $(3,3)$.
\smallskip\noindent\emph{Heights.}
For a pair $\{0,j\}$ with quadratic annihilating character of
conductor $c$, let $y=q_0+q_j$ and let $B$ be the sum of the
ordinary Fourier amplitude bounds for the other two characters.
Since $p_0\le q_0\le y/2$, Fourier inversion gives
$|Q|\ge t-2y$ and $|P|\ge s-2y-B$. A tangent to the mixed bound,
with parameter $r>0$, now gives
\[
 y\ge\frac{\sqrt c\,t+2cr(s-B)-cr^2-2}{2\sqrt c+4cr}.
\tag{R1}
\]
The separate bound gives $2y\ge t-\sqrt{(c,4)/c}$. For an adjacent
pair in a cyclic field of faithful conductor $J$ replace the
subtracted term by $\sqrt{2(J,4)/J}$. Thus, apart from the quadratic
characters of conductors $8,12$ and adjacent pairs at $J=16,20$,
$2y\ge t-1/\sqrt5$. There are no other even faithful conductors
below $40$ and the least odd one is at least $13$. If every pair has
this bound, $q_j\ge y/2$ gives
$A<1+4s/(t-1/\sqrt5)\le21+8\sqrt5<39$.

For the exceptional quadratic characters outside fields containing
$\sqrt5$, the mass lemma and the preceding local argument give
$s\ge1/2$ except possibly at $c=12$, $s=1/3$, $t\ge2/3$; the latter
case gives $A\le15$ at this pair. For $s\ge1/2$ set $r=1/2-B$
and
\[
 \ell=\frac{\sqrt c/2+cr^2-2}{2\sqrt c+4cr}.
\]
When $r,\ell>0$, (R1) and $t\ge s\ge1/2$ give $y\ge2s\ell$.
Thus the trace cap at this pair is $A<1+1/\ell$. The maximum cap
over all three pairs suffices. Here is the complete infinite-range
check. For $\mathbb Q(\sqrt2,\sqrt d)$ take $d$ odd squarefree;
$B=3/(2\sqrt d)$ for $d\equiv1\pmod4$ and
$B=(1/\sqrt2+1/2)/\sqrt d$ otherwise. For
$\mathbb Q(\sqrt3,\sqrt d)$ take $d$ squarefree prime to three;
$B=\sqrt{(c_d,2)/c_d}+\sqrt{(c_{3d},2)/c_{3d}}\le3/(2\sqrt d)$,
where $c_n=\operatorname{disc}\mathbb Q(\sqrt n)$. Compute (R1)
up to $d=54,34$ respectively; beyond these bounds substitute
$3/(2\sqrt{55})$, $3/(2\sqrt{35})$ and use monotonicity. In cyclic
fields with quadratic subfield $\mathbb Q(\sqrt2)$ the faithful
conductor is $16m$ for odd squarefree $m$, and $B=2\sqrt{2/J}$.
The \codecheck{quartic_rational_frontier/higher_degree_quartic_bounds.py}{conductor comparison} checks $J\le176$ and then uses monotonicity. This gives
exactly the displayed twenty-four exceptions; the caps elsewhere in
these three families are respectively $264,437,206$. When the second
smallest character conductor is at least $1000$, both other amplitudes
at an exceptional character total at most $2\sqrt{2/1000}$, giving
$\ell>1/38$ and $A\le38$.

\smallskip\noindent\emph{The finite field inventory.}
Write $a\le b\le c$ for the three conductors. The relative
discriminant bound from Lemma~\ref{lem:relative-trace-discriminant}
gives $A\ge2\sqrt{\sqrt{bc/a}+4}$ and so $b\le\sqrt a(A^2/4-4)$.
For $a>16$ every pair has effective separate amplitude at most
$2/\sqrt a$ (for a cyclic adjacent pair use $J\ge a$ when $J$ is odd and $J\ge2a$ when $4\mid J$). Thus
$A<1+2/(1/2-2/\sqrt a)$. At $a\ge88$ this gives $A<8$, whereas
the relative bound gives $A>7$. For the positive fundamental
$a<88$ and $b\ge1000$ test $5\le A\le38$, imposing
$a(A-5)^2<16(A-1)^2$ if $a>16$, and then the bound on $b$.
This gives $b\le1471$. The field inventory may also use
$ab\le25007$ outside the stated exceptions: if $a=8,12$ and
$b<1000$ use $b\le999$; otherwise the same bound $A\le38$
applies unless $K$ is already listed. In the biquadratic case
loop over the two smallest fundamental discriminants and multiply
their quadratic characters to recover the third. Keep the
field-specific height and relative discriminant tests above.

For a cyclic field write $J$ for its conductor and $d$ for the
radicand of the quadratic subfield. The local components of a
primitive even quartic character enumerate $(a,J)$: the odd part of
$J$ is squarefree, $v_2(J)\in\{0,2,3,4\}$, and the odd local orders
are two or four. The quadratic conductor $a$ is the product of the
odd primes of local order four, times eight if $v_2(J)=4$.
The local parities multiply to one.

\smallskip\noindent\emph{Enumerating the integral ratios.}
Use Lemma~\ref{lem:quartic-trace-box} for the 710 biquadratic
fields. For each cyclic pair choose $h,k>0$ with $d=h^2+k^2$.
All cyclic quartics with this quadratic subfield are rational twists
of the same field:
\[
 K=\mathbb Q(\sqrt d,Y),\quad
 Y^2=m(d+h\sqrt d),\quad Z^2=m(d-h\sqrt d).
\]
Indeed the quadratic radicand has norm $d$ times a square
\cite[pp.~145--146]{HudsonWilliams}. Its quotient by $d+h\sqrt d$
has square norm $r^2$ and equals a rational multiple of
$(1+z/r)^2$ if the quotient is $z$ and $r>0$.
At odd primes outside $d$, the two factors $d\pm h\sqrt d$
have even valuations since their product is $dk^2$ and they are
relatively prime there. The squarefree twist $m$ therefore has the
same odd prime factors as $J$ outside $d$. Trying the divisors of
$2d$ for the remaining primes, up to the twist by $d$, covers every
field.

The trace Gram diagonal of $1,\sqrt d,Y,Z$ is $(4,4d,4md,4md)$,
giving index $I=\sqrt{256m^2d^3/(aJ^2)}$ in the maximal order.
Requiring $I$ integral leaves 232 presentations. Relative traces
write $x=(A+u\sqrt a)/4+(vY+wZ)/I$. Positivity bounds the lattice
by $au^2<A^2$ and $16md(v^2+w^2)<I^2(A^2+au^2)$.

In both enumerations retain primitive classes with monic integral
polynomial $X^4-AX^3+BX^2-CX+H$, four positive roots, least root
below one, and $H\mid A^4$, $H^2\mid B^4$, $H^3\mid C^4$. The
biquadratic norm must be a square or twice a square: odd inertia
indices divide two and the ideal is invariant. A dyadic odd norm
valuation requires ramification index four. The respective counts
are $96\,433\,626$ and $7\,581\,015$ canonical tuples, leaving 258
and 281 ratios or outer presentations.

For the final exclusions Corollary~\ref{cor:galois-columns} gives
$h_j=\sum_i q_i x_{i+j}\ge s$.

Use the probability variables $z_i=(p_i+q_i)/(s+t)$; then
$q_i/(s+t)=z_i/(1+x_i)$, $p_i/(s+t)=x_i z_i/(1+x_i)$ and
$1/4\le s/(s+t)\le1/2$. The separate and mixed Fourier estimates
continue to hold after division by $s+t\ge1$. For each candidate
impose these inequalities, the two minimum conditions, and
$h_j\ge s$. Outward rational root intervals and supporting
half-spaces give a system $Bz\le d$. The companion
\codefile{quartic_rational_frontier/verify_rational_frontier.py} rebuilds both inventories,
enumerates the boxes and checks nonnegative integer weights $w$
with $\min_i(B^{\mathsf T}w)_i>d^{\mathsf T}w$ for all $258+281$
candidates. This contradicts $\sum z_i=1$ and proves the theorem.
\end{proof}

\subsection{The exceptional fields with rational complementary codegrees}
\begin{theorem}\label{thm:half-ten-grams}
Let $K=\mathbb Q(D)$ have one of the following ten values:
\[
 \mathbb Q(\sqrt2,\sqrt d),\quad d\in\{15,19,21,23,29,33,37,41\},
 \qquad \mathbb Q\!\left(\sqrt{m(2+\sqrt 2)}\right),\quad m\in\{5,7\}.
\]
Suppose the two ordinary and the two signed complementary codegrees
are all $4$. After Galois transport maximizing $E$,
\[
                    \operatorname{Tr}_K(E/D)\le19.
\]
\end{theorem}
\begin{proof}
Normalize $p_i=i(2/D)$ and $q_i=i(2/E)$ to have sum one. Label the
involution fixing $\sqrt2$ by $j=2$ and the remaining nonidentity
elements by $1,3$. Put $x_i=p_i/q_i$. At common codegree four the
complements contribute equally to the ordinary and signed squared
norms of the simultaneous column.
Lemmas~\ref{lem:column-comparison} and~\ref{lem:mass-moments} therefore give
\begin{equation}\label{eq:quartic-half-ten-grams}
 \sum_l q_{l+i}^2/q_l\le\sum_l p_{l+i}^2/p_l,\qquad
 K_j=\sum_l q_l x_l x_{l+j}\in 2\mathbb Z+1 .
\end{equation}
The normalized Fourier caps at conductor $c$ are
\[
 c|P|^2\le4(c,2),\qquad c|Q|^2\le4(c,4),\qquad
 c|P|^2+2\sqrt c\,|Q|\le8 .
\]

Here and in the next height estimate we use the following test for a
small $q_j$. Put $a=p_0/q_0$, $b=p_0/q_j$; then $0<b\le a\le1$.
For $i=1,3$ the spectral bound on the columns and their squared-norm
bound imply
\begin{align*}
 aq_i,\ bq_{i+j}&\le p_i,\\
 aq_i^2+bq_{i+j}^2&\le p_i^2+
 p_0\left(\delta+\sum_{l\ne0}p_{l+i}^2/p_l\right).
\end{align*}
Here $\delta=0$ in \eqref{eq:quartic-half-ten-grams}; more generally
one adds $\delta$ when the complementary terms give that correction
to the first inequality. On a box in $p,q$ these are linear bounds
on $a,b$ and bound the positive moment
\[
  K_j=\sum_l p_l p_{l+j}/q_l=(a+b)p_j+p_1p_3(1/q_1+1/q_3).
\]
Bisect the box after applying the Fourier constraints. Discard a
box when its upper bound for $K_j$ is below one or when its trace
upper bound $1+\sum_{i\ne0}\sup(p_i)/\inf(q_i)$ is below the
proposed trace. The trace test is used only for positive denominators.

Among the eight biquadratic fields the two other conductors are
bounded below by $(60,120)$ when the first is divisible by four,
and by $(21,168)$ when it is odd. In the two cyclic fields the
faithful conductor is at least $80$. Within each group the separate
and mixed Fourier caps decrease, so it suffices to check the three
least conductor tuples. For $i=1,3$ the separate bounds give
$q_i>9/100$: use $4q_i\ge 2(q_0+q_i)\ge1-|Q|$ in the
biquadratic fields, and $1-\sqrt2|Q|$ for a cyclic adjacent pair.
If the trace were at least $20$, we would have $p_0\le q_0\le
q_j<1/19$, since $x_0\le1$. The
\codecheck{quartic_more/ten_half_height.py}{three interval checks}
exhaust precisely these initial boxes using the
\codecheck{quartic_more/half_bound.py}{two-variable bound} above.
This proves the assertion.
\end{proof}

\begin{corollary}\label{cor:quartic-half-ten-exclusion}
A rank-six fusion category in the ten fields of
Theorem~\ref{thm:half-ten-grams} with both complementary codegree pairs
equal to $(4,4)$ has a pseudo-unitary Galois conjugate.
\end{corollary}
\begin{proof}
By Theorem~\ref{thm:all-proper-quartic} take $x=E/D$ primitive.
Write $P=X^4-AX^3+BX^2-CX+H$ for its polynomial. By
Lemma~\ref{lem:mass-moments}, $A,B,C$ are even and $H$ is odd.
In the trace boxes of Lemma~\ref{lem:quartic-trace-box} retain the
positive primitive $d$-numbers with these parities and least root
below one. The \codecheck{verify_quartic_half_ten.py}{lattice enumeration}
finds no polynomial with $A\le19$ (already its seven candidates
up to trace $39$ start at $A=20$).
\end{proof}

For the neighboring fields we will need the weights at somewhat
larger traces. Given a candidate $P$, let $x_i$ be its four roots
in increasing order, with $q_i$ permuted accordingly, and put
$m_r=\sum q_i x_i^r$. At equal half mass $m_0=m_1=1$ and every
$m_r$ is a positive odd integer by Lemma~\ref{lem:mass-moments}.
Thus $3\le m_2<A$ and $m_2^2\le m_3<A m_2$.
For each possible pair $(m_2,m_3)$ one has
\begin{equation}\label{eq:quartic-recovered-weights}
 q_i=\frac{x_i^3+(1-A)x_i^2+(B-A+m_2)x_i+
                        m_3-Am_2+B-C}{P'(x_i)} .
\end{equation}
Indeed the numerator is the polynomial part of
$P(X)\sum_{r=0}^3m_rX^{-r-1}$, and the residues of its quotient
by $P$ are the weights. Retain positive weights with $q_0$ least
and all $\sum_i q_ix_{i+j}$ odd in the correct Galois ordering.

\begin{theorem}\label{thm:quartic-three-rational}
Let $\mathcal C$ be a rank-six fusion category with quartic dimension
field $K$ and rational signed complementary codegrees. If
\[
 K=\mathbb Q(\sqrt3,\sqrt d),\qquad d=5,7,10,11,13,
\]
then $\mathcal C$ admits a pseudo-unitary Galois conjugate.
At signed mass $t=1/2$ the same conclusion holds for
$K=\mathbb Q(\sqrt2,\sqrt d)$, $d=7,11,13,17$, and for
$K=\mathbb Q(\sqrt{3(2+\sqrt2)})$.
\end{theorem}
\begin{proof}
Use the three mass ranges of Lemma~\ref{lem:quartic-rational-mass-strata}
and the integral shifts of Corollary~\ref{cor:galois-columns}.
At $t=1/2$ we have $s=t$, and all complementary signed evaluations
of the ordinary columns are one.

Use the normalized probabilities $p_i=i(2/D)$, $q_i=i(2/E)$.
Their shifted two-column moments are positive odd integers. If the
ordinary complementary codegrees are rational, both are fixed under
Galois and the simultaneous-column Gram inequalities hold without
change. Otherwise their quadratic codegrees have $u=1/2$. The last two lines
of Lemma~\ref{lem:quartic-short-pairs} give
$f,f'=n\pm\sqrt{n(n-4)}$ with $(c,n)=(5,5),(8,8)$.

If the automorphism
fixes $f$ the Gram inequality still holds as stated. If it
exchanges $f,f'$, its ordinary complementary squared-norm contribution
is $f/f'^2+f'/f^2=(n-3)/2$ and its signed contribution is $1/2$.
Thus in every case
\[
 \sum_l q_{l+i}^2/q_l\le \sum_l p_{l+i}^2/p_l+\delta,
 \qquad \delta=0\ \text{or}\ n-4.
\]

Take $j$ fixing the named quadratic subfield, of conductor $8$ or
$12$. For $i=1,3$ the other two pair sums give $q_i>1/64$:
their lower bounds for $2(q_0+q_i)$ are $1-2\sqrt{(c,4)/c}$ in
a biquadratic field and $1-\sqrt{32/48}$ for a cyclic adjacent
pair. Apply the small-$q_j$ test in Theorem~\ref{thm:half-ten-grams}
to $p_0,q_0,q_j\le1/64$, with the preceding correction $\delta$.

For all five fields containing $\sqrt2$ an irrational ordinary
pair must have $(c,n)=(8,8)$. Then $c O_{\chi_8}=2$ in \eqref{eq:orbit-projection}, which forces $T_{\chi_8}=0$ and gives $q_0+q_j=1/2$, hence $q_j\ge1/4$ directly.
For fields containing $\sqrt3$ there is no irrational ordinary
pair unless the field contains $\sqrt5$, where $n-4=1$.
Thus $\delta=0$ suffices except in $\mathbb Q(\sqrt3,\sqrt5)$,
where one suffices. The interval calculation permits larger
corrections on several rows.

The \codecheck{quartic_more/half_bound.py}{interval calculation}
gives $K_j<1$ throughout these boxes, hence $q_j>1/64$ and $A<65$.
Apply the same lattice enumeration and
\eqref{eq:quartic-recovered-weights} with $A\le64$.

Exactly one tuple survives them, at $(8,17,136)$ with
$(A,B,C,H)=(24,78,24,1)$. For its possible $m_2=7,m_3=111$
one weight exceeds $7/10$. But the conductor-$136$ signed Fourier
amplitude is at most $4/\sqrt{136}<2/5$; positivity of the other
three weights gives $2\max q_i-1\le4/\sqrt{136}$. This contradiction
excludes the last tuple.

Suppose $t>1/2$ in one of the five fields containing $\sqrt3$.
A quadratic common codegree gives $t\ge3/5$. For rational rows
of codegrees $g\le h$, $gh$ is a square or twice a square:
at each odd prime the inertia index divides two, so $N(E)$ has
even valuation, as does the biquadratic discriminant. Apply
Lemma~\ref{lem:trace-order-conductor}. For $t<3/5$ we must have
$g=3,4$. The only pair with $1/2<t<3/5$ in these square classes is
$(3,12)$, giving $t=7/12$. At $d=5,11$ three is inert
in $\mathbb Q(\sqrt d)$, and Lemma~\ref{lem:quartic-three-induction}
excludes this possibility.
All nontrivial signed Fourier amplitudes in these five fields are
at most $1/\sqrt3$. In unnormalized variables $q_i=i(1/E)$,
$p_i=i(1/D)$ it follows that $q_i\ge(t-1/\sqrt3)/4$ for
$i\ne0$. Using $s\le t$ gives
$A<1+4t/(t-1/\sqrt3)$. Its integer cap is $106$ at
$t\ge3/5$ and $390$ at $t\ge7/12$. The same biquadratic
lattice enumeration, now without the half-mass parity condition,
gives respectively $168,321,108,16,139$ candidates in the order
$d=5,7,10,11,13$. For each candidate impose the coordinate minima, $1/3\le s\le t\le1$,
the stated lower bound on $t$, and $\sum q_ix_{i+j}\ge s$ at each
shift. If $\sqrt5\notin K$ also impose $s+2t\ge5/3$ from the mass strata. Use the separate and mixed
Fourier inequalities of Lemma~\ref{lem:sextic-matrix-compression},
including the joint projection in directions $(u,v)$,
$-3\le u,v\le3$. Substitution of $p_i=x_iq_i$ gives linear
inequalities in the four weights. The $752$ integer certificates in
\codefile{quartic_more/verify.py} exclude every primitive ratio.
\end{proof}

\begin{theorem}\label{thm:quartic-fifteen-rational}
Suppose a rank-six fusion category has rational signed complementary
codegrees. It has a pseudo-unitary Galois conjugate if its quartic
dimension field is one of the following fifteen fields:
\[
 \mathbb Q(\sqrt2,\sqrt d),\qquad
 d=7,11,13,15,17,19,21,23,29,33,37,41,\quad \mathbb Q\!\left(\sqrt{m(2+\sqrt2)}\right),\quad m=3,5,7 .
\]
\end{theorem}
\begin{proof}
We prove two mass gaps, now using both the ordinary and the signed
algebras. In any displayed field $K$, all odd valuations of the
norm of a $d$-number are even, since the ramification index is at
most two at odd primes. For the cyclic fields this follows since
the quadratic subfield is unramified there. Also the field
discriminant has only even odd-prime valuations. The regular-trace
discriminant of the integral character order with rational
complementary axes of codegrees $g,h$ therefore shows that $gh$ is a
square or twice a square. For a quadratic complementary factor $F$
with common rational codegree $h$, the same identity shows that
$|\operatorname{disc}F|$ has even odd-prime valuations. Thus
$F=\mathbb Q(\sqrt2),\mathbb Q(i)$ or $\mathbb Q(\sqrt{-2})$.
Lemma~\ref{lem:quadratic-row-difference} gives $h\ge8$ in the
real case. In the imaginary case duality exchanges just two
characters and hence just two basis elements, by invariance of
trace. The same lemma gives $h=|\operatorname{disc}F|b^2$. A quadratic pair with common rational
codegree has $h=4$ or $h\ge8$.

Three and five are inert in $\mathbb Q(\sqrt2)$.
Lemma~\ref{lem:quartic-three-induction} excludes codegree three;
Lemma~\ref{lem:quartic-inert-congruence} excludes $(5,5),(5,10)$
at five and $(6,6)$ at three. In either sector a rational-codegree
mass is thus $1/2$, with pair $(4,4)$, or at least $5/8$:
the square classes above leave only $(3,12),(5,5)$ in the
intervening interval, and mass below $1/2$ would require codegree three.
For a distinct irrational ordinary pair let $u=1-s$ and
$c=\operatorname{disc}\mathbb Q(f)$. If $1/2\le s\le5/8$,
then $u\ge3/8$ and $c\le(1+2u)/u^2\le112/9$.
Among the quadratic subfields here this leaves $c=8$.
Both pairs in Lemma~\ref{lem:quartic-short-pairs} saturate
$c(1/f-1/f')^2=1$; their masses $u$ are $1/2,3/8$.
Thus $cO_{\chi_8}=2$ and the mixed bound
in \eqref{eq:orbit-projection} gives $\widehat q(\chi_8)=0$.

We have thus shown in either sector that the mass equals $1/2$
or is at least $5/8$. At half mass rational complementary
codegrees must be $(4,4)$ since $(3,6)$ was excluded.

The five half-mass fields in
Theorem~\ref{thm:quartic-three-rational} are already settled.
For the other ten at half mass, ordinary rational codegrees give
Corollary~\ref{cor:quartic-half-ten-exclusion}. An irrational
ordinary pair saturates at conductor eight, giving normalized
$q_0+q_j=1/2$ for the involution fixing $\sqrt2$. Also $q_1,q_3>9/100$ as in Theorem~\ref{thm:half-ten-grams},
and $q_j\ge1/4$ gives $A<1+100/9<13$. The empty trace-$19$
inventory of Corollary~\ref{cor:quartic-half-ten-exclusion}
requires only the signed equal-mass parity and the trace cap.

Suppose $t>1/2$, and use unnormalized $p_i=i(1/D)$,
$q_i=i(1/E)$. If $s\ge5/8$, then either $t\ge23/36$ or
$t=s=5/8$. Indeed for rational rows $g,h\ge4$ with $gh$ square or twice a
square there is no mass in $(5/8,23/36)$, and a common codegree
in this interval is impossible.
Write $y=q_0+q_j$ for the
conductor-eight pair and let $B$ be the sum of the other two
ordinary Fourier amplitudes, or twice the faithful amplitude in
the cyclic case. At this pair use the tangent estimate (R1) with $c=8$.

For $s\ge s_0$, $t\ge t_0$ appearing below, use $r=s_0-B>0$.
The right side divided by $s$ is minimized at $s=s_0,t=t_0$
since $\sqrt8\,t_0-16rB-8r^2-2<0$. Once this minimum
$\ell$ is positive, $q_j\ge s\ell/2$. The other
coordinates satisfy $q_i\ge(t-a_i)/4$ where $a_i$ is the
quadratic amplitude $\sqrt{\gcd(c,4)/c}$ or the adjacent cyclic
amplitude $\sqrt{8/J}$. These give
$A<1+\max(2/\ell,\max_i4t/(t-a_i))$.
Apply the formula with $(s_0,t_0)=(5/8,5/8)$, except at $d=7$,
where we use $(5/8,23/36)$ and treat $s=t=5/8$ below. If instead
$s=1/2$, $t\ge5/7$, apply it with $(1/2,5/7)$.
The resulting integer caps fit in the final table below.

It remains to treat $s=1/2$, $5/8\le t<5/7$, and the one case
$s=t=5/8$ at $d=7$. In the first range the signed pair has
separate rational characters. If their signed evaluations on $V$
are $b,e$, then $b/g+e/h=1/2$. For $4\le g\le h$ the range gives $g\le6$. At $g=4$ positivity
forces $b=1$ and $h=4e<28$; the square classes leave
$(h,e)=(8,2),(16,4)$. At $g=5,6$ only the three excluded pairs
$(5,5),(5,10),(6,6)$ satisfy the bounds and square classes.
For $K_j=\sum_l p_l p_{l+j}/q_l$, the unit coefficient gives
$K_j+b^2/g+e^2/h\in\mathbb Z$; hence $K_j\ge1/4,3/4$ respectively.
At $s=t=5/8$, Lemma~\ref{lem:mass-moments} gives $K_j\ge5/8$.
If the ordinary pair is irrational, the saturation proved above
already bounds every coordinate. Otherwise its simultaneous
columns have both complementary evaluations one, as do the signed
ones. The unsigned Gram bound thus becomes
\[
 \sum_l q_{l+i}^2/q_l\le\sum_l p_{l+i}^2/p_l+t-s .
\]

The \codecheck{quartic_more/nonhalf_bound.py}{interval bounds} give
$q_j>1/256$. For $(4,8)$ they use $(s,t,K_{\min},\delta)=
(1/2,5/8,1/4,1/8)$, and for $(4,16)$ its weaker enclosing
range $(1/2,[5/8,11/16],1/2,1/4)$. At $d=7$ also use
$(5/8,5/8,5/8,0)$. Apply the two-variable argument of
Theorem~\ref{thm:half-ten-grams} to the unnormalized masses.
There is one extra linear constraint: by the shifted first moment,
\[
 \frac ba p_j +a q_j + q_k p_m/q_m+q_m p_k/q_k
       \in s+\mathbb Z_{\ge0},
 \qquad a=p_0/q_0,\quad b=p_0/q_j .
\]
Here $b$ denotes the variable of the two-variable program, distinct
from the integer character evaluation above. Interval evaluation
of the last three terms bounds $b/a$ above and below using the
permitted integers; the resulting linear bounds are included before
taking the objective maximum. The checks visit at most $17239$
boxes per row. The other two coordinates exceed $1/256$ by their
separate Fourier bounds. Thus $A\le128$ at $s=1/2$ or $A\le160$
at $s=5/8$.

The final integer caps and counts of primitive $d$-number candidates
are as follows:
\[
\setlength{\arraycolsep}{3.4pt}\begin{array}{c|rrrrrrrrrrrr|rrr}
 d&7&11&13&15&17&19&21&23&29&33&37&41&m=3&m=5&m=7\\\hline
 A\le&397&128&230&128&128&128&128&128&128&128&128&128&128&128&128\\
 \#&283&9&49&40&7&5&81&9&2&4&0&0&29&12&26
\end{array}
\]
Use the mass constraints $1/2\le s\le t$, $5/8\le t\le1$, the
nonnegative integral shifted differences, coordinate minima, and
the separate and mixed Fourier inequalities as in
Theorem~\ref{thm:quartic-three-rational}. The file
\codefile{quartic_more/two_nonhalf.py} verifies positive integer
dual gaps for each of the $556$ candidates. It generates the
cyclic fields in the integral basis from
Corollary~\ref{cor:quartic-half-ten-exclusion}, valid also at
$m=3$. This eliminates every primitive ratio and proves the theorem.
\end{proof}

\subsection{The remaining fields containing the golden ratio}
\begin{theorem}\label{thm:quartic-cyclic-golden-family}
Let $K$ be a real cyclic quartic field with quadratic subfield $\mathbb Q(\sqrt5)$, and suppose $K\ne\mathbb Q(\zeta_{20})^+$. If $K=\mathbb Q(D)$ and the signed complementary pair has distinct irrational codegrees, then the category has a pseudo-unitary Galois conjugate.
\end{theorem}
\begin{proof}
The pair has field $\mathbb Q(\sqrt5)$. With the notation of Theorem~\ref{thm:quartic-high-quadratic}, either $t\ge16/25$, or $s=t=1/2$ and the signed amplitude at conductor five is at most $\sqrt5/10$. The conductor $J$ is a multiple of five, and if even it is a multiple of twenty. A real quartic field cannot have conductor five. For $J\ne20,40$, the adjacent-pair amplitude is at most $\sqrt{2/15}$, since an odd conductor is at least fifteen and an even one is at least sixty. Thus the three pair bounds again give $A\le22$. For $J=40$ one may use the weaker bound $1/\sqrt5$ for all three amplitudes. It gives $A\le21$ in the first case and $A\le38$ in the second.

The \codecheck{generic_quartic/cyclic_five.py}{quadratic resolvent test} makes this enumeration small. There is an integer root $r$ of the cubic resolvent for which $A^2-4(B-r)=5v^2$, $v\in\mathbb Z$: the two pairs are the orbits of the involution, and the difference of their sums is an anti-invariant integer in $\mathbb Q(\sqrt5)$. For each $A,H,B,v$, putting $r=B-(A^2-5v^2)/4$ in the resolvent gives a quadratic equation for $C$. At trace at most $22$, the positivity, primitivity and $d$-number tests give $53$ cyclic polynomials with the required quadratic subfield. At trace at most $38$ a second list retains only those for which $\operatorname{disc}(P)/(5\cdot40^2)$ is an integer square; it has $62$ entries. These are necessary conductor conditions, and overlap of the two searches is harmless.

Except for two polynomials, the linear constraints of Theorem~\ref{thm:quartic-high-quadratic} exclude every entry of both lists. With the minimum root at zero there are three choices for the opposite root, and interchange of the two adjacent roots preserves the bounds. For the first list the constants of odd conductor fifteen weaken both separate bounds and the mixed bound at the faithful character, since $J$ is odd and at least fifteen or a multiple of twenty at least sixty. Use the actual constants of conductor forty for the second list. The two choices for the signed mass give $678$ rational separating certificates, with three more for the short-pair branch of the last case below.

The polynomial $X^4-8X^3+19X^2-12X+1$ has root $2-(\zeta_{20}+\zeta_{20}^{-1})$ and generates the field excluded in the statement. For the remaining polynomial put
\[
           P=X^4-7X^3+14X^2-8X+1,\qquad x=E/D .
\]
Choose $b_i=\zeta_{15}^{2^i}+\zeta_{15}^{-2^i}$, $i\in\mathbb Z/4$.
Then $x=2-b_0$ at its least embedding. The $b_i$ have trace one and
circulant Gram matrix with first row $(9,-1,-6,-1)$. Its determinant
$1125=5\cdot15^2$ is the field discriminant, so they form an integral
basis. The involution $\tau$ sends $b_0$ to $b_2$ and fixes the
complementary field. Reduction modulo $P$ gives
\[
 (\tau x-x)P'(x)(x^3-7x^2+13x-4)=15 .
\]
These \codecheck{generic_quartic/c15_lattice.py}{period identities}
and Lemma~\ref{thm:primitive-divisibility} give
$z=15/E=\sum c_i b_i$ with integral $c_i$.

Write $T=\sum c_i$, $U=c_0-c_1+c_2-c_3$, $V=c_0-c_2$, $W=c_1-c_3$.
The Gram eigenvalues at the quadratic and faithful characters are
$5,15$, respectively. The squared Fourier amplitudes of $z$ there
are $5U^2$ and $15(V^2+W^2)$. The separate bounds for $q=z/15$
therefore give $|U|\le3$ and $V^2+W^2\le1$.
The half-mass case already passed the separating test, so
$T=15t\ge15(16/25)$ and $T\ge10$. The same Gram matrix gives
\[
 15s=\operatorname{Tr}(xz)=(7T-5U-30V)/4<15 .
\]
Hence $V>0$: at $T=10$ use also $U\equiv T\pmod2$, and at
$T\ge11$ use $U\le3$. This gives $V=1,W=0$. But then
$z-\tau z=b_0-b_2>0$, contradicting minimality of $q$.

\end{proof}

\begin{theorem}\label{thm:quartic-three-five}
A complex fusion category of rank six with $\mathbb Q(D)=\mathbb Q(\sqrt3,\sqrt5)$ has a pseudo-unitary Galois conjugate.
\end{theorem}
\begin{proof}
The rational complementary case is Theorem~\ref{thm:quartic-three-rational}. For an irrational complementary pair let $c_0=5,12$ or $60$ be its discriminant. If $c_0=60$, then $t\ge3/4$ by the proof of Corollary~\ref{cor:quartic-five-unbounded}. All three quadratic amplitudes are at most $1/\sqrt3$, so $A\le24$ by the pair bound. For $c_0=12$ we have $A\le94$ by Proposition~\ref{prop:quartic-primitive-norm-height} and $t\ge5/9>11/20$ by Lemma~\ref{lem:quartic-short-pairs}.

At $c_0=5$ a pair other than $5\pm\sqrt5$ gives $t\ge16/25$, and the same pair bound gives $A\le64$. For the short pair, Lemma~\ref{thm:quartic-five-to-three} gives
$s=t=1/2$ and the positive odd normalized moment of
Theorem~\ref{thm:half-ten-grams}, since the involution fixing
$\sqrt3$ exchanges the two complementary codegrees. Set
$\varphi=(1+\sqrt5)/2$.

The signed complementary contribution to the squared norm of the simultaneous column is $1/2$ for an automorphism fixing the pair, and $\operatorname{Tr}(\varphi^4/f)=1$ if it interchanges it. Both are at least the $1/2$ subtracted in the Gram inequality of Theorem~\ref{thm:quartic-three-rational}. The ordinary correction is therefore at most one, by the same calculation there. The already verified row $(12,5,60)$ of that theorem's height calculation yields $A\le64$; its inputs are the half-mass equalities, the Fourier bounds, a Gram correction at most $16$, and $K_j\ge1$.

The biquadratic coefficient enumeration at trace $94$ produces $143$ primitive $d$-number candidates. For each test the four alternatives just described, retaining the corresponding trace cap, the lower bound for $t$, and in the short case $s=t=1/2$ and $|\widehat q(\chi_5)|\le\sqrt5/10$. Every test has a strict rational separating certificate as in Theorem~\ref{thm:quartic-high-quadratic}. The verifier is \codefile{generic_quartic/three_five.py}. The proper ratios are covered by Theorem~\ref{thm:all-proper-quartic}.
\end{proof}

\begin{theorem}\label{thm:quartic-cyclic-twenty-all}
Every rank-six fusion category with dimension field
$K=\mathbb Q(\zeta_{20}+\zeta_{20}^{-1})$ has a pseudo-unitary
Galois conjugate.
\end{theorem}
\begin{proof}
Let $g$ generate its cyclic group in the order $0,1,2,3$ and
write $p_i=g^i(1/D)$, $q_i=g^i(1/E)$ after transport maximizing
$E$. The quartic and quadratic character conductors are $20$ and
$5$. In particular $q_2\ge(t-1/\sqrt5)/4$ and
$q_1,q_3\ge(t-\sqrt{2/5})/4$.

Here $s\ge1/2$ in both signed cases. Indeed the ordinary part of the
proof of Lemma~\ref{lem:quartic-rational-mass-strata} leaves only
$(3,3)$ below half mass, and three is inert in $\mathbb Q(\sqrt5)$.
Apply Lemma~\ref{lem:quartic-three-induction}.

For rational signed codegrees we have $s\le t$. If $t\ge23/36$ then
Lemma~\ref{lem:quartic-pair-height} gives $A\le398$. Otherwise
$1/2\le t<23/36$. For separate rational rows of codegrees $g_0,h$,
$g_0h$ is a square or five times a square: only five has inertia index
four and the field discriminant is $5\cdot20^2$. For the same reason
a quadratic common-codegree field has square class $-1,5$ or $-5$,
by Lemma~\ref{lem:trace-order-conductor}. A common codegree in this
mass range is four or five. For separate rows $4\le g_0\le h$,
the square classes leave $(4,4),(4,5),(5,5)$. At five $K$ has
a single prime, so two distinct residue rows cannot both be assigned
a quartic branch. Thus in the last pair the rows agree modulo five,
but their squared distance is ten. This excludes $(5,5)$.

In each remaining rational case $s=t$. Indeed the complementary
values of the ordinary induction class are totally positive integers
on separate rational axes, or conjugate algebraic integers on a
common-codegree orbit. Their weighted sum $1-s$ is at most $1/2$.
At $(4,4)$ and $(4,5)$ this forces both values to be one. At common
five their sum is at most $5/2$, hence two by positivity and
integrality of their norm. The three rational low masses are
$1/2,11/20,3/5$.

For distinct signed codegrees Lemma~\ref{thm:quartic-five-to-three}
gives either $t\ge2/3$, whence $A\le117$, or the short golden
pair, with $s=t=1/2$. In the former case the involution $g^2$ fixes
the signed complementary field, so
$\sum_i q_i x_{i+2}\ge s$ by Lemma~\ref{thm:primitive-divisibility}.

In all four low cases
$K_j=\sum_i p_i p_{i-j}/q_i\ge t$. For rational codegrees this
follows from Lemma~\ref{lem:mass-moments}; for the short golden
pair use Lemma~\ref{thm:quartic-five-to-three}, which also gives
all three shifted first-moment bounds. Its signed quadratic
amplitude is at most $\sqrt5/10$ by Lemma~\ref{lem:quadratic-compression}.
We claim $A\le77$. The opposite weight $q_2$ already exceeds $1/128$;
by reflection, consider $p_0,q_0,q_1\le1/128$. With
$a=p_0/q_0,\ b=p_0/q_1$ the spectral inequalities are
$0\le b\le a\le1$, $a q_i,bq_{i+1}\le p_i$, and
\[
 K_1=a p_3+b p_1+p_1p_2/q_2+p_2p_3/q_3.
\]
At half mass the ordinary irrational pair can only have
$(c,n)=(5,5)$ by Lemma~\ref{lem:quartic-short-pairs}, so its
ordinary Gram contribution is at most $(n-3)/2=1$.
The signed contribution is at least $1/2$ also for the short
golden pair, as in Theorem~\ref{thm:quartic-three-five}. Thus the
unnormalized Gram correction is at most $1/2$.
The \codecheck{quartic_more/c20_bound.py}{interval calculation}
with these bounds and the separate and joint Fourier inequalities
gives $K_1<t$ on every such box. Hence $A<1+128s<78$.

For $y=\zeta_{20}+\zeta_{20}^{-1}$ one has $y^4-5y^2+5=0$ and
$\mathcal O_K=\mathbb Z[y]$: write each cyclotomic integer as
$u(y)+v(y)\zeta_{20}$ using $\zeta_{20}^2=y\zeta_{20}-1$ and take
complex conjugation.
Write $x=a+by+cy^2+dy^3$. The inverse trace Gram matrix of this
basis has diagonal $(3/2,1,1/5,1/10)$, so $\operatorname{Tr}(x^2)<A^2$
bounds $b^2,5c^2,10d^2<A^2$. Conjugating to have $0<x<1$ at a
fixed embedding puts $a$ in an interval of length one. The norm
polynomial of each lattice point is tested for positivity and the
$d$-number divisibilities, leaving $721$ polynomials at $A\le398$.
For each ratio impose the coordinate minima and Fourier bounds with
$p_i=x_iq_i$. In the high rational case use $t\ge23/36$, $s\le t$
and all three shifted inequalities; in the high irrational case use
$t\ge2/3$, $s\ge1/2$ and the shift by two. Both have $s+t\ge1$.
At the four low profiles use their mass equalities and
$\sum_iq_ix_{i+j}\ge s$, $\sum_iq_ix_ix_{i+j}\ge t$ for every
shift, with the improved golden amplitude or the rational half-mass
parity where applicable.
The \codecheck{quartic_more/cyclic20.cpp}{lattice enumeration},
\codecheck{quartic_more/c20_candidates.py}{rational inequalities} and
\codecheck{quartic_more/c20_duals.json}{integer certificates} exclude
all $721$ high rational and $183$ high irrational candidates. At the rational low masses,
Lemma~\ref{lem:mass-moments} gives $P(X)\equiv(X-1)^4$ modulo
$2,20,5$, leaving $24,3,12$ respective candidates. The short
golden pair has $108$. Thus the four low cases use $147$ certificates.
\end{proof}
\begin{proof}[Proof of Theorem~\ref{thm:quartic-final}]
Theorem~\ref{thm:all-proper-quartic} treats proper dimension ratios.
For a primitive ratio, two cases exhaust the two complementary signed
characters. Rational codegrees leave twenty-four fields by
Theorem~\ref{thm:quartic-rational-frontier}; Theorems~\ref{thm:quartic-three-rational},
\ref{thm:quartic-fifteen-rational} and~\ref{thm:quartic-cyclic-twenty-all}
exclude twenty-one. For irrational codegrees, Theorems~\ref{thm:quartic-high-quadratic}
and~\ref{thm:quartic-silver-fields} leave the three stated fields,
$\mathbb Q(\sqrt3,\sqrt5)$ and the cyclic quartics containing $\sqrt5$.
The latter two cases are covered by Theorems~\ref{thm:quartic-three-five},
\ref{thm:quartic-cyclic-golden-family} and~\ref{thm:quartic-cyclic-twenty-all}.
\end{proof}

\section{The quartic fields of conductors sixteen and twenty-four}\label{sec:guide-1624}
\begin{theorem}\label{thm:quartic-four-complete}
Every complex fusion category of rank six whose global Frobenius--Perron dimension generates $\mathbb Q(\sqrt2,\sqrt3)$ or $\mathbb Q(\zeta_{16})^+$ admits a spherical structure.
\end{theorem}
Failure of sphericality would force both complementary pairs to have
codegrees $4,4$. The dimension orbit would then carry half the
reciprocal mass. Its odd integral moments determine a short list of
possibilities over $\mathbb R$; reduction modulo small primes excludes
each of them.
\subsection{Small complementary codegrees}

\begin{lemma}\label{lem:quartic-four-rows}
Suppose $K=\mathbb Q(D)$ is quartic with one prime over two, and
the two complementary ordinary characters are rational of codegree four.
There is an invertible $g$ of order two acting by two transpositions.
If the signed complementary characters are conjugate over a real
quadratic field, their rows have the form
$(1,-1,y,-y,0,0)$ and its conjugate, with $yy'=-1$. Both quartic
dimension orbits then have reciprocal mass $1/2$.
\end{lemma}
\begin{proof}
Neither rational branch splits off at two, since its unit coefficient
would be $1/4$, and their union cannot split off either, since its
coefficient would be $1/2$. Thus the two rows agree modulo two.
Each has four entries $\pm1$. By orthogonality their sum is supported
at $1,g$ with entries $2,\pm2$. Orthogonality to the Perron row
gives $d_g=1$ and the entry $-2$. Multiplication by $g$ has eigenvalues
$1$ four times and $-1$ twice, so $g$ is an involution with two
transpositions. Its positive lift has the same permutation and is $1$
on the quartic orbit. The two conjugate characters take the same sign
on $g$, necessarily $-1$ by the zero unit coefficient. Their rows
therefore vanish at the two fixed points and alternate on each
transposition. Orthogonality gives $yy'=-1$ and hence
$1/(2+2y^2)+1/(2+2y'^2)=1/2$.
\end{proof}

\subsubsection{\texorpdfstring{The pair $8\pm4\sqrt2$ in $\mathbb Q(\sqrt2,\sqrt3)$}{The conjugate pair at conductor twenty-four}}
\begin{lemma}\label{prop:quartic-twenty-four-fold}
At either conductor, suppose the signed complementary codegrees are
$f=4(2+\sqrt2)$ and $f'$. Then $s=t=1/2$. Put $q=2/E$, $x=E/D$ and
$L(B)=\operatorname{Tr}_K(qB)$. Every monomial in conjugates of $x$
has positive integral $L$-value, and each $L(\sigma x)$ is odd.
The squared-norm inequality in \eqref{eq:quartic-half-ten-grams} holds.
\end{lemma}
\begin{proof}
Write $F=\mathbb Q(\sqrt2)$ and $\beta=a+b\sqrt2$ for the
complementary value of the signed Perron induction class. Its
contribution to the unit coefficient is
$\operatorname{Tr}_F(\beta/f)=(a-b)/2$. It lies strictly between
zero and one, so $a-b=1$ and $s=1/2=t$. Total positivity then gives
$b=0,1,2$. By Corollary~\ref{cor:galois-columns} the integral
translates have coordinates $(\sigma x,\sigma\beta)$. For any
product with coordinates $(B,\gamma)$, its doubled unit coefficient
is $L(B)+2\operatorname{Tr}_F(\gamma/f)$. The second summand is
integral for every $\gamma\in\mathcal O_F$, and is odd for either
conjugate of each of the three possible $\beta$. This proves the
moment assertions.

By Lemma~\ref{lem:quartic-short-pairs} the ordinary complementary
codegrees are rational or are the same pair $f,f'$: neither field
contains $\sqrt5$. A rational pair contributes $1/2$ to every
simultaneous-column squared norm. The signed pair contributes $1/2$
when fixed and $f/f'^2+f'/f^2\ge1/2$ when exchanged. Subtract the
complementary contributions in Lemma~\ref{lem:column-comparison}.
\end{proof}
\begin{theorem}\label{thm:quartic-twenty-four-silver}
A rank-six complex fusion category with quartic dimension field
$\mathbb Q(\sqrt2,\sqrt3)$ and complementary signed codegrees
$8\pm4\sqrt2$ admits a spherical structure.
\end{theorem}
\begin{proof}
It suffices to exclude the primitive ratio $x=E/D$ by
Theorem~\ref{thm:all-proper-quartic}. Use the normalized probabilities
$p_i=i(2/D)$, $q_i=i(2/E)$ as in Theorem~\ref{thm:half-ten-grams};
their sums are one by Lemma~\ref{prop:quartic-twenty-four-fold}.
Write $P_c,Q_c$ for their character transforms. Keeping the signed
complementary pair in the conductor-eight projection gives
$2P_8^2+Q_8^2\le1$ by \eqref{eq:orbit-projection}; also
$|Q_8|\le2\sqrt2/3$ by Lemma~\ref{lem:quadratic-compression}.
Let $j$ fix $\sqrt3$. The two-variable test of
Theorem~\ref{thm:half-ten-grams} applies with
$K_j=L(xj(x))\ge1$ by Lemma~\ref{prop:quartic-twenty-four-fold}.
If $q_j\le1/320$, the strengthened conductor-eight bounds give
$K_j<1$ in all \codecheck{generic_quartic/short24/height.py}{1077 interval boxes}.
The bounds on $Q_8,Q_{24}$ give $q_i>1/320$ for the other
nonidentity indices. Since $x_0\le1$ and $\sum_i p_i=1$,
$\operatorname{Tr}(x)\le320$.
The \codecheck{generic_quartic/short24/verify.py}{biquadratic enumeration}
of Theorem~\ref{thm:quartic-silver-fields} at this cap gives
444 primitive totally positive $d$-number polynomials with least
root below one. Impose the preceding Fourier bounds on
$p_i=x_iq_i$, $\sum p_i=\sum q_i=1$, with $p_0,q_0$ least.
The \codecheck{generic_quartic/short24/short24_duals.json}{separating weights}
exclude 443 polynomials.

The single remaining minimum embedding is
\[
 x=(52-18\sqrt8+14\sqrt{12}-10\sqrt{24})/4.
\]
Put $M_l=L(lx)$. By Lemma~\ref{prop:quartic-twenty-four-fold},
these are positive odd integers and $M'=L(x^2)$ is integral.
Moreover $M_0=1$ and $\sum_l M_l=52$.
In conductor order $8,12,24$ put
$T_c=\sum_i\chi_c(i)M_i$. Since $T_c=Q_c\widehat x_c$,
Fourier inversion gives
\[
 4M'=\operatorname{Tr}(x^2)+\tfrac{148}{3}T_8+
             \tfrac{362}{7}T_{12}+\tfrac{256}{5}T_{24}.
\]
The coefficients are $\widehat{x^2}_c/\widehat x_c$. Testing
positive odd $M_1,M_2$ with $M_3=51-M_1-M_2>0$ leaves just
$(1,25,25),(7,39,5),(25,11,15)$ by integrality. In each case
\[
 2(q_0-q_2)=Q_8+Q_{24}
   =-T_8/(18\sqrt8)-T_{24}/(10\sqrt{24})>0,
\]
contrary to minimality of $q_0$. This excludes all primitive ratios.
\end{proof}

\begin{lemma}\label{lem:quartic-large-mass-height}
Suppose the quartic dimension field is $K=\mathbb Q(\zeta_c)^+$,
$c=16,24$. Let $p_i=i(1/D)$,
$q_i=i(1/E)$, $t=\sum q_i$, after maximizing $E$ by transport.
If $t\ge71/100$, then $A=\sum p_i/q_i$ is less than $740$ at
$c=16$ and $758$ at $c=24$.
\end{lemma}
\begin{proof}
We have $p_0/q_0\le1$, $q_0\le q_i$, and $\sum p_i<1$. If every
$q_i$, $i\ne0$, exceeds $1/400$ then $A<401$. Otherwise choose
$j\ne0$ with $q_j\le1/400$. The pair imbalance
$|2(q_0+q_j)-t|$ is at least $7/10$. At conductor $24$
this is a Fourier modulus $|Q_d|$, and the separate bounds force
$d=8$. At conductor $16$ it is $|Q_2|$ for the opposite pair
or at most $\sqrt2|Q_1|$ for an adjacent pair (cyclic indices).
In each case the separate and mixed bounds of
Lemma~\ref{lem:sextic-matrix-compression} give
\[
 q_j\ge(t-1/\sqrt2)/4,\qquad
 p_j<(1+51/1000)/2 .
\]
Indeed the corresponding imbalance for $p$ is bounded by $|P_8|$,
$|P_2|$ or $\sqrt2|P_1|$, respectively, each below $51/1000$.

At $c=24$ the other two nonidentity entries of $q$ are at least
$(t-1/\sqrt3)/4$ by the separate pair bounds. At $c=16$ they are
greater than $1/10$: their sum is at least $.705$; for an opposite
pair their difference is at most $|Q_1|\le1/2$, and for an adjacent
pair each is at most $1/2+1/400$, by comparison with the opposite
small entry. Thus $A$ is less than
\[
 1+\frac{2(1+51/1000)}{71/100-1/\sqrt2}
    +\begin{cases}10,&c=16,\\4/(71/100-1/\sqrt3),&c=24.\end{cases}
\]
This gives the stated bounds.
\end{proof}

\begin{theorem}\label{thm:quartic-twenty-four-large-mass}
Let $K=\mathbb Q(\sqrt2,\sqrt3)$ be the quartic dimension field of a
rank-six fusion category. If $\operatorname{Tr}_K(1/E)\ge71/100$,
the category admits a spherical structure. This includes all cases
where the two complementary signed codegrees are distinct conjugates
in $\mathbb Q(\sqrt6)$.
\end{theorem}
\begin{proof}
Let $p_i=i(1/D)$, $q_i=i(1/E)$ and $s=\sum p_i$, $t=\sum q_i$.
Lemma~\ref{lem:quartic-large-mass-height} gives
$\operatorname{Tr}(E/D)<758$.

Use the cap $760$ in the same biquadratic integer enumeration as
above. It gives $1061$ primitive polynomials. The complete list is
excluded by $q_i,p_i\ge0$, $q_0\le q_i$, $p_0\le p_i$, $p_i=x_iq_i$,
$s,t\le1$, $t\ge71/100$, and the separate and coupled Fourier
estimates at conductors $8,12,24$. The verifier
\codefile{generic_quartic/short24/large_mass.py} regenerates the list and checks
integer weights $w\ge0$ for each outward-rounded system $Rq\le b$,
with $\min(0,\min_i(w^{\mathsf T}R)_i)>w^{\mathsf T}b$. The
weights are in \codefile{generic_quartic/short24/large_mass_duals.json}.
Corollary~\ref{cor:quartic-five-unbounded} proves $t>71/100$ for
quadratic complementary field of discriminant twenty-four in its
explicit bounds on the complementary mass. Finally apply
Theorem~\ref{thm:all-proper-quartic}.
\end{proof}

\subsection{The reciprocal mass of the quartic factor}

\subsubsection{Quadratic codegrees and small weights}
\begin{lemma}\label{thm:cyclic16-irrational-finite}
Let $K=\mathbb Q(\zeta_{16})^+$ and suppose that the signed
complementary codegrees $f,f'$ are distinct and irrational. Write
$s=\operatorname{Tr}_K(1/D)$, $t=\operatorname{Tr}_K(1/E)$,
and $v=1/f-1/f'$. Then $s\ge1/2$. For the Fourier coefficients
of $1/D,1/E$ at the faithful and quadratic characters one has
\[
    4|Q_1|+16|P_1|^2\le2,\qquad 2P_2^2+Q_2^2+2v^2\le1/2 .
\]
If the two complementary ordinary characters are rational of codegree
four, there is an involutive invertible acting by two transpositions.
In this case $s=t=1/2$ and $f=2+2y^2$, where
$y\in\mathbb Z[\sqrt2]$ has norm $-1$.
\end{lemma}
\begin{proof}
Write $F=\mathbb Q(\sqrt2)$ for the unique quadratic subfield and put $u=1/f+1/f'=1-t$.
An ordinary character of codegree two would force a rational signed
character of codegree two, by the invertible-object argument of
Lemma~\ref{lem:quartic-rational-mass-strata}. The signed algebra
$K\times F$ has no rational character. A complementary codegree three
is excluded by Lemma~\ref{lem:quartic-three-induction}. Hence rational
ordinary complementary codegrees leave $s\ge1/2$.
For a distinct irrational ordinary pair the same bound $s\ge1/2$
is Lemma~\ref{thm:quadratic-pair-mass}. The two Fourier inequalities
are \eqref{eq:orbit-projection} and the mixed bound of
Lemma~\ref{lem:sextic-matrix-compression}, retaining the additional
signed quadratic pair at conductor eight.
The last assertion is Lemma~\ref{lem:quartic-four-rows}, since two
is totally ramified in $K$.
\end{proof}

\begin{lemma}\label{lem:round2-quartic24-norm}
Let $K=\mathbb Q(\sqrt2,\sqrt3)$. The norm of every totally positive $d$-number in $K$ is a rational square. In either character algebra of a rank-six category with dimension field $K$, the quadratic field of a pair of distinct complementary codegrees, if present, is $\mathbb Q(\sqrt2)$ and the norm of its codegree is twice a square.
\end{lemma}
\begin{proof}
For an invariant ideal the norm valuation at a prime of ramification
index $e$ is a multiple of $4/e$. Only two can contribute an odd
valuation, so the positive norm is a square or twice a square. The
latter is impossible: it would be a norm from $\mathbb Q(\sqrt3)$,
but $a^2-3b^2=2c^2$ has no nonzero integral solution by descent
modulo three.

For distinct irrational complementary codegrees $f,f'$ write $F=\mathbb Q(f)=\mathbb Q(\sqrt d)$, with $d=2,3$ or $6$. The rational algebra of the corresponding quotient is $K\times F$. By its order discriminant identity, $N_K(H)ff'$ has the square class of $\operatorname{disc}(K)\operatorname{disc}(F)$, where $H=D$ or $E$. The first part and the square discriminant of the biquadratic field imply $N_F(f)=d h^2$ for some $h\in\mathbb Q^\times$. Consequently $f/(h\sqrt d)$ has norm $-1$. This is impossible for $d=3$ or $6$: clearing denominators in a putative norm equation gives $a^2+c^2=d b^2$, which forces $3\mid a,c$ and then $3\mid b$. This leaves $d=2$.
\end{proof}

\begin{lemma}[A small weight in the presence of a quadratic pair]\label{lem:1624-weight-cut}
Suppose the signed complementary codegrees $f,f'$ are distinct and
lie in $\mathbb Q(\sqrt2)$. At either conductor the unnormalized
weights $q_i=i(1/E)$ satisfy $t-|Q_8|>2/94$. Now exclude the pair
$8\pm4\sqrt2$ and take $q_0$ least. At conductor $24$ let $j$ fix
$\sqrt3$. At conductor $16$ let $j$ generate the cyclic group and
assume $t<71/100$. Then
\[
 \begin{array}{c|c|c}
 c&q_j\text{ at most}&\text{consequence}\\\hline
 24&1/188&s<2/3,\quad(1/f-1/f')^2>5/36\\
 16&1/1000&s<5/8 .
 \end{array}
\]
\end{lemma}
\begin{proof}
Put $u=1-t$, $v=1/f-1/f'$ and $N=ff'$, with $f>f'$.
The projection gives $Q_8^2+2v^2\le1/2$. For $N\ge99$,
$t^2+2u^2\ge2/3$ and $v^2=u^2-4/N$ give
$t-|Q_8|\ge1/12-4/N>2/94$. For $N\le98$ the signed row of
Lemma~\ref{lem:quartic-short-pairs} gives $(f+f',N)=(24,72)$ or
$(16,32)$. In the former case $t-|Q_8|\ge2/3-\sqrt{7/18}>2/94$;
in the latter use $|Q_8|\le\sqrt2/3$ from
Lemma~\ref{lem:quadratic-compression} and $t=1/2$.

The unit
$f/f'=\eta^n$, $\eta=3+2\sqrt2$, gives
$v^2/u^2=((\eta^n-1)/(\eta^n+1))^2$.
For $n=1$ one has $f=k(2+\sqrt2)$ with $k$ integral, and
$8\mid f+f'=4k$ makes $k$ even. Thus $u=2/k$.
For $n\ge2$ we have $v^2\ge8u^2/9$, and for $n\ge3$
we have $v^2\ge49u^2/50$.

At conductor $24$, Lemma~\ref{lem:round2-quartic24-norm} gives
$ff'=2h^2$. The square root $f/(h\sqrt2)$ of $f/f'$ is a unit of
norm $-1$ in $\mathbb Q(\sqrt2)$, hence $n$ is odd.
Write $e=1/188$ and suppose $q_j\le e$. Then
$t\le1/\sqrt3+4e<3/5$ by the conductor-twelve bound, while
$t\ge1/2$. This rules out $n=1$ except for the excluded $k=4$.
In particular $v^2\ge49u^2/50>5/36$.
The identities $P_{12}=2(p_0+p_j)-s$ and
$P_8+P_{24}=2(p_0-p_j)$ give
$s\le4e+|P_8|+|P_{24}|+|P_{12}|$. The separate and mixed bounds
give $|P_{24}|\le1/\sqrt{12}$ and
\[
 |P_8|^2\le1/4-v^2,\qquad
 |P_{12}|^2\le 1/6-(t-4e)\sqrt3/6 .
\]
Use $(a+b)^2\le(4/3)(a^2+3b^2)$ for these last two moduli.
Completing the square in $u=1-t$ gives
\[
 s\le4e+1/\sqrt{12}
  +\sqrt{\tfrac43\bigl(\tfrac34-(1-4e)\sqrt3/2+75/392\bigr)}
 <2/3 .
\]

At conductor $16$ put $e=1/1000$ and suppose $q_j\le e$.
The adjacent imbalance gives $\sqrt2|Q_1|\ge t-4e$.
From $s\le4e+2|P_1|+|P_2|$ and the two Fourier inequalities
of Lemma~\ref{thm:cyclic16-irrational-finite} we obtain
\begin{equation}\label{eq:round2-c16-cut}
 s\le4e+\sqrt{1/2-(t-4e)/\sqrt2}+\sqrt{1/4-v^2}.
\end{equation}
If $n=1$, the mass range leaves only $k=6$, so $t=2/3$ and
$v^2=1/18$; the two radicals are less than $.179$ and $.441$.
Otherwise $(a+b)^2\le2(a^2+b^2)$ and $v^2\ge8u^2/9$ give
\[
 s\le4e+\sqrt{2\bigl(3/4-(1-4e)/\sqrt2+9/64\bigr)}<5/8 .
\]
All strict numerical bounds follow by squaring with
$1/\sqrt2>.7071$, $\sqrt3/2>.86602$ and $1/\sqrt{12}<.28868$.
\end{proof}

\subsubsection{Conductor twenty-four with a quadratic complementary field}
\begin{theorem}\label{prop:quartic-twenty-four-silver-general}
Let $K=\mathbb Q(\sqrt2,\sqrt3)$ be the dimension field, with distinct
complementary signed codegrees $f,f'$ in $\mathbb Q(\sqrt2)$.
Then the category admits a spherical structure.
\end{theorem}
\begin{proof}
Write $p_i=1/i(D)$, $q_i=1/i(E)$, with sums $s,t$.
Put $N=ff'$, $L=f+f'$, $v=1/f-1/f'$, so $v^2=(1-t)^2-4/N$.
The short pair $L=16,N=32$ is covered by
Theorem~\ref{thm:quartic-twenty-four-silver}.

If both ordinary characters are rational of codegree four,
Lemma~\ref{lem:quartic-four-rows} gives $s=t=1/2$ and
$f=2+2y^2$ for a norm-minus-one unit $y\in\mathbb Z[\sqrt2]$.

The fundamental unit is $1+\sqrt2$. Apart from the short pair,
$|y|$ or $|y'|$ is at least $(1+\sqrt2)^3$; hence
$N=4(2+y^2+y'^2)\ge800$ and $v^2\ge49/200$.

In this paragraph normalize $p,q$ by doubling both. The
conductor-eight projection estimate used in Theorem~\ref{thm:quartic-twenty-four-silver}
gives $2P_8^2+Q_8^2\le2-8v^2\le1/25$.
The simultaneous-column norm estimate there still holds: the ordinary
pair contributes $1/2$ and the signed contribution is at least $1/2$.
The complementary contribution to the signed coefficient of the product of the Perron
column and its integral Galois translate at $j$ fixing $\sqrt3$ is
$N_F(\beta)/2$,
as $j$ interchanges $f,f'$. Thus the same positive integral
moment and the same height and separation certificates apply. Their last
ratio has both conjugates on the positive $\chi_8$ half less than $2/5$,
as read from its displayed radical expression. This gives an ordinary
mass there at most $(2/5)(1+1/5)/2=6/25$, while
$|P_8|\le1/\sqrt{50}$ requires mass exceeding $2/5$.

Return to unnormalized coordinates. We claim that every remaining
ordinary complement has $s\ge63/100$. Each rational row has even
codegree: on the four signed dimension labels its normalized entries
are one, on the two signed complementary labels they are equal
integers, and their sum is the codegree. For split rational
characters of codegrees $g\le h$ the discriminant identity makes
$gh$ a square by Lemma~\ref{lem:round2-quartic24-norm}. Unless
$(g,h)=(4,4)$, already treated, these even denominators give
$s\ge2/3$: for $g=4$ the next even square is $h=16$, and for
$g\ge6$ the reciprocal sum is at most $1/3$. Codegrees two and
three are excluded by Lemmas~\ref{lem:quartic-rational-mass-strata}
and~\ref{lem:quartic-three-induction}.
A nonsplit character pair of common rational codegree is impossible:
all simples are self-dual since the signed algebra is totally real,
whereas the discriminant identity would make the positive
quadratic field discriminant a square.

A distinct irrational ordinary pair lies in $\mathbb Q(\sqrt2)$ by
Lemma~\ref{lem:round2-quartic24-norm}, and $1-s\le1/2$ by
Lemma~\ref{thm:quadratic-pair-mass}. If $1-s\ge37/100$, both possibilities in
Lemma~\ref{lem:quartic-short-pairs} saturate the conductor-eight
ordinary bound, contradicting $v\ne0$ in the mixed projection.
This proves the claim. Moreover, writing $w$ for the ordinary
reciprocal gap, its unit ratio gives $w^2\ge(1-s)^2/2$.
Equation~\eqref{eq:orbit-projection} gives $4w^2+2v^2\le1/2$.
Thus $v^2>5/36$ forces $s>2/3$ in this last case as well.

Lemma~\ref{lem:1624-weight-cut} now gives $q_j>1/188$ for $j$
fixing $\sqrt3$. For the other two nonidentity indices use
$t-|Q_8|>2/94$ from the same lemma
and $t-|Q_{24}|\ge1/2-1/\sqrt6>2/94$. Since $q_0$ is least,
both weights are at least $1/188$. Thus $\operatorname{Tr}(x)<189$.

It remains to test primitive ratios through trace $188$ by
Theorem~\ref{thm:all-proper-quartic}. Put $p_i=x_iq_i$ and impose
$s\ge63/100$, $t\ge1/2$, the two common minima and the separate and
mixed Fourier bounds. In particular
\[
 2P_8^2+Q_8^2+2(1-t-1/12)^2\le1/2+1/72,
\]
since $L\ge24$ gives $8/N\le(1-t)/3$. The shift fixing $\sqrt2$
also satisfies $\operatorname{Tr}(q\sigma x)\ge s$ by
Corollary~\ref{cor:galois-columns}, with $q=1/E$. The
\codecheck{generic_quartic/general24/verify.py}{exact enumeration}
gives $247$ ratios; nonnegative rational combinations of
the outward-rounded inequalities exclude each of them.
\end{proof}
\subsubsection{Large reciprocal mass at conductor sixteen}
\begin{theorem}\label{thm:cyclic-sixteen-large}
Let $K=\mathbb Q(\zeta_{16})^+$ be the global Frobenius--Perron dimension field of a complex fusion category of rank six. If $t=\operatorname{Tr}_K(1/E)\ge71/100$, a Galois conjugate is pseudo-unitary.
\end{theorem}
\begin{proof}
With the notation of Lemma~\ref{lem:quartic-large-mass-height},
$A=\operatorname{Tr}(E/D)<740$. Take $\sigma$ a cyclic generator,
put $s=\sum p_i$ and write $P_k,Q_k$ for the cyclic transforms. By
Theorem~\ref{thm:all-proper-quartic} we may assume $x=E/D$ primitive.

At cap $740$ the cyclic enumeration yields $695$ ratios. More explicitly write $y=\sqrt{2+\sqrt2}$, $z=\sqrt{2-\sqrt2}$. The integral basis $1,\sqrt2,y,z$ follows from $\mathcal O_K=\mathbb Z[y]$, since $y^2=2+\sqrt2$ and $y^3=3y+z$. The generator $\sigma$ sends $y$ to $z$ and $z$ to $-y$. The presentation input for the cyclic enumerator is $(8,16,2,1,1,1,1,740)$. Use the inequalities for the separate and coupled projections, the two coordinate minima, $s\le1$ and $71/100\le t\le1$. Inserting $p_i=x_iq_i$ gives a strict rational separating certificate for all but one polynomial. The remaining ratio has $(A,u,v,w)=(460,162,106,44)$ in the presentation $A/4+u\sqrt8/4+v y+w z$, up to conjugacy.

Here $s\ge1/2$ by Lemmas~\ref{lem:quartic-rational-mass-strata} and~\ref{lem:quartic-three-induction} for rational signed complements and by the ordinary-pair argument of Lemma~\ref{thm:cyclic16-irrational-finite} otherwise. The involution $\tau=\sigma^2$ fixes the complementary codegrees, so
Corollary~\ref{cor:galois-columns} gives
$\operatorname{Tr}(q\tau(x))\ge s$. Four additional separating certificates for this one ratio yield
\[
 t<711/1000,\qquad 63/100<s<67/100,\qquad 2(q_0+q_2)<2/94.
\]
The last inequality is incompatible with a distinct irrational signed pair in $\mathbb Q(\sqrt2)$ by Lemma~\ref{lem:1624-weight-cut}, since $t-|Q_2|\le2(q_0+q_2)$. In the rational case the complementary codegrees are integers at least four. In the interval $71/100\le t<711/1000$, the reciprocal mass $1-t$ forces the pair $(4,25)$: the smaller denominator is at most six, and five and six permit no integral partner in this interval. The two signed evaluations $a,b$ of the ordinary FP induction column are then positive integers, satisfying $s=1-a/4-b/25$. The displayed interval forces $a=1$, and contains no number of the form $3/4-b/25$. The verifier \codefile{generic_quartic/sixteen.py} checks the finite enumerations and all $698$ strict certificates.
\end{proof}

\subsubsection{The nonhalf rational pairs at both conductors}
\begin{lemma}\label{lem:quartic24-fast}
Let $K=\mathbb Q(\sqrt2,\sqrt3)$, with $U_1=3+2\sqrt2$,
$U_2=2+\sqrt3$, $U_3=5+2\sqrt6$. Every totally positive integral
$d$-number in $K$ has a unique expression
$n z\prod_{i=1}^3 U_i^{k_i}$, with $n\in\mathbb Z_{>0}$,
$k_i\in\mathbb Z$, and $4z=A+a\sqrt8+b\sqrt{12}+c\sqrt{24}$
from the following table:
\[
\begin{array}{rrrr|r}
 A&a&b&c&N(z)\\\hline
 4&0&0&0&1\\
 8&-2&0&0&4\\
 8&-1&0&-1&1\\
 12&-3&-2&1&9\\
 12&0&-2&0&36\\
 12&0&0&-2&9\\
 12&2&-2&-2&4\\
 24&-6&4&-4&36
\end{array}
\]
\end{lemma}
\begin{proof}
First remove the largest positive rational integer dividing $z$.
The remaining ideal has norm $2^e9^f$, with $0\le e<4$ and
$0\le f<2$: the ramification indices at two and three are four and
two, and the residue degree at three is two. Lemma~\ref{lem:round2-quartic24-norm}
makes $e$ even.

In any totally real biquadratic field, let $z$ be a totally positive
$d$-number of norm $m^2$. Its relative norm to each quadratic
subfield $L_i$ has invariant ideal of norm $m^2$, hence ideal $(m)$:
the norm determines the common valuation above every prime. If $U_i$
generates the totally positive units of $L_i$, it follows that
\begin{equation}\label{eq:biquadratic-relative-norms}
 N_{K/L_i}(z)=mU_i^{r_i},\qquad z^2=m\prod_i U_i^{r_i}.
\end{equation}
The second equality follows by multiplying the three relative norms.

Here $m=2^\alpha3^\beta$ with $\alpha,\beta\in\{0,1\}$, and the
quadratic Pell equations give the three units in the statement.
Multiplying by $U_i$ changes only $r_i$, by two, so take
$r_i\in\{-1,0\}$. The positive square roots of $U_i$ are
$1+\sqrt2$, $(\sqrt2+\sqrt6)/2$ and $\sqrt2+\sqrt3$.
Their signs at the four embeddings are those of $\sqrt2,\sqrt6,\sqrt3$,
respectively. Thus total positivity of
$\sqrt m\prod_i(\sqrt{U_i})^{r_i}$ is equivalent to
\[
 r_1+r_2\equiv\alpha,\qquad r_2+r_3\equiv\beta\pmod2.
\]
The two solutions for each $m$ give the eight entries displayed.
They are integral in the basis $1,\sqrt2,\sqrt3,(\sqrt2+\sqrt6)/2$,
which in the table coordinates means
$4\mid A$, $2\mid b$ and $a\equiv c\pmod2$. Their norms and the
unique primes over two and three give invariant ideals.

If $\operatorname{Tr}(z)\le C$ then both embeddings of each relative
norm are at most $C^2$, and their product is at least one.
Equation~\eqref{eq:biquadratic-relative-norms} therefore bounds the
three exponents in the stated decomposition by
$|k_i|\le\log C/\log U_i+1/2$. The enumeration in
\codefile{quartic24_fast.py} uses this box, the eight seeds and
positive rational integers. It keeps one representative per Galois
orbit with a conjugate below one.
\end{proof}

\begin{theorem}\label{thm:quartic24-rational-nonhalf}
Let $K=\mathbb Q(\zeta_c)^+$ with $c=16$ or $24$ be the dimension
field of a rank-six fusion category. If its signed complementary
codegrees are $(4,h)$ with $h=8,9,16$ or $18$, the category has a
pseudo-unitary Galois conjugate.
\end{theorem}
\begin{proof}
By Theorem~\ref{thm:all-proper-quartic} we may assume
$x=E/D$ primitive after transport maximizing $E$. Write $h$ for
the second codegree. The dyadic completion of $K$ is a field.
Thus if $h$ is even, the rational rows have the same mod-two
reduction: otherwise both local factors would need the unique field
branch, as their idempotents have nonintegral trace. The four-row
has four odd entries, so $4\mid h$, excluding eighteen.

Put $s=\operatorname{Tr}_K(1/D)$, $t=\operatorname{Tr}_K(1/E)$.
Here $s\ge1/2$ by Lemma~\ref{lem:quartic-rational-mass-strata}:
$s=1/3$ gives ordinary
complementary codegrees three, excluded by
Lemma~\ref{lem:quartic-three-induction}.
The two rational evaluations of the signed FP column are positive
integers, so the first is one, and if the second is $b$ then
$s=3/4-b/h$. Congruence of the two rows makes $b$ odd at even $h$.
The only alternatives to $b=1$ are $(h,b)=(9,2),(16,3)$.
Their ordinary complementary masses are $17/36,7/16$ respectively.
Neither is the sum of reciprocals of two integers $\ge4$: the
smaller would be four and the other would be $9/2,16/3$.
Ordinary codegrees two and three have already been excluded. Thus an
irrational ordinary pair of mass $a/d=17/36$ or $7/16$ would be
required. Lemma~\ref{lem:fixed-quadratic-mass} gives a squared
reciprocal gap at least $(a^2-4)/d^2>1/8$ in both cases. This exceeds
the ordinary projection bound at each possible conductor $8,12,24$.
Therefore $b=1$ and $s=t\ge5/8$.

Normalize $p_i=i(1/D)/s$ and $q_i=i(1/E)/s$ to have sum one.
By Lemma~\ref{lem:mass-moments}, the shifted first and mixed moments
$L(jx),L(xj(x))$ are at least one, where $L(v)=\operatorname{Tr}(qv)$.
For each character of conductor $d$, the normalized Fourier
coefficients satisfy
\[
 d s^2|P|^2\le(d,2),\qquad d s^2|Q|^2\le(d,4),\qquad
 d s^2|P|^2+\sqrt d\,s|Q|\le2.
\]
Use the weaker common value $s=5/8$ in these bounds.

A small coordinate $q_j$, $j\ne0$, would contradict a mixed moment.
Write the group additively and put $a=p_0/q_0$, $b'=p_0/q_j$.
Simultaneous-column dominance gives $0<b'\le a\le1$ and
$a q_i\le p_i$, $b' q_i\le p_{i-j}$. Reindexing the four terms
of the appropriate mixed moment gives
\begin{equation}\label{eq:round2-cyclic-reverse-product}
  1\le a p_{-j}+b'p_j+
       \sum_{i\notin\{0,j\}}p_i p_{i-j}/q_i .
\end{equation}
The other two denominators have uniform lower bounds.
At conductor $24$, pair inversion and minimality of $q_0$ give
$q_i\ge(1-8/(5\sqrt3))/4>1/60$ outside the kernel of the
conductor-eight character. Thus only its nonidentity element needs
testing. At conductor $16$, if $q_j\le1/100$ then the other two
coordinates sum to at least $.98$. The normalized faithful inequality
$|Q_1|\le.8$ bounds their difference for an opposite pair and each
coordinate by $.81$ for an adjacent pair. Both therefore exceed $1/20$.

The following two finite tests complete the argument in each field.
First impose $q_j\le\delta$, the probability and Fourier bounds,
and the column inequalities above. Bisection of the three remaining
$p$ intervals and two $q$ intervals, rounding outwards to
$2^{-20}\mathbb Z$, makes the right side of
\eqref{eq:round2-cyclic-reverse-product} strictly less than one on
every feasible box. The resulting lower bounds on $q_i$ for $i\ne0$,
together with $x_0\le1$, give the indicated integral trace caps.
Then enumerate the primitive ratios at that cap. For fixed $x$
substitute $p_i=x_iq_i$ in the probability, minimum, Fourier and
moment inequalities. Nonnegative rational combinations of the
outward linear bounds contradict $\sum q_i=1$ for every candidate.
\[
\begin{array}{c|c|r|r|r}
 c&\delta&\text{boxes}&\operatorname{Tr}(x)\text{ cap}&
   \text{primitive ratios}\\\hline
 24&1/50&3117&60&62\\
 16&1/100&125+1379&100&86
\end{array}
\]
The complete checks are linked separately for conductor
\codecheck{quartic24_nonhalf_height.py}{twenty-four}
(\codecheck{quartic24_nonhalf_check.py}{separating weights}) and
\codecheck{field_checks/sixteen/nonhalf_height.py}{sixteen}
(\codecheck{field_checks/sixteen/nonhalf_verify.py}{separating weights}).
\end{proof}

\subsection{Reduction to common complementary codegrees four}
\begin{corollary}\label{cor:round2-quartic24-irrational}
If $K=\mathbb Q(\sqrt2,\sqrt3)$ and the complementary signed codegrees are distinct irrational conjugates, the category admits a spherical structure.
\end{corollary}
\begin{proof}
Combine Lemma~\ref{lem:round2-quartic24-norm} with Theorem~\ref{prop:quartic-twenty-four-silver-general}.
\end{proof}

\begin{proposition}\label{prop:round2-quartic24-ordinary}
At either conductor, if both signed complementary codegrees are four
and the ordinary complementary codegrees are distinct irrational
conjugates, some Galois conjugate is pseudo-unitary.
\end{proposition}
\begin{proof}
We may assume $x=E/D$ primitive by Theorem~\ref{thm:all-proper-quartic}.
Lemmas~\ref{lem:quartic-rational-mass-strata} and
\ref{lem:quartic-three-induction} give $s=t=1/2$.
The ordinary pair must therefore be $8\pm4\sqrt2$ by
Lemma~\ref{lem:quartic-short-pairs}. Its reciprocal gap saturates
the conductor-eight projection. Thus $P_8=Q_8=0$ by
\eqref{eq:orbit-projection}.

Double $p,q$ to give them sum one and let $j$ fix $\sqrt2$.
For each $i$ the pair $i,ij$ has total weight $1/2$ in both
vectors. Hence $p_i=x_iq_i$ determines the weights:
\[
             q_i=\frac{x_{ij}-1}{2(x_{ij}-x_i)} .
\]
Primitivity ensures a nonzero denominator. It remains to check
positivity and minimality of $q_0$ at the least conjugate of $x$.
For finiteness we bound the other three $q_i$ from below.

At conductor $24$ let $b,c$ fix $\sqrt3,\sqrt6$, respectively.
We have $q_j\ge1/4$ and $q_c\ge(1-\sqrt{2/3})/4>1/64$.
If $q_b\le1/64$, then $p_0\le q_0\le q_b$,
$|Q_{12}|\ge15/16$, and $P_{12}+P_{24}=4p_0-1$.
The Fourier bounds would give
\[
 15/16\le |P_{12}|+|P_{24}|
 \le\sqrt{(8-2\sqrt{12}\,(15/16))/12}+1/\sqrt3<15/16.
\]
Thus $\operatorname{Tr}(x)\le64$.
At conductor $16$ write $j=2$ in cyclic indices. If $q_1$ or
$q_3$ were at most $.01$, the opposite half sums would give
$|Q_1|\ge .48\sqrt2$, $|P_1|\ge .48$, contrary to
$4|P_1|^2+2|Q_1|\le2$. Since $q_2\ge1/4$,
$\operatorname{Tr}(x)\le100$.

Substitute the primitive ratios at these caps into the displayed
formula. The two exact checks give
\[
 \begin{array}{c|r|r|r}
 \text{conductor}&\text{ratios}&q>0&q>0,\ q_0\text{ least}\\\hline
 \codecheck{field_checks/twentyfour/ordinary_silver.py}{24}&69&13&0\\
 \codecheck{field_checks/sixteen/half_silver.py}{16}&86&20&0 .
 \end{array}
\]
\end{proof}

\begin{lemma}\label{lem:1624-rational-shortlist}
At either conductor, suppose the two signed complementary codegrees
are rational and $t<71/100$. Then they are both four, or the
category has a pseudo-unitary Galois conjugate.
\end{lemma}
\begin{proof}
Codegrees two and three are excluded by
Lemmas~\ref{lem:quartic-rational-mass-strata} and~\ref{lem:quartic-three-induction}.
For two rational characters of codegrees $g\le h$, the discriminant
identity gives $gh$ square or twice a square. Indeed $N_K(E)$ is a
square at conductor $24$ by Lemma~\ref{lem:round2-quartic24-norm};
at conductor $16$ its odd valuations are multiples of four because
the ideal $(E)$ is Galois invariant and unramified at odd primes.
The two quartic discriminants have squareclasses one and two.
Since $1/g+1/h>29/100$, the pairs are
\[
 (4,4),(4,8),(4,9),(4,16),(4,18),(5,5),(5,10),(6,6).
\]
For $g=5,6$ take $p=5,3$ respectively. The quartic branches contain
the unramified quadratic field at $p$, so the two rows coincide
modulo $p$ by Lemma~\ref{lem:silver-local-partner}; their squared
distance $g+h$ is not divisible by $p^2$. Theorem~\ref{thm:quartic24-rational-nonhalf}
handles $(4,h)$ for $h>4$.

For a quadratic character orbit the common codegree $g$ is $4,5$
or $6$. A real field of discriminant $d$ needs $d\mid2g$ and
$d\le g$ by Lemma~\ref{lem:quadratic-row-difference}. Only $d=g=5$
survives, contrary to the discriminant identity at five. In the
imaginary case there is one non-self-dual pair, since duality has a
one-dimensional negative eigenspace in $K\times F$. Hence $g=|d|b^2$
by the same lemma, so $g=4$.
\end{proof}

\begin{theorem}\label{thm:round2-quartic24-finite}
If $K=\mathbb Q(\sqrt2,\sqrt3)$, failure of sphericality forces both pairs of complementary formal codegrees to be $4,4$.
\end{theorem}
\begin{proof}
Corollary~\ref{cor:round2-quartic24-irrational} and
Theorem~\ref{thm:quartic-twenty-four-large-mass} reduce the signed
complement to rational codegrees and $t<71/100$.
Apply Lemma~\ref{lem:1624-rational-shortlist}.
The signed pair is thus $4,4$ and $s=1/2$. An irrational ordinary complementary pair is excluded by Proposition~\ref{prop:round2-quartic24-ordinary}. The remaining ordinary codegrees are both at least four and their reciprocals sum to $1/2$, so both equal four.
\end{proof}

\subsubsection{Reduction to two pairs of codegree four at conductor sixteen}
\begin{theorem}\label{thm:round2-sixteen}
Let $K=\mathbb Q(\zeta_{16})^+$ be the global Frobenius--Perron dimension field of a complex fusion category of rank six. Unless both ordinary and signed complementary codegrees are $4,4$, the category has a pseudo-unitary Galois conjugate.
\end{theorem}
\begin{proof}
Choose the embedding maximizing $E$ and take $x=E/D$ primitive by Theorem~\ref{thm:all-proper-quartic}. Throughout, $p_i=\sigma^i(1/D)$ and $q_i=\sigma^i(1/E)$ with respective sums $s,t$, where $\sigma$ generates the cyclic Galois group and indices run modulo four. We may assume $t<71/100$ by Theorem~\ref{thm:cyclic-sixteen-large}. Write $P_k,Q_k$ for the Fourier transforms at frequencies $k=1,2$ of conductors $16,8$. We have $s\ge1/2$, by Lemmas~\ref{lem:quartic-rational-mass-strata} and~\ref{lem:quartic-three-induction} in the rational case and by the ordinary-pair argument in Lemma~\ref{thm:cyclic16-irrational-finite} otherwise. Also $t\ge1/2$. If $x_i=p_i/q_i$, the separate, mixed and simultaneous-column bounds are
\[
 c|P_k|^2\le(c,2),\quad c|Q_k|^2\le(c,4),\quad
 c|P_k|^2+\sqrt c|Q_k|\le2,\qquad p_0q_{l+i}\le p_iq_l.
\]
Here $(c,d)=\gcd(c,d)$. We distinguish the two possibilities for the signed complement.

\smallskip\noindent\emph{Rational formal codegrees.}
Lemma~\ref{lem:1624-rational-shortlist} leaves two signed codegrees
four, so $s=t=1/2$. Rational ordinary codegrees, all at least four,
then also equal four.
Proposition~\ref{prop:round2-quartic24-ordinary} excludes the
distinct irrational ordinary case.

\smallskip\noindent\emph{Distinct irrational signed codegrees.}
Return to the unnormalized vectors and set $u=1-t$, $v=1/f-1/f'$, $L=f+f'$, $N=ff'$. This field has just one quadratic subfield, $\mathbb Q(\sqrt2)$. Lemmas~\ref{lem:1624-weight-cut} and~\ref{thm:cyclic16-irrational-finite} give
\[
 t-|Q_2|\ge 2/94,\qquad 2P_2^2+Q_2^2+2v^2\le1/2.
\]
As $q_0$ is least, $q_2\ge1/188$. Let $e=1/1000$. For a nonshort signed pair, if $q_1\le e$ or
$q_3\le e$, Lemma~\ref{lem:1624-weight-cut} gives $s<5/8$.

For $s<5/8$ the ordinary pair must consist of two rational four-rows. Indeed split rational characters have even codegrees (the equal entries at the two signed complementary labels give $g=4+2k$), and their product is a square or twice a square. Reciprocal sum greater than $3/8$ then leaves just $(4,4)$. The signed algebra is totally real, so all simples are self-dual. A quadratic ordinary pair of common rational codegree four or five is excluded by the squared row-difference test, using also the order-discriminant identity at five. If the ordinary pair is irrational, its field is $\mathbb Q(\sqrt2)$.
Lemma~\ref{lem:quartic-short-pairs} leaves only the ordinary codegrees
$8\pm4\sqrt2$ at reciprocal mass greater than $3/8$.
Their projection equality contradicts $v\ne0$.

With two rational ordinary four-rows there is an actual invertible with two transpositions, by the proof of Lemma~\ref{thm:cyclic16-irrational-finite}. Hence $s=t=1/2$ and $f=2+2y^2$ for $y\in\mathbb Z[\sqrt2]$ with $yy'=-1$. Outside the short pair, $N\ge800$ and $v^2\ge49/200$. Thus $|P_2|\le1/\sqrt{200}$. Under the assumed collapse, $p_2-p_0\ge (s-|P_2|)/2-2e>21/100$, and $|Q_1|\ge(t-4e)/\sqrt2$. These bounds violate $4|Q_1|+16|P_1|^2\le2$. This proves the claim for every nonshort pair and gives $\operatorname{Tr}_K(x)\le1000$.

For the short pair, Lemma~\ref{prop:quartic-twenty-four-fold} gives
$s=t=1/2$, Gram dominance, positive integral mixed moments and odd
first moments. Normalize $p,q$ to sum one. We have
$2P_2^2+Q_2^2\le1$, $|Q_2|\le2\sqrt2/3$ and
$4|P_1|^2+2|Q_1|\le2$. If $q_1\le e$, then $p_0,q_0\le e$ and
$q_2,q_3>1/40$ by the quadratic Fourier bound.

For clarity we describe the interval contradiction certifying $q_1>e$. With $a=p_0/q_0$, $b=p_0/q_1$, impose $0\le b\le a\le1$ and
\[
 a q_i\le p_i,\quad b q_{i+1}\le p_i,\quad
 a q_i^2+bq_{i+1}^2\le p_i^2+p_0\sum_{l\ne0}p_{l+i}^2/p_l
 \qquad (i=1,2,3).
\]
The shifted moment is $(b/a)p_1+q_1p_2/q_2+q_2p_3/q_3+q_3a$. Outward interval bounds for its last three terms and the odd integral values in its range give two further linear inequalities on $a,b$. On each box, two-variable linear duality then bounds $ap_3+bp_1$ in \eqref{eq:round2-cyclic-reverse-product}. The verifier \codefile{field_checks/sixteen/short_height.py} rules out all $697$ boxes by Fourier or shifted-moment infeasibility or by $k_{-1}<1$. The same holds after reversing the generator; the quadratic pair bound gives $q_2>e$. The trace cap $1000$ follows in this case as well.

\smallskip\noindent\emph{Primitive ratios through trace $1000$.}
Work again with unnormalized masses. In addition to the projection bounds, $1/2\le s$, $1/2\le t<71/100$, and the two lower bounds $2(q_0+q_2),2(q_1+q_3)\ge2/94$, the involution $\tau=\sigma^2$ fixes the complementary character field. Thus $\sum q_i(x_{i+2}-x_i)\ge0$, by Corollary~\ref{cor:galois-columns}. Since $f/f'$ is a nontrivial totally positive unit, $v^2\ge u^2/2$. At the $915$ ratios furnished by the cyclic enumerator, these inequalities and $2P_2^2+Q_2^2+u^2\le1/2$ leave twelve. For those twelve use the following four cases.
\begin{enumerate}
\item At the short pair both masses are $1/2$, all normalized first moments are at least one, and the normalized mixed moments are at least one. For even shifts the latter follows from integrality of $2\operatorname{Tr}(\beta^2/f)$ for the three possible $\beta$. Here $2P_2^2+Q_2^2\le1/4$.
\item At an ordinary four-pair and a longer signed unit $y$, both masses are $1/2$ and $2P_2^2+Q_2^2\le1/100$. All normalized first moments are positive integers: multiply each shifted FP column divided by its simultaneous unit by $1+g$ to isolate the dimension factor.
\item In the remaining cases $s\ge5/8$ by the ordinary-pair argument. If $f/f'=3+2\sqrt2$, then $f=m(2+\sqrt2)$ with $m$ integral and even by $8\mid L$. The mass interval leaves only $m=4,6$. The first was treated above; the second gives $t=2/3$, $2P_2^2+Q_2^2\le1/2-1/9$.
\item For the remaining powers of the fundamental positive unit, $v^2\ge8u^2/9$, giving $2P_2^2+Q_2^2+16u^2/9\le1/2$.
\end{enumerate}
Using rational supporting half-spaces of these ellipsoids excludes every ratio except one in the last case. The sole polynomial and a representative are
\[
 X^4-24X^3+36X^2-16X+2,\qquad x=6+4\sqrt2+5y+2z,
\]
with $y=\sqrt{2+\sqrt2}$, $z=\sqrt{2-\sqrt2}$. The same tests give $51/100<t<7/10$ for this ratio. The verifier \codefile{field_checks/sixteen/irr_verify.py} checks the $952$ strict certificates by outward radical enclosures and integer arithmetic.

The remaining obstruction is a trace duality argument that will also
be used at conductor forty. Let $S=\mathbb Z[Gx]$, and suppose
$\tau\in G$ fixes the complementary evaluations of the FP column.
By Corollary~\ref{cor:galois-columns} its translated columns have
coordinates $(gx,g\beta)$. Their differences and products show that
the signed order contains $(J,0)$, where
\begin{equation}\label{eq:quartic-difference-ideal}
 J=\sum_{g\in G}S\,g(\tau x-x),\qquad
 \operatorname{Tr}_K(q_0J)\subseteq\mathbb Z .
\end{equation}

Here $[\mathcal O_K:\mathbb Z[x]]=17$ by the polynomial discriminant,
while $1,x,\sigma x,\tau x$ have determinant $232$ in the
basis $1,\sqrt2,y,z$. Hence $S=\mathcal O_K$.
Now $\tau x-x=-2(5y+2z)$ has norm $32$, since
$(5y+2z)^2=58+41\sqrt2$. Both its ideal and the different (of
norm $2^{11}$) are supported at the unique dyadic prime, so their
product is $16\mathcal O_K$. Equation~\eqref{eq:quartic-difference-ideal}
gives $16q_0\in\mathcal O_K$. All integral traces are multiples of
four in the displayed basis, forcing $t\in\frac14\mathbb Z$,
contrary to $51/100<t<7/10$.
\end{proof}

\subsection{The common four-codegree case}
\begin{lemma}[A lower bound for the trace]\label{thm:half-million}
At either conductor, suppose both signed complementary codegrees equal
four. Unless a Galois conjugate is pseudo-unitary, transport maximizing
$E$ gives $\operatorname{Tr}_K(E/D)>10^6$.
\end{lemma}
\begin{proof}
By Lemmas~\ref{lem:quartic-rational-mass-strata} and
\ref{lem:quartic-three-induction}, $s=t=1/2$. Put $x=E/D$, $q=2/E$,
$R_x=\mathbb Z[Gx]$, with $G=\operatorname{Gal}(K/\mathbb Q)$.
Theorem~\ref{thm:all-proper-quartic} and Lemma~\ref{lem:mass-moments}
allow us to take $x$ primitive with characteristic polynomial
$(X-1)^4$ modulo two. The functional $L(v)=\operatorname{Tr}(qv)$
is integral on $R_x$, odd on its monomials, and $L(1)=L(x)=1$.
The least conjugate of $x$ is below one and minimizes both $q$ and $qx$.

\emph{Bounding the denominator.} Write $q=z/T$ with $z$ a primitive
integral lattice vector and $T=\operatorname{Tr}(z)$. Then $T\mid M$
whenever $Mq$ is integral. The following bounds come from the trace
dual of explicit elements of $R_x$.

At conductor $24$, write $4x=A+\sum c_i\sqrt{d_i}$ for
$(d_1,d_2,d_3)=(8,12,24)$. Character sums give $c_i\sqrt{d_i}$ in
$R_x$. Their pairwise products give $g_i\sqrt{d_i}\in R_x$, where
$g_i=\gcd(c_i,\sqrt{d_jd_k/d_i}\,c_jc_k)$ for $\{i,j,k\}=\{1,2,3\}$.
Each $g_i$ is nonzero by primitivity. Pairing with $q$ gives
$T\mid4\operatorname{lcm}_i(d_i g_i)$.

At conductor $16$, take $y=\sqrt{2+\sqrt2}$ and $w=\sqrt{2-\sqrt2}$
with $x=a+b\sqrt2+cy+dw$. The conjugate sums and differences give
$4b\sqrt2$, $2(cy+dw)$ and $2(cw-dy)$ in $R_x$. Products and
integer combinations give $H\sqrt2$, $2Ry$ and $2Rw$, where
$H=4\gcd(b,c^2-d^2+2cd,c^2-d^2-2cd)$ and $R=c^2+d^2$.
The integral basis $1,\sqrt2,y,w$ has trace Gram diagonal $(4,8,8,8)$,
so $T\mid\operatorname{lcm}(4,8H,16R)$.

\emph{Enumerating the two integral $d$-numbers.} At conductor $24$
use Lemma~\ref{lem:quartic24-fast}. At conductor $16$ the stripped
ideal, after division by its largest rational integer factor, has norm
$1,2,4$ or $8$: the only ramified prime is two, and it is totally
ramified. Use the positive units
$U=(1+\sqrt2)^2$, $V=(1+y)^2$, $W=(1+w)^2$. Their logarithms span
the trace-zero hyperplane: $\log U$ alternates in cyclic order and
the opposite half-differences of $\log V,\log W$ are
$(c_0,d_0),(d_0,-c_0)$ with $c_0,d_0>1$.
Every embedding of $U^{\pm1}$ is below six and of $V^{\pm1},W^{\pm1}$
below nineteen. Rounding all three logarithmic coordinates to the
nearest integer therefore puts a representative of every stripped
orbit at trace at most $4\cdot8^{1/4}\sqrt{6\cdot19^2}<320$.
The \codecheck{field_checks/half/seeds16.cpp}{integral basis enumeration}
at this bound retains the local trace minima under multiplication by
the three units and their inverses. Every orbit has such a minimum
in the box. Reducing these minima gives one representative
per norm:
\[
 1,\qquad 2-y,\qquad 2-\sqrt2,\qquad 4-2\sqrt2-y+w.
\]
For the exponent search at trace $C$, the centered log vector lies in
the simplex with vertices the permutations of $(-3,1,1,1)\log C$.
The inverse half-difference matrix bounds the $V,W$ coordinates by
$2\max(c_0,d_0)\log C/(c_0^2+d_0^2)<\log C$. Their alternating
components are half of $\log U$ up to sign, giving $2\log C$ for
the $U$ coordinate. Add the same bounds for the representative at
$C=320$.

List $x$ through trace $10^6$, then the stripped $z$ at the trace
divisors just obtained. Impose $L(x)=1$, oddness of $L(\sigma x)$,
the common minimum, oddness of $L(x^i)$ for $2\le i\le5$, and the
Fourier bounds, in that order. The
\codecheck{field_checks/half/verify.py}{exact enumeration} gives
\[
\begin{array}{c|r|r|r|r|r|r|r}
 c&x&(x,z)&L(x)=1&\text{shifts}&\text{minima}&\text{powers}&\text{Fourier}\\\hline
 24&113412&23270270&3311&2615&3&1&0\\
 16&38115&5032847&1597&1031&1&1&0
\end{array}
\]
\end{proof}
\subsubsection{Odd trace identities}
At either conductor, the nonspherical alternative has two complementary codegrees four in each sector and $\operatorname{Tr}(E/D)>10^6$. We use the normalized quantities $p=2/D$, $q=2/E$, $x=p/q$ and $L(v)=\operatorname{Tr}(qv)$. The normalization gives $L(1)=L(x)=1$. At conductor sixteen the signed complement is split or $\mathbb Q(i)$.
Lemma~\ref{lem:quadratic-row-difference} excludes a real quadratic
pair of common codegree four, and its imaginary analogue gives
$|\operatorname{disc} F|=4$. In either case the absolute discriminant
of the complement is a square. The order discriminant formula then
gives
\begin{equation}\label{eq:16-common-four-norm}
       N_K(q)\in2\mathbb Q^{\times2},\qquad \mathbb Q(q)=K,
\end{equation}
since $\operatorname{disc}(K)=2^{11}$ and an element of a proper
subfield of a cyclic quartic field has square norm.
\begin{lemma}\label{lem:four-opposite-recovery}
For either quartic field let $j$ be an involution, let $q$ and $p=qx$
be the two probability vectors, and put $L(v)=\operatorname{Tr}(qv)$.
Use the notation $b=q_j$, $z=p_j$, $a=x_0$, $r=q_0/q_j$,
$s=q_m/q_l\le1$ for the two other indices $l,m$. Assume that
$b\le1/250000$, $0<r<.003$, and that the entries of
\[
 (h,k,M,v,N)=
 \bigl(L(jx),L(xj(x)),L((jx)^2),L(W^2),L(jx/W)\bigr),
 \qquad W=j(q)/q,
\]
are odd positive integers. In the two indicated ranges the remaining
two moments are uniquely determined.
\[
\begin{array}{cc|ccc|cc}
 h&k&z&s&a&M&v\\ \hline
 3&1&[.228,.332]&[.138,.185]&(0,.270]&11&5\\
 7&3&[.228,.233]&[.054,.058]&(0,.409]&53&17
\end{array}
\]
\end{lemma}
\begin{proof}
Write $A=x_l$, $B=x_m$, $S=s+s^{-1}$, $t=r/b$ and
\[
 C=q_l+q_m=1-b(1+r),\qquad
 D_1=q_l A+q_m B=1-z-arb,\qquad H=q_lB+q_m A=h-zr-ba .
\]
Then $CAB=k-za(1+r)$. Grouping the four trace terms into the two
orbits of $j$ gives
\begin{align*}
 M&=H(H+D_1)/C-CAB+z^2t+ba^2,\\
 v&=C(S-1)+1/t+br^2,\\
 N&=HS-D_1+zr^2+a/t .
\end{align*}
For example, $q_lB^2+q_mA^2=H(A+B)-CAB$ and
$A+B=(H+D_1)/C$. These equalities isolate the only unbounded
quantity, $t$. Substituting the bounds in either row yields
\begin{align*}
 \big|M-[h(h+1)-k-z(h-a)+z^2t]\big|&<.02,\\
 \big|v-(S-1+1/t)\big|&<.001,\\
 \big|N-hv-(h-1)-z-(a-h)/t\big|&<.02 .
\end{align*}
For example, put $e=1-C$, $d=1-z-D_1=arb$ and $y=h-H=zr+ba$.
In both rows $e<.00000402$, $d<5\cdot10^{-9}$, $y<.001$ and $S<19$.
The first error is at most
$[eh(h+1)+y(2h+1)+hd]/C+zar+ba^2$, the second at most
$e(S-1)+br^2$, and the third at most
$yS+he(S-1)+d+(z+hb)r^2$. These give the displayed bounds.
If $t\le4$, the first expression places the odd integer $M$
in $(9,11)$ or $(51,53)$, respectively. Hence $t>4$. The last
expression now puts the even integer $N-hv-(h-1)$ in $(-2,2)$.
Its vanishing gives $7<t<15$ and $26<t<34$ in the two rows.
Inserting these intervals in the first two expressions gives
$9<M<13$, $4<v<7$ in the first row and
$51<M<55$, $16<v<18$ in the second. Oddness proves the table.
\end{proof}

\begin{lemma}\label{lem:cyclic16-small-adjacent}
Let $K=\mathbb Q(\zeta_{16})^+$ in the two complementary four-codegree
case, and let $p=2/D$, $q=2/E$ at an embedding maximizing $E$.
For a generator $g$ of $\operatorname{Gal}(K/\mathbb Q)$, write
$p_i=g^i(p)$, $q_i=g^i(q)$, $x_i=p_i/q_i$. Neither $q_1$ nor $q_3$
can be at most $4\cdot10^{-6}$. In particular, if
$\operatorname{Tr}(E/D)>10^6$, its largest conjugate is $x_2$.
\end{lemma}
\begin{proof}
We may reverse the generator and assume $e=q_1\le4\cdot10^{-6}$.
Put $r=q_0/e$, $a=x_0$, $Q=q_2$, $B=q_3$, $z=p_1$, $v=p_2$,
$w=p_3$, $b=v/Q$, and $c=w/B$. We have $0<a,r\le1$ and
$a\le b,c$. The faithful Fourier moduli are
\[
 P=\sqrt{(p_0-v)^2+(z-w)^2},\qquad
 T=\sqrt{(q_0-Q)^2+(e-B)^2}.
\]
The two mixed Fourier inequalities are
\[
 16P^2+8T\le8,\qquad
 8(2(p_0+v)-1)^2+2\sqrt8\,|2(q_0+Q)-1|\le8.
\]
Since $T\ge(1-4e)/\sqrt2>0.707$, we have
$v<0.4$ and $v+|z-w|\le p_0+\sqrt2 P<0.55$.
Thus $z,w>1/5$. If $\min(Q,B)\le1/25$, then
$T>0.9599$ and $|2(q_0+Q)-1|>0.9199$.
It follows that $v<p_0+0.142$; the left side of the second
inequality is greater than $8(0.715)^2+2\sqrt8(0.9199)>8$.
Consequently $1/25<Q,B<24/25$.

By Lemma~\ref{lem:mass-moments}, $H=L(gx)$ and $J=L(xg(x))$ are positive odd integers, where $L(v)=\operatorname{Tr}(qv)$.
The expansion
\[
  J=rza+zb+Qbc+Bca
\]
gives $1\le J\le(2+c)b\le27b$. In particular $v>1/1000$.
There is an additional quadratic unit
\[
    U=\frac{q_0q_2}{q_1q_3}\in\mathbb Q(\sqrt2).
\]
Indeed the conjugate ratios of the $d$-number $E$ are units,
and translation by $g^2$ fixes this product while $g$ inverts it.
Thus $U=(3+2\sqrt2)^n$ for $n\in\mathbb Z$. Since $U=rQ/B<24$,
one has $n\le1$. For fixed $U$ and $Q$, the identities
\[
 r=\frac{U(1-Q-e)}{Q+Ue},\quad
 B=1-Q-e(1+r),\quad
 H=rz+eb+\frac{Q(1-z-v)}{B}
       +a\left(B-\frac{Qer}{B}\right)
\]
determine $r,B$ and then $a$ once the odd integer $H$ is chosen.
The coefficient of $a$ is positive since $B>1/25$.

The \codecheck{field_checks/moments/adjacent16.cpp}{interval check}
starts from $(Q,z,v)\in[1/25,24/25]\times[1/5,4/5]\times[1/1000,2/5]$
and $0\le e\le4\cdot10^{-6}$. For each admissible odd $H$ recover
$r,B,a$ as above. Impose the two Fourier inequalities,
$0\le a\le\min(1,b,c)$, $w=1-z-v-aer\ge1/5$, and odd positive $J$.
Multiplying Gram dominance by $p_0$ and retaining two terms on the
right gives the last test, for $i=1,2,3$,
\[
 p_i^2+p_0\sum_{k=1}^3p_{k+i}^2/p_k\ge a(q_i^2+r q_{1+i}^2).
\]
Bisect the widest input interval, treating $-5\le n\le1$
separately and $0\le U\le(3+2\sqrt2)^{-6}$ together.
Outward operations on $10^{-9}\mathbb Z$ and integer square-root
comparisons discard every box. Since $x_j>250000$ implies
$q_j<4\cdot10^{-6}$, the trace assertion follows.
\end{proof}
\begin{proposition}\label{prop:k24-four-small-moments}\label{prop:k16-four-small-moments}
In either common-four case assume $\operatorname{Tr}(x)>10^6$
and let $j$ maximize $j(x)$. It is the
involution fixing $\sqrt2$. Put $W=j(q)/q$. Then
\[
 \bigl(L(jx),L(xj(x)),L((jx)^2),L(W^2)\bigr)
   =(3,1,11,5)
\]
at conductor $24$; at conductor $16$ the other possibility is
$(7,3,53,17)$.
\end{proposition}
\begin{proof}
In the cyclic case Lemma~\ref{lem:cyclic16-small-adjacent} identifies
$j$; in the biquadratic case every nonidentity element is an involution.
In the notation of Lemma~\ref{lem:four-opposite-recovery}, put
$e=1/250000$, $C=q_l+q_m=1-b(1+r)$, $A=x_l$, $B=x_m$.
Lemma~\ref{thm:half-million} gives $b=q_j<e$. The first two
moments $h,k$ and the remaining ones in that lemma are positive odd
integers by Lemma~\ref{lem:mass-moments}. Pairing indices gives
\[
 h=zr+ba+q_lB+q_mA,\qquad k=za(1+r)+CAB.
\]
There are two further inequalities common to both fields. First put
$S_r=r+r^{-1}$, $S_s=s+s^{-1}$. These are the two conjugates of
$W+W^{-1}$ over the rationals, and
$L(W^2)+1=(1-C)S_r+CS_s$. If $S_r\ne S_s$, the algebraic integer
$L(W^2)+1-S_r$ has nonzero norm, giving
\[
       (1-C)C(S_r-S_s)^2\ge1.
\]
Second, Gram dominance for the simultaneous columns, multiplied by
$p_0$ and retaining two terms on the signed side, gives for $i\ne0$
\[
 p_i^2+e(1/z+1/p_l+1/p_m)\ge a(q_i^2+r q_{i+j}^2).
\]

For clarity, the only field-dependent inputs are the Fourier bounds.
At conductor $24$ these are $dP^2+2\sqrt d\,|Q|\le8$,
$dQ^2\le16$ at each character of conductor $d$. The character
trivial on $j$ has $|Q|>1-4e$, so its conductor is $8$ or $12$;
also $1/5<z<4/5$. The conductor-$24$ inequality gives
$|q_l-q_m|\le\sqrt{2/3}+e$, hence $q_l,q_m>1/12$.
Using $a\le\min(A,B)$ in the expression for $k$ gives
$k\le48\min(p_l,p_m)$. The bounds
$p_lp_m<(4/5)^2/4$ and $q_lq_m>(1-2e-1/12)/12$ give $k<5$;
the formula for $h$ gives $h<12$.

At conductor $16$ the quadratic and faithful mixed inequalities give
\[
 8(2(z+p_0)-1)^2+2\sqrt8(1-2(q_0+b))\le8,\qquad
 16((z-p_0)^2+(p_l-p_m)^2)+8(q_l-q_m)\le8.
\]
Thus $.22<z<.8$ and $q_l,q_m>1/25$. The same moment formulas yield
$k<7$, $h<25$, and $k\le100\min(p_l,p_m)$.
Here $S_r\ne S_s$: equality would give $[\mathbb Q(W):\mathbb Q]\le2$,
so $jW=W$ and $W=W^{-1}=1$. Then $jq=q$, contradicting \eqref{eq:16-common-four-norm}.
Since $S_s<26$, the norm inequality already gives $r<.003$:
otherwise $(1-C)C(S_r-S_s)^2<2e\cdot335^2<1$.

For each permissible conductor assignment and each odd $h,k$ in
these ranges, first enclose $r$ using $0<ba<e$ in the equation for
$h$, then enclose $a$ using the equation for $k$. The interval variables are
$(z,q_l-q_m,p_l-p_m)$, initially in
$[.2,.8]\times[0,.82]\times[-1,1]$ at conductor $24$ and in
$[.2,.8]\times[0,1]\times[-1,1]$ at conductor $16$.
With $p_0,q_0,b\in[0,e]$, impose the displayed bounds and
$0<r\le1$, $0<a\le\min(1,A,B)$. Bisection of the longest edge,
rounding outwards to $10^{-9}\mathbb Z$, leaves the following
enclosures. The full tests are linked for conductors
\codecheck{field_checks/moments/moments24.py}{twenty-four} and
\codecheck{field_checks/moments/moments16.py}{sixteen}.
\[
\begin{array}{c|c|cc|ccc}
 c&K^{\langle j\rangle}&h&k&z&s&a\\\hline
 24&\mathbb Q(\sqrt2)&3&1&[.228,.332]&[.138,.185]&(0,.232]\\
 16&\mathbb Q(\sqrt2)&3&1&[.228,.323]&[.138,.183]&(0,.270]\\
 16&\mathbb Q(\sqrt2)&7&3&[.228,.233]&[.054,.058]&(0,.409]
\end{array}
\]
In the first row equality $S_r=S_s$ would make $W$ a unit of degree
at most two inverted by $j$. Its smaller conjugate $s<1$ would be a
positive power of $2-\sqrt3$ or $5-2\sqrt6$, neither possible in
the displayed interval. Thus the same norm inequality gives $r<.003$
here as well. Lemma~\ref{lem:four-opposite-recovery} now gives the
asserted moments.
\end{proof}

\begin{lemma}\label{lem:quartic-local-moments}
Let $K=\mathbb Q(\zeta_c)^+$ for $c=16$ or $24$, and let $j$ be the
nontrivial automorphism over $\mathbb Q(\sqrt2)$. Let $z,x$ be totally
positive integral $d$-numbers.
Suppose $N(x)$ is odd and $q=z/\operatorname{Tr}(z)$ satisfies
$\operatorname{Tr}(qx)=1$. At $c=16$ assume also $N(z)\in2\mathbb Q^{\times2}$.
Write $L(v)=\operatorname{Tr}(qv)$ and $W=jz/z$.
The vector $(L(j(x)),L(xj(x)),L(j(x)^2),L(W^2))$ cannot be
$(3,1,11,5)$ at either conductor, or $(7,3,53,17)$ at conductor sixteen.
\end{lemma}
\begin{proof}
The unit-orbit descriptions in Lemma~\ref{lem:quartic24-fast} and
Lemma~\ref{thm:half-million} reduce the test to finite rings.
Remove the rational integer factor of $z$, which cancels in $q$.

\emph{Conductor twenty-four.}
In the integral basis $1,\sqrt2,\sqrt3,(\sqrt2+\sqrt6)/2$, write
$[v]_0=\operatorname{Tr}(v)/4$. For $(H,K_0,M,V)=(3,1,11,5)$
the five moment identities say
\[
 [z(x,j(x),xj(x),j(x)^2,W^2)]_0=[z]_0(1,H,K_0,M,V)
\]
coordinatewise. Reduce these equations modulo $8$ and $11$ independently.
For each of the eight seeds $z_*$, multiply by the residue subgroup
generated by $U_1,U_2,U_3$ from Lemma~\ref{lem:quartic24-fast}.
One obtains $W$ even for a nonunit seed by starting with
$j(z_*)/z_*$ and multiplying by $j(U_i)/U_i$ at each step.
For $x$ use the four seeds of odd norm, the same units, and all odd
scalar multipliers. The numbers of residue classes of $z$ admitting
an $x$ are, in seed order,
\[
\begin{array}{c|rrrrrrrr}
8 &0&8&0&0&0&0&0&16\\
11&14&0&12&12&32&14&32&0 .
\end{array}
\]
The \codecheck{field_checks/moments/quartic24_two_primes.py}{two residue calculations}
have disjoint seed supports. A
\codecheck{field_checks/moments/quartic24.py}{four-modulus check}
also gives the exclusion for arbitrary integral $x$.

\emph{Conductor sixteen.}
Take $y=\sqrt{2+\sqrt2}$, $w=\sqrt{2-\sqrt2}$. Modulo $17$ the
four evaluations starting with $(\sqrt2,y,w)=(6,5,8)$ identify
$\mathcal O_K/17$ with $\mathbb F_{17}^4$. The involution $j$ shifts
by two. Normalize unit vectors to first coordinate one; the three
generators $3+2\sqrt2,(1+y)^2,(1+w)^2$ give a group $G$ of order
$128$ (all generator orders divide eight). The norm classes allow
only the seeds $2-y$, $4-2\sqrt2-y+w$ for $z$, and $1$ for $x$.
For $Z$ in either normalized seed translate, put $T=\sum Z_i$.
Retain $\sum_i Z_{i+2}^2/Z_i=VT$. All surviving $Z$ have $T\ne0$;
for $U\in G$ the only possible $X$ proportional to $U$ is then
$TU/\sum_iZ_iU_i$, omitting zero denominators.
The \codecheck{field_checks/moments/quartic16.py}{residue calculation}
gives $44$ possible $Z$ for $V=5$ and $12$ for $V=17$.
None admits an $X$ satisfying the remaining three scalar equations.
\end{proof}

\begin{proof}[Proof of Theorem~\ref{thm:quartic-four-complete}]
If there were no spherical structure, Theorems~\ref{thm:round2-sixteen},
\ref{thm:round2-quartic24-finite} and Lemma~\ref{thm:half-million} would
give the common four-rows and $\operatorname{Tr}(x)>10^6$.
Clear the denominators of $q$ to write $q=z/\operatorname{Tr}(z)$
with $z$ a totally positive integral $d$-number. Lemma~\ref{lem:mass-moments} gives $N(x)$ odd. At conductor sixteen \eqref{eq:16-common-four-norm} gives
$N(z)\in2\mathbb Q^{\times2}$.
Proposition~\ref{prop:k24-four-small-moments} now forces one of the
moment vectors excluded by Lemma~\ref{lem:quartic-local-moments}.
\end{proof}

\section{The quartic field of conductor forty}
\begin{theorem}\label{thm:forty-closed}
Every complex rank-six fusion category with $K=\mathbb Q(\sqrt2,\sqrt5)$ admits a spherical structure.
\end{theorem}
Write $F=\mathbb Q(\sqrt5)$. The Fourier bounds leave a signed
complement $F$ and two rational ordinary characters. Induction then
identifies certain square equations modulo four in the two orders.
Their solution counts rule out the remaining codegrees.
\begin{lemma}\label{lem:forty-unit-param}
The totally positive $d$-numbers $z$ in the ring of integers of
$K=\mathbb Q(\sqrt2,\sqrt5)$ are precisely the totally positive elements
\[
 h\sqrt d\, (1+\sqrt2)^i\phi^j(3+\sqrt{10})^l,\qquad
 h\in\mathbb Z_{>0},\quad d\in\{1,2,5,10\},\quad i,j,l\in\mathbb Z ,
 \qquad \phi=(1+\sqrt5)/2 .
\]
Here $N_K(z)=h^4d^2$. If $\operatorname{Tr}_K(z)\le1001$, the three
exponents have absolute values at most $7,12,3$, respectively.
\end{lemma}
\begin{proof}
All ramification indices are at most two, so an invariant ideal has
square norm, say $m^2$. The quadratic Pell equations give totally
positive unit generators $(1+\sqrt2)^2,\phi^2,(3+\sqrt{10})^2$.
Equation~\eqref{eq:biquadratic-relative-norms} now writes $z$ as
$\sqrt m$ times the asserted product of units. Since $\sqrt m\in K$,
we have $m=dh^2$ with $d\in\{1,2,5,10\}$. Conversely, every positive
element of the displayed form is integral and has invariant ideal.
Each real conjugate of a relative norm is at most
$(\operatorname{Tr}_K(z)/2)^2$ by the arithmetic--geometric mean
inequality. Applying this to each generator and its conjugate gives
the stated exponent bounds.
\end{proof}

\begin{proposition}\label{prop:forty-trace-form}
Let $K=\mathbb Q(\sqrt2,\sqrt5)$ be the Frobenius--Perron dimension
field in rank six. Suppose one of the two character algebras has
complementary codegrees $f,f'$ generating $F=\mathbb Q(\sqrt5)$.
Write $H$ for its dimension-orbit codegree, $D$ in the ordinary sector
and $E$ in the signed sector. Then
\[
  ff'=5k^2,\qquad k=a^2+b^2\quad(a,b\in\mathbb Z),\qquad
  \operatorname{Tr}_{K/\mathbb Q}(z^2/H)\simeq_{\mathbb Q}
  \langle1,1,1,1\rangle .
\]
Choose $f>f'$ and put $\phi=(1+\sqrt5)/2$.
Then $f=k\sqrt5\,\phi^m$ with $m$ positive and odd, and
$\operatorname{Tr}_F(1/f)=F_m/k$, where $F_m$ is the $m$th
Fibonacci number. In particular $ff'=45$ and $180$ are impossible.
At the prime over five, normalize $v(\sqrt5)=1$. If $v(H)$ is even,
$5^{-v(H)/2}H$ has nonsquare residue in $\mathbb F_{25}$.
\end{proposition}
\begin{proof}
Both character factors are totally real, so all simples are self-dual.
Their orthonormal basis for the unit-coefficient pairing gives
$A\perp B\simeq6\langle1\rangle$, with
\[
 A(z)=\operatorname{Tr}_K(z^2/H),\qquad B(y)=\operatorname{Tr}_F(y^2/f).
\]
The determinant of $A$ is a square, since both $\operatorname{disc}(K)$
and $N(H)$ are squares (Lemma~\ref{lem:forty-unit-param}). Thus
$\det B=5/N(f)$ modulo squares gives $N(f)=5k^2$.

Put $w=(5+\sqrt5)/2=\sqrt5\,\phi$. The invariant ideals $(f)$
and $(kw)$ coincide by equality of their norms. Their quotient is
therefore a totally positive unit, an even power of $\phi$.
This gives the asserted $f$ and its reciprocal trace by Binet's
formula. In the basis $1,\phi$ the form
$\operatorname{Tr}_F(y^2/w)$ has identity Gram matrix,
so $B\simeq k\langle1,1\rangle$.

At each odd $p\ne5$ the algebra $K\otimes\mathbb Q_p$ is unramified.
Remove the common rational $p$-power of $H$. Such scaling preserves
the Hasse invariant of a rank-four form of square determinant:
the factor is $(d,d)_p^6(d,\det A)_p^3=1$. The resulting form is
unimodular, hence has Hasse invariant one. Since $A\perp B$ is
standard, $(k,k)_p=(k,-1)_p=1$. Thus each $p\equiv3\pmod4$ divides
$k$ to even order, so $k=a^2+b^2$. Multiplication by $a+bi$
identifies $B$ with $\langle1,1\rangle$, and Witt cancellation gives
the assertion about $A$ \cite[Chapters~III--IV]{SerreArithmetic}.

Finally put $L=\mathbb Q_5(\sqrt2)$, the unramified quadratic extension
of $\mathbb Q_5$, so $K\otimes\mathbb Q_5=L(\sqrt5)$. Rational scaling
does not change the isometry class of $A$ at five, by its square
determinant and the scaling calculation above. If $v(H)$ is even, take
$H$ to be a unit. Choose a unit $h\in L$ with the same residue.
Hensel's lemma makes $H/h$ a square. Thus $A$ has the form $B_0\perp5B_0$,
where $B_0(y)=\operatorname{Tr}_{L/\mathbb Q_5}(2y^2/h)$ is unimodular
of rank two. Its determinant modulo squares is $2/N_{L/\mathbb Q_5}(h)$.
The Hasse invariant at five of $B_0\perp5B_0$ is the quadratic
character of this determinant: in a diagonal unit basis use
$(5,5)_5=(5,-1)_5=1$, while a Hilbert symbol of two units is one.
It must equal one. Since two is a nonsquare modulo five, the residue norm
of $h$ is a nonsquare in $\mathbb F_5^\times$, equivalently the residue
of $h$ is a nonsquare in $\mathbb F_{25}^\times$.
\end{proof}

\smallskip
\begin{theorem}\label{thm:quartic-silvergold-reduction}
Let $K=\mathbb Q(\sqrt2,\sqrt5)$ be the dimension field of a
rank-six category. Unless it has a pseudo-unitary Galois conjugate,
the signed complementary codegrees are distinct irrational
conjugates in $\mathbb Q(\sqrt5)$ and
\[
 6837/10000< \operatorname{Tr}_K(1/E)<143/200,\qquad
 \operatorname{Tr}_K(E/D)>94
\]
after transport maximizing $E$.
\end{theorem}
\begin{proof}
Theorem~\ref{thm:all-proper-quartic} permits us to assume $x=E/D$
primitive. Write $p=1/D$, $q=1/E$, $s=\operatorname{Tr}(p)$,
$t=\operatorname{Tr}(q)$ and let $j$ fix $\sqrt2$. Only the pair $q_0,q_j$ can be simultaneously small:
the other two pair bounds are $q_i\ge(t-1/\sqrt5)/4$ in
unnormalized coordinates.

Suppose first the signed complementary codegrees are rational.
Lemma~\ref{lem:forty-unit-param} gives square norm. Thus the
order-discriminant identity makes $gh$ a square for two rational axes
of codegrees $g,h$. A quadratic factor of common rational codegree
must instead be $\mathbb Q(i)$ with codegree $4n^2$. For $t<18/25$
the smaller denominator is at most seven. The local congruences at
three and five from Theorem~\ref{thm:quartic-fifteen-rational}
eliminate $(5,5),(6,6)$; Lemma~\ref{lem:quartic-three-induction}
excludes three in either sector.
The square condition leaves
\[
 (4,4),\quad(4,9),\quad(4,16),\quad(4,25),\quad(7,7).
\]

At half mass one has $s=t=1/2$. At $(4,h)$ with $h=9,16,25$,
the signed evaluations at $V_0$ must have the form $(1,b)$,
$1\le b\le h/4$. Indeed $s\ge1/2$ by Lemmas~\ref{lem:quartic-rational-mass-strata}
and~\ref{lem:quartic-short-pairs}, so the first entry is one.
Thus $s=3/4-b/h$ and $t=3/4-1/h$.

The ordinary complementary masses $u=17/36,7/16,41/100$ are
impossible. None is a sum of two reciprocals of integers at least
three. For a quadratic pair of reduced mass $u=a/b$,
Lemma~\ref{lem:fixed-quadratic-mass} gives squared gap at least
$(a^2-4)/b^2>1/8$ in these three cases, while
\eqref{eq:orbit-projection} bounds it by $1/10,1/8,1/40$ at
conductors $5,8,40$.
At signed half mass the first shift lies in
$1/2+\mathbb Z_{\ge0}$ and the mixed moment is at least $1/2$.
First dispose of an ordinary short pair in $\mathbb Q(\sqrt2)$.
Its squared reciprocal gap is $1/8$, so at conductor eight
$8O_\eta\ge2$. The mixed bound in \eqref{eq:orbit-projection}
forces $T_\eta=0$, hence $Q_8=0$ and
$q_0+q_j=t/2=1/4$. In this case $q_j\ge1/8$ already.
Otherwise Lemma~\ref{lem:quartic-short-pairs} leaves rational ordinary
codegrees or $5\pm\sqrt5$. At every shift the former give zero
Gram correction. The golden pair contributes at most $(5-3)/2=1$
to the ordinary squared length, the signed contribution is at least
$1/2$, and hence the correction is at most $1/2$.
The \codecheck{quartic_more/verify.py}{two-variable bound} of
Theorem~\ref{thm:quartic-fifteen-rational}, with this correction,
gives $q_j>1/76$ in $19473$ boxes, hence $A\le38$.
Of the primitive quartic polynomials only two satisfy parity at
this cap: their $(A,B,C,H)$ are $(16,62,64,1)$ and
$(32,254,368,1)$. The positive odd-moment vectors in their outer
boxes number $15$ and $197$; none has least weight at the least root
(\codecheck{quartic_more/silvergold.py}{moment calculation}).

If $t=23/36$ or $11/16$, omit the Gram correction entirely and
use only the two linear column bounds on $a=p_0/q_0,b'=p_0/q_j$.
Use $h_j\in s+\mathbb Z_{\ge0}$ as before and
$K_j\ge1+\lfloor1/4+b^2/h\rfloor-(1/4+b^2/h)$, since the
complementary contribution is $1/4+b^2/h$. At each allowed $s$
the same box argument gives $q_j>1/300$, hence $A\le207$.
At $t=71/100$ we will use $A\le550$. If $q_j\ge1/400$ this is
immediate. Otherwise $|Q_8|>7/10$ gives $|P_8|<51/1000$, so
$p_j<(s+51/1000)/2\le381/1000$ while
$q_j\ge(t-1/\sqrt2)/4$. The two other nonidentity coordinates
are at least $(t-1/\sqrt5)/4$. Substitution gives $A<550$.
For $t\ge5/7$ the separate estimates give
$A<1+4t/(t-1/\sqrt2)<399$.

At trace cap $550$, impose the minima, separate and mixed Fourier
bounds, $2/5\le s\le t$, $t\ge23/36$, $s+t\ge1$, and the three
first-shift inequalities. Of $594$ polynomials only
$(64,62,16,1),(166,167,34,1)$ survive. For each impose $t\ge18/25$,
the $(7,7)$ masses $t=5/7$, $s=4/7,5/7$, or the displayed
$(4,h)$ masses and mixed moments. The
\codecheck{quartic_more/silvergold.py}{separation test} eliminates both.
Suppose instead the signed codegrees form a quadratic pair of
discriminant $c$. Again $s\ge1/2$: an ordinary codegree two would
force a rational signed codegree two, and three is excluded by
Lemma~\ref{lem:quartic-three-induction}. Apply
Lemmas~\ref{lem:quartic-rational-mass-strata} and~\ref{lem:quartic-short-pairs}.
For $c=5$ the short-pair lemma leaves $5\pm\sqrt5$ or $t\ge2/3$.
For the longer pairs add the first-shift inequality fixing the
complementary field; at $c=8$ also impose $t-|Q_8|\ge2/94$.
The \codecheck{quartic_more/silvergold_duals.json}{separating weights}
exclude all ratios in the following ranges:
\[
\begin{array}{c|c|r|r}
 c&t\text{ at least}& A\text{ through}&\text{polynomials}\\\hline
 8&1/2&94&75\\ 40&3/4&94&75\\ 5&2/3&94&75\\ 5&143/200&510&553
\end{array}
\]
At $c=8$ the cap $94$ is Proposition~\ref{prop:quartic-primitive-norm-height};
at $c=40$ the initial row bound of Theorem~\ref{thm:quartic-all-fields}
gives $t\ge3/4$ and the same cap. In the last row
$A<1+4/(143/200-1/\sqrt2)<511$.

For $f,f'=5\pm\sqrt5$ one has $s=t=1/2$ and the short evaluation
is $\beta=1$ or $\phi^2$. In either case $K_j\ge1/2$ for the shift
exchanging $\sqrt5$ (Theorem~\ref{thm:quartic-cyclic-twenty-all}).
For $\beta=1$ the first shift is half-integral and the preceding
trace-$38$ calculation applies. For $\beta=\phi^2$ it is integral.
Again $q_j\ge1/8$ for an ordinary pair in $\mathbb Q(\sqrt2)$.
In all remaining cases the Gram correction at every shift is at most
$1/2$, and
$|Q_5|\le\sqrt5/10$ by Lemma~\ref{lem:quadratic-compression}.
The \codecheck{quartic_more/verify.py}{two-variable bound} now gives
$q_j>1/48$ in $46225$ boxes. Since $|Q_{40}|\le1/\sqrt{10}$,
the two remaining nonidentity weights also exceed $1/48$.
Thus $A\le24$. The shifted complementary polynomial $(X^2-3X+1)^2$
makes $A,C$ even and $B,H$ odd; $2\operatorname{Tr}(\beta^r/f)=2,5$
makes the normalized moments of orders two and three even and odd,
respectively.

Only $(A,B,C,H)=(14,47,26,1),(18,43,18,1)$ remain. The same
\codecheck{quartic_more/silvergold.py}{moment calculation} gives
$15$ and $23$ positive weight vectors in their outer boxes. Only
one has least weight at the least root, at moments $(6,81)$ in the
second polynomial. Its largest weight exceeds $0.85$, whereas Fourier
inversion gives $\max_i 2q_i\le(2t+\sqrt5/10+1/\sqrt{10})/2<0.78$.

It remains to exclude $2/3\le t\le6837/10000$. By
Proposition~\ref{prop:forty-trace-form}, $u=1-t=F_m/k$, with $m$
positive and odd and $k$ a sum of two squares. The signed conductor-five
bound $u^2\le1/10+4/(5k^2)$ gives
\[
  10F_m^2-k^2\le8,\qquad 3163k/10000\le F_m\le k/3 .
\]
Together these imply $k\le132$ and $F_m\le44$. The odd-indexed
Fibonacci numbers in this range are $1,2,5,13,34$. For the last
three the first inequality requires $k\ge16,42,108$, respectively,
while the second gives $k\le15,41,107$. The first two force $k=3,6$,
neither a sum of two squares.
\end{proof}

\begin{theorem}\label{prop:forty-divergent-complements}
In the remaining case of Theorem~\ref{thm:quartic-silvergold-reduction},
$\operatorname{Tr}_K(E/D)>1001$ and the ordinary complement consists
of two rational characters with codegrees $(4,4n^2)$, $n\ge2$, or $(8,8)$.
\end{theorem}
\begin{proof}
Put $A=\operatorname{Tr}_K(E/D)$ and let $j$ fix $\sqrt2$.
First suppose $A>1001$. Since $p_0/q_0\le1$ and $s<1$, the
pair bounds $q_i\ge(t-1/\sqrt5)/4$ for $i\ne0,j$ give
$q_0\le q_j<1/1000$, and $p_0<1/1000$. Then
$|Q_8|=t-2(q_0+q_j)>679/1000$. Consequently the mixed projection
and the separate conductor-forty estimate yield
\[
 |P_8|\le\sqrt{(2-\sqrt8\,|Q_8|)/8}<1/10,\qquad
 |P_{40}|\le1/\sqrt{20}<224/1000.
\]
Since $P_8=2(p_0+p_j)-s$ and $P_5+P_{40}=2(p_0-p_j)$,
we obtain $|P_5|>s-328/1000$. The separate bound
$|P_5|\le1/\sqrt5$ gives $s<776/1000$.

Suppose there is a distinct irrational ordinary pair $g,g'$, and put
$w=1/g-1/g'$ and $N_o=gg'$. Its quadratic field has conductor
$c=5,8$ or $40$. Write $U>1$ for the least totally positive unit.
Since $g/g'$ is a nontrivial power of $U$ up to conjugation,
$w^2\ge[(U-1)/(U+1)]^2(1-s)^2$. For $c=8$ the multiplier is $1/2$,
whereas the twist-fixed projection gives
$w^2\le1/8-Q_8^2/4<1/8-(679/1000)^2/4$, strictly below
$(1/2)(224/1000)^2$. At $c=40$ the multiplying factor is $9/10$
from $U=19+6\sqrt{10}$. Then $w^2\le1/40$ again contradicts the unit-ratio bound.

At $c=5$, \eqref{eq:orbit-projection} gives $P_5^2+2w^2\le1/5$. Since $s\ge1/2$ and
$w^2=(1-s)^2-4/N_o$, this implies
\[
 8/N_o\ \ge\ \tfrac23(1-328/1000)^2-1/5>1/10.
\]
Here we minimized $(s-328/1000)^2+2(1-s)^2$ over the real line.
Thus $N_o<80$. In Proposition~\ref{prop:forty-trace-form} the
only possible $k$ at $5k^2<80$ are one and two, as three is excluded.
The row trace $L_o=g+g'$ is a multiple of five, at least ten and
at most $N_o/2$. Thus only $L_o=10,N_o=20$ remains, with
$g,g'=5\pm\sqrt5$.

Let $e$ be the idempotent of this ordinary quadratic factor.
Its different is $(\sqrt5)$ and $(g)=(2\sqrt5)$, so the conductor
formula of Lemma~\ref{lem:trace-order-conductor} gives $2e\in R$.
The involution $1-2e$ has squared length one and unit coefficient
$1-2(1/g+1/g')=0$, hence is a signed nontrivial basis element.
Its Frobenius--Perron value is one, so it is an actual invertible.
Its positive lift is one on the four dimension characters and
has a constant sign on the conjugate pair. Its zero unit coefficient
would give $t=1/2$, outside the permitted interval.

An ordinary quadratic pair of common rational codegree is likewise
impossible. All simples are self-dual since the signed algebra is
totally real. The ordinary field discriminant would be a positive
quadratic field discriminant times a square, whereas $N(D)g^2$ is
a square. Thus both ordinary characters are rational and their
codegrees $g\le h$ are even: in the $g$-row the four signed
dimension induction evaluations equal one and the two remaining
signed evaluations are equal integers, giving $g=4+2a$ by the
unit-coefficient trace and symmetry. The product $gh$ is a square
by the discriminant calculation. Since $1/g+1/h>224/1000$ and $g\ge4$, one gets $(4,4n^2),(6,6)$ or $(8,8)$.
For $(6,6)$ the two rows must pair at three by
Lemma~\ref{lem:silver-local-partner}: all quartic residue branches
contain the unramified quadratic field from $\sqrt2$. Their squared
row difference twelve would then be divisible by nine. Lemma~\ref{lem:quartic-four-rows} excludes $(4,4)$, since it forces $t=1/2$.

It remains to exclude $94<A\le1001$. Write $N=ff'=5k^2$ and
$u=1-t=F_m/k$ by Proposition~\ref{prop:forty-trace-form}.
If $k<16$, then $285/1000<F_m/k<3163/10000$ gives $F_m\le4$.
Since $m$ is odd, only $F_m=1,2$ remain, and the interval forces
$F_m=2$, $k=7$, impossible for a sum of two squares. Consequently
$N\ge1280$ and the amplified signed projection gives
\[
  |Q_5|^2\le1/5-2(u^2-4/N)<(21/100)^2.
\]
The unit bounds in Lemma~\ref{lem:forty-unit-param} give $1065$
primitive ratios through trace $1001$, including $75$ through trace $94$. For the others use the
separate and mixed projections, $s\ge1/2$, $t>6837/10000$, the last
conductor-five bound, and $\operatorname{Tr}(q\sigma x)\ge s$ for
$\sigma$ fixing $F$. Only the following two tuples lack a strict
rational separating certificate, with $4x=A+u_8\sqrt8+u_5\sqrt5+u_{40}\sqrt{40}$
up to conjugacy:
\[
 (A,u_8,u_5,u_{40})=(166,58,74,26),\qquad(564,199,252,89).
\]
Apply \eqref{eq:quartic-difference-ideal} with $\sigma$ fixing $F$.
Starting with the four conjugates of $\sigma x-x$, adjoin their
products by $gx$ until the Hermite basis stabilizes. In the basis
$1,\sqrt2,\phi,\sqrt2\phi$ this gives the two diagonal basis
matrices $\operatorname{diag}(16,4,16,4)$ and
$\operatorname{diag}(4,2,4,2)$, respectively.
The second tuple forces $4t\in\mathbb Z$, incompatible with the permitted
interval for $t$. The first gives $16t\in\mathbb Z$ and therefore
$t=11/16$, together with $4\sqrt2\,Q_8\in\mathbb Z$. At this fixed
mass the same linear system supplies two strict certificates giving
$-7/10<Q_8<-54/100$. This interval has no integral multiple of
$1/(4\sqrt2)$. The file \codefile{field_checks/forty/verify.py} regenerates the
census, checks the two Hermite bases and verifies all separating weights.
\end{proof}

\begin{corollary}\label{cor:forty-dyadic-scalar}\label{prop:forty-golden-period}
Both dyadic orders are local with residue field $\mathbb F_2$, and
$4\mid k$ in $ff'=5k^2$. Write
\[
 H=h_H\sqrt{d_H}\,\xi^{i_H}\phi^{j_H}\theta^{l_H}\quad(H=D,E),
 \qquad f=k\sqrt5\,\phi^m,
 \qquad \xi=1+\sqrt2,\quad\theta=3+\sqrt{10}.
\]
Then $3\mid j_D,j_E,m$ and $1-t=F_m/k$ with $m\equiv3\pmod6$,
on choosing $f>f'$.
\end{corollary}
\begin{proof}
Neither ordinary rational axis splits off at two, since its idempotent
has nonintegral trace. Their sum cannot split off either, for the same
reason, and the quartic completion is a field. Thus the ordinary
order is local with residue field $\mathbb F_2$; its reduction
modulo two gives the same conclusion for the signed order.
The signed $F$-row $\psi$ agrees with an ordinary short row $r$
modulo $2\mathcal O_F$, as $F$ is unramified. Hence
$f=\sum\psi_i^2\equiv\sum r_i^2\equiv0\pmod{4\mathcal O_F}$.

Let $j$ negate $\sqrt5$ and fix $\sqrt2$. The simultaneous column
for $j$ is a unit (its coordinates are ratios of conjugate
$d$-numbers), hence reduces to one. In the quartic coordinates
$H/jH$ the radical contributes a sign, $j$ fixes $\xi$, and
$\theta$ reduces to one. In the complementary coordinate
$f/jf=\phi^{2m}$. Since $j\phi=-\phi^{-1}$ and $\bar\phi$ has
order three in $\mathbb F_4^\times$, we get
$3\mid j_D,j_E,m$. The remaining assertions are
Proposition~\ref{prop:forty-trace-form}.
\end{proof}

\subsection{A congruence from the second indicator}

\begin{lemma}\label{lem:induction-fs-four}
In the remaining conductor-forty case choose the lift of positive second Frobenius--Schur indicator for each simple. Let $\Phi:R\to T$ be the isometry of underlying lattices obtained by matching these basis elements. For every induction constituent $C$ one has
\[
      \Phi([FC]_R)=[FC]_T\pmod {4T}.
\]
For every $a,b\in R$ one also has
\[
      \Phi(a^2 b)=\Phi(a)^2\Phi(b)\pmod {4T}.
\]
\end{lemma}
\begin{proof}
The two character algebras are commutative and totally real. Thus
all simple lifts are self-dual: every character takes the same value
on a lift and on its dual. The pivotal maps on the two lifts differ by a sign, so the pivotal
transpositions and the indicators have opposite signs. In a symmetric self-dual object the multiplicity
space of an indicator-$-1$ simple has a nondegenerate alternating
form and has even dimension. Its ordinary and signed multiplicities
therefore agree modulo four in the stated basis.

The pivotal fusion category $\mathcal S$ is unimodular, so
$I(\mathbf1)$ is symmetric Frobenius in $Z(\mathcal S)$
\cite[Lemma~7.1]{SchweigertWoike}. Each constituent $A_\chi$ is
self-dual: the rows $S_{A_\chi^*,Y}=S_{A_\chi,Y^*}$ and $S_{A_\chi,Y}$
agree by the character formula and $[FY^*]=[FY]$. Its multiplicity
is one, so the symmetric pairing restricts nondegenerately to it.
The pivotal forgetful functor preserves this self-duality.

For the product assertion, let $j$ be the pivotal map and
$f:Y\to Y^*$ a symmetric self-duality. Then
$j_X\otimes f\otimes1_{X^*}$ is a symmetric self-duality on
$XYX^*$: pivotal transposition fixes it by $f^*j_Y=f$ and
$j_X^*j_{X^*}=1$. Commutativity and self-duality identify its class
with $[X]^2[Y]$. This proves the assertion on basis elements.
Extend linearly in $b$ and expand the square in $a$: the mixed
terms have coefficient two, and products agree modulo two.
\end{proof}
Let $C=FA$ for an ordinary rational character label $A$ of codegree $h$.
\begin{lemma}\label{lem:forty-rational-row}
The signed coordinates of $C$ are $(E/h,(h/2-2)f/h)$. In particular
every $(z,w)\in T$ satisfies
\begin{equation}\label{eq:forty-trace-transfer}
       \operatorname{Tr}_K z+(h/2-2)\operatorname{Tr}_F w\in h\mathbb Z.
\end{equation}
On the ordinary dimension orbit and the two complementary labels,
with $A$ listed first among the latter, the row $2hS_{A,-}$ takes
integral values $\gamma,m,b$, respectively. They satisfy
\begin{equation}\label{eq:forty-row-bound}
 \gamma>0,\qquad 4\gamma+m+b=h,\qquad
 4\gamma^2+m^2+b^2\le h(h+m)/2 .
\end{equation}
\end{lemma}
\begin{proof}
On the signed dimension orbit all four entries of $2hS_{A,-}$ are one.
The two remaining entries agree by Galois invariance. Since the six
entries sum to $h$, their common value is $h/2-2$. Symmetry gives the
signed coordinates of $C$, and the integral unit coefficient of
$C(z,w)$ gives \eqref{eq:forty-trace-transfer}.

On the ordinary side, rationality and Galois invariance again give
integral entries constant on each orbit. Symmetry identifies
$\gamma$ with $h\operatorname{FPdim}(C)/D$, and the unit coefficient
gives the sum in \eqref{eq:forty-row-bound}. For the inequality take the $A,A$ entry of $S\Theta S$.
It equals $S_{A,A}$ since $A$ has twist one and the center has
trivial anomaly. The row of this self-dual label is real.
Its squared norm is equally divided between the two sign sectors
by orthogonality with the $JA$ row. The signed twist terms cancel
in $J$-pairs. Of the ordinary mass $1/2$, at least
$M=(4\gamma^2+m^2+b^2)/(2h^2)$ has twist one, counting the
$J$-translates. Hence $2M-1/2\le S_{A,A}=m/(2h)$.
\end{proof}

\begin{proposition}\label{prop:forty-four-column}\label{prop:forty-eight-columns}
The forgetful image of an induction label of smallest ordinary
complementary codegree has the following coordinates:
\[
\begin{array}{c|c|c|l}
\text{ordinary codegrees}& K\times\mathbb Q^2 & K\times F&\text{parameters}\\\hline
(4,4n^2)&(D/4,m,-n^2m)&(E/4,0)&m=0,2\\
(8,8)&(\gamma D/8,m,m)&(E/8,f/4)&(\gamma,m)=(1,2),(2,0)
\end{array}
\]
In the first row all signed quartic traces, and all ordinary quartic
traces when $m=0$, are divisible by four. In the second row both
induction labels have isomorphic forgetful images, their ordinary
$16S$-submatrix has every entry $m$, and $8\mid k$.
\end{proposition}
\begin{proof}
The signed columns are given by Lemma~\ref{lem:forty-rational-row}.
For $h=4$, \eqref{eq:forty-row-bound} gives $\gamma=1$, $b=-m$,
$-1\le m\le2$. The signed complementary value is zero and the
residue character is common, so $m$ is even. Symmetry gives the
other evaluations; pairing the column with an arbitrary integral
element gives the trace assertions.

For $h=8$ let $m_i,c$ be the diagonal and off-diagonal entries of the
symmetric $16S$-submatrix. The row bound becomes
\[
 4\gamma_i+m_i+c=8,\qquad (m_i-2)^2+c^2\le36-4\gamma_i^2 .
\]
Thus $\gamma_i=1,2$: the value three would require $m_i=2,c=0$,
contrary to the first equation. The ordinary rows coincide modulo
two on the FP induction column, so the $\gamma_i$ agree, as do the
$m_i=m$.
The two rational character rows $r_1,r_2$ each have four odd
entries $\pm1$, one $\pm2$ and a zero. If they agreed modulo four
their nonzero difference would be supported at the single position
with entries $\pm2$, contradicting orthogonality to the Perron row.
Thus one coordinate of $(r_1-r_2)/2$ is odd. The signed induction
columns agree, so by Lemma~\ref{lem:induction-fs-four} their ordinary
difference belongs to $4R$. Character inversion gives its coefficient
vector $(m-c)(r_1-r_2)/8$, forcing $16\mid m-c$. The two displayed
bounds now give exactly the second row of the table. The images
agree in both quotients and hence in the sphericalization.
Finally $f/4\equiv m\equiv0\pmod{2\mathcal O_F}$ by the common
residue character, giving $8\mid k$.
\end{proof}

\subsection{The valuations of the two global dimensions}
Write $v_2$ for the rational $2$-adic valuation. On $K$ use the integral
valuation at the prime over two, normalized by $v(\sqrt2)=1$.
Thus $v$ restricts to $2v_2$ on $\mathbb Q$. The induction columns just
computed supply both the scale and the parity information in the
following comparison.

\begin{proposition}\label{prop:forty-exponent-parity}
In the setting of Corollary~\ref{prop:forty-golden-period}, write
$H=h_H\sqrt{d_H}\,\xi^{i_H}\phi^{3r_H}\theta^{l_H}$ for $H=D,E$.
Then $i_D\equiv i_E$ and $l_D\equiv l_E$ modulo two.
In the case $\gamma=2$ of Proposition~\ref{prop:forty-eight-columns}
all four exponents are even.
In particular $E/D$ is an integer times a unit times either $1$ or
$\sqrt5$ in the parametrization of Lemma~\ref{lem:forty-unit-param}.
\end{proposition}
\begin{proof}
For the simultaneous column $W_\sigma$, pair $W_\sigma-1$ with the
column of the smallest ordinary codegree $h=4,8$. The ordinary complementary evaluations of $W_\sigma$
are one. In the signed complement they are $f/\sigma f$. Thus
Lemma~\ref{lem:induction-fs-four} and
Proposition~\ref{prop:forty-four-column} give a single congruence
\[
 \gamma\bigl(\operatorname{Tr}_K(D/\sigma D)-4\bigr)
 \equiv\operatorname{Tr}_K(E/\sigma E)-4
 +(h/2-2)\bigl(\operatorname{Tr}_F(f/\sigma f)-2\bigr)
 \pmod{4h}.
\]
The last term vanishes modulo $4h$. Only $h=8$ and a nontrivial
action on $F$ require checking. There the trace difference is
$(\phi^m-\phi^{-m})^2$, divisible by $16$ because the rational
part of $(2+\sqrt5)^{m/3}$ is even.

Negating $\sqrt5$ alone gives
$\operatorname{Tr}_K(H/\sigma H)/4\equiv3^{l_H}\pmod8$:
it is the product of the rational parts of $\phi^{6r_H}$ and
$\theta^{2l_H}$, with $\phi^6=9+4\sqrt5$, $\theta^2=19+6\sqrt{10}$.
Negating both radicals gives $3^{i_H}$ instead, since
$\xi^2=3+2\sqrt2$. Thus for $a=i,l$ the two congruences are
$\gamma(3^{a_D}-1)\equiv3^{a_E}-1\pmod h$. These give matching
parities at $\gamma=1$ and even exponents at $\gamma=2$, $h=8$.
The sign conditions on $E/D$ now exclude a $\sqrt2$ component in
its radical factor.
\end{proof}

\begin{proposition}\label{prop:forty-marked-scales}
In the residual conductor-forty case write
$H=h_H\sqrt{d_H}\xi^i\phi^{3J}\theta^l$ for $H=D,E$ as in
Corollary~\ref{prop:forty-golden-period}.
Set $s=\operatorname{Tr}(1/D)$, $b=v_2(k)$, and let
$\delta_H=1$ if $d_H=5$ and $l_H$ is odd, and zero otherwise.
Then $d_H\in\{2,5,10\}$ and
\[
 v_2(h_D)=2+\delta_D-v_2(s),\qquad v_2(h_E)=1+\delta_E+b .
\]
Whenever $\gamma=1$ the two dimensions have the same dyadic
valuation, $\delta_D=\delta_E$, and
\[
                         b=1-v_2(s).
\]
For $(4,4n^2)$ with $n$ even, the ordinary complementary
evaluations of the codegree-four induction column both vanish.
\end{proposition}
\begin{proof}
For $H=E$ its quartic trace form is a sum of four rational squares by
Proposition~\ref{prop:forty-trace-form}. For $H=D$ the ordinary
complement is a sum of two squares in either case: in the eight case
each weight is $1/8$, and $\langle2,2\rangle\simeq\langle1,1\rangle$.
Witt cancellation gives the same quartic form. At five the argument
in Proposition~\ref{prop:forty-trace-form} therefore applies to both.
If $d_H=1$ the unit remaining at five after removing the scalar
has square residue in $\mathbb F_{25}$: $\xi$ has norm $-1$,
a square modulo five, and $\phi,\theta$ have rational residue.
This excludes $d_H=1$.

Write $\xi^{-i}=a+c\sqrt2$, $(2+\sqrt5)^{-J}=u+w\sqrt5$,
$\theta^{-l}=v+y\sqrt{10}$. Here $a,v$ are odd and
$(c,y,u,w)\equiv(i,l,1-J,J)\pmod2$; also $4\mid w$ for even $J$
since $(2+\sqrt5)^2=9+4\sqrt5$. Positivity gives
$i+l\equiv[2\mid d_H]$ and $J+l\equiv[5\mid d_H]\pmod2$.
Extract the rational coefficient of $h_H/H$:
\[
\begin{array}{c|ccc}
 d_H&2&5&10\\\hline
 h_H\operatorname{Tr}(1/H)/4&cuv+5awy&awv+2cuy&auy+cwv .
\end{array}
\]
In each column its valuation is $\delta_H$. Reversing the exponents
also gives $v_2(\operatorname{Tr}H)=2+v_2(h_H)+v_2(d_H)+\delta_H$.
For $H=D$ the first scale formula follows immediately. For $H=E$
use $t=1-F_m/k$ with $4\mid k$, $m\equiv3\pmod6$, and
$F_m\equiv2\pmod4$ by the Fibonacci recurrence.

The signed FP induction column has quartic value $E/D$ and residue
$\gamma$. At $\gamma=1$ this is a unit, so $d_D,d_E$ have the same
parity and $v_2(h_D)=v_2(h_E)$. The only odd seed is $5$ and
$l_D\equiv l_E\pmod2$ by Proposition~\ref{prop:forty-exponent-parity}.
Thus $\delta_D=\delta_E$ and $b=1-v_2(s)$.

For the codegrees $(4,4n^2)$ write $m_C\in\{0,2\}$ for the
ordinary evaluation of the four-column. Comparing the two squared
lengths modulo eight by Lemma~\ref{lem:induction-fs-four} gives
\begin{equation}\label{eq:forty-four-norm-congruence}
 \operatorname{Tr}E-\operatorname{Tr}D
                 \equiv 4m_C^2(1+n^2)\pmod{128}.
\end{equation}
For even $n$, $v_2(s)=-2-2v_2(n)$, so both traces are divisible
by $128$. Thus $m_C=0$.
\end{proof}
\begin{proposition}\label{prop:forty-eight-norm}
In the ordinary $(8,8)$ case of Proposition~\ref{prop:forty-eight-columns} one has $\gamma=1$ and $v_2(k)=3$. The dyadic valuations of $D$ and $E$ agree (normalized by $v_{\mathfrak p}(\sqrt2)=1$) and belong to $\{9,10\}$.
\end{proposition}
\begin{proof}
Put $b=v_2(k)\ge3$ and
$P_H=\operatorname{Tr}_K H\,\operatorname{Tr}_K(1/H)/16$ for $H=D,E$.
When $d_H=5$, write
\[
  \xi^i=a+b'\sqrt2,\quad (2+\sqrt5)^r=u+v\sqrt5,\quad
  \theta^l=c+d'\sqrt{10}.
\]
Multiplying these expressions and their inverses gives
$P_H=5(-1)^l\big((avc)^2-(2b'ud')^2\big)$.
In particular $P_H\equiv5\pmod8$ when $l$ is even: then $i$ is
even, $r$ is odd and $a,v,c$ are odd.

At $\gamma=1$, Proposition~\ref{prop:forty-marked-scales} gives $b=3$
and a common valuation $8$ ($d=5$, $l$ even), $9$ ($d=2,10$) or
$10$ ($d=5$, $l$ odd). At $\gamma=2$,
Proposition~\ref{prop:forty-exponent-parity} makes $i,l$ even, whence
$d_D=d_E=5$. The signed evaluation $E/D$ of the ordinary
Frobenius--Perron induction class reduces to $\gamma$, hence has
positive valuation. The trace formulas in the proof of
Proposition~\ref{prop:forty-marked-scales} give
$v_2(h_D)=4$ and $v_2(h_E)=1+b$, so $b\ge4$.

By Lemma~\ref{lem:induction-fs-four} the two squared lengths of
$C=FA_i$ agree modulo eight. Write $k=2^bK_0$ and $F_m=2A$ with
$K_0,A$ odd. The coordinates in Proposition~\ref{prop:forty-eight-columns}
give the signed squared length
\[
  \frac{\operatorname{Tr}_K E}{64}
    +\frac{\operatorname{Tr}_F f}{16}
   =2^{b-3}K_0\left(\frac{P_E}{2^{b-1}K_0-A}+5A\right).
\]
At $\gamma=2$ the bracket vanishes modulo four, since
$P_E\equiv A^2\equiv1\pmod4$. The ordinary length squared is
$4P_D/3\equiv4\pmod8$, a contradiction. At $\gamma=1$ with
valuation eight the ordinary length squared is $P_D/3+1\equiv0\pmod8$.
The signed length squared is
$K_0(5(4K_0-A)+5A)\equiv4\pmod8$, since odd residues modulo eight
are self-inverse. This leaves only valuations nine and ten.
\end{proof}

\subsection{Counting square roots in the two local orders}

By Lemma~\ref{lem:induction-fs-four} the lattice isometry identifies
the solutions in $R/2R$ and $T/2T$ of each marked equation
$bx^2=y\pmod4$. Here $C$ is a forgotten induction column of smallest ordinary
complementary codegree, $b=1,C$, and $y$ lies in the integral span
of the unit and the columns $C,W_j,W_l,W_{jl}$.
\begin{proposition}\label{prop:forty-local-orders}
For $(4,4n^2)$ with $n$ odd and $v_2(k)=2$, the common valuation
$m=v(D)=v(E)$ is neither six, seven, nor eight. For $(8,8)$ with
$v_2(k)=3$ and $\gamma=1$, it cannot be nine.
\end{proposition}
\begin{proof}
Let $F_2=\mathbb Q_2(\sqrt5)$, $\pi=\sqrt2$, and write $A$ for either
completed order. Its normalization is
$O=\mathcal O_{F_2}[\pi]\times B$, with $B=\mathbb Z_2^2$ in the
ordinary sector and $B=\mathcal O_{F_2}$ in the signed sector.
Let $c=2$ for $(4,4n^2)$ and $c=3$ for $(8,8)$.
Since all simples are self-dual, $A$ is self-dual for the
unit-coefficient trace. The quartic weight is $1/H$ ($H=D,E$), and the
complementary weights have denominator $2^c$ up to units. The
different of $F_2(\pi)$ is $(2\pi)$, since $F_2$ is unramified.
Lemma~\ref{lem:trace-order-conductor} thus gives
\[
 I=O^\#=\pi^{m-3}\mathcal O_{F_2}[\pi]\times2^cB
 \subseteq A\subseteq O,\qquad [O:A]=[A:I]=2^{m-3+c}.
\]
All square tests below take place in $O/4I$.

\smallskip\noindent\emph{The first valuation can be counted by hand.}
When $m=6$, put $\overline O=\mathbb F_4[\pi]/(\pi^2)\times S$,
where $S=\mathbb F_2^2$ or $\mathbb F_4$, and
$U=(A+2O)/2O$, $V=\{\bar z:2z\in A\}$. Here $4O\subset A$.
The reduction of $4\tau$ on $(X+Y\pi,s)$ is
$\ell=\operatorname{Tr}_{\mathbb F_4/\mathbb F_2}(X)+
\operatorname{Tr}_{S/\mathbb F_2}(s)$, with the sum trace in the split
case. Indeed $8/H$ is one modulo $\pi$ by
Corollary~\ref{prop:forty-golden-period} and Proposition~\ref{prop:forty-marked-scales},
and the complementary weights of $4\tau$ are one modulo two.
Thus $V\subset\ker\ell$ and the index gives $\dim U+\dim V=7$.
Locality gives $U=\mathbb F_2\mathbf1+(\pi W,0)$ with
$W\subset\mathbb F_4$. Squaring a lift of $(\pi e,0)$ shows
$(e^2,0)\in V$, hence $\operatorname{Tr}(e^2)=0$.
Consequently $W=\mathbb F_2$ and $V=\ker\ell$.

A solution of $x^2=1\pmod{4A}$ has the form $x=1+2y$ with $y\in O$.
Indeed both factors of $(x-1)(x+1)$ would have valuation one in the
quartic component if $v(x-1)=1$, and the other components are
unramified. Membership is then equivalent to $\bar y\in V$ and
$y^2+y\in A$. Every $b\in V$ with $b^2+b\in U$ has such a lift:
changing $y$ by $2z$ changes $y^2+y$ by $2z$ modulo $4O$.
Two successful lifts of the same $b$ differ by an element of $A$.
The number of solutions is therefore
$|\{b\in V:b^2+b\in U\}|/|U|$. Write $b=(x_0+x_1\pi,s)$. One has $x_1\in\mathbb F_2$ in both
sectors. In the ordinary sector $x_0\in\mathbb F_2$ and $s=(t,t)$.
In the signed sector $x_0,s\in\mathbb F_4$ have the same trace.
There are thus $8$ and $16$ admissible $b$, respectively. Division
by $|U|=4$ gives two ordinary and four signed square roots.

\smallskip\noindent\emph{Valuation eight.}
Here $n$ is odd, $d_H=5$, $v_2(h_H)=4$ and $i_H,l_H$ are odd
by Proposition~\ref{prop:forty-marked-scales}. Thus $h=H/16$ is
one modulo $2\mathcal O_{K,2}$: the two units $\xi,\theta$ each
reduce to $1+\pi$, and $\sqrt5,\phi^{3r}$ reduce to one. Also
$128\mid\operatorname{Tr}H$ in both sectors. Equation~\eqref{eq:forty-four-norm-congruence}
and $v_2(1+n^2)=1$ give $m_C=0$. Consequently $C=(4h,0)$
in either normalization.

We compare $x^2=1+C\pmod{4A}$. Any solution has $x=1+2y$ in
$O$, as at valuation six. Write $[a+b\pi]_0=a$ for $a,b\in\mathcal O_{F_2}$,
and $\operatorname{tr}$ for the trace from $F_2$ to $\mathbb Q_2$
or its reduction from $\mathbb F_4$ to $\mathbb F_2$. In the ordinary
order $2y\in A$ forces $\operatorname{tr}(\bar y_K)=0$, where the bar
denotes reduction modulo $\pi$: the quartic contribution to $4\tau(2y)$
is $\operatorname{tr}([y_K/h]_0)$ and both rational contributions are even.
The residues of $y^2+y-C/4$ are therefore one on the quartic branch
and zero on the rational branches. By locality this element is outside
$A$, contradicting the square equation.

In the signed order put $u=4/f\equiv1\pmod2$. Hensel's lemma gives
$y_K=2r$ with $y_K^2+y_K=1+h$ and $y_F$ with $y_F^2+y_F=1$;
the latter has residue of trace one. We claim $\operatorname{tr}(\bar r)=1$.
Write $h=1+2e$. The first equation gives $\bar r=1+\bar e$.
Normalization gives
\[
 8=8\tau(1)=\operatorname{tr}([h^{-1}]_0)+2\operatorname{tr}(u).
\]
Since $h^{-1}\equiv1-2e\pmod4$ and $\operatorname{tr}(u)$ is even,
$\operatorname{tr}(\bar e)=1$, proving the claim. For $(w,z)\in A$,
locality gives $\bar w=\bar z\in\mathbb F_2$. Therefore
\[
 2\tau\big((4r,2y_F)(w,z)\big)
 =\operatorname{tr}([rw/h]_0)+\operatorname{tr}(u y_Fz)
 \equiv \bar w\big(\operatorname{tr}(\bar r)+\operatorname{tr}(\bar y_F)\big)=0\pmod2.
\]
Self-duality gives $2y=(4r,2y_F)\in A$. Now $x=1+2y$ satisfies
$x^2=1+C+4\mathbf1$, a signed solution to the same equation.

\smallskip\noindent\emph{Valuations seven and nine.}
Let $j$ negate $\sqrt5$ and $l$ negate $\sqrt2$, each fixing the other
radical. We count solutions of the following equations
in $A/2A$, with all right sides read modulo $4A$:
\[
\begin{array}{c|lll}
m&\multicolumn{3}{l}{\text{equations}}\\ \hline
7&x^2=2C-W_j-W_l-W_{jl}&x^2=2-W_l&Cx^2=2C\\
9&x^2=C-W_l&x^2=-1-C&Cx^2=2W_j+2W_{jl}-C .
\end{array}
\]
In each row the sets of possible counts are disjoint:
\[
\begin{array}{c|l|l}
m&R&T\\ \hline
7&(0,0,16),(0,2,0),(2,0,0),(2,2,0)&(0,0,0),(0,2,16)\\
9&(0,0,8),(1,0,8),(1,0,16),(1,1,0)&(0,0,0),(0,0,16),(0,1,8)
\end{array}
\]

Use coordinates $1,\phi,\pi,\pi\phi$ and the two coordinates of
$B$, with $\phi^2=1+\phi$. Work modulo $16$ at $m=7$ and $32$
at $m=9$; these moduli kill $O/4I$. Write
$v=\pi^m/H$. Modulo the indicated coordinate modulus its possible
values have the form
\[
 a(\sqrt5)^e\xi^i\phi^h\theta^u,\qquad a\text{ odd},\quad
 e=0,1,\quad i+u\equiv1,\quad h+u+e\equiv0\pmod2,\quad 3\mid h.
\]
These are the signs and the residue condition of
Lemma~\ref{lem:forty-unit-param} and Corollary~\ref{prop:forty-golden-period}
applied to $1/H$.
At every valuation the signed weight $2^c/f$ ranges through odd
multiples of $\sqrt5\phi^h$ with $h\equiv3\pmod6$. On the ordinary
side run over odd $n$ when $c=2$. In each loop all powers can be
taken through their periods in the finite unit group. The four
columns are obtained from the ratios of Galois conjugates and from
Proposition~\ref{prop:forty-four-column}.
Retain the normalization $\tau(1)=1$ to the available precision.

The remaining search is intrinsic to $O/I$. Enumerate its unital
subrings of order $2^{m-3+c}$ containing the four columns, with
residue field $\mathbb F_2$ and $\tau(A)\subset\mathbb Z_2$.
Here $I=2^cO$. Use the filtration
$U_d=(2^{-d}(A\cap2^dO)+2O)/2O$. The ambient residue algebra $O/2O$ has the perfect pairing
$\lambda(xy)$, where
$\lambda(X+Y\pi,s)=\operatorname{Tr}(X+Y)+\operatorname{Tr}(s)$,
since $\pi^m/H\equiv1+\pi$ and $2^c/f\equiv1$ modulo two.
Self-duality gives $U_{c-1-d}=U_d^\perp$, and locality and
orthogonality to $1$ give $U_0=\langle1\rangle$ or
$\langle1,(\pi,0)\rangle$. For $c=3$ the intermediate $U_1$ is
a totally isotropic three-space containing $U_0$ and the
characteristic vector $(\pi,1_S)$. Lift a basis of $U_d/U_{d-1}$
at each level $d$, taking higher digits at level $a$ modulo $U_a$.
This gives every lattice; retain those closed under products.
For $m=9$ the signed side must also satisfy
$\tau(C^2)\equiv3\pmod4$ by Lemma~\ref{lem:induction-fs-four}.

Each retained lattice has just $2^6$ classes in $A/2A$. Substitution
of these classes in the displayed equations gives the table, as
verified for
\codecheck{field_checks/forty/quadratic40.py}{valuation seven} and
\codecheck{field_checks/forty/primitive9.py}{valuation nine}.
\end{proof}

\begin{corollary}\label{cor:odd-forty-excluded}
In the ordinary complementary pair $(4,4n^2)$ the integer $n$ is even. In the pair $(8,8)$ both global dimensions have valuation ten at the prime above two, normalized by $v_{\mathfrak p}(\sqrt2)=1$.
\end{corollary}
\begin{proof}
If $n$ is odd, Proposition~\ref{prop:forty-marked-scales} gives $v_2(k)=2$ and the same valuation of $D$ and $E$. It is six or eight for the seed $d_D=5$, and seven for $d_D=2,10$. Proposition~\ref{prop:forty-local-orders} excludes all three valuations. For $(8,8)$, Proposition~\ref{prop:forty-eight-norm} leaves nine or ten, and Proposition~\ref{prop:forty-local-orders} excludes nine.
\end{proof}

\subsection{The last two ordinary pairs}

\begin{theorem}\label{thm:forty-even-excluded}
The ordinary complementary codegrees $4,4n^2$ with even $n$ do not
occur in the residual conductor-forty case.
\end{theorem}
\begin{proof}
Write $H=h_H\sqrt{d_H}\xi^{i_H}\phi^{3r_H}\theta^{l_H}$
for $H=D,E$. Put $e=v_2(n)\ge1$, $s=\operatorname{Tr}_K(1/D)$
and $t=\operatorname{Tr}_K(1/E)$. The ordinary complementary terms
give $s=(3n^2-1)/(4n^2)$, and hence $v_2(s)=-2-2e$.
Let $j$ negate $\sqrt5$ and fix $\sqrt2$, and let $u$ negate
$\sqrt2$ and fix $\sqrt5$. Apply \eqref{eq:forty-trace-transfer} with $h=4n^2$ to the signed simultaneous
units for $j$ and $ju$. Their complementary value $\phi^{2m}$ is the
odd $(m/3)$th power of $9+4\sqrt5$, so its trace is two modulo
sixteen. Therefore the quartic traces are four modulo sixteen.
The trace calculations in Proposition~\ref{prop:forty-exponent-parity}
give
\[
 \tfrac14\operatorname{Tr}_K(E/jE)\equiv3^{l_E},\qquad
 \tfrac14\operatorname{Tr}_K(E/juE)\equiv3^{i_E}\pmod4.
\]
Both exponents are even, and the same holds for $D$ by that proposition.
The exclusion of $d=1$ in Proposition~\ref{prop:forty-marked-scales}
and total positivity now force $d_D=d_E=5$, $r_D,r_E$ odd. Furthermore
the same proposition gives $b=v_2(k)=3+2e$, $v_2(t)=-2-2e$ and
$v_2(h_D)=v_2(h_E)=4+2e$.

Let $W_j$ be the simultaneous induction unit for $j$ and put
\[
       P=W_j+W_j^{-1}-2,\qquad q=2^{2e-2}.
\]
We compare the unit coefficients of $qP$ in the two orders. They must
agree modulo two because the based reductions and the images of
$W_j$ agree there.

The ordinary complementary values of $P$ vanish. Its signed
complementary value $(\phi^m-\phi^{-m})^2$ equals $16$ times an
odd square: an odd power of $2+\sqrt5$ has rational part two
modulo four and has norm $-1$. Thus it is $16\pmod{128}$.
On a quartic factor set $A=\phi^{3r_H}$, $B=\theta^{l_H}$ and
$\Lambda=\mathbb Z_2[\sqrt2,\sqrt5]$. The value $P_H$ of $P$
is $(AB-A^{-1}B^{-1})^2$. As $r_H$ is odd, $(A-A^{-1})/4$
is an odd integer and $A^{-1}\equiv\sqrt5\pmod{2\Lambda}$.
As $l_H$ is even and $\theta^2=19+6\sqrt{10}$, we have
$B\equiv1\pmod{2\Lambda}$ and $(B-B^{-1})/4\in\sqrt{10}\mathbb Z_2$.
Consequently
\[
 \frac{AB-A^{-1}B^{-1}}4
   =\frac{A-A^{-1}}4 B+A^{-1}\frac{B-B^{-1}}4
     \equiv 1+c\sqrt2\pmod{2\Lambda}
\]
for some $c\in\mathbb F_2$. Squaring gives $P_H\in16(1+2\Lambda)$.

Since $h_H/H\in\Lambda$ and $\operatorname{Tr}_K\Lambda\subset4\mathbb Z_2$,
the error $q\operatorname{Tr}_K((P_H-16)/H)$ has valuation at least
$(2e-2)+5+2-(4+2e)=1$. The signed complementary error is also even,
since $v_2(1-t)=-2-2e$. Thus the ordinary unit coefficient of $qP$
is $16qs$ modulo two, a dyadic unit. The signed
coefficient is $16qt+16q(1-t)=16q$ modulo two. This contradiction
holds for every $e\ge1$.
\end{proof}

\begin{proposition}\label{prop:forty-ten}
The codegree pair $(8,8)$ cannot occur.
\end{proposition}
\begin{proof}
By Propositions~\ref{prop:forty-eight-norm} and~\ref{prop:forty-local-orders}
only the common ramified valuation ten remains. Then $v=32/D$ is an
integral dyadic unit, and $v\equiv1\pmod{2\mathcal O_{K,2}}$.
Indeed in the notation of Proposition~\ref{prop:forty-eight-norm} one has
$d=5$, both $i,l$ odd, and the units $\xi,\theta$ are $1+\epsilon$
modulo two, where $\epsilon=\sqrt2$ and $\epsilon^2=0$ modulo two.
The factors $\sqrt5,\phi^{3r}$ both reduce to one.

Write $v=v_0+\epsilon v_1$ with $v_i\in\mathcal O_{F,2}$ and
$F=\mathbb Q(\sqrt5)$. For an element
$x=(a+b\phi+\epsilon y,r,s)$ of the ordinary completed order we have
$r\equiv s\equiv a\pmod2$ by locality. Write $\operatorname{tr}$
for the trace on $F$. Normalization of the unit coefficient gives
$\operatorname{tr}(v_0)=12$, since $\tau(1)=\operatorname{tr}(v_0)/16+1/4$.
It follows that
\[
 16\tau(x)=a\operatorname{tr}(v_0)
    +b\operatorname{tr}(v_0\phi)+2\operatorname{tr}(v_1y)+2(r+s)
    \equiv b\operatorname{tr}(v_0\phi)\pmod4.
\]
The displayed trace factor is odd because $v_0\equiv1$ and
$\operatorname{tr}(\phi)=1$, so $4\mid b$. Now the codegree-eight
induction class is $C=(D/8,2,2)$, whence
$4\tau(Cx)=2a+b+r+s$. Thus $2a+r+s\equiv0\pmod4$, and as
$r\equiv a\pmod2$ this gives $s-r\equiv0\pmod4$. Both rational
coordinates would agree modulo four on every basis element, contradicting
the primitivity of $(r_1-r_2)/2$ in the proof of
Proposition~\ref{prop:forty-eight-columns}.
\end{proof}

\begin{proof}[Proof of Theorem~\ref{thm:forty-closed}]
In the absence of a spherical structure Theorem~\ref{prop:forty-divergent-complements} gives two possibilities for the complementary ordinary codegrees. Corollary~\ref{cor:odd-forty-excluded} and Theorem~\ref{thm:forty-even-excluded} exclude $(4,4n^2)$ for all $n$. Proposition~\ref{prop:forty-ten} excludes $(8,8)$.
\end{proof}

\begin{thebibliography}{99}

\bibitem{BJ} M.~Bischoff and C.~Jones, \href{https://doi.org/10.1007/s00029-021-00725-3}{\emph{Computing fusion rules for spherical $G$-extensions of fusion categories}}, Selecta Math. (N.S.) \textbf{28} (2022), Paper No.~26.

\bibitem{BN} S.~Burciu and S.~Natale, \href{https://doi.org/10.1063/1.4774293}{\emph{Fusion rules of equivariantizations of fusion categories}}, J. Math. Phys. \textbf{54} (2013), 013511.

\bibitem{DLN} C.~Dong, X.~Lin and S.-H.~Ng, \href{https://doi.org/10.2140/ant.2015.9.2121}{\emph{Congruence property in conformal field theory}}, Algebra Number Theory \textbf{9} (2015), 2121--2166.

\bibitem{DCW} J.~Dong, G.~Chen and Z.~Wang, \href{https://doi.org/10.1016/j.jalgebra.2022.02.020}{\emph{Fusion categories containing a fusion subcategory with maximal rank}}, J. Algebra \textbf{604} (2022), 107--127.

\bibitem{DPP} J.~Dong, S.~Palcoux and A.~Plessis, \emph{Grothendieck rings of pivotal fusion categories of rank four}, \href{https://arxiv.org/abs/2609.33272}{arXiv:2609.33272}.

\bibitem{EGNO} P.~Etingof, S.~Gelaki, D.~Nikshych and V.~Ostrik, \emph{Tensor categories}, Mathematical Surveys and Monographs, vol.~205, AMS, 2015.

\bibitem{ENO} P.~Etingof, D.~Nikshych and V.~Ostrik, \href{https://doi.org/10.4007/annals.2005.162.581}{\emph{On fusion categories}}, Ann. of Math. (2) \textbf{162} (2005), 581--642.

\bibitem{GSFields} T.~Gannon and A.~Schopieray, \href{https://doi.org/10.4310/CNTP.241202233757}{\emph{Algebraic number fields generated by dimensions in fusion rings}}, Commun. Number Theory Phys. \textbf{18} (2024), 705--743.

\bibitem{HaagGS} P.~Grossman and N.~Snyder, \emph{Quantum subgroups of the Haagerup fusion categories}, Comm. Math. Phys. \textbf{311} (2012), 617--643. \href{https://arxiv.org/abs/1102.2631}{arXiv:1102.2631}.

\bibitem{Rank4Sphere} \emph{Grothendieck rings of complex fusion categories: rank 4}, contributed by S.~Palcoux, Hexagon (2026). \href{https://hexagonmath.org/2611.00645v1}{hexagon:2611.00645v1}.

\bibitem{Class6Rank5} \emph{Grothendieck rings of complex fusion categories: rank 5}, contributed by S.~Palcoux, Hexagon (2026). \href{https://hexagonmath.org/2611.00986v1}{hexagon:2611.00986v1}.

\bibitem{HudsonWilliams} R.~H.~Hudson and K.~S.~Williams, \href{https://doi.org/10.1216/RMJM/1181073167}{\emph{The integers of a cyclic quartic field}}, Rocky Mountain J. Math. \textbf{20} (1990), 145--150.

\bibitem{KMRS} V.~Koppen, V.~Mulevi\v{c}ius, I.~Runkel and C.~Schweigert, \href{https://doi.org/10.1007/s00220-022-04489-5}{\emph{Domain walls between 3d phases of Reshetikhin--Turaev TQFTs}}, Comm. Math. Phys. \textbf{396} (2022), 1187--1220.

\bibitem{MuleviciusOrb} V.~Mulevi\v{c}ius, \emph{Condensation inversion and Witt equivalence via generalised orbifolds}, Theory Appl. Categ. \textbf{41} (2024), 1203--1292. \href{https://tac.mta.ca/tac/volumes/41/36/41-36abs.html}{Journal article}.

\bibitem{Narkiewicz} W. Narkiewicz, \emph{Elementary and Analytic Theory of Algebraic Numbers}, 3rd ed., Springer Monographs in Mathematics, Springer, Berlin, 2004.

\bibitem{NTV} S.~Nasseh, S.~Sather-Wagstaff, R.~Takahashi and K.~VandeBogert, \href{https://doi.org/10.1016/j.jpaa.2018.06.006}{\emph{Applications and homological properties of local rings with decomposable maximal ideals}}, J. Pure Appl. Algebra \textbf{223} (2019), 1272--1287.

\bibitem{NgSchInd} S.-H.~Ng and P.~Schauenburg, \href{https://doi.org/10.1007/s00220-010-1096-6}{\emph{Congruence subgroups and generalized Frobenius--Schur indicators}}, Comm. Math. Phys. \textbf{300} (2010), 1--46.

\bibitem{Ostrik2} V.~Ostrik, \emph{Fusion categories of rank~2}, Math. Res. Lett. \textbf{10} (2003), 177--183. \href{https://arxiv.org/abs/math/0203255}{arXiv:math/0203255}.

\bibitem{OstrikCodegrees} V.~Ostrik, \href{https://doi.org/10.4310/MRL.2009.v16.n5.a11}{\emph{On formal codegrees of fusion categories}}, Math. Res. Lett. \textbf{16} (2009), 895--901.

\bibitem{Ostrik3} V.~Ostrik, \href{https://doi.org/10.17323/1609-4514-2015-15-2-373-396}{\emph{Pivotal fusion categories of rank~3}}, Mosc. Math. J. \textbf{15} (2015), 373--396.

\bibitem{SchopExt} A.~Schopieray, \href{https://doi.org/10.1112/jlms.12789}{\emph{Categorification of integral group rings extended by one dimension}}, J. Lond. Math. Soc. (2) \textbf{108} (2023), 1617--1641.

\bibitem{SchopNP} A.~Schopieray, \emph{Non-pseudounitary fusion}, J. Pure Appl. Algebra \textbf{226} (2022), 106927. \href{https://arxiv.org/abs/2010.02958}{arXiv:2010.02958}.

\bibitem{SchweigertWoike} C.~Schweigert and L.~Woike, \href{https://doi.org/10.1016/j.aim.2023.109006}{\emph{Homotopy invariants of braided commutative algebras and the Deligne conjecture for finite tensor categories}}, Adv. Math. \textbf{422} (2023), article 109006.
\bibitem{SerreArithmetic} J.-P.~Serre, \href{https://doi.org/10.1007/978-1-4684-9884-4}{\emph{A course in arithmetic}}, Graduate Texts in Mathematics, vol.~7, Springer, 1973.

\end{thebibliography}
\end{document}
